
\documentclass[11pt]{article}
\usepackage[a4paper,margin=25mm]{geometry}
\usepackage[T1]{fontenc}
\usepackage[utf8]{inputenc}
\usepackage[english]{babel}
\usepackage{lmodern}
\usepackage{amsmath,amssymb,amsthm,mathtools,array,booktabs,graphicx,xcolor,enumitem,pdflscape}
\usepackage{microtype}
\usepackage[colorlinks=true,urlcolor=black,linkcolor=black]{hyperref}
\setlength{\parindent}{1.2em}
\setlength{\parskip}{0.15em}
\emergencystretch=3em
\allowdisplaybreaks
\hbadness=10000
\newcommand{\grp}{\mathfrak{H}}
\newcommand{\Hgrp}{\mathfrak{H}}
\newcommand{\Ggrp}{\mathfrak{G}}
\newcommand{\Egrp}{\mathfrak{E}}
\newcommand{\zchar}{\mathfrak{z}}
\newcommand{\Hgr}{\mathfrak{H}}
\newcommand{\Ggr}{\mathfrak{G}}
\newcommand{\Dgr}{\mathfrak{D}}
\newcommand{\Fgr}{\mathfrak{F}}
\newcommand{\FrobP}{\mathfrak{P}}
\newcommand{\Char}{\chi}
\newcommand{\eps}{\varepsilon}
\newcommand{\ten}{\mathrm{ten}}
\newcommand{\te}{\mathrm{te}}
\newcommand{\Class}[1]{#1^{\mathrm{te}}}
\newcommand{\class}[1]{\textsuperscript{#1}}
\newcommand{\Cl}[1]{#1^{\mathrm{ten}}}
\DeclareMathOperator{\rang}{Rang}
\newcommand{\resizeMath}[1]{\resizebox{\textwidth}{!}{$#1$}}
\newcommand{\half}{\tfrac{1}{2}}
\newcommand{\frobtitle}[1]{\begin{center}{\Large #1}\\[0.75em]By G. \textsc{Frobenius}.\end{center}\vspace{0.5em}}
\newcommand{\sign}{\operatorname{sign}}
\newcommand{\coeff}[1]{\left[#1\right]}
\newcommand{\CharTwo}[4]{\left(\begin{matrix}#1&\cdots&#2\\#3&\cdots&#4\end{matrix}\right)}
\newcommand{\ch}[2]{\left(\begin{smallmatrix}#1\\#2\end{smallmatrix}\right)}
\newcommand{\vtheta}{\vartheta}
\newcommand{\fnum}[1]{\textup{(\ensuremath{#1}.)}\quad}
\newcommand{\tr}{\operatorname{tr}}
\newcommand{\rank}{\operatorname{Rang}}
\newcommand{\grS}{\mathfrak{S}}
\newcommand{\modeta}{\pmod{\eta}}
\newcommand{\frobtag}[2]{\[\begin{array}{@{}r@{\quad}l@{}}\textup{(#1)} & \displaystyle #2\end{array}\]}
\newcommand{\frobtagc}[2]{\[\begin{array}{@{}r@{\quad}c@{}}\textup{(#1)} & \displaystyle #2\end{array}\]}
\begin{document}
\begin{center}{\huge Frobenius: Selected Works on Group Theory and Hypercomplex Quantities}\end{center}
\bigskip

% ===== BEGIN WORK 053: 053_ber_Gruppencharaktere =====
\begin{center}
{\Large \textbf{On group characters.}}\\[0.8em]
{\large By G. \textsc{Frobenius}.}\\[0.6em]
\rule{6em}{0.4pt}
\end{center}

In the proof of the theorem that every linear function of one variable represents infinitely many prime numbers, if its coefficients are relatively prime integers, \textsc{Dirichlet} for the first time used certain systems of roots of unity. These also occur in the closely related question concerning the number of ideal classes in a cyclotomic field (compare \textsc{Dedekind}'s remark in \textsc{Dirichlet}'s lectures on number theory, 4th ed., p. 625), as well as in the generalization of that theorem to quadratic forms and in the investigations concerning their division into genera. According to \textsc{Dedekind}, the characteristic property of these expressions consists in this: they are quantities $\chi(n)$ depending on a variable positive integer $n$, have only a finite number of values, and satisfy the condition
\[
\chi(m)\chi(n)=\chi(mn).
\]
As he carries this out in a purely abstract form, to the elements $A,B,C,\ldots$ of every finite group $\grp$ of permutable elements (an \textsc{Abel}ian group) one can assign such roots of unity $\chi(A),\chi(B),\chi(C),\ldots$ which satisfy the equations
\[
\chi(A)\chi(B)=\chi(AB),
\]
and which, following \textsc{Gauss}, he called the \emph{characters of the group}.

By a \emph{character} of a quadratic form, \textsc{Gauss}, \emph{Disqu. arithm.} Art. 230, understands a relation of the numbers representable by the form to the odd prime numbers $p$ (or to 4 or 8) entering into its determinant. He expresses that relation by the signs $Rp$ and $Np$. \textsc{Dirichlet}, \emph{Recherches sur diverses applications de l'analyse infinitésimale à la théorie des nombres}, \S~3 (\textsc{Crelle}'s Journal, vol. 19), replaces these symbols by the \textsc{Legendre} (and \textsc{Jacobi}) symbol $\left(\frac{m}{p}\right)$, which, next to \textsc{Lagrange}'s resolvent, probably offers the oldest example of the use of characters of commutative groups


known here. The advantage of this transformation consists in the fact that the genus characters of \textsc{Gauss} are only relations, whereas those of \textsc{Dirichlet} are numbers with which one can calculate. Thus multiplication of these characteristic numbers replaces the composition of the characters corresponding to the composition of the genera (Arts. 246 to 248)\footnote{I take the remarks contained in this paragraph from a letter of \textsc{Dedekind}'s dated 8 July 1896.}.

The number $h$ of characters of an \textsc{Abel}ian group $\grp$ is equal to the order of the group. One can obtain the $h$ characters by representing the elements of $\grp$ by means of a basis of independent elements and assigning to the basis elements arbitrary roots of unity whose degree is equal to their order. They can be assigned to the elements of the group (in several ways), and therefore can be denoted by $\chi_B(A)$. Since the product of two characters is again a character, they form a group, and this group is isomorphic with $\grp$. Their relations to the subgroups of $\grp$ were discussed most fully by \textsc{Weber} (\emph{Theorie der \textsc{Abel}'schen Zahlkörper}, I. \S~3, IV. \S~2 and 3, Acta Math. vols. 8 and 9).

In April of this year \textsc{Dedekind} communicated to me a problem which had occurred to him in the year 1880, and which, because it belonged both to group theory and to determinant theory, would in his opinion probably interest me, whereas a closer examination of it would draw him himself too far away from his arithmetical investigations. Its solution, which I hope soon to be able to communicate, led me to a generalization of the concept of characters to arbitrary finite groups. I wish here to develop this concept, in the belief that through its introduction group theory may receive an essential promotion and enrichment. The theory of characters acquires a further special interest through its remarkable relations to the theory of complex quantities formed from several principal units.

\section*{\S~1.}

Two elements $A$ and $B$ of a finite group $\grp$ are called \emph{conjugate} (with respect to $\grp$) if there is in $\grp$ an element $T$ satisfying the condition $T^{-1}AT=B$. If two elements are conjugate to a third, then they are also conjugate to one another. Hence the $h$ elements of $\grp$ can be divided into \emph{classes of conjugate elements}, a division of which I have repeatedly made advantageous use, especially in my new proof


of \textsc{Sylow}'s theorem, \textsc{Crelle}'s Journal vol. 100. The principal element $E$ forms a class by itself, the \emph{principal class}. It shall be denoted by $(0)$, and the remaining classes by $(1),(2),\ldots,(k-1)$, where $k$ is the number of classes. If $A$ is any element of the $\alpha$th class, then the elements of $\grp$ that are permutable with $A$ form a group contained in $\grp$. If $\frac{h}{h_\alpha}$ is its order, then $h_\alpha$ is the number of different elements of the $\alpha$th class; thus $h_0=1$. Since every element of $\grp$ belongs to one and only one of these $k$ classes, we have
\begin{equation*}\tag{1.}
\sum_\alpha h_\alpha=h.
\end{equation*}

The class number $k$ can also be obtained in the following way. Let each of the two variable elements $R$ and $S$ run independently through the $h$ elements of $\grp$, and count how often $SR=RS$ is true. If for $R$ one puts a fixed element of the $\alpha$th class, then there are $\frac{h}{h_\alpha}$ elements $S$ permutable with $R$. Thus, if one puts for $R$ in turn each of the $h_\alpha$ elements of the $\alpha$th class, one obtains $\frac{h}{h_\alpha}h_\alpha$ solutions of that equation. Since this number is the same for each of the $k$ classes, $hk$ is the number of solutions of the equation $SR=RS$.

As $A$ runs through the $h_\alpha$ elements of the $\alpha$th class, $A^{-1}$ also runs through all elements of a class. This class shall be called the \emph{inverse class} of $(\alpha)$ and denoted by $(\alpha')$. Inverse classes contain equally many elements,
\begin{equation*}\tag{2.}
h_{\alpha'}=h_\alpha.
\end{equation*}
If, as for example with the principal class, $(\alpha')=(\alpha)$, then the $\alpha$th class is called \emph{two-sided}.

Let $(\alpha),(\beta),(\gamma)$ be any three different or equal classes. Let $A$ run through the $h_\alpha$ different elements of the $\alpha$th class, $B$ through the $h_\beta$ elements of the $\beta$th, and $C$ through the $h_\gamma$ elements of the $\gamma$th. The number $h_{\alpha\beta\gamma}$ (which may also be zero) shall indicate how many of the $h_\alpha h_\beta h_\gamma$ elements $ABC$ are equal to the principal element, that is, satisfy the equation
\begin{equation*}\tag{3.}
ABC=E.
\end{equation*}
Since then $AB=C^{-1}$, $h_{\alpha\beta\gamma}$ also states how many of the $h_\alpha h_\beta$ elements $AB$ belong to the class $(\gamma')$. Equation (3.) is identical with $BCA=E$ and $CAB=E$. Hence $h_{\alpha\beta\gamma}$ of the $h_\beta h_\gamma$ products $BC$ are contained in $(\alpha')$, and $h_{\alpha\beta\gamma}$ of the $h_\gamma h_\alpha$ products $CA$ are contained in $(\beta')$. Therefore $h_{\alpha\beta\gamma}$ is not greater than the smallest of the three numbers $h_\beta h_\gamma$, $h_\gamma h_\alpha$, and $h_\alpha h_\beta$.


But the two elements $AB$ and
\begin{equation*}\tag{4.}
BA=A^{-1}(AB)A=B(AB)B^{-1}
\end{equation*}
are conjugate. Therefore they either both belong to the group $(\gamma')$, or both do not. Hence $h_{\beta\alpha\gamma}=h_{\alpha\beta\gamma}$. In connection with the above remarks it follows that the number $h_{\alpha\beta\gamma}$ remains unchanged under every interchange of the three indices. Since equation (3.) also implies $C^{-1}B^{-1}A^{-1}=E$, we have
\begin{equation*}\tag{5.}
h_{\alpha'\beta'\gamma'}=h_{\alpha\beta\gamma}.
\end{equation*}

If in equation (3.) one puts for $A$ a fixed element of the $\alpha$th class, while $B$ and $C$ remain variable as above, suppose that it has $m$ solutions. If $A'=T^{-1}AT$ is any other fixed element of the $\alpha$th class, and one puts $B'=T^{-1}BT$ and $C'=T^{-1}CT$, then also $A'B'C'=E$. Now $B'$ runs simultaneously with $B$ through the $h_\beta$ elements of the $\beta$th class, and $C'$ simultaneously with $C$ through the $h_\gamma$ elements of the $\gamma$th class. Therefore this equation also has $m$ solutions. Thus, if in (3.) $A$ is taken successively to be the $h_\alpha$ elements of the $\alpha$th class, it has altogether $h_\alpha m=h_{\alpha\beta\gamma}$ solutions. Likewise that equation has $\frac{h_{\alpha\beta\gamma}}{h_\beta}$ solutions if $B$ is a fixed element of the $\beta$th class, while $A$ and $C$ are variable elements of the $\alpha$th and $\gamma$th classes. Hence $h_{\alpha\beta\gamma}$ is divisible by each of the three numbers $h_\alpha,h_\beta,h_\gamma$, and therefore also by their least common multiple.

If $A$ is taken successively as the $h_\alpha$ elements of the $\alpha$th class and $B$ as the $h_\beta$ elements of the $\beta$th class, one obtains $h_\alpha h_\beta$ elements $AB$, which need not be distinct. Each of them belongs either to the class $(0)$ or to the class $(1)\ldots$ or to the class $(k-1)$. Therefore
\begin{equation*}\tag{6.}
\sum_\gamma h_{\alpha\beta\gamma}=h_\alpha h_\beta.
\end{equation*}

The consideration made for three classes can be carried out in the same way for arbitrarily many classes. Let, for example, $(\alpha),(\beta),(\gamma),(\delta)$ be four different or equal classes, and let $h_{\alpha\beta\gamma\delta}$ be the number of solutions of the equation
\begin{equation*}\tag{7.}
ABCD=E,
\end{equation*}
when $A$ (respectively $B,C,D$) runs through the $h_\alpha$ (respectively $h_\beta,h_\gamma,h_\delta$) elements of the class $(\alpha)$ (respectively $(\beta),(\gamma),(\delta)$). Then also $BCDA=CDAB=DABC=E$, so that $h_{\alpha\beta\gamma\delta}$ remains unchanged under cyclic permutation of the indices. Further, $BCD=A^{-1}$, $CDA=B^{-1}$, and so on. Thus $h_{\alpha\beta\gamma\delta}$ tells how many of the $h_\beta h_\gamma h_\delta$ elements $BCD$ belong to the class $(\alpha')$, or how many of the $h_\gamma h_\delta h_\alpha$ elements $CDA$ belong to the class $(\beta')$. Consequently $h_{\alpha\beta\gamma\delta}$ is not greater than the smallest of the four numbers $h_\beta h_\gamma h_\delta$, $h_\gamma h_\delta h_\alpha$,


$h_\delta h_\alpha h_\beta$, $h_\alpha h_\beta h_\gamma$. If $B$ is a fixed element of the $\beta$th class, while $A,C,D$ run through the $\alpha$th, $\gamma$th, and $\delta$th classes, then equation (7.) has only $\frac{h_{\alpha\beta\gamma\delta}}{h_\beta}$ solutions. Accordingly $h_{\alpha\beta\gamma\delta}$ is divisible by each of the four numbers $h_\alpha,h_\beta,h_\gamma,h_\delta$. Finally,
\begin{equation*}\tag{8.}
\sum_\delta h_{\alpha\beta\gamma\delta}=h_\alpha h_\beta h_\gamma.
\end{equation*}

If $A,B,C,D$ is a solution of equation (7.), and one puts
\[
B^{-1}AB=A',\qquad BA'B^{-1}=A,
\]
then $BA'CD=E$. Since $A'$ belongs to the $\alpha$th class just as $A$ does, one can in this way put the solutions of this equation and of equation (7.) into one-to-one correspondence. Hence $h_{\beta\alpha\gamma\delta}=h_{\alpha\beta\gamma\delta}$. In this way one recognizes that $h_{\alpha\beta\gamma\delta}$ remains unchanged under every interchange of the four indices. Therefore also
\begin{equation*}\tag{9.}
h_{\alpha'\beta'\gamma'\delta'}=h_{\alpha\beta\gamma\delta}.
\end{equation*}

For $n$ classes $(\alpha),(\beta),(\gamma),\ldots,(\nu)$ the number $h_{\alpha\beta\gamma\ldots\nu}$ has the same meaning and analogous properties. For the case $n=1$ this notation is not to be used, but it is still to be used for $n=2$. Thus $h_{\alpha\beta}$ is the number of solutions of the equation $AB=E$ or $B=A^{-1}$. Hence, if $(\alpha)$ and $(\beta)$ are not inverse classes, then $h_{\alpha\beta}=0$. But if $(\beta)=(\alpha')$, then
\begin{equation*}\tag{10.}
h_{\alpha\alpha'}=h_\alpha=h_{\alpha'}.
\end{equation*}
If $\nu=0$, then $h_{\alpha\beta\gamma\ldots\mu0}$ is the number of solutions of the equation $ABC\ldots ME=E$. Therefore
\begin{equation*}\tag{11.}
h_{\alpha\beta\gamma\ldots\mu0}=h_{\alpha\beta\gamma\ldots\mu}.
\end{equation*}
In particular,
\begin{equation*}\tag{12.}
h_{\alpha\beta0}=h_{\alpha\beta},
\end{equation*}
that is, zero unless $(\beta)=(\alpha')$, but then equal to $h_\alpha$.

\section*{\S~2.}

The most important property of the numbers $h_{\alpha\beta\gamma}$ is obtained by decomposing the equation $ABCD=E$ into the two equations
\[
AB=L^{-1},\qquad CD=L.
\]
If for $L$ one puts a fixed element of the $\lambda$th class, then the first equation has $\frac{h_{\alpha\beta\lambda}}{h_\lambda}$, and the second $\frac{h_{\lambda'\gamma\delta}}{h_\lambda}$, solutions. Thus there arise $\frac{h_{\alpha\beta\lambda}h_{\lambda'\gamma\delta}}{h_\lambda^2}=m_\lambda$ solutions of equation (7.) of \S~1. If $L$ is taken successively as each of the $h_\lambda$ elements of the $\lambda$th class, one obtains




$h_\lambda m_\lambda$ solutions, and if finally one lets $(\lambda)$ be each of the $k$ classes, one sees that in all it has
\begin{equation*}\tag{1.}
h_{\alpha\beta\gamma\delta}=\sum_\lambda \frac{1}{h_\lambda}h_{\alpha\beta\lambda}h_{\lambda'\gamma\delta}
\end{equation*}
solutions. In the same way one obtains
\begin{equation*}\tag{2.}
h_{\alpha\beta\ldots\xi\eta\ldots\sigma\tau}=\sum_\lambda \frac{1}{h_\lambda}h_{\alpha\beta\ldots\xi\lambda}h_{\lambda'\eta\ldots\sigma\tau}.
\end{equation*}
Taking into account also equations (10.) and (12.) of \S~1, it follows that all the numbers $h_{\alpha\beta\gamma\ldots}$ can be composed from the numbers $h_{\alpha\beta\gamma}$.

Consequently the sum on the right-hand side of equation (1.) remains unchanged, not only when $\alpha$ is interchanged with $\beta$, or $\gamma$ with $\delta$, but also when $\beta$ is interchanged with $\delta$, or when any permutation of the four quantities $\alpha,\beta,\gamma,\delta$ is made. The same conclusions are used by \textsc{Gauss}, \emph{Theoria residuorum biquadraticorum}, Comm. prima, \S~17 (Collected Works vol. II, p. 81). Thus, if for a moment one sets
\[
\frac{1}{h_\alpha}h_{\alpha\beta\gamma}=a_{\alpha\beta\gamma},
\]
then
\begin{equation*}\tag{3.}
a_{\alpha\beta\gamma}=a_{\alpha\gamma\beta},\qquad
\sum_\lambda a_{\alpha\lambda\gamma}a_{\lambda\beta\delta}=\sum_\lambda a_{\alpha\lambda\delta}a_{\lambda\beta\gamma}.
\end{equation*}
Therefore one can apply to these quantities the theorems developed by \textsc{Weierstrass} and \textsc{Dedekind} in their papers \emph{On the theory of complex quantities formed from $n$ principal units}, Göttingen Nachrichten 1884 and 1885, and which I, since they form the basis of this investigation, derived again and generalized in my paper \emph{On permutable matrices} (p. 601 of this volume): if the determinant of degree $k$ formed from the $k^2$ quantities
\[
p_{\alpha\beta}=\sum_{\varkappa,\lambda}a_{\varkappa\lambda\alpha}a_{\lambda\varkappa\beta}=\sum_{\varkappa,\lambda}a_{\varkappa\varkappa\lambda}a_{\lambda\alpha\beta}
\]
is different from zero, then the equations
\[
r_\beta r_\gamma=\sum_\alpha a_{\alpha\beta\gamma}r_\alpha
\]
have exactly $k$ distinct solutions $r_\alpha=r_\alpha^{(\varkappa)}$, and the determinant of degree $k$ formed from these solutions is different from zero. If $x_0,x_1,\ldots,x_{k-1}$ are variables, and if one puts $a_{\alpha\beta}=\sum_\gamma a_{\alpha\beta\gamma}x_\gamma$, then the determinant of degree $k$ is
\[
\left|a_{\alpha\beta}-r\epsilon_{\alpha\beta}\right|=\prod_\varkappa\left(r_0^{(\varkappa)}x_0+\cdots+r_{k-1}^{(\varkappa)}x_{k-1}-r\right).
\]
If $(s_\alpha^{(\varkappa)})$ is the complementary system to $(r_\alpha^{(\varkappa)})$, then the ratios of the $k$ quantities
\[
s_0=s_0^{(\varkappa)},\qquad s_1=s_1^{(\varkappa)},\ldots\qquad s_{k-1}=s_{k-1}^{(\varkappa)}
\]


are completely determined by the linear equations
\[
s_\alpha r=\sum_\beta a_{\alpha\beta}s_\beta,
\]
provided one puts $r=r_0^{(\varkappa)}x_0+\cdots+r_{k-1}^{(\varkappa)}x_{k-1}$ in them (pp. 613 and 614).

In the present case,
\[
p_{\alpha\beta}=\sum_{\varkappa,\lambda}\frac{h_{\varkappa\lambda\alpha}h_{\lambda\varkappa\beta}}{h_\varkappa h_\lambda}.
\]
Since $(\varkappa')$ runs through all $k$ classes at the same time as $(\varkappa)$, one may also replace $(\varkappa)$ by $(\varkappa')$ in this sum, and then, by (5.) of \S~1, obtains
\begin{equation*}\tag{4.}
p_{\alpha\beta'}=\sum_{\varkappa,\lambda}\frac{h_{\varkappa\lambda\alpha}h_{\varkappa\lambda\beta}}{h_\varkappa h_\lambda},
\end{equation*}
and hence
\begin{equation*}\tag{5.}
p_{\alpha\beta'}=p_{\beta\alpha'}=p_{\alpha'\beta}=p_{\beta'\alpha}.
\end{equation*}
The determinant of degree $k$, $|p_{\alpha\beta'}|$, is obtained from $|p_{\alpha\beta}|$ by an interchange among the columns, and therefore differs from it at most by sign. Now consider the $k$ quantities
\[
\frac{h_{\varkappa\lambda0}}{\sqrt{h_\varkappa h_\lambda}},\qquad
\frac{h_{\varkappa\lambda1}}{\sqrt{h_\varkappa h_\lambda}},\ldots\qquad
\frac{h_{\varkappa\lambda,k-1}}{\sqrt{h_\varkappa h_\lambda}}.
\]
Putting for $\varkappa$ and for $\lambda$ the values $0,1,\ldots,k-1$, one obtains a system of $k^3$ quantities arranged in $k^2$ rows and $k$ columns. From this system one can form a determinant of degree $k$ by choosing any $k$ rows, hence in $\binom{k^2}{k}$ ways. By the general multiplication theorem of determinant theory, the sum of the squares of all these determinants is, according to equation (4.), equal to $|p_{\alpha\beta'}|$. Put $\varkappa=0$ and $\lambda'=0,1,\ldots,k-1$; then one obtains a definite system of $k$ rows. In the determinant formed from their elements, the diagonal elements $\frac{h_{0\lambda'\lambda}}{\sqrt{h_\lambda}}=\sqrt{h_\lambda}$ are different from zero, and the remaining elements are zero. Consequently the determinant of degree $k$, $|p_{\alpha\beta}|$, is different from zero.

Let $f$ be a proportionality factor, provisionally left undetermined, and let
\[
r_\alpha=\frac{h_\alpha\chi_\alpha}{f}.
\]
Then the equations
\begin{equation*}\tag{6.}
h_\beta h_\gamma\chi_\beta\chi_\gamma=f\sum_\alpha h_{\alpha'\beta\gamma}\chi_\alpha
\end{equation*}
have $k$ different systems of solutions
\begin{equation*}\tag{7.}
\chi_\alpha=\chi_\alpha^{(\varkappa)},\qquad f=f^{(\varkappa)}\qquad(\varkappa=0,1,\ldots,k-1),
\end{equation*}


and the determinant of degree $k$
\begin{equation*}\tag{8.}
\left|\chi_\alpha^{(\varkappa)}\right|
\end{equation*}
is different from zero. Later I shall make a definite provision concerning the choice of the factor $f$, that is, of the $k$ factors $f^{(\varkappa)}$, and shall then call the quantities $\chi_\alpha^{(\varkappa)}$ the $k$ \emph{characters of the group $\grp$}. A character $\chi$ is a system of $k$ numbers $\chi_0,\chi_1,\ldots,\chi_{k-1}$ corresponding to the $k$ classes $(0),(1),\ldots,(k-1)$. The $\varkappa$th character $\chi^{(\varkappa)}$ is the system of the $k$ numbers
\[
\chi_0=\chi_0^{(\varkappa)},\qquad \chi_1=\chi_1^{(\varkappa)},\ldots\qquad \chi_{k-1}=\chi_{k-1}^{(\varkappa)}\qquad(\varkappa=0,1,\ldots,k-1).
\]
The meaning of equations (6.) follows from the following consideration. Let $x_0,x_1,\ldots,x_{k-1}$ and $y_0,y_1,\ldots,y_{k-1}$ be $2k$ independent variables, and let
\begin{equation*}\tag{9.}
h_\gamma z_\gamma=\sum_{\alpha,\beta}h_{\alpha\beta\gamma}x_\alpha y_\beta
\end{equation*}
be $k$ bilinear functions of them. Then
\[
f\sum_\gamma h_\gamma\chi_\gamma z_\gamma
=f\sum_{\alpha,\beta,\gamma}h_{\alpha\beta\gamma}\chi_\gamma x_\alpha y_\beta
=\sum_{\alpha,\beta}h_\alpha\chi_\alpha x_\alpha\,h_\beta\chi_\beta y_\beta,
\]
and hence
\begin{equation*}\tag{10.}
f\left(\sum_\gamma h_\gamma\chi_\gamma z_\gamma\right)=
\left(\sum_\alpha h_\alpha\chi_\alpha x_\alpha\right)
\left(\sum_\beta h_\beta\chi_\beta y_\beta\right).
\end{equation*}
I now define a linear function $\xi^{(\varkappa)}$ of $x_0,x_1,\ldots,x_{k-1}$ corresponding to the character $\chi^{(\varkappa)}$ by the equation
\begin{equation*}\tag{11.}
\sum_\alpha h_\alpha\chi_\alpha^{(\varkappa)}x_\alpha=f^{(\varkappa)}\xi^{(\varkappa)}\qquad(\varkappa=0,1,\ldots,k-1).
\end{equation*}
Then the determinant of degree $k$ is
\begin{equation*}\tag{12.}
\left|\sum_\gamma h_{\alpha\beta'\gamma}x_\gamma-h_{\alpha\beta'}r\right|=
\prod_\varkappa h_\varkappa\left(\xi^{(\varkappa)}-r\right).
\end{equation*}
If one puts $x_1=\cdots=x_{k-1}=0$ and $x_0=1$, then consequently
\[
(1-r)^k=\prod_\varkappa\left(\frac{\chi_0^{(\varkappa)}}{f^{(\varkappa)}}-r\right)
\]
and therefore
\begin{equation*}\tag{13.}
\chi_0^{(\varkappa)}=f^{(\varkappa)},\qquad \chi_0=f,
\end{equation*}
since, in equations which hold for every value of the upper index $\varkappa$, I shall also omit it. (Compare \textsc{Dedekind}'s discussions of many-valued systems of quantities, loc. cit., p. 144.) Thus equation (12.) is no more general than the equation
\begin{equation*}\tag{14.}
\left|\sum_\gamma h_{\alpha\beta'\gamma}x_\gamma\right|=\prod_\varkappa h_\varkappa\xi^{(\varkappa)},
\end{equation*}
because it follows from this equation by replacing $x_0$ by $x_0-r$. Writing the former in the form
\[
\left|\sum_\gamma \frac{h_{\alpha\beta'\gamma}}{h_\alpha}x_\gamma-\frac{h_{\alpha\beta'}}{h_\alpha}r\right|=\prod_\varkappa(\xi^{(\varkappa)}-r),
\]


we have $\frac{h_{\alpha\beta'}}{h_\alpha}$ equal to 1 in the diagonal elements and 0 in the remaining elements. Further, the quantities $\frac{h_{\alpha\beta'\gamma}}{h_\alpha}$ are integers. Therefore if all variables $x_\gamma=0$ except one, it shows that the quantities
\begin{equation*}\tag{15.}
\frac{h_\alpha\chi_\alpha}{f}
\end{equation*}
are algebraic integers.

\section*{\S~3.}

The elements of each row of the system $(s_\alpha^{(\varkappa)})$ complementary to $(r_\alpha^{(\varkappa)})$ are determined, up to a common factor, by the equations
\[
\sum_\beta \frac{h_{\alpha'\beta\gamma}}{h_\alpha}s_\beta=\frac{h_\gamma\chi_\gamma}{f}s_\alpha.
\]
By virtue of the symmetry properties of the quantities $h_{\alpha\beta\gamma}$, these equations are satisfied, according to (6.) of \S~2, by the values $s_\beta=\chi_{\beta'}$. Hence $s_\beta=\frac{e}{h}\chi_{\beta'}$, where $e$ is a new proportionality factor. Accordingly the two systems
\begin{equation*}\tag{1.}
\left(\frac{h_\alpha\chi_\alpha^{(\varkappa)}}{f^{(\varkappa)}}\right),\qquad
\left(\frac{e^{(\varkappa)}\chi_{\alpha'}^{(\varkappa)}}{h}\right)
\end{equation*}
are complementary. Thus the equations
\begin{equation*}\tag{2.}
\sum_\alpha h_\alpha\chi_\alpha^{(\varkappa)}\chi_{\alpha'}^{(\lambda)}=0,
\end{equation*}
hold if $\varkappa$ and $\lambda$ are different, while
\begin{equation*}\tag{3.}
\sum_\alpha h_\alpha\chi_\alpha^{(\varkappa)}\chi_{\alpha'}^{(\varkappa)}=\frac{h f^{(\varkappa)}}{e^{(\varkappa)}}.
\end{equation*}
Further,
\begin{equation*}\tag{4.}
\sum_\varkappa\frac{e^{(\varkappa)}}{f^{(\varkappa)}}\chi_\alpha^{(\varkappa)}\chi_\beta^{(\varkappa)}=\frac{h h_{\alpha\beta}}{h_\alpha h_\beta},
\end{equation*}
and hence, by (13.) of \S~2, for $\beta=0$,
\begin{equation*}\tag{5.}
\sum_\varkappa e^{(\varkappa)}\chi_\alpha^{(\varkappa)}=0,
\end{equation*}
if $\alpha$ is different from 0, while for $\alpha=0$
\begin{equation*}\tag{6.}
\sum_\varkappa e^{(\varkappa)}f^{(\varkappa)}=h.
\end{equation*}
I also write equation (3.) in the simpler form
\begin{equation*}\tag{7.}
\sum_\alpha h_\alpha\chi_\alpha\chi_{\alpha'}=\frac{hf}{e}
\end{equation*}
and equation (2.) in the form
\begin{equation*}\tag{8.}
\sum_\alpha h_\alpha\chi_\alpha\psi_{\alpha'}=0,
\end{equation*}
where $\psi_\alpha$ denotes a character different from $\chi_\alpha$.


By equation (5.) of \S~1 and (6.) of \S~2,
\[
h_\alpha h_\gamma\chi_\alpha\chi_{\gamma'}=f\sum_\beta h_{\alpha\beta\gamma'}\chi_\beta=f\sum_\beta h_{\alpha'\beta\gamma}\chi_\beta.
\]
Multiplying by $\chi_\gamma$ and summing with respect to $\gamma$, one obtains from (7.) the formula
\begin{equation*}\tag{9.}
\frac{h}{e}h_\alpha\chi_\alpha=\sum_{\beta,\gamma}h_{\alpha'\beta\gamma}\chi_\beta\chi_\gamma,
\end{equation*}
which for $\alpha=0$ agrees with (7.). In the same way one finds the formula
\begin{equation*}\tag{10.}
\sum_{\beta,\gamma}h_{\alpha\beta\gamma}\chi_\beta\psi_\gamma=0,
\end{equation*}
which for $\alpha=0$ passes into (8.).

Eliminating from equations (6.) of \S~2 and (9.) the products $\chi_\beta\chi_\gamma$, one obtains the linear equations
\begin{equation*}\tag{11.}
\frac{h}{ef}h_\alpha\chi_\alpha=\sum_\beta p_{\alpha\beta'}\chi_\beta.
\end{equation*}
Therefore the $k$ quantities $g^{(\varkappa)}=e^{(\varkappa)}f^{(\varkappa)}$ are the roots of the equation of degree $k$
\begin{equation*}\tag{12.}
\left|g p_{\alpha\beta'}-h h_{\alpha\beta'}\right|=0.
\end{equation*}
By equation (4.) of \S~2,
\begin{equation*}\tag{13.}
\sum p_{\alpha\beta'}x_\alpha x_\beta=\sum_{\varkappa,\lambda}\frac{1}{h_\varkappa h_\lambda}\left(\sum_\alpha h_{\varkappa\lambda\alpha}x_\alpha\right)^2.
\end{equation*}
This quadratic form, whose determinant is different from zero, is therefore a positive definite form. Hence the $k$ roots $g$ of equation (12.) are all real positive quantities, and the elementary divisors of that determinant are all linear. For an $m$-fold root $g=ef$, therefore, its subdeterminants of degree $(k-m+1)$ and higher also vanish. Accordingly the rank of the system of linear equations (11.) is equal to $k-m$, and they possess $m$ independent solutions. The $m$ characters $\chi_\alpha$ corresponding to an $m$-fold root $g=ef$ can then first be represented as linear combinations of such $m$ independent solutions, and, when these are substituted in equations (6.) of \S~2, one finds the coefficients of these combinations by solving an equation of degree $m$. The equation of degree $k$ originally found for the determination of the $k$ characters, namely (12.) of \S~2, therefore decomposes, after equation (12.) has been solved, into as many factors as this equation has distinct roots; and to an $m$-fold root of equation (12.) there corresponds a factor of degree $m$.

As regards equation (12.) itself, however, it can further be shown that its roots are not merely real and positive, but


that they are integers (\S~6), and even squares of integers that divide $h$. Up to now, however, I have not succeeded in proving these theorems by considerations as simple as the preceding results.

The factors $f^{(\varkappa)}$ have so far been assumed arbitrary, but different from zero. By means of formulas (3.) the factors $e^{(\varkappa)}$ are obtained from them. Since $g^{(\varkappa)}=e^{(\varkappa)}f^{(\varkappa)}$ is a real positive quantity, $f^{(\varkappa)}$ too shall be chosen real and positive. The same then holds for $e^{(\varkappa)}$. One arrives at the simplest formulas by setting $f=\sqrt{g}$. Then $e=f=\chi_0$. As already mentioned, $e$ is then a positive integer dividing $h$, and it can be shown that $\chi_\alpha$ is a sum of $f$ roots of unity of degree $n$, if $n$ denotes the order of the elements of the $\alpha$th class.

By equation (1.) and (6.) of \S~1 and (4.) of \S~2,
\begin{equation*}\tag{14.}
\sum_\beta p_{\alpha\beta'}=h h_\alpha.
\end{equation*}
Now, by (6.) of \S~1, equations (6.) of \S~2 are satisfied by setting all quantities $\chi_\alpha=f$. The corresponding value of $ef$ is 1. Hence one may set $e=f=1$, and therefore $\chi_\alpha=1$. This character may be called the principal character. If, in equation (8.), this character is chosen for $\psi$, one sees that every other character satisfies the equation
\begin{equation*}\tag{15.}
\sum_\alpha h_\alpha\chi_\alpha=0.
\end{equation*}

Since the determinant (8.) of \S~2 is different from zero, equation (12.) of \S~2 has no multiple roots. Nor, if $\chi$ and $\psi$ are two different characters, can the quantities $\chi_0,\chi_1,\ldots,\chi_{k-1}$ be proportional to the quantities $\psi_0,\psi_1,\ldots,\psi_{k-1}$. Likewise every class $(\alpha)$ is completely determined by the $k$ corresponding values $\chi_\alpha^{(\varkappa)}$; if $\alpha$ and $\beta$ are different, it cannot be that $\chi_\alpha^{(\varkappa)}=\chi_\beta^{(\varkappa)}$ for all values of $\varkappa$. Since the coefficients of equation (12.) of \S~2 are real, then, if $x_0,x_1,\ldots,x_{k-1}$ are real variables, to every complex root there corresponds a conjugate complex root. Thus, if $\chi_\alpha$ and $\chi'_\alpha$ are conjugate complex quantities, then to every character $\chi_\alpha$ there corresponds a conjugate complex character $\chi'_\alpha$. But it must be that $\chi'_\alpha=\chi_{\alpha'}$. For first, by (6.) of \S~2,
\[
h_\beta h_\gamma\chi_\beta\chi_\gamma=f\sum_\alpha h_{\alpha\beta\gamma}\chi_\alpha=f\sum_\alpha h_{\alpha'\beta\gamma}\chi_{\alpha'},
\]
and consequently $\psi_\alpha=\chi_{\alpha'}$ is a character. Were this different from $\chi'_\alpha$, then by (8.) one would have $\sum h_\alpha\chi'_{\alpha}\psi_{\alpha'}=0$, hence $\sum h_\alpha\chi_\alpha\chi'_{\alpha}=0$, whereas each term, as the product of two conjugate complex quantities,


is positive. In every character, therefore, inverse classes correspond (if $f$ is assumed real) to conjugate complex values $\chi_\alpha$ and $\chi_{\alpha'}$, and hence to a two-sided class there corresponds a real value. From every complex character $\chi_0,\chi_1,\ldots,\chi_{k-1}$ one obtains a conjugate complex character, and this is $\chi_{0'},\chi_{1'},\ldots,\chi_{(k-1)'}$. A real character is also called \emph{two-sided}; two conjugate complex characters are also called \emph{inverse} (compare \textsc{Weber}, \emph{Proof of the theorem that every properly primitive quadratic form is capable of representing infinitely many prime numbers}. Math. Ann. vol. 20, p. 308).

\section*{\S~4.}

If one multiplies equation (6.) of \S~2,
\[
h_\beta h_\gamma\chi_\beta^{(\varkappa)}\chi_\gamma^{(\varkappa)}=f^{(\varkappa)}\sum_\lambda h_{\lambda\beta\gamma}\chi_\lambda^{(\varkappa)}
\]
by $\frac{e^{(\varkappa)}\chi_\alpha^{(\varkappa)}}{f^{(\varkappa)2}}$ and then sums with respect to $\varkappa$, one obtains
\begin{equation*}\tag{1.}
\frac{h h_{\alpha\beta\gamma}}{h_\alpha h_\beta h_\gamma}=\sum_\varkappa \frac{e^{(\varkappa)}}{f^{(\varkappa)2}}\chi_\alpha^{(\varkappa)}\chi_\beta^{(\varkappa)}\chi_\gamma^{(\varkappa)},
\end{equation*}
and likewise, in general, for $n$ indices $\alpha,\beta,\gamma\ldots\nu$, the formula holds
\begin{equation*}\tag{2.}
\frac{h h_{\alpha\beta\gamma\ldots\nu}}{h_\alpha h_\beta h_\gamma\cdots h_\nu}
=\sum_\varkappa \frac{e^{(\varkappa)}}{f^{(\varkappa)n-1}}\chi_\alpha^{(\varkappa)}\chi_\beta^{(\varkappa)}\chi_\gamma^{(\varkappa)}\cdots\chi_\nu^{(\varkappa)},
\end{equation*}
which for $n=2$ passes into (4.) of \S~3. Hence, by (11.) of \S~2,
\begin{equation*}\tag{3.}
h\sum_{\alpha,\beta,\gamma,\ldots,\nu}h_{\alpha\beta\gamma\ldots\nu}x_\alpha x_\beta x_\gamma\cdots x_\nu=\sum_\varkappa e^{(\varkappa)}f^{(\varkappa)}\xi^{(\varkappa)n}.
\end{equation*}
By (4.) of \S~2 and (5.) of \S~1,
\[
p_{\alpha\beta}=\sum_{\varkappa,\lambda}\frac{h_{\lambda\varkappa\alpha}h_{\varkappa'\lambda'\beta}}{h_\varkappa h_\lambda},
\]
and therefore, if the summation with respect to $\varkappa$ is carried out by means of formula (1.) of \S~2,
\begin{equation*}\tag{4.}
p_{\alpha\beta}=\sum_\lambda \frac{1}{h_\lambda}h_{\alpha\beta\lambda'},
\end{equation*}
and consequently, by (3.) of \S~3 and (2.),
\begin{equation*}\tag{5.}
\frac{p_{\alpha\beta}}{h_\alpha h_\beta}=\sum_\varkappa \frac{\chi_\alpha^{(\varkappa)}\chi_\beta^{(\varkappa)}}{f^{(\varkappa)2}}.
\end{equation*}
Thus
\begin{equation*}\tag{6.}
\sum_{\alpha,\beta}p_{\alpha\beta}x_\alpha x_\beta=\sum_\varkappa(\xi^{(\varkappa)})^2
\end{equation*}
and
\begin{equation*}\tag{7.}
\sum_{\alpha,\beta}p_{\alpha\beta'}x_\alpha x_\beta=\sum_\varkappa \xi^{(\varkappa)}\xi^{(\varkappa)'}.
\end{equation*}



Since $\xi^{(\varkappa)}$ and $\xi^{(\varkappa)'}$ are conjugate complex quantities, this formula again shows that this form is positive.

The coefficients of equation (12.), \S~3, which serves for determining the quantities $g=ef$, can also be computed by determining the power sums of its roots. If one regards $\frac{h}{g}$ as the unknown, the sum of its roots is
\[
\begin{aligned}
\sum_\varkappa \frac{h}{e^{(\varkappa)}f^{(\varkappa)}}
&=\sum_\alpha \frac{p_{\alpha\alpha'}}{h_\alpha}
=\sum_{\alpha,\beta}\frac{h_{\alpha\alpha'\beta\beta'}}{h_\alpha h_\beta} \\
&=\sum_{\alpha,\beta,\gamma}\frac{h_{\alpha\beta\gamma}h_{\alpha'\beta'\gamma'}}{h_\alpha h_\beta h_\gamma}
=\sum_{\alpha,\beta,\gamma}\frac{h_{\alpha\beta\gamma}^{2}}{h_\alpha h_\beta h_\gamma}.
\end{aligned}
\]
If one composes the system $\frac{p_{\alpha\beta'}}{h_\alpha}$ with itself, one obtains
\[
\sum_\varkappa \frac{p_{\alpha\varkappa'}p_{\varkappa\beta'}}{h_\alpha h_\varkappa}
=\sum_{\varkappa,\lambda,\mu}\frac{h_{\alpha\varkappa\lambda\lambda'}h_{\varkappa\beta'\mu\mu'}}{h_\alpha h_\varkappa h_\lambda h_\mu}
=\sum_{\lambda,\mu}\frac{h_{\alpha\beta\lambda\lambda'\mu\mu'}}{h_\alpha h_\lambda h_\mu}
=\frac{1}{h_\alpha}\sum_{\varkappa,\lambda,\mu}\frac{h_{\alpha\varkappa\lambda\mu}h_{\beta\varkappa\lambda\mu}}{h_\varkappa h_\lambda h_\mu}.
\]
Therefore the sum of the squares of the roots is
\[
\sum_\varkappa\left(\frac{h}{e^{(\varkappa)}f^{(\varkappa)}}\right)^2
=\sum_{\alpha,\beta,\gamma}\frac{h_{\alpha\alpha'\beta\beta'\gamma\gamma'}}{h_\alpha h_\beta h_\gamma}
=\sum_{\alpha,\beta,\gamma,\delta}\frac{h_{\alpha\beta\gamma\delta}^{2}}{h_\alpha h_\beta h_\gamma h_\delta}.
\]
By composing the system $n-2$ times with itself, one finds in the same way, for the sum of the $(n-2)^{\text{nd}}$ powers of its roots,
\begin{equation*}\tag{8.}
\sum_\varkappa\left(\frac{h}{e^{(\varkappa)}f^{(\varkappa)}}\right)^{n-2}
=\sum_{\alpha,\beta,\gamma\ldots\mu}\frac{h_{\alpha\alpha'\beta\beta'\ldots\mu\mu'}}{h_\alpha h_\beta\cdots h_\mu}
=\sum_{\alpha,\beta,\gamma\ldots\mu,\nu}\frac{h_{\alpha\beta\cdots\mu\nu}^{2}}{h_\alpha h_\beta\cdots h_\mu h_\nu},
\end{equation*}
where $n$ is the number of the indices $\alpha,\beta,\ldots\mu,\nu$. (Compare \textsc{Borchardt}, Crelle's Journal vol.~30, p.~38.) One can also reduce the $k^2$ quantities $p_{\alpha\beta}$ to the $k$ quantities
\begin{equation*}\tag{9.}
p_{\alpha0}=p_\alpha=p_{\alpha'}=\sum_\lambda\frac{h_{\alpha\lambda\lambda'}}{h_\lambda}=h_\alpha\sum_\varkappa\frac{\chi_\alpha^{(\varkappa)}}{f^{(\varkappa)}}
\end{equation*}
(denoted by \textsc{Dedekind}, loc. cit. p.~147 (25.), as $\sigma_r$). For, by (1.), \S~2 and (4.),
\[
p_{\alpha\beta}=\sum_\lambda\frac{1}{h_\lambda}\left(\sum_\gamma\frac{1}{h_\gamma}h_{\alpha\beta\gamma}h_{\gamma\lambda\lambda'}\right),
\]
and hence
\begin{equation*}\tag{10.}
p_{\alpha\beta}=\sum_\gamma\frac{1}{h_\gamma}h_{\alpha\beta\gamma}p_\gamma.
\end{equation*}
Thus equation (12.), \S~3 passes into equation (12.), \S~2 when one sets $x_\gamma=\frac{p_\gamma}{h_\gamma}$ and $\xi=\frac{h}{g}$. Indeed, by (11.), \S~3 and (13.), \S~2,
\begin{equation*}\tag{11.}
\sum_\alpha p_\alpha\chi_\alpha=\frac{h}{e}.
\end{equation*}
Thus if $f$ is chosen so that $\chi_0,\chi_1,\ldots\chi_{k-1}$ are algebraic integers,


then $e$ is a divisor of $h$. Formula (3.) gives, for these values of the variables,
\begin{equation*}\tag{12.}
\sum_\varkappa\left(\frac{h}{e^{(\varkappa)}f^{(\varkappa)}}\right)^{n-1}
=\sum_{\alpha,\beta,\gamma\ldots\nu}h_{\alpha\beta\gamma\ldots\nu}\frac{p_\alpha p_\beta p_\gamma\cdots p_\nu}{h_\alpha h_\beta h_\gamma\cdots h_\nu}.
\end{equation*}
The number $p_\alpha$ has the following meaning: If $A$ is a fixed element of the $\alpha^{\text{th}}$ class, and if each of the two variable elements $R$ and $S$ runs independently of the other through the $h$ elements of $\grp$, then $\frac{hp_\alpha}{h_\alpha}$ is the number of solutions of the equation
\begin{equation*}\tag{13.}
SR=RSA.
\end{equation*}
For, if $A$ is a fixed element of the $\alpha^{\text{th}}$ class, if $R^{-1}$ first runs through the $h_\lambda$ elements of the class $(\lambda)$, and if $R'$ runs through the $h_\lambda$ elements of the inverse class $(\lambda')$, then the equation $R^{-1}R'A=E$ has exactly $\frac{h_{\alpha\lambda\lambda'}}{h_\alpha}$ solutions.

If, however, for $R'$ one puts each element of the class $(\lambda')$ $\frac{h}{h_\lambda}$ times, then it has $\frac{h h_{\alpha\lambda\lambda'}}{h_\alpha h_\lambda}$ solutions. This is achieved by setting, for each $R$, $R'=S^{-1}RS$ and then letting $S$ run through all $h$ elements of $\grp$. Finally, if $\lambda$ also moves from $0$ to $k-1$, then
\[
\sum_\lambda\frac{h h_{\alpha\lambda\lambda'}}{h_\alpha h_\lambda}=\frac{hp_\alpha}{h_\alpha}
\]
is the number of solutions of the equation $R^{-1}(S^{-1}RS)A=E$, that is, of equation (13.). In particular $p_0=k$, and therefore $hk$ is the number of solutions of the equation $SR=RS$.

Similarly, more generally, $hp_{\alpha\beta}$ is the number of solutions of the equation
\begin{equation*}\tag{14.}
SR=RSAB,
\end{equation*}
when $A$ runs through the $h_\alpha$ elements of the $\alpha^{\text{th}}$ class, $B$ through the $h_\beta$ elements of the $\beta^{\text{th}}$ class, and $R$ and $S$ through all $h$ elements of $\grp$. If, however, one takes for $A$ a fixed element of the $\alpha^{\text{th}}$ class, then that equation has $\frac{hp_{\alpha\beta}}{h_\alpha}$ solutions. With the help of formulas (5.) and (6.), \S~3, one can also prove this and similar propositions by calculation.

\section*{\S~5.}

For the further development of this theory it is advantageous to replace the symbol $\chi_\alpha$, when $A$ is an element of the $\alpha^{\text{th}}$ class, by $\chi_A$, or also, using $\chi$ as a function symbol, by $\chi(A)$. Then, since this quantity has the same value for all elements of the $\alpha^{\text{th}}$ class,
\begin{equation*}\tag{1.}
\chi(B^{-1}AB)=\chi(A),
\end{equation*}


or, if one replaces $A$ by $BA$,
\begin{equation*}\tag{2.}
\chi(AB)=\chi(BA).
\end{equation*}
In particular $\chi(E)=f$. I likewise denote the variables $x_\alpha$ and $y_\alpha$ by $x_A$ and $y_A$, and the constants $h_\alpha$ and $p_\alpha$ by $h_A$ and $p_A$. For all these quantities the equation analogous to (2.) holds. Further, $h_{A^{-1}}=h_A$, while $\chi(A^{-1})$ and $\chi(A)$ are conjugate complex quantities.

Equation (6.), \S~2, from which the ratios of the $\chi_\alpha$ are obtained, can then be written in the form
\begin{equation*}\tag{3.}
h_B\chi(A)\chi(B)=f\sum_S'\chi(AS)
\end{equation*}
where $S$ runs through the $h_B$ elements conjugate with $B$. For if $(\alpha)$ and $(\beta)$ are the classes of $A$ and $B$, then for $S$ one must put the $h_\beta$ elements of the $\beta^{\text{th}}$ class. Since $A$ is a fixed element of the $\alpha^{\text{th}}$ class, it occurs $\frac{h_{\alpha\beta\gamma'}}{h_\alpha}$ times that $AS$ belongs to the $\gamma^{\text{th}}$ class, and hence $\chi(AS)=\chi_\gamma$. Thus $\sum_S'\chi(AS)=\frac{1}{h_\alpha}\sum_\gamma h_{\alpha\beta\gamma'}\chi_\gamma$.

Since $\chi(E)=f$ and $\chi(S)=\chi(B)$, one can also put equation (3.) in the form
\begin{equation*}\tag{4.}
\sum_S'\left(\chi(E)\chi(AS)-\chi(A)\chi(S)\right)=0
\end{equation*}
where $S$ runs through the elements of a class. Put $S=R^{-1}BR$ and let $R$ run through all $h$ elements of $\grp$; then $S$ becomes equal to every element of the $\beta^{\text{th}}$ class, each exactly $\frac{h}{h_\beta}$ times. Hence
\begin{equation*}\tag{5.}
h\chi(A)\chi(B)=f\sum_R\chi(AR^{-1}BR).
\end{equation*}
By repeated application of this relation one obtains
\begin{equation*}\tag{6.}
\left(\frac{h}{f}\right)^{n-1}\chi(A_1)\chi(A_2)\cdots\chi(A_n)=\sum'\chi(A_1R_1A_2R_2\cdots A_nR_n),
\end{equation*}
where $R_1,R_2,\ldots R_n$ run through all systems of $n$ elements of the group $\grp$ which satisfy the equation
\[
R_1R_2\cdots R_n=E.
\]
In the same way one obtains the formula
\begin{equation*}\tag{7.}
p_{\alpha\beta}=\sum_{R,S}'\frac{p_{RS}}{h_{RS}},
\end{equation*}
where $R$ runs through the $h_\alpha$ elements of the $\alpha^{\text{th}}$ class and $S$ through the $h_\beta$ elements of the $\beta^{\text{th}}$.

If $R$ and $S$ run through all elements that satisfy the condition $RS=A$, then by (9.) and (10.), \S~3,
\begin{equation*}\tag{8.}
\frac{h}{e}\chi(A)=\sum\chi(R)\chi(S),\qquad 0=\sum\chi(R)\psi(S)\qquad(RS=A).
\end{equation*}
For, if $A$ is a fixed element of the $\alpha^{\text{th}}$ class, then while $R$


runs through the $\beta^{\text{th}}$ class and $S$ the $\gamma^{\text{th}}$ class, it occurs $\frac{h_{\alpha'\beta\gamma}}{h_\alpha}$ times that $RS=A$. This can also be written
\begin{equation*}\tag{9.}
\frac{h}{e}\chi(AB^{-1})=\sum_R\chi(AR^{-1})\chi(RB^{-1}),\qquad 0=\sum_R\chi(AR^{-1})\psi(RB^{-1}),
\end{equation*}
where $R$ runs through all elements of $\grp$. In particular,
\begin{equation*}\tag{10.}
\frac{hf}{e}=\sum_R\chi(R^{-1})\chi(R),\qquad 0=\sum_R\chi(R^{-1})\psi(R).
\end{equation*}
The $h$ variables $x_R$ (respectively $y_R$) reduce, as a consequence of the conditions
\begin{equation*}\tag{11.}
x_{AB}=x_{BA},\qquad y_{AB}=y_{BA}
\end{equation*}
to only $k$ independent variables each; they have the same value for all elements $R$ which belong to the same class. From them I form the two systems, each of $h^2$ elements,
\begin{equation*}\tag{12.}
(x_{PQ^{-1}}),\qquad (y_{PQ^{-1}}).
\end{equation*}
The $h$ rows of the first system are obtained by putting for $P$ the $h$ elements of $\grp$ in any order, and the $h$ columns by putting for $Q$ the same elements in the same order. The system has certain symmetry properties determined by $\grp$, in consequence of which there are only $k$ distinct elements among its $h^2$ elements. The system formed by composing the two systems (12.) has the same symmetry properties and is moreover independent of the order in which the composition is carried out. It is
\begin{equation*}\tag{13.}
z_{PQ^{-1}}=\sum_R x_{PR^{-1}}y_{RQ^{-1}}=\sum_S y_{PS^{-1}}x_{SQ^{-1}}.
\end{equation*}
For if one sets $PR^{-1}=SQ^{-1}$, then $S=PR^{-1}Q$; hence $S$ runs, simultaneously with $R$, through all $h$ elements of $\grp$, and $S^{-1}P=Q^{-1}R$, so $y_{PS^{-1}}=y_{S^{-1}P}=y_{Q^{-1}R}=y_{RQ^{-1}}$. If, moreover, one replaces $R$ by $RP$, one sees that the first sum depends only on the product $PQ^{-1}$. At first one would have been allowed to denote it only by $z_{P,Q}$, but now may set it equal to a quantity $z_R=z_{PQ^{-1}}$ depending on only one element $R=PQ^{-1}$. Finally,
\[
z_{AB}=\sum x_{AR^{-1}}y_{RB}=\sum y_{AS^{-1}}x_{SB}=\sum x_{BS}y_{S^{-1}A}=\sum x_{BR^{-1}}y_{RA}=z_{BA}.
\]
For $R=S^{-1}$ runs, simultaneously with $S$, through all $h$ elements of $\grp$.

Any two systems (12.) which possess the symmetry properties presupposed here are commutable with one another. If one puts $x_\varkappa=1$ and the remaining $k-1$ of the variables $x_0,x_1,\ldots x_{k-1}$ equal to zero, one obtains $k$ special systems of this kind for $\varkappa=0,1,\ldots k-1$. Of these $k$ systems, therefore, any two are commutable. From them the most general system again arises by multiplying them by the $k$ variables


$x_0,x_1,\ldots x_{k-1}$ and adding. Consequently the determinant of degree $h$
\begin{equation*}\tag{14.}
\left|x_{PQ^{-1}}\right|=\Theta(x_0,x_1,\ldots x_{k-1})=\Theta((x))
\end{equation*}
is a product of $h$ linear functions of the $k$ variables (\emph{Über vertauschbare Matrizen}, p.~601 of this volume). In this determinant the $h$ elements on the diagonal, and only these, are equal to $x_0$; hence $x_0^h$ has coefficient 1. Therefore one may stipulate that also in each linear factor of $\Theta$
\begin{equation*}\tag{15.}
\xi=\frac{1}{f}\sum h_\alpha\chi_\alpha x_\alpha
\end{equation*}
the coefficient of $x_0$ is equal to 1. Here $h_\alpha$ is to have the same meaning as before, $f$ is to be different from zero, and $\chi_0,\chi_1,\ldots\chi_{k-1}$ are to denote the still unknown coefficients. Accordingly $\chi_0=f$. But now, by (13.),
\[
\left|z_{PQ^{-1}}\right|=\left|x_{PQ^{-1}}\right|\left|y_{PQ^{-1}}\right|.
\]
Consequently each linear factor of $\Theta((z))$ must also split into the product of a linear function of $x_0,x_1,\ldots x_{k-1}$ and one of $y_0,y_1,\ldots y_{k-1}$,
\[
\frac{1}{f}\sum h_\gamma\chi_\gamma z_\gamma=\left(\sum a_\alpha x_\alpha\right)\left(\sum b_\beta y_\beta\right),
\]
where $a_0=b_0=1$. If one sets $y_0=y_E=1$ and $y_1=\cdots=y_{k-1}=0$, then $z_R=x_R$, and therefore $\frac{1}{f}\sum h_\gamma\chi_\gamma x_\gamma=\sum a_\alpha x_\alpha$. Hence
\begin{equation*}\tag{16.}
f\sum h_\gamma\chi_\gamma z_\gamma=\left(\sum h_\alpha\chi_\alpha x_\alpha\right)\left(\sum h_\beta\chi_\beta y_\beta\right).
\end{equation*}
But equation (13.), or
\begin{equation*}\tag{17.}
z_C=\sum x_Ry_S\qquad(RS=C)
\end{equation*}
is identical with
\begin{equation*}\tag{18.}
h_\gamma z_\gamma=\sum_{\alpha,\beta}h_{\alpha\beta\gamma'}x_\alpha y_\beta.
\end{equation*}
For, if $C$ is an element of the $\gamma^{\text{th}}$ class, then, while $R$ runs through the elements of the $\alpha^{\text{th}}$ class and $S$ through those of the $\beta^{\text{th}}$, it occurs $\frac{h_{\alpha\beta\gamma'}}{h_\gamma}$ times that $RS=C$. Comparing the coefficients of $x_\alpha y_\beta$ therefore gives from (16.)
\begin{equation*}\tag{19.}
h_\alpha h_\beta\chi_\alpha\chi_\beta=f\sum_\gamma h_{\alpha\beta\gamma'}\chi_\gamma,
\end{equation*}
and consequently the value system $\chi_0,\chi_1,\ldots\chi_{k-1}$ is equal to one of the $k$ distinct characters. Equation (15.) can be put in the form
\begin{equation*}\tag{20.}
f\xi=\sum\chi(R)x_R.
\end{equation*}


Conversely, if $\chi$ is any one of the $k$ characters, then $\xi$ is a linear factor of the determinant $\Theta$. Namely, if in equations (17.) and (18.) one sets $y_\beta=\chi_{\beta'}$, then
\[
z_C=\sum_S x_{CS^{-1}}y_S=\sum_R x_{CR}\chi(R),\qquad h_\gamma z_\gamma=\sum_{\alpha,\beta}h_{\alpha\beta\gamma'}x_\alpha\chi_{\beta'}=\frac{1}{f}\sum_\alpha h_\alpha h_\gamma\chi_\alpha\chi_{\gamma'}x_\alpha,
\]
and so
\begin{equation*}\tag{21.}
\sum_R \chi(R)x_{CR}=\chi(C^{-1})\xi,
\qquad
\sum_R \chi(CR)x_R=\chi(C)\xi.
\end{equation*}
In determinant (14.), multiply the elements of the first row $(P=E)$ by the non-zero number $\chi(E)=f$, the elements of the row characterized by the index $P$ by $\chi(P)$, and then add all rows to the first. Then any element of the first row becomes
\[
\sum_P\chi(P)x_{PQ^{-1}}=\chi(Q)\xi,
\]
and consequently $\Theta$ is divisible by $\xi$. Therefore
\begin{equation*}\tag{22.}
\Theta=\left|x_{PQ^{-1}}\right|=\prod_\varkappa\left(\xi^{(\varkappa)}\right)^{g^{(\varkappa)}}
=\prod_\varkappa\left(\frac{1}{f^{(\varkappa)}}\sum h_\alpha\chi_\alpha^{(\varkappa)}x_\alpha\right)^{g^{(\varkappa)}},
\end{equation*}
where $g^{(\varkappa)}$ is a non-zero integer.

With the help of the theorems on matrices (linear systems) that I developed in my paper \emph{Über lineare Substitutionen und bilineare Formen}, Crelle's Journal vol.~84, a second proof of this theorem is obtained, from which the meaning of the numbers $g^{(\varkappa)}$ also becomes clear. The system of $h$ rows and columns so far denoted by $(x_{PQ^{-1}})$ I shall denote more briefly by $(x)$. If $\varepsilon_R=0$ when $R$ is different from $E$, but $\varepsilon_E=1$, then $(\varepsilon_{PQ^{-1}})=(\varepsilon)$ is the identity system obtained from $(x)$ by setting $x_0=1$, $x_1=x_2=\cdots=x_{k-1}=0$. Then the content of equations (5.) and (6.), \S~3 is expressed by the formula
\begin{equation*}\tag{23.}
\sum_\varkappa\left(\frac{1}{h}e^{(\varkappa)}\chi^{(\varkappa)}\right)=(\varepsilon)
\end{equation*}
and that of equations (9.) by the formulas
\begin{equation*}\tag{24.}
\left(\frac{1}{h}e^{(\varkappa)}\chi^{(\varkappa)}\right)^2=\left(\frac{1}{h}e^{(\varkappa)}\chi^{(\varkappa)}\right),
\qquad
\left(\frac{1}{h}e^{(\varkappa)}\chi^{(\varkappa)}\right)\left(\frac{1}{h}e^{(\lambda)}\chi^{(\lambda)}\right)=0.
\end{equation*}
If in the first of these one suppresses the index $\varkappa$, it shows: the equation of lowest degree satisfied by the system $\left(\frac{e\chi}{h}\right)$ is $\psi\left(\left(\frac{e\chi}{h}\right)\right)=0$, where $\psi(r)=r(r-1)$. Therefore the characteristic determinant $\varphi(r)$ of this system can vanish only for $r=0$ and $r=1$; and because $\psi(r)$ has no repeated linear factor, its elementary divisors are all of first degree. Thus if
\begin{equation*}\tag{25.}
\varphi(r)=\left|r(\varepsilon)-\left(\frac{e\chi}{h}\right)\right|=r^{h-g}(r-1)^g,
\end{equation*}


then for $r=0$ all subdeterminants of $\varphi(0)$ of degree $h,h-1,\ldots g+1$ vanish, while not all of degree $g$ vanish. Consequently $g$ is the rank of the system $(\chi)$.

By solving the linear equations (11.), \S~2 one obtains
\begin{equation*}\tag{26.}
hx_\alpha=\sum_\varkappa \chi_\alpha^{(\varkappa)}e^{(\varkappa)}\xi^{(\varkappa)}.
\end{equation*}
Hence
\begin{equation*}\tag{27.}
(x)=\sum_\varkappa \xi^{(\varkappa)}\left(\frac{1}{h}e^{(\varkappa)}\chi^{(\varkappa)}\right)
\end{equation*}
and consequently
\[
\prod_\varkappa\left(r(\varepsilon)-\xi^{(\varkappa)}\left(\frac{1}{h}e^{(\varkappa)}\chi^{(\varkappa)}\right)\right)=r^h(\varepsilon)-r^{h-1}(x),
\]
because in the expansion of this symbolic product all other terms vanish by (24.). Thus the determinants of these two systems are also equal:
\[
\prod_\varkappa\left(r^{h-g^{(\varkappa)}}(r-\xi^{(\varkappa)})^{g^{(\varkappa)}}\right)=r^{h(h-1)}\left|r(\varepsilon)-(x)\right|.
\]
Therefore, first, the powers of $r$ on both sides of this equation must have equal exponents. If one cancels them and then sets $r=0$, equation (22.) is obtained. Consequently in it $g^{(\varkappa)}$ is the rank of the system
\begin{equation*}\tag{28.}
\left(\chi^{(\varkappa)}(PQ^{-1})\right).
\end{equation*}

\section*{\S~6.}

In the determinant of degree $h$, $\left|x_{P,Q}\right|$, let the $h^2$ elements $x_{P,Q}$ first be independent variables. Let $T$ be a fixed element of $\grp$, and put
\[
T^{-1}PT=P',\qquad T^{-1}QT=Q'.
\]
Then $P'$ runs, together with $P$, through the $h$ elements of $\grp$, and $Q'$ through the same elements in the same order. Therefore
\begin{equation*}\tag{1.}
\Theta=\left|x_{P,Q}\right|=\left|x_{P',Q'}\right|,
\end{equation*}
and also the subdeterminant $\Theta_{A,B}$ complementary to the element $x_{A,B}$ in the first determinant is equal to the subdeterminant which is complementary, in the second determinant, to the same element $x_{A,B}$, though it stands there in a different position.

I now restrict the variability of the $h^2$ elements of the determinant $\Theta$ at first by setting $x_{P,Q}=x_{PQ^{-1}}$. Then it depends only on $h$ independent variables; each set of $h$ of the $h^2$ elements are equal to one another, and in every row all $h$ distinct variables occur, likewise in every column. The dif-


ferent rows differ from one another only by the order of the variables, which is determined by the constitution of the group $\grp$. I now assert that the subdeterminant $\Theta_{A,B}$ also does not depend on $A$ and $B$ separately, but only on $AB^{-1}=C$. In other words: the variable $x_C$ occurs in $\Theta$ in $h$ places. At each place the same subdeterminant is complementary to it. First, for a variable on the diagonal, $\Theta_{A,A}=\left|x_{R,S}\right|$, where $R$ runs through the elements of $\grp$ except $A$, and $S$ through the same $h-1$ elements in the same order. If $R=UA$, $S=VA$, then
\[
x_{R,S}=x_{UA,VA}=x_{UAA^{-1}V^{-1}}=x_{UV^{-1}}=x_{U,V},
\]
so $\Theta_{A,A}=\left|x_{U,V}\right|$, where $U$ and $V$ run through the elements of $\grp$ except $E$. Hence $\left|x_{U,V}\right|=\Theta_{E,E}$.

If $A$ and $B$ are two fixed elements of $\grp$, and $C=AB^{-1}$, then $CP$ runs through the $h$ elements of $\grp$ simultaneously with $P$, only in a different order. Hence
\[
\pm\left|x_{P,Q}\right|=\left|x_{CP,Q}\right|=\left|y_{P,Q}\right|,
\]
provided one sets $x_{CPQ^{-1}}=y_{PQ^{-1}}=y_{P,Q}$. Now $\Theta_{B,B}(y)=\Theta_{E,E}(y)$ and hence
\begin{equation*}\tag{2.}
\Theta_{A,B}=\Theta_{C,E}\qquad(C=AB^{-1}).
\end{equation*}
From this theorem follows the relation
\begin{equation*}\tag{3.}
\Theta_{A,B}=\frac{1}{h}\frac{\partial\Theta}{\partial x_{AB^{-1}}}.
\end{equation*}
The theory of these and still more general determinants, whose basis is formed by the present investigation, I treat in the paper announced at the beginning. Now, however, I further restrict the variability of the $h$ quantities $x_R$ by the conditions (11.), \S~5; thus I put $x_R=x_S$ when $R$ and $S$ are conjugate. If $C$ belongs to the $\gamma^{\text{th}}$ class, then the variable $x_C=x_\gamma$ now occurs in every row $h_\gamma$ times, and in the whole determinant $h h_\gamma$ times. But here too the subdeterminant $\Theta_{A,B}$ has the same value for each of these $h h_\gamma$ places.

Under the assumption just made, the two determinants (1.) not only have the same value, but the variables standing in corresponding positions are also equal, $x_{P,Q}=x_{P',Q'}$. Therefore the subdeterminant complementary to the variable $x_{P,Q}$ in the first is also equal to the subdeterminant complementary to the variable $x_{P',Q'}$ in the second; but this, as noted above, is equal to the subdeterminant complementary to the variable $x_{P',Q'}$ in the first. In connection with (2.) this proves the assertion
\begin{equation*}\tag{4.}
\Theta_{P,Q}=\Theta_{R,S},\quad \text{if } RS^{-1}=T^{-1}PQ^{-1}T.
\end{equation*}



For it is $RS^{-1}=P'Q'^{-1}$, hence
\[
\Theta_{P,Q}=\Theta_{P',Q'}=\Theta_{R,S}.
\]
Thus now
\begin{equation*}\tag{5.}
\Theta_{A,B}=\frac{1}{h h_{AB^{-1}}}\frac{\partial\Theta}{\partial x_{AB^{-1}}}.
\end{equation*}
From the identical equations between the elements of a determinant and the complementary minors belonging to them there follows therefore
\[
\sum_R \frac{1}{h_R}\frac{\partial\log\Theta}{\partial x_R}x_{RA^{-1}}=h\varepsilon_A,
\]
and so, if one substitutes for $\Theta$ its expression (22.), \S~5,
\[
\sum_R\sum_\varkappa \frac{g^{(\varkappa)}}{f^{(\varkappa)}\xi^{(\varkappa)}}\chi_R^{(\varkappa)}x_{RA^{-1}}=h\varepsilon_A,
\]
and, since by (21.), \S~5
\[
\sum_R \chi_R^{(\varkappa)}x_{RA^{-1}}=\chi^{(\varkappa)}(A)\xi^{(\varkappa)}
\]
is,
\[
\sum_\varkappa \frac{g^{(\varkappa)}}{f^{(\varkappa)}}\chi^{(\varkappa)}(A)=h\varepsilon_A,
\]
and since by (5.) and (6.), \S~3 also
\begin{equation*}\tag{6.}
\sum_\varkappa e^{(\varkappa)}\chi^{(\varkappa)}(A)=h\varepsilon_A
\end{equation*}
is true,
\[
\sum_\varkappa\left(\frac{g^{(\varkappa)}}{f^{(\varkappa)}}-e^{(\varkappa)}\right)\chi_\alpha^{(\varkappa)}=0.
\]
Since determinant (8.), \S~2 is different from zero, it follows from this that
\begin{equation*}\tag{7.}
g^{(\varkappa)}=e^{(\varkappa)}f^{(\varkappa)},
\end{equation*}
and with this it is proved that the roots $g=ef$ of equation (12.), \S~3 are integers.

One obtains the characteristic function of $\Theta$ by replacing $x_0$ by $x_0-r$, and therefore also $\xi^{(\varkappa)}$ by $\xi^{(\varkappa)}-r$. Hence
\[
\Theta(x_0-r,x_1,\ldots x_{k-1})=\prod_\varkappa (\xi^{(\varkappa)}-r)^{g^{(\varkappa)}}
\]
and consequently by (3.), \S~4
\begin{equation*}\tag{8.}
\begin{aligned}
\frac{1}{h}\frac{\partial\log\Theta(x_0-r,x_1,\ldots x_{k-1})}{\partial r}
&=\frac{1}{h}\sum_\varkappa \frac{e^{(\varkappa)}f^{(\varkappa)}}{r-\xi^{(\varkappa)}} \\
&=r^{-1}+x_0r^{-2}+\left(\sum_{\alpha,\beta}h_{\alpha\beta}x_\alpha x_\beta\right)r^{-3}
+\left(\sum_{\alpha,\beta,\gamma}h_{\alpha\beta\gamma}x_\alpha x_\beta x_\gamma\right)r^{-4}\\
&\quad+\left(\sum_{\alpha,\beta,\gamma,\delta}h_{\alpha\beta\gamma\delta}x_\alpha x_\beta x_\gamma x_\delta\right)r^{-5}+\cdots .
\end{aligned}
\end{equation*}
Accordingly this expression is the generating function of all the numbers $h_{\alpha\beta\gamma\ldots}$ defined in \S~1. After the numbers $h_{\alpha\beta\gamma}$ have been obtained from this expansion, determinant (14.), \S~2 supplies the product of the $k$ different linear factors of $\Theta$. The relation between these two determinants of degrees $h$ and $k$ is surely one of the most remarkable results of the theory developed here.

It follows further from equation (5.) that, if $\xi^g$ is a factor of $\Theta$, the minors of degree $(h-1)$ of $\Theta$ are all divisible by $\xi^{g-1}$. Therefore the minors of degree $(h-2)$ are all divisible by $\xi^{g-2}$, \ldots{} those of degree $(h-g+1)$ are all divisible by $\xi$, but those of degree $(h-g)$ are no longer all divisible by $\xi$. Thus, if $x_0,x_1,\ldots x_{k-1}$ are independent variables, $\Theta$ has only linear elementary divisors.

If each two elements of $\grp$ are interchangeable, then $k=h$; the characters determined by the conditions
\begin{equation*}\tag{9.}
\chi(AB)=\chi(A)\chi(B)
\end{equation*}
are all of degree $f=1$; the determinant $\Theta$ of degree $h$, which becomes identical with determinant (14.), \S~2 of degree $k$, contains every linear factor in the power $e=1$, and its decomposition is easy to carry out directly with the aid of relations (9.). For cyclic groups, for which $\xi$ passes into the resolvent of \textsc{Lagrange}, it was already given in 1853 by \textsc{Spottiswoode}, \emph{Elementary Theorems relating to Determinants}, \textsc{Crelle}'s Journal vol.~51, p.~375. Other cases have been treated by \textsc{Nöther}, \emph{Notiz über eine Classe symmetrischer Determinanten}, Math. Ann. vol.~16; \textsc{Gegenbauer}, \emph{Über eine specielle symmetrische Determinante}, Wiener Ber. 1880; \textsc{Puchta}, \emph{Ein neuer Satz aus der Theorie der Determinanten}, Wiener Denkschriften vol.~43. The general formula for the determinant of a commutative group was obtained by \textsc{Dedekind} in 1880, on the occasion of the investigation on group determinants mentioned at the beginning, by multiplication of the determinants $\left|x_{PQ^{-1}}\right|$ and $\left|\chi^{(\varkappa)}(R)\right|$. As I gather from the 25th volume of the \emph{Fortschritte der Mathematik}, p.~220, \textsc{Burnside}, in the journal \emph{Messenger of Math.} (2) XXIII, p.~112, which is not accessible to me, also derived this formula in a paper \emph{On a property of certain determinants}.

\section*{\S~7.}

The theorems and formulae developed so far all retain their validity if the concept of class on which they rest is taken more broadly. If the group $\grp$ is an invariant subgroup of another group $\grp'$, then two elements $R$ and $R'$ of $\grp$ shall be called conjugate (with respect to $\grp'$) if there is an element $T$ in $\grp'$ which satisfies the condition
\begin{equation*}\tag{1.}
R'=T^{-1}RT
\end{equation*}

Let $T$ be a fixed element of $\grp'$. Then, on the basis of equation (1.), to each element $R$ of $\grp$ there corresponds an element $R'$ of $\grp$. If $R$ runs through all elements of $\grp$, then $R'$ also runs through the same elements, only in another order. In this way one obtains an isomorphism of $\grp$ into itself, and indeed every possible one, by choosing the group $\grp'$ and within it the element $T$ suitably (\emph{Über endliche Gruppen} \S~5, Sitzungsber. 1895, p.~22).

Let $A$, $B$, and $P$ be elements of $\grp$, and let $B=P^{-1}AP$ be conjugate with $A$ with respect to $\grp$. Then $A'$ and $B'=P'^{-1}A'P'$ are also conjugate with respect to $\grp$, and conversely. Through that isomorphism, therefore, the elements of a class $(\alpha)$ pass into the elements of the same or of another class $(\beta)$. For two such classes, which I shall call conjugate, one therefore has $h_\alpha=h_\beta$. If one now unites all elements of $\grp$ which are conjugate to one another with respect to $\grp'$ into a class, then the number of the new classes is $k'\le k$, and each new class $(\rho)'$ arises by uniting a certain number $r$ of conjugate old classes $(\alpha)$, $(\beta)$, $(\gamma)$, \ldots. Since $h_\alpha=h_\beta=h_\gamma=\cdots$, it contains
\begin{equation*}\tag{2.}
h'_\rho=h_\alpha+h_\beta+h_\gamma+\cdots=rh_\alpha
\end{equation*}
elements of $\grp$.

The $k$ quantities $\chi_0,\chi_1,\ldots \chi_{k-1}$ which form a character $\chi$ are determined up to a common factor $f$ by the fact that
\begin{equation*}\tag{3.}
f\xi=\sum h_\alpha\chi_\alpha x_\alpha
\end{equation*}
is a linear factor of determinant $\Theta$. Likewise every new character $\chi'$ is determined by the fact that
\[
f'\xi'=\sum h'_\rho\chi'_\rho x'_\rho
\]
is a linear factor of an analogous determinant $\Theta'$. This latter arises from $\Theta$ by setting, for each $r$ conjugate classes $(\alpha)$, $(\beta)$, $(\gamma)$, \ldots{} which unite to form a class $(\rho)'$, $x_\alpha=x_\beta=x_\gamma=\cdots=x'_\rho$. Hence
\[
\frac{1}{f'}h'_\rho\chi'_\rho=\frac{1}{f}(h_\alpha\chi_\alpha+h_\beta\chi_\beta+h_\gamma\chi_\gamma+\cdots)
\]
or, by (2.),
\begin{equation*}\tag{4.}
\frac{f}{f'}\chi'_\rho=\frac{1}{r}(\chi_\alpha+\chi_\beta+\chi_\gamma+\cdots).
\end{equation*}
If $R$ is an element of one of those $r$ classes, then this equation can also be replaced by
\begin{equation*}\tag{5.}
\frac{f}{f'}\chi'(R)=\frac{1}{h'}\sum_U'\chi(U^{-1}RU),
\end{equation*}
where $U$ runs through all $h'$ elements of $\grp'$. For

in this sum of $h'$ terms, each of the $r$ terms of the sum (4.) occurs, and each equally often.

By setting certain complexes of the variables of $\Theta$ equal to one another, however, two different linear factors $\xi$ and $\eta$ of $\Theta$ can also become equal to one another. Thus two different characters $\chi$ and $\psi$ can produce the same character $\chi'$. This occurs in the following case:

If $T$ is a fixed element of $\grp'$, then by (1.) to each element $R$ of $\grp$ there corresponds an element $R'$ of $\grp$. If $S$ runs through the elements of a class $(\alpha)$, then $S'$ also runs through the elements of a class $(\beta)$. Therefore, by equation (5.), \S~5,
\[
\sum_S\bigl(\chi(E)\chi(A'S')-\chi(A')\chi(S')\bigr)=0.
\]
If therefore one sets
\[
\chi(R')=\chi(T^{-1}RT)=\psi(R),
\]
then, since $A'S'=(AS)'$,
\[
\sum_S\bigl(\psi(E)\psi(AS)-\psi(A)\psi(S)\bigr)=0.
\]
From this relation, which is identical with equation (6.), \S~2, it follows that $\psi(R)$ is a character. Such two characters shall be called conjugate (with respect to $\grp'$). If the substitution $T$ does not carry each of the old classes into itself, then one can choose the character $\chi$ so that $\psi$ is different from it. For if, for instance, $A$ and $A'$ represent different classes, then it cannot be the case that $\chi^{(\varkappa)}(A)=\chi^{(\varkappa)}(A')$ for every value of $\varkappa$. If one chooses $\varkappa$ so that this equation does not hold, and then sets $\chi^{(\varkappa)}=\chi$, then $\psi(A)=\chi(A')$ is different from $\chi(A)$; hence $\psi$ is another character than $\chi$.

Two conjugate characters differ only by the arrangement of the $k$ quantities $\chi_0,\chi_1,\ldots \chi_{k-1}$. But by formula (10.), \S~5,
\begin{equation*}\tag{6.}
\sum_R \chi(R^{-1})\chi(R)=\frac{hf}{e},
\end{equation*}
and therefore also
\[
\sum_R \psi(R^{-1})\psi(R)=\sum_{R'}\chi(R'^{-1})\chi(R')=\frac{hf}{e},
\]
since $R'$ runs at the same time as $R$ through all elements of $\grp$. For two conjugate characters, therefore, because $f=\chi(E)=\psi(E)$, the number $g=ef$, the exponent of the linear factor $\xi$ in $\Theta$, also has the same value.

Two conjugate characters $\chi$ and $\psi$ produce the same new character $\chi'$. For
\begin{equation*}\tag{7.}
\frac{h'f}{f'}\chi'(R)=\sum_U'\chi(U^{-1}RU)=\sum_U'\psi(U^{-1}RU).
\end{equation*}

The last sum is namely equal to $\sum_U'\chi((UT)^{-1}R(UT))$; it therefore differs from the first only by the order of the summands, since $UT$ runs at the same time as $U$ through all elements of $\grp'$.

Conversely, two characters $\chi$ and $\psi$ which produce the same character $\chi'$ must always be conjugate. Thus, if equation (7.) holds, there must be an element $T$ in $\grp'$ such that $\psi(R)=\chi(T^{-1}RT)$ for every element $R$ of $\grp$. For from that equation there follows
\[
\sum_U'\sum_R \chi(U^{-1}RU)\psi(R^{-1})=\sum_U'\sum_R \psi(U^{-1}RU)\psi(R^{-1}).
\]
Let $U$ be a fixed element of $\grp'$. If then $\psi(U^{-1}RU)=\psi(R)$ for every element $R$ of $\grp$, then by (6.) $\sum_R\psi(U^{-1}RU)\psi(R^{-1})$ is a positive quantity different from zero. Every element $U$ of $\grp'$ which is contained in $\grp$ satisfies this condition. But if it is not fulfilled, then $\psi(U^{-1}RU)=\vartheta(R)$ is a character different from $\psi(R)$, and hence this sum vanishes by equation (10.), \S~5. Thus the right-hand side of the last equation has a value different from zero, and hence also the sum on the left-hand side. But now, if again $U$ is a fixed element of $\grp'$ and $R$ is a variable element of $\grp$, $\chi(U^{-1}RU)=\vartheta(R)$ is a character. If this is different from $\psi(R)$, then $\sum_R\chi(U^{-1}RU)\psi(R^{-1})=0$. Thus $\vartheta$ cannot be different from $\psi$ for every $U$; hence there must be an element $U=T$ for which $\chi(T^{-1}RT)=\psi(R)$, and consequently $\chi$ and $\psi$ are conjugate characters.

In the sum (5.) all the characters conjugate with $\chi$ occur, and each equally often. For those elements $U$ of $\grp'$ which satisfy the equation $\chi(U^{-1}RU)=\chi(R)$ for every element $R$ of $\grp$ form a group contained in $\grp'$ which contains $\grp$. Therefore the number $s$ of the characters $\chi^{(\varkappa)},\chi^{(\lambda)},\chi^{(\mu)},\ldots$ which are conjugate with $\chi$ is a divisor of $\frac{h'}{h}$, and one can also put that equation into the form
\begin{equation*}\tag{8.}
\frac{f}{f'}\chi'=\frac{1}{s}(\chi^{(\varkappa)}+\chi^{(\lambda)}+\chi^{(\mu)}+\cdots).
\end{equation*}
The formulae (4.) and (8.) exhibit the two different ways in which one can think of the origin of the new characters, which I shall call the relative characters of $\grp$ with respect to $\grp'$, from the old ones.

If, in the passage from $\Theta$ to $\Theta'$, several different linear factors $\xi,\eta,\ldots$ of $\Theta$ become equal to one another, then they correspond to conjugate characters $\chi,\psi,\ldots$ and conversely. Hence the powers of $\xi,\eta,\ldots$ which divide $\Theta$ all have the same exponent $g$.

Consequently the exponent $g'$ of the power in which the corresponding linear factor $\xi'$ is contained in $\Theta'$ is
\begin{equation*}\tag{9.}
g'=sg.
\end{equation*}
Thus one may put, for example, $f'=sf$, $e'=e$. While, as mentioned above, $g$ is a square, $g'$ need not be a square. This explains why the proof of that theorem cannot be carried out by the means used so far alone.

The relative characters of $\grp$ have a simple meaning for the group $\grp'$ itself. The new classes into which I have divided the elements of $\grp$ are also classes of $\grp'$, because $\grp$ is an invariant subgroup of $\grp'$. Since I no longer make any use of the old classes and the function $\Theta$, I shall change the notation throughout, and shall denote by $k$ and $\Theta$ what I have until now called $k'$ and $\Theta'$, and the new classes themselves by $(0),(1)\ldots(k-1)$, the variables corresponding to them by $x_0,x_1,\ldots x_{k-1}$. On the other hand, let $k'$ now be the number of classes of $\grp'$. Then the classes $(k),(k+1),\ldots(k'-1)$ contain no element of $\grp$. Let the determinant corresponding to the group $\grp'$, namely (14.), \S~5 of degree $h'$, be
\[
\Theta'=\left|x_{P,Q}\right|=\left|x_{PQ^{-1}}\right|,
\]
where $P$ and $Q$ run through the $h'=nh$ elements of $\grp'$ in the same order. If $\grp'$ decomposes (mod. $\grp$) into the $n$ complexes
\[
\grp'=\grp+\grp T+\grp U+\cdots,
\]
then I arrange the rows and columns of that determinant so that $P$ first runs through the $h$ elements of $\grp$, then the $h$ elements of $\grp T$, then those of $\grp U$, and so on. If one now sets in $\Theta'$ $x_k=\cdots=x_{k'-1}=0$, only the $k$ variables $x_R$ remain, whose index $R$ is an element of $\grp$.

In the first $h$ rows only the elements of the first $h$ columns remain, and these form the determinant of degree $h$
\[
\Theta=\left|x_{R,S}\right|=\left|x_{RS^{-1}}\right|,
\]
where $R$ and $S$ run through the elements of $\grp$. In the following $h$ rows only the elements $x_{RT,ST}=x_{RT(ST)^{-1}}=x_{RS^{-1}}$ remain, which stand in columns $h+1$ through $2h$, and these again form the determinant $\Theta$, and so on. Therefore, if one sets $x_k=\cdots=x_{k'-1}=0$, one obtains
\begin{equation*}\tag{10.}
\Theta'=\Theta^n.
\end{equation*}
But the coefficients of the linear factors of $\Theta$ are the relative characters $\chi_0,\chi_1,\ldots \chi_{k-1}$ of $\grp$, and those of the linear factors of $\Theta'$ are the characters of $\grp'$. Thus, in every character $\chi_0,\ldots \chi_{k-1},\chi_k,\ldots \chi_{k'-1}$ of $\grp'$, the first $k$ values $\chi_0,\ldots \chi_{k-1}$ form, apart from a common

factor $f$, a relative character of $\grp$, and conversely every relative character of $\grp$, $\chi_0,\ldots \chi_{k-1}$, can be completed in one or more ways to a character of $\grp'$ by adjoining suitable values $\chi_k,\ldots \chi_{k'-1}$.

\section*{\S~8.}

I now wish to illustrate the theory of group characters by several examples. The even permutations of 4 symbols form a group $\grp$ of order $h=12$. Its elements fall into 4 classes; the elements of order 2 form a two-sided class (1), those of order 3 form two inverse classes (2) and $(3)=(2')$. Let $\rho$ be a primitive cubic root of unity.

\begin{center}
Tetrahedron. $h=12$.
\[
\begin{array}{c|rrrr|r}
 & \chi^{(0)} & \chi^{(1)} & \chi^{(2)} & \chi^{(3)} & h_\alpha\\ \hline
\chi_0 & 1 & 3 & 1 & 1 & 1\\
\chi_1 & 1 & -1 & 1 & 1 & 3\\
\chi_2 & 1 & 0 & \rho & \rho^2 & 4\\
\chi_3 & 1 & 0 & \rho^2 & \rho & 4
\end{array}
\]
\end{center}

The values of $\chi_0$ are at the same time those of $f=e$.

All permutations of 4 symbols form a group $\grp'$ of order $h'=24$, of which $\grp$ is an invariant subgroup. The classes (2) and (3) are conjugate and unite to a class $(2)'$. Likewise the characters $\chi^{(2)}$ and $\chi^{(3)}$ are conjugate. The relative characters are
\[
\begin{array}{c|rrr|r}
 & \chi^{(0)} & \chi^{(1)} & \chi^{(2)} & h_\alpha\\ \hline
\chi_0 & 1 & 3 & 2 & 1\\
\chi_1 & 1 & -1 & 2 & 3\\
\chi_2 & 1 & 0 & -1 & 8
\end{array}
\]
To obtain this table, add according to formula (8), \S~7, the two last columns in the preceding one, and keep only one of the two last rows, which then become equal to one another. Or, according to formula (4), \S~7, take the arithmetic mean from the elements of the two last rows, and keep only one of the two last columns. To obtain integers, multiply the elements of the last column further by $f'=2$. For the first two characters one has $e=f=\chi_0$; for the last, however, $e=1$, $f=2$, and therefore $g=2$ is not a square.

In the group $\grp'$ the even permutations of order 2 form a class (1), the odd ones a class (2), and the permuta-

tions of orders 3 and 4 each form a class (3) and (4). Each of the 5 classes is two-sided. The classes (2) and (4) contain the odd permutations; the classes (0), (1), and (3) contain the even permutations and together form the group $\grp$.

\begin{center}
Octahedron. $h=24$.
\[
\begin{array}{c|rrrrr|r}
 & \chi^{(0)} & \chi^{(1)} & \chi^{(2)} & \chi^{(3)} & \chi^{(4)} & h_\alpha\\ \hline
\chi_0 & 1 & 3 & 3 & 2 & 1 & 1\\
\chi_1 & 1 & -1 & -1 & 2 & 1 & 3\\
\chi_2 & 1 & 1 & -1 & 0 & -1 & 6\\
\chi_3 & 1 & 0 & 0 & -1 & 1 & 8\\
\chi_4 & 1 & -1 & 1 & 0 & -1 & 6
\end{array}
\]
\end{center}

For every character $e=f=\chi_0$.

2. The permutations of 5 symbols form a group $\grp'$ of order $h'=120$, while the even ones alone form a group $\grp$ of order $h=60$. The elements of $\grp'$ fall into 7 classes. The elements of orders 3, 4, 5, and 6 each form a class (3), (4), (5), (6). Of the permutations of order 2, the even ones form class (1), the odd ones class (2). Classes (2), (4), and (6) contain the odd permutations; the remaining ones contain the even permutations and together form the subgroup $\grp$. Each of the 7 classes is two-sided. If one considers $\grp$ by itself, then class (5) of $\grp'$ splits into 2 classes. Therefore the 5 classes of $\grp$ shall be denoted by (0), (1), (3), (4), (5), where (4) and (5) contain the elements of order 5. Both in $\grp$ and in $\grp'$, each class is two-sided.

\begin{center}
Icosahedron. $h=60$.
\[
\begin{array}{c|rrrrr|r}
 & \chi^{(0)} & \chi^{(1)} & \chi^{(3)} & \chi^{(4)} & \chi^{(5)} & h_\alpha\\ \hline
\chi_0 & 1 & 5 & 4 & 3 & 3 & 1\\
\chi_1 & 1 & 1 & 0 & -1 & -1 & 15\\
\chi_3 & 1 & -1 & 1 & 0 & 0 & 20\\
\chi_4 & 1 & 0 & -1 & \frac{1}{2}(1+\sqrt{5}) & \frac{1}{2}(1-\sqrt{5}) & 12\\
\chi_5 & 1 & 0 & -1 & \frac{1}{2}(1-\sqrt{5}) & \frac{1}{2}(1+\sqrt{5}) & 12
\end{array}
\]
\end{center}

Here always $e=f=\chi_0$. The relative characters with respect to $\grp'$ are, however,
\[
\begin{array}{c|rrrr|r}
 & \chi^{(0)} & \chi^{(1)} & \chi^{(3)} & \chi^{(5)} & h_\alpha\\ \hline
\chi_0 & 1 & 5 & 4 & 6 & 1\\
\chi_1 & 1 & 1 & 0 & -2 & 15\\
\chi_3 & 1 & -1 & 1 & 0 & 20\\
\chi_5 & 1 & 0 & -1 & 1 & 24
\end{array}
\]
Here $e^{(5)}=3$, $f^{(5)}=6$, hence $g^{(5)}=18$ is not a square. Finally, for the group $\grp'$ one obtains:

\begin{center}
$h=120$
\[
\begin{array}{c|rrrrrrr|r}
 & \chi^{(0)} & \chi^{(1)} & \chi^{(2)} & \chi^{(3)} & \chi^{(4)} & \chi^{(5)} & \chi^{(6)} & h_\alpha\\ \hline
\chi_0 & 1 & 5 & 5 & 4 & 4 & 6 & 1 & 1\\
\chi_1 & 1 & 1 & 1 & 0 & 0 & -2 & 1 & 15\\
\chi_2 & 1 & 1 & -1 & 2 & -2 & 0 & -1 & 10\\
\chi_3 & 1 & -1 & -1 & 1 & 1 & 0 & 1 & 20\\
\chi_4 & 1 & -1 & 1 & 0 & 0 & 0 & -1 & 30\\
\chi_5 & 1 & 0 & 0 & -1 & -1 & 1 & 1 & 24\\
\chi_6 & 1 & 1 & -1 & -1 & 1 & 0 & -1 & 20
\end{array}
\]
\end{center}

For the symmetric group $\grp$ of degree $n$ all characters are also rational integers. The reason lies in the fact that $R$ and $R^a$ are always conjugate in $\grp$ whenever $a$ is relatively prime to the order of $R$.

\section*{\S~9.}

To carry out a still more general example, I choose the group $\grp$ of order
\begin{equation*}\tag{1.}
h=\frac{1}{2}p(p^2-1),
\end{equation*}
which is formed by the linear substitutions
\begin{equation*}\tag{2.}
y=\frac{\gamma+\delta x}{\alpha+\beta x}\qquad(\operatorname{mod}.p)
\end{equation*}
of determinant
\begin{equation*}\tag{3.}
\alpha\delta-\beta\gamma\equiv 1\qquad(\operatorname{mod}.p)
\end{equation*}
provided that $p$ denotes an odd prime number. The properties which come into consideration here are derived most fully by \textsc{Gierster}, \emph{Die Untergruppen der \textsc{Galois}'schen Gruppe der Modulargleichungen für den Fall eines primzahligen Transformationsgrades}, Math. Ann. vol.~18, p.~319. From this interesting and important work I take the following results:

The $h$ elements of the group $\grp$ fall into
\begin{equation*}\tag{4.}
k=\frac{1}{2}(p-1)+3
\end{equation*}
classes of conjugate elements. The principal class formed by the element $E$ shall not, as in the general theory, be denoted by (0), but by $(\lambda)$. The elements of order $p$ form two classes, which I shall also denote not by numerals, but by the letters $(\mu)$ and $(\nu)$. Let $P$ and $Q$ be two non-conjugate elements of order $p$. If $a$ is a quadratic residue of $p$ and $b$ a non-residue, then $P^a$ is conjugate with $P$, while $P^b$ is not. Therefore, if $p\equiv3\pmod{4}$, then $(\mu)$ and $(\nu)$ are inverse classes, and one may choose $Q=P^{-1}$. If however $p\equiv1\pmod{4}$, then $P$ and $P^{-1}$ are conjugate, as likewise are $Q$ and $Q^{-1}$; thus each of these two classes is a two-sided one.

The order of every other element of the group is a divisor of $\frac12(p-1)$ or of $\frac12(p+1)$. There is an element $R$ of order $\frac12(p-1)$ and an element $S$ of order $\frac12(p+1)$. Among the powers of $R$, any two with opposite exponents $R^a$ and $R^{-a}$ are conjugate, and likewise among the powers of $S$. On the other hand, according as $p\equiv1$ or $3\pmod{4}$,
\begin{equation*}\tag{5.}
\begin{array}{c}
R,\ R^2,\ldots,\ R^{\frac{p-1}{4}}\quad\left(\text{resp. }R^{\frac{p-3}{4}}\right),\\[0.35em]
S,\ S^2,\ldots,\ S^{\frac{p-1}{4}}\quad\left(\text{resp. }S^{\frac{p+1}{4}}\right),
\end{array}
\end{equation*}
represent $\frac12(p-1)$ distinct classes. If the three elements $E$, $P$, and $Q$ are added to these, one has one and only one representative for each of the $k$ classes. Every class is two-sided. Only when $p\equiv3\pmod{4}$ are $(\mu)$ and $(\nu)$ inverse classes. After $R$ and $S$ have been chosen in some definite way, let the number $\pm a$ $\left(\operatorname{mod}.\frac12(p-1)\right)$ be called the index of the class represented by $R^a$ or $R^{-a}$, and similarly the number $\pm b$ $\left(\operatorname{mod}.\frac12(p+1)\right)$ the index of the class represented by $S^b$ or $S^{-b}$.

For any two conjugate substitutions (2.) the number
\begin{equation*}\tag{6.}
\varkappa\equiv \frac12(\alpha+\delta)\qquad(\operatorname{mod}.p)
\end{equation*}
has, apart from sign, the same value, and therefore so do the two roots of each of the two congruences
\begin{equation*}\tag{7.}
x^2\pm2\varkappa x+1\equiv0\qquad(\operatorname{mod}.p).
\end{equation*}
I therefore call $\pm\varkappa$ the invariant of the class to which the substitution (2.) belongs. Conversely, if two substitutions have the same invariant $\pm\varkappa$, then they are conjugate, except when $\varkappa=\pm1$. The substitutions of invariant $\varkappa=\pm1$ split into three classes according to the quadratic character $(\operatorname{mod}.p)$ of the numbers represented by the form
\begin{equation*}\tag{8.}
\frac12(\alpha+\delta)\bigl(\beta x^2+(\alpha-\delta)xy-\gamma y^2\bigr).
\end{equation*}
If all these are divisible by $p$, then $\beta=\gamma=0$, $\alpha=\delta=\pm1$, and we obtain the principal class $(\lambda)$. If they are not all divisible by $p$, then the numbers representable by the form and prime to $p$ are either all residues or all non-residues of $p$. In the first case $\varkappa\beta$ is a residue, and, if $\beta=0$, $-\varkappa\gamma$ is a residue. This class shall be called $(\mu)$. In the second case $\varkappa\beta$ is a non-residue, and, if $\beta=0$, $-\varkappa\gamma$ is a non-residue. This class shall be called $(\nu)$. But if $\varkappa$ is different from $\pm1$, then there is only one class of invariant $\pm\varkappa$; it shall be denoted by $(\varkappa)$ or also by $(+\varkappa)$ $\left(\varkappa=0,2,3,\ldots,\frac{p-1}{2}\right)$. Put
\begin{equation*}\tag{9.}
\varepsilon_\varkappa=\left(\frac{\varkappa^2-1}{p}\right),\qquad
\varepsilon_0=\varepsilon=\left(\frac{-1}{p}\right),\qquad
\varepsilon_1=0.
\end{equation*}
Then, for the number of elements in each of these $k$ classes, one obtains the formulae
\begin{equation*}\tag{10.}
h_\lambda=1,\qquad
h_\mu=h_\nu=\frac12(p^2-1),\qquad
h_0=\frac12p(p+\varepsilon),\qquad
h_\varkappa=p(p+\varepsilon_\varkappa).
\end{equation*}

The elements of the classes $(\mu)$ and $(\nu)$ all have order $p$. The class $(0)$ consists of all elements of $\grp$ whose order is equal to $2$. Its somewhat different behaviour with respect to the other classes $(\varkappa)$ is explained as follows: if $A$ is an element of the class $(\varkappa)$, then $A$ and $A^{-1}$ are conjugate but distinct. If however $\varkappa=0$, then these two elements are equal. If $R^a$ does not have order $2$ (or $1$), then the $\frac12(p-1)$ powers of $R$ are the only elements of $\grp$ that are commutable with $R^a$. But with the group of these powers there is also commutable an element $T$ of order $2$ that satisfies the condition $T^{-1}RT=R^{-1}$. The analogous assertion holds for $S^b$. If however $R^a$ (respectively $S^b$) has order $2$, then, besides the powers of $R$, $T$ is also commutable with $R^a$.

If $\varkappa^2-1$ is a quadratic residue of $p$, hence $\varepsilon_\varkappa=+1$, then the order of an element of the class $(\varkappa)$ is the exponent to which $(\varkappa+\sqrt{\varkappa^2-1})^2$ belongs $(\operatorname{mod}.p)$, and so is a divisor of $\frac12(p-1)$. If however $\varkappa^2-1$ is a non-residue, hence $\varepsilon_\varkappa=-1$, then this order is the exponent to which $x^2$ belongs $(\operatorname{mod}.p,\ x^2-2\varkappa x+1)$, and so is a divisor of $\frac12(p+1)$; consequently in both cases it is a divisor of $\frac12(p-\varepsilon_\varkappa)$. Conversely, therefore, for the class to which $R^a$ belongs one has $\varepsilon_\varkappa=+1$, and for the class to which $S^b$ belongs one has $\varepsilon_\varkappa=-1$. By (5.) the elements whose order divides
\begin{equation*}\tag{11.}
\frac12(p-1),\quad \text{fall into}\quad \frac14(p-\varepsilon)-\frac12(1-\varepsilon)
\end{equation*}
classes, and those whose order divides
\begin{equation*}\tag{12.}
\frac12(p+1),\quad \text{fall into}\quad \frac14(p-\varepsilon)
\end{equation*}
classes. Here the principal element is excluded. Or, summarized differently: the elements whose order divides
\begin{equation*}\tag{13.}
\frac12(p-\varepsilon),\quad \text{fall into}\quad \frac14(p-\varepsilon)
\end{equation*}
classes, and those whose order divides
\begin{equation*}\tag{14.}
\frac12(p+\varepsilon),\quad \text{fall into}\quad \frac14(p-\varepsilon)-\frac12(1-\varepsilon)
\end{equation*}
classes. If $\pm\varkappa$ is the invariant of a class, I shall now define its index as follows. Let $r$ and $s$ be primitive roots of the congruences
\begin{equation*}\tag{15.}
r^{p-1}\equiv1,\qquad s^{p+1}\equiv1\qquad(\operatorname{mod}.p).
\end{equation*}
Then $r$ is real, and although $s$ is imaginary, $s^b+s^{-b}$ is real. If now $\varepsilon_\alpha=+1$, then each of the two congruences (7.) has two real roots, $r^{\pm a}$, $r^{\frac12(p-1)\pm a}$. I then call the number $\pm a$, taken modulo $\frac12(p-1)$, the index of the class $(\pm\alpha)$. If however $\varepsilon_\beta=-1$, then each has two imaginary roots $s^{\pm b}$, $s^{\frac12(p+1)\pm b}$. I then call the number $\pm b$, taken modulo $\frac12(p+1)$, the index of the class $(\pm\beta)$. This definition passes over into the preceding one if for $R$ one takes the substitution $y=\frac{rx}{r^{-1}}$, and for $S$ chooses a substitution which, in imaginary form, is $y=\frac{sx}{s^{-1}}$. (\textsc{Gierster}, loc. cit. \S~3.) For if $r$ and $r^{-1}$ are the roots of the characteristic equation (7.) for the substitution $R$, then $r^a$ and $r^{-a}$ are those for the substitution $R^a$.

\section*{\S~10.}

Three (distinct or equal) classes shall be called \emph{concordant} if between their invariants $\alpha$, $\beta$, $\gamma$ the relation
\begin{equation*}\tag{1.}
\alpha^2+\beta^2+\gamma^2\pm2\alpha\beta\gamma\equiv1\qquad(\operatorname{mod}.p)
\end{equation*}
holds, and otherwise \emph{discordant}. If, for example, $\gamma=\pm1$, then this relation reduces to $\beta=\pm\alpha$. If $\gamma=0$, it reads $\alpha^2+\beta^2=1$. I now exclude the case in which one of the invariants is equal to $\pm1$. Writing equation (1.) in the form
\[
(\alpha^2-1)(\beta^2-1)\equiv(\alpha\beta\pm\gamma)^2,
\]
one obtains from it
\begin{equation*}\tag{2.}
\varepsilon_\alpha=\varepsilon_\beta=\varepsilon_\gamma.
\end{equation*}
If $\alpha$ and $\beta$ are given, then concordant with them are
\begin{equation*}\tag{3.}
\pm\gamma\equiv\alpha\beta+\sqrt{(\alpha^2-1)(\beta^2-1)},\qquad
\pm\delta\equiv\alpha\beta-\sqrt{(\alpha^2-1)(\beta^2-1)},
\end{equation*}
and $\varepsilon_\alpha=\varepsilon_\beta=\varepsilon_\gamma=\varepsilon_\delta$. If, for example, $\beta=\pm\alpha$, then $\pm\gamma=2\alpha^2-1$ and $\pm\delta=1$, and consequently always $\varepsilon_\alpha=\varepsilon_{(2\alpha^2-1)}$. The orders of the elements of three concordant classes all divide $\frac12(p-1)$, or all divide $\frac12(p+1)$. Thus, if $a$, $b$, $c$ are their indices, these all refer to the same primitive root $r$ (or $s$). If
\[
2\alpha\equiv r^a+r^{-a},\qquad
2\beta\equiv r^b+r^{-b},\qquad
2\gamma\equiv r^c+r^{-c},
\]
then congruence (1.) (for the upper sign) passes over into
\[
(r^{a+b+c}+1)(r^{a-b-c}+1)(r^{-a+b-c}+1)(r^{-a-b+c}+1)\equiv0,
\]
and is therefore identical with the condition
\begin{equation*}\tag{4.}
a\pm b\pm c\equiv0\qquad\left(\operatorname{mod}.\frac12(p-1),\ \text{resp. }\frac12(p+1)\right).
\end{equation*}
Consequently the indices of the two classes (3.) are
\begin{equation*}\tag{5.}
c\equiv a+b,\qquad d\equiv a-b.
\end{equation*}
I also put, for abbreviation,
\begin{equation*}\tag{6.}
2\eta_\varkappa=
\left(\frac{-2}{p}\right)
\left(\left(\frac{1+\varkappa}{p}\right)+\left(\frac{1-\varkappa}{p}\right)\right),
\qquad
\eta_0=\left(\frac{-2}{p}\right),\qquad
2\eta_1=\varepsilon.
\end{equation*}
If $\varepsilon_\beta=-\varepsilon$, then $\eta_\beta=0$. If however $\varepsilon_\alpha=+\varepsilon$, then
\begin{equation*}\tag{7.}
\eta_\alpha=
\left(\frac{-2}{p}\right)\left(\frac{1+\alpha}{p}\right)
=
\left(\frac{-2}{p}\right)\left(\frac{1-\alpha}{p}\right)
=(-1)^a\varepsilon.
\end{equation*}
For if $\varepsilon=+1$ and $2\alpha\equiv r^a+r^{-a}$, then $2(\alpha+1)\equiv(r^{-a}+1)^2r^a$, hence
\[
\left(\frac{2}{p}\right)\left(\frac{1+\alpha}{p}\right)=\left(\frac{r^a}{p}\right)=(-1)^a.
\]
If however $\varepsilon=-1$ and $2\alpha\equiv s^a+s^{-a}$, then $2\alpha+2\equiv s^a+s^{-a}+2$. Thus, if $a=2b$ is even, then $2\alpha+2\equiv(s^b+s^{-b})^2$, hence a quadratic residue, since $s^b+s^{-b}$ is real. Conversely, if $2\alpha+2=4\beta^2$ is a quadratic residue, then $2\alpha-2=4(\beta^2-1)$ is a non-residue, since $\varepsilon_\alpha=-1$, hence $\alpha^2-1=(\alpha-1)(\alpha+1)$ is a non-residue. If then $s^b$ and $s^{-b}$ are the roots of the congruence $x^2-2\beta x+1\equiv0$, then $2\beta\equiv s^b+s^{-b}$, hence $s^a+s^{-a}+2\equiv(s^b+s^{-b})^2$ or $s^a+s^{-a}\equiv s^{2b}+s^{-2b}$. Consequently $a\equiv\pm2b$ $(\operatorname{mod}.p+1)$, so that $a$ is even.

Let now $\alpha$, $\beta$, $\gamma$ be numbers from $2$ to $\frac12(p-1)$. Then
\[
h_{000}=\frac12h,\qquad
h_{00\alpha}=h,\qquad
h_{0\alpha\beta}=2h,\qquad
h_{\alpha\beta\gamma}=4h,
\]
if the three classes are discordant. If however they are concordant, then
\[
h_{0\alpha\beta}=2h+\varepsilon p^2(p+\varepsilon),\qquad
h_{\alpha\beta\gamma}=4h+\varepsilon_\alpha p^2(p+\varepsilon_\alpha).
\]
Furthermore,
\[
\begin{array}{lll}
h_{\lambda\lambda\lambda}=1, &
h_{\lambda\mu\mu}=h_{\lambda\nu\nu}=h_\mu\frac12(1+\varepsilon), &
h_{\lambda\mu\nu}=h_\mu\frac12(1-\varepsilon),\\[0.35em]
h_{00\lambda}=\frac12p(p+\varepsilon), &
h_{\alpha\alpha\lambda}=p(p+\varepsilon_\alpha), &
h_{\alpha\beta\lambda}=0,\\[0.35em]
h_{\mu\mu\mu}=h_{\nu\nu\nu}=h_\mu\left(\frac34(p-\varepsilon)+\frac12(p-1)\right),&
\multicolumn{2}{l}{
h_{\mu\mu\nu}=h_{\mu\nu\nu}=h_\mu\left(\frac14(p-\varepsilon)-\frac12(1-\varepsilon)\right),}\\[0.35em]
h_{00\mu}=h_{00\nu}=h\frac12(1+\varepsilon),&
h_{0\alpha\mu}=h_{0\alpha\nu}=h,&
h_{\alpha\beta\mu}=h_{\alpha\beta\nu}=2h,\\[0.35em]
h_{\alpha\alpha\mu}=h_{\alpha\alpha\nu}=2h+\varepsilon_\alpha h,&&\\[-0.2em]
h_{0\mu\mu}=h_{0\nu\nu}=h\frac12(1+\varepsilon\eta),&
h_{0\mu\nu}=h\frac12(1-\varepsilon\eta),&\\[0.35em]
h_{\alpha\mu\mu}=h_{\alpha\nu\nu}=h(1+\varepsilon\eta_\alpha),&
h_{\alpha\mu\nu}=h(1-\varepsilon\eta_\alpha).&
\end{array}
\]
I shall only indicate briefly the way by which I have calculated these numbers. To obtain $\frac{h_{0\alpha\beta}}{h_0}$, take a definite substitution of the class $(0)$ and compose it with all substitutions of the class $(\alpha)$,
\[
\begin{pmatrix}0&1\\-1&0\end{pmatrix}
\begin{pmatrix}\xi&\eta\\ \zeta&2\alpha-\xi\end{pmatrix}
=
\begin{pmatrix}\zeta&2\alpha-\xi\\ -\xi&-\eta\end{pmatrix}.
\]
In order that the composite substitution belong to the class $(\beta)'=(\beta)$, it is necessary that
\[
\zeta-\eta=\pm2\beta,\qquad
\xi(2\alpha-\xi)-\eta\zeta\equiv1,
\]
hence
\[
(\xi-\alpha)^2+(\eta\pm\beta)^2\equiv\alpha^2+\beta^2-1.
\]
If $\beta$ is different from zero, then the number of solutions of these two congruences is twice as large as that of the congruence $\xi^2+\eta^2\equiv\omega$, where $\omega\equiv\alpha^2+\beta^2-1$. If $\omega$ is different from zero, this number is $p-\varepsilon$. But if $\omega=0$, so that $(0)$, $(\alpha)$, $(\beta)$ are concordant classes, it is $p-\varepsilon+\varepsilon p$. Since now $h_0=\frac12p(p+\varepsilon)$, it follows that $h_{0\alpha\beta}=2h$, except when $\alpha^2+\beta^2=1$; in that case, however, $h_{0\alpha\beta}=2h+\varepsilon p^2(p+\varepsilon)$.

The group $\grp$ can be enlarged to a group of order $2h$ in which $P$ and $Q$ are conjugate elements. It therefore has an automorphism which carries the classes $(\mu)$ and $(\nu)$ into each other. Hence the same relations are obtained for the class $(\mu)$ as for $(\nu)$. For example, $h_{\mu\nu\nu}=h_{\nu\mu\mu}$. By formulae (10.), \S~9, $h_{\lambda\mu\nu}=h_\mu\frac12(1-\varepsilon)$. Now to find $\frac{1}{h_\mu}(h_{\mu\nu\mu}+h_{\mu\nu\nu}+h_{\mu\nu\lambda})$, take a definite substitution of the class $(\mu)$ and compose it with all substitutions of $(\nu)$,
\[
\begin{pmatrix}0&1\\-1&2\end{pmatrix}
\begin{pmatrix}\xi&\eta\\ \zeta&2-\xi\end{pmatrix},
\]
and put the invariant of the composite substitution equal to $\pm1$. This gives the equations
\[
\zeta-\eta+4-2\xi=\pm2,\qquad
2\xi-\xi^2-\eta\zeta\equiv1.
\]
Taking the lower sign, one obtains by elimination of $\xi$
\[
(\eta+\zeta+4)^2\equiv16\eta.
\]
This congruence has no admissible solution. For $\eta$ is required to be a non-residue, and, if $\eta=0$, then $-2$ is required to be a non-residue. If the upper sign is taken, one obtains $(\eta+\zeta)^2\equiv0$, hence $\eta=-\zeta$. Therefore any non-residue may be put for $\eta$, and one obtains $\frac12(p-1)$ solutions. Consequently $2h_{\mu\mu\nu}+h_\mu\frac12(1-\varepsilon)=h_\mu\frac12(p-1)$, and therefore
\[
h_{\mu\mu\nu}=h_\mu\left(\frac14(p-\varepsilon)-\frac12(1-\varepsilon)\right).
\]

To determine the characters of the group $\grp$, one accordingly has to solve the following equations. I put, for abbreviation,
\[
x=\chi_0+2\sum\chi_\varkappa+\chi_\mu+\chi_\nu,\qquad
y=\eta_0\chi_0+2\sum\eta_\varkappa\chi_\varkappa+\eta_1(\chi_\mu+\chi_\nu),
\]
where $\varkappa=2,3,\ldots,\frac12(p-1)$. Let $\alpha$ and $\beta$ be numbers in the sequence $0,2,3,\ldots,\frac12(p-1)$.

\[
\begin{array}{ll}
1)&\text{If }\alpha\text{ is different from }\beta\text{ and }\varepsilon_\beta=-\varepsilon_\alpha,\\[0.25em]
& p\chi_\alpha\chi_\beta=fx.\\[0.65em]
2)&\text{If however }\varepsilon_\beta=+\varepsilon_\alpha,\\[0.25em]
& \displaystyle \frac{p(p+\varepsilon_\alpha)}{f}\chi_\alpha\chi_\beta
=x(p-\varepsilon_\alpha)+\varepsilon_\alpha p(\chi_\gamma+\chi_\delta),
\end{array}
\]
where $\gamma$ and $\delta$ are determined by (3.).
\[
\begin{array}{ll}
3)&\displaystyle \frac{(p+\varepsilon_\alpha)}{f}\chi_\alpha\chi_\mu
=x+\varepsilon_\alpha\chi_\alpha+\eta_\alpha(\chi_\mu-\chi_\nu),\\[0.75em]
&\displaystyle \frac{(p+\varepsilon_\alpha)}{f}\chi_\alpha\chi_\nu
=x+\varepsilon_\alpha\chi_\alpha+\eta_\alpha(\chi_\nu-\chi_\mu).\\[0.75em]
4)&\displaystyle \frac{p(p+\varepsilon)}{f}\chi_0^2
=x(p-\varepsilon)+\varepsilon(p-\varepsilon)(\chi_\mu+\chi_\nu)+2f,\\[0.75em]
&\displaystyle \frac{p(p+\varepsilon_\alpha)}{f}\chi_\alpha^2
=x(p-\varepsilon_\alpha)+\frac12\varepsilon_\alpha(p-\varepsilon_\alpha)(\chi_\mu+\chi_\nu)+f+\varepsilon_\alpha p\chi_{(2\alpha^2-1)},\\[0.75em]
5)&\displaystyle \frac{p^2-1}{pf}\chi_\mu^2
=x+\varepsilon y+\frac{f}{p}(1+\varepsilon)-\frac1p\left(1+\frac{\varepsilon}{2}\right)(\chi_\mu+\chi_\nu)+\frac{\varepsilon(p-\varepsilon)}{p}(\chi_\mu-\chi_\nu),\\[0.75em]
&\displaystyle \frac{p^2-1}{pf}\chi_\nu^2
=x+\varepsilon y+\frac{f}{p}(1+\varepsilon)-\frac1p\left(1+\frac{\varepsilon}{2}\right)(\chi_\mu+\chi_\nu)+\frac{\varepsilon(p-\varepsilon)}{p}(\chi_\nu-\chi_\mu),\\[0.75em]
&\displaystyle \frac{p^2-1}{pf}\chi_\mu\chi_\nu
=x-\varepsilon y+\frac{f}{p}(1-\varepsilon)-\frac1p\left(1-\frac{\varepsilon}{2}\right)(\chi_\mu+\chi_\nu).
\end{array}
\]

I. I first investigate whether there are solutions in which $\chi_\mu$ and $\chi_\nu$ are different. From 3) it follows that $(p+\varepsilon_\alpha)\chi_\alpha=2f\eta_\alpha$. If $\varepsilon_\alpha=-\varepsilon$, then $\eta_\alpha=0$, and therefore also $\chi_\alpha=0$. Thus, putting the provisionally arbitrary proportionality factor $f=\frac12(p+\varepsilon)$, one has
\[
\chi_\alpha=\eta_\alpha\qquad\left(\alpha=0,2,3,\ldots,\frac12(p-1)\right).
\]
By 1), $x=0$. By equation (2.), $\varepsilon_\alpha=\varepsilon_{(2\alpha^2-1)}$. Thus, if in 4) $\varepsilon_\alpha=-\varepsilon$, one obtains $\chi_\mu+\chi_\nu=\varepsilon$. By 5) one then has $y=\frac12(p-\varepsilon-1)$ and $4\chi_\mu\chi_\nu=1-\varepsilon p$, hence
\[
\chi_\mu=\frac12(\varepsilon\pm\sqrt{\varepsilon p}),\qquad
\chi_\nu=\frac12(\varepsilon\mp\sqrt{\varepsilon p}).
\]
That the two solutions so obtained really satisfy all the equations, for example the equations 2),
\begin{equation*}\tag{8.}
2\eta_\alpha\eta_\beta=\varepsilon_\alpha(\eta_\gamma+\eta_\delta),
\end{equation*}
is most simply verified by means of formula (7.), \S~9, and (5.).

II. For every other solution $\chi_\mu=\chi_\nu$. The equations thereby simplify considerably; for example, 5) is to be replaced by
\[
5^*)\qquad
\frac{p^2-1}{f}\chi_\mu^2=px+f-2\chi_\mu,\qquad
py=\chi_\mu-f.
\]

I now investigate whether there are solutions in which $x$ is different from zero. Then by 1), $\chi_0$, $\chi_2$, $\chi_3,\ldots,\chi_{\frac12(p-1)}$ are different from zero; moreover all $\chi_\alpha$ are equal to each other for which $\varepsilon_\alpha$ has the same value. Let $\varepsilon_\alpha=+\varepsilon$ and $\varepsilon_\delta=-\varepsilon$. Then
\[
y=\eta_0\chi_0+2\chi_0\sum\eta_\alpha+\varepsilon\chi_\mu,
\]
and hence, since $\eta_0+2\sum\eta_\alpha=-\varepsilon$,
\[
y=\varepsilon(\chi_\mu-\chi_0)=\frac{\chi_\mu-f}{p}.
\]
Further, in the sum $x=\chi_0+2\chi_\mu+2\sum(\chi_\alpha+\chi_\delta)$, the number of the characters $\chi_\alpha$ is equal to $\frac14(p-\varepsilon)-1$, while that of the characters $\chi_\delta$ is equal to $\frac14(p-\varepsilon)-\frac12(1-\varepsilon)$. Therefore
\[
x=2\chi_\mu+\left(\frac12(p-\varepsilon)-1\right)\chi_0+
\left(\frac12(p+\varepsilon)-1\right)\chi_\delta.
\]
By 1) and 4),
\[
fx=p\chi_0\chi_\delta,\qquad
\frac{p(p-\varepsilon)}{f}\chi_\delta^2=x(p+\varepsilon)-2\varepsilon p\chi_\delta,
\]
and hence
\[
(p-\varepsilon)\chi_\delta=(p+\varepsilon)\chi_0-2\varepsilon f,\qquad
px=(p^2+1)\chi_\mu-f.
\]
Substitution of these values in $5^*)$ gives
\[
\chi_\mu(\chi_\mu-f)=0,
\]
therefore either $\chi_\mu=f$ or $\chi_\mu=0$. In the first case let $f=1$. Then $\chi_\mu=\chi_\nu=1$, $\chi_0=1$, $\chi_\delta=1$, $\chi_\alpha=1$ $(x=p,\ y=0)$. In the second case let $f=p$. Then $x=-1$, $y=-1$, and $\chi_0=\varepsilon$, $\chi_\delta=-\varepsilon$, $\chi_\alpha=\varepsilon$, or in general $\chi_\varkappa=\varepsilon_\varkappa$; also $\chi_\mu=\chi_\nu=\varepsilon_1=0$.

III. For all other solutions $x=0$. By 1), therefore, either all $\chi_\alpha=0$ for which $\varepsilon_\alpha=-1$, or all those for which $\varepsilon_\alpha=+1$. First let $\chi_\alpha=0$ if $\varepsilon_\alpha=-1$, and let $f=p+1$. By 3), if $\varepsilon_\alpha=+1$, one has $\chi_\alpha\chi_\mu=\chi_\alpha$, and since not all $\chi_\alpha$ can vanish, $\chi_\mu=+1$. If $\varepsilon_\alpha=\varepsilon_\beta=+1$, and therefore also $\varepsilon_\gamma=\varepsilon_\delta=+1$, then by 2) and 4)
\[
\chi_\alpha\chi_\beta=\chi_\gamma+\chi_\delta,\qquad
\chi_\alpha^2=\chi_{(2\alpha^2-1)}+2.
\]
Thus, if $a$ and $b$ are the indices of the classes $(\alpha)$ and $(\beta)$, put $\chi_\alpha=\xi_a$, $\chi_\beta=\xi_b$, and $\xi_0=2$. Then by (5.)
\[
\xi_a\xi_b=\xi_{a+b}+\xi_{a-b},
\]
even when $b=a$. If $\rho$ is a new unknown and $\xi_1=\rho+\rho^{-1}$, then from $\xi_1\xi_1=\xi_2+\xi_0$ it follows that $\xi_2=\rho^2+\rho^{-2}$; then from $\xi_1\xi_2=\xi_3+\xi_1$ it follows that $\xi_3=\rho^3+\rho^{-3}$, and in general that $\xi_a=\rho^a+\rho^{-a}$. From $\xi_{\frac12(p-1)}=\xi_0$ it then follows that
\begin{equation*}\tag{9.}
\rho^{\frac12(p-1)}=1.
\end{equation*}
Thus one reaches the solutions
\[
\begin{array}{ll}
\chi_\lambda=p+1,\quad \chi_\mu=\chi_\nu=1,\quad
\chi_\alpha=\rho^a+\rho^{-a},&
\text{if }\varepsilon_\alpha=+1\text{ and }2\alpha\equiv r^a+r^{-a},\\
\chi_\beta=0,&
\text{if }\varepsilon_\beta=-1.
\end{array}
\]
The values so obtained also satisfy all the equations, except that when $\rho=1$ they do not satisfy the equation $x=0$, and when $\frac12(p-1)$ is even and $\rho=-1$ they do not satisfy the equation $y=-1$. $(5^*)$. Their number is therefore
\begin{equation*}\tag{10.}
n=\frac14(p-\varepsilon)-1.
\end{equation*}

IV. Finally, let $\chi_\alpha=0$ if $\varepsilon_\alpha=+1$, and let $f=p-1$. Then $\chi_\beta\chi_\mu=-\chi_\beta$ and consequently $\chi_\mu=-1$. If $\varepsilon_\alpha=\varepsilon_\beta=-1$, then
\[
\chi_\alpha\chi_\beta=-\chi_\gamma-\chi_\delta,\qquad
\chi_\alpha^2=-\chi_{(2\alpha^2-1)}+2.
\]
Then, if $a$ and $b$ are the indices of the classes $(\alpha)$ and $(\beta)$, put $\chi_\alpha=-\xi_a$, $\chi_\beta=-\xi_b$ and $\xi_0=2$. This gives
\[
\xi_a\xi_b=\xi_{a+b}+\xi_{a-b}.
\]
Thus, if $\xi_1=\sigma+\sigma^{-1}$, then $\xi_b=\sigma^b+\sigma^{-b}$; and since $\xi_{\frac12(p+1)}=\xi_0$,
\begin{equation*}\tag{11.}
\sigma^{\frac12(p+1)}=1.
\end{equation*}
Thus one reaches the solutions
\[
\begin{array}{ll}
\chi_\lambda=p-1,\quad \chi_\mu=\chi_\nu=-1,\quad
\chi_\beta=-\sigma^b-\sigma^{-b},&
\text{if }\varepsilon_\beta=-1\text{ and }2\beta\equiv s^b+s^{-b},\\
\chi_\alpha=0,&
\text{if }\varepsilon_\alpha=+1.
\end{array}
\]
For $\sigma$ the value $+1$ is inadmissible, and if $\frac12(p+1)$ is even, so is the value $-1$. The number of these solutions is therefore
\begin{equation*}\tag{12.}
n=\frac14(p-\varepsilon)-\frac12(1-\varepsilon).
\end{equation*}
Thus all $k=\frac12(p-1)+3$ characters have been determined. For each there remains to determine a number $e$ by means of the formula
\[
\frac{hf}{e}=\sum h_\alpha\chi_\alpha\chi_\alpha'.
\]
One finds that in all $k$ cases $e=f$. The proportionality factors have therefore been chosen in accordance with the rule stated in \S~3.

I shall collect the $k$ characters once more, using the notation introduced in \S~5:
\[
\begin{array}{c|ccccc}
n & 1 & 1 & 2 & \frac14(p-\varepsilon)-1 & \frac14(p-\varepsilon)-\frac12(1-\varepsilon)\\ \hline
\chi(E) & 1 & p & \frac12(p+\varepsilon) & p+1 & p-1\\
\chi(P) & 1 & 0 & \frac12(\varepsilon\pm\sqrt{\varepsilon p}) & 1 & -1\\
\chi(Q) & 1 & 0 & \frac12(\varepsilon\mp\sqrt{\varepsilon p}) & 1 & -1\\
\chi(R^a) & 1 & 1 & \frac12(1+\varepsilon)(-1)^a & \rho^a+\rho^{-a} & 0\\
\chi(S^b) & 1 & -1 & -\frac12(1-\varepsilon)(-1)^b & 0 & -\sigma^b-\sigma^{-b}
\end{array}
\]
For $p=3$ and $p=5$ this gives the characters of the group of the tetrahedron and of the icosahedron stated in \S~8.

\begin{center}
\rule{11em}{0.4pt}\\[-0.2em]
Issued on 27 August.\\[-0.2em]
\rule{11em}{0.4pt}

{\small Berlin, printed in the Reich Printing Office.}
\end{center}
% ===== END WORK 053 =====

\clearpage

% ===== BEGIN WORK 054: 054_ber_die_Primfactoren_der_Gruppendetermina =====
\begin{center}
{\Large On the Prime Factors of the Group Determinant.}\\[0.65em]
By G. Frobenius.\\[0.6em]
\rule{0.22\linewidth}{0.4pt}
\end{center}

The theory of the characters of a group, whose foundations I developed in my last paper, requires for its further development the investigation of a determinant whose degree is equal to the order of the group. Following Dedekind, who first recognized its significance for the theory of groups and directed my attention to it, I call it the \emph{group determinant} corresponding to the group. I use the $h$ elements $A,B,C,\ldots$ of the group $\Hgrp$ as indices for $h$ independent variables $x_A,x_B,x_C,\ldots$. In choosing this notation, I make the convention that, if $L=MN$, then also $x_L=x_{MN}$. From these $h$ quantities, which are distinguished from one another by one index, I form $h^2$ quantities furnished with two indices, by setting $x_{P,Q}=x_{PQ^{-1}}$. If $G_1,G_2,\ldots,G_h$ are the $h$ elements of $\Hgrp$ in some definite order, I consider the matrix $(x_{P,Q})=(x_{PQ^{-1}})$, whose $h$ rows are obtained by successively putting for $P$ the $h$ elements $G_1,G_2,\ldots,G_h$, and whose $h$ columns are obtained by putting for $Q$ the same $h$ elements in the same order. This matrix has certain symmetry properties conditioned by the constitution of the group $\Hgrp$. In every row the $h$ variables all occur, and likewise in every column. The various rows (columns) differ from one another only by the arrangement of the variables. The group determinant corresponding to the group $\Hgrp$ is the determinant of this matrix
\[
\Theta=|x_{P,Q}|=|x_{PQ^{-1}}|.
\]

If to the elements of the first row one adds those of all the other rows, then those elements all become equal to $\sum x_R=\xi$. Hence the integral function $\Theta$ of degree $h$ in the $h$ variables $x_A,x_B,x_C,\ldots$ is divisible by the linear function $\xi$. Thus, apart from the trivial case $h=1$, $\Theta$ always decomposes into factors of lower degree. The number $k$ of the distinct irreducible factors, or prime factors, of $\Theta$ is equal to the number of classes of conjugate elements into which the elements of $\Hgrp$ split. If $f$ is the degree of such a prime factor $\Phi$, then $\Theta$ is divisible by the $f$th power, and by no higher power, of $\Phi$. The degree $f$ is a divisor of the order $h$. By a linear substitution $\Phi$ can be transformed into a function of $f^2$, but not of fewer, variables; and if each prime factor of
\[
\Theta=\Pi \Phi^f
\]
is transformed in this way, then the $\sum f^2=h$ new variables are all independent of one another. If one always sets equal those variables $x_R$ whose indices are elements of the same class, then
\[
\Phi=\xi^f
\]
becomes the $f$th power of a linear function $\xi$ of $k$ independent variables, and the $k$ linear functions that thus arise from the $k$ prime factors of $\Theta$ are linearly independent. From the coefficients of the linear function $\xi$ there results a character $\chi$ of the group $\Hgrp$, and from its $k$ values $\chi(R)$ all the coefficients of the prime function $\Phi$ can be calculated. The theory of the general group determinant, in which the $h$ quantities $x_R$ are unrestricted variables, is thus reduced to the theory of the special group determinant, in which the variability of these quantities is restricted by the conditions $x_{BA}=x_{AB}$. But the calculation of this determinant of degree $h$,
\[
\Theta=\Pi \xi^{f^2},
\]
can be reduced to that of a determinant of degree $k$,
\[
\left|\sum_\gamma \frac{1}{h_\alpha}h_{\alpha\beta\gamma}x_\gamma\right|=\Pi\xi,
\]
in which the linear factor $\xi$, which occurs in $\Theta$ raised to the power $f^2$, is contained only simply. The definition of the positive integers $h_{\alpha\beta\gamma}$ and the development of their relations to the characters form the principal content of my paper \emph{On Group Characters}, Sitzungsberichte 1896, which in what follows I shall cite as \emph{Ch.}

Quite analogous properties to those of such a prime factor $\Phi$ of a group determinant are possessed by the integral function of degree $2^p$ in $2^{2p}$ variables, which I investigated in my paper \emph{On Theta Functions of Several Variables} (Crelle's Journal, vol. 96), although there I was led to it not by the consideration of the group of relations existing among the theta functions, but by the addition theorem valid for these functions. Otherwise the group determinant had hitherto been investigated only for the case of commutative groups, where all its prime factors are linear. For some particularly simple non-commutative groups Dedekind in the year 1886 found, by calculation, that the determinant $\Theta$ decomposes into prime factors; and his interesting results, which he recently communicated to me, have induced me to investigate generally the decomposition of the group determinant into prime factors for an arbitrarily given group.

\section*{§ 1.}

I denote the determinant of the matrix of degree $h$
\frobtagc{1.}{(x_{P,Q})=(x_{PQ^{-1}})=(x)}
by
\frobtagc{2.}{|x_{P,Q}|=|x_{PQ^{-1}}|=\Theta(x_E,x_A,x_B,x_C,\ldots)=\Theta(x_R)=\Theta(x)=\Theta.}
By $E$ I always mean the principal element. In the symbol $\Theta(x_R)$, $R$ denotes a variable element, for which the $h$ elements $E,A,B,C,\ldots$ of the group $\Hgrp$ are to be put. In the use of the notation $(x)$ or $\Theta(x)$, $x$ is an empty symbol, which receives a meaning only from the fact that the indices $E,A,B,C,\ldots$ are attached to it.

Now let $y_E,y_A,y_B,y_C,\ldots$ be a second system of $h$ independent variables. From them I form the matrix
\frobtagc{3.}{(y_{P,Q})=(y_{PQ^{-1}})=(y).}
Its rows (columns) are obtained by putting for $P$ ($Q$) the $h$ elements $G_1,G_2,\ldots,G_h$ of $\Hgrp$ in the same order in which they were used in forming the matrix (1.).

From those two systems, each of $h$ variables $x_R$ and $y_R$, I form a third system $z_R$ by setting
\frobtagc{4.}{z_A=\sum x_Ry_S\qquad (RS=A).}
In this sum the $h$ elements of $\Hgrp$ are to be put for $R$, and each element $R$ is to be joined with the element $S$ $(=R^{-1}A)$ that satisfies the condition $RS=A$ (not $SR=A$), so that $S$ too runs through the $h$ elements of $\Hgrp$, each time joined with $R=AS^{-1}$. Then, by the composition of the two matrices (1.) and (3.), which have the symmetry properties assumed here and conditioned by the constitution of the group $\Hgrp$, there results the matrix
\frobtagc{5.}{(z_{P,Q})=(z_{PQ^{-1}})=(z)=(x)(y)}
with the same symmetry properties. For it is
\[
z_{P,Q}=\sum_R x_{P,R}y_{R,Q}=\sum_R x_{PR^{-1}}y_{RQ^{-1}}.
\]
If in this sum one sets $R=SQ$, then $S$ runs simultaneously with $R$ through the $h$ elements of $\Hgrp$, only in another order, and by (4.) one obtains
\[
z_{P,Q}=\sum_S x_{PQ^{-1}S^{-1}}y_S=z_{PQ^{-1}}.
\]

The same theorem holds when arbitrarily many such matrices are composed. If, for example, $z_E,z_A,z_B,z_C,\ldots$ are $h$ arbitrary quantities, and one sets
\[
v_A=\sum x_Ry_Sz_T,\qquad (RST=A),
\]
then the matrix $(v_{PQ^{-1}})=(v)=(x)(y)(z)$ is composed from the three matrices $(x)$, $(y)$, and $(z)$ in this order. If the matrix $(x)$ is composed with itself $n$ times, let the matrix
\[
(x_{P,Q})^n=(x)^n=(x_{PQ^{-1}}^{(n)})=(x^{(n)})
\]
result. Then
\frobtagc{6.}{x_R^{(n)}=\sum x_{R_1}x_{R_2}\cdots x_{R_n}\qquad (R_1R_2\cdots R_n=R).}

Consequently every function of the matrix $(x)$ also has the symmetry properties assumed here; for example the matrix adjugate to $(x)$ (\emph{On Linear Substitutions and Bilinear Forms}, Crelle's Journal, vol. 84, p. 7), and likewise the principal matrix (identity matrix)
\[
(x)^0=(\varepsilon_{P,Q})=(\varepsilon_{PQ^{-1}})=(\varepsilon),
\]
where $\varepsilon_R=0$ except when $R=E$, and $\varepsilon_E=1$.

Now let $\Phi,\Phi',\Phi'',\ldots$ be the different indecomposable functions (prime functions) occurring in
\frobtagc{7.}{\Theta=\Phi^e\Phi'^{e'}\Phi''^{e''}\cdots=\Pi\Phi^e,}
and let $f,f',f'',\ldots$ be the degrees of these integral homogeneous functions of the $h$ variables $x_E,x_A,x_B,x_C,\ldots$. In the group determinant $\Theta$, the diagonal elements, and only these, are equal to $x_E$. Hence $\Theta$ is reduced to $x_E^h$ when all variables except $x_E$ are set equal to zero, and consequently $\Phi$ is then also reduced to the $f$th power of $x_E$. Hence one can choose the still undetermined constant factor of $\Phi$ so that in this function $x_E^f$ has coefficient 1.

By the multiplication theorem for determinants, the relation
\frobtagc{8.}{\Theta(z)=\Theta(x)\Theta(y)}
follows from equations (2.) and (5.). From this one obtains, for every prime factor of $\Theta$,
\[
\Phi=\Phi(x_E,x_A,x_B,x_C,\ldots)=\Phi(x_R)=\Phi(x),
\]
the analogous relation
\frobtagc{9.}{\Phi(z)=\Phi(x)\Phi(y),\quad \text{when }(z)=(x)(y),}
by which the function $\Phi$ can be characterized independently of its relation to the group determinant $\Theta$. For, if the right-hand side of equation (8.) is decomposed into prime factors, it follows that $\Phi(z)$ splits into the product of a function $\Lambda(x)$ of the $h$ variables $x_R$ alone and a function $M(y)$ of the $h$ variables $y_R$ alone. If one then sets $y_R=\varepsilon_R$ in the equation $\Phi(z)=\Lambda(x)M(y)$, then $z_R=x_R$, and hence $\Phi(x)=\Lambda(x)M(\varepsilon)$. Likewise $M(y)=\Lambda(\varepsilon)\Phi(y)$ and $\Lambda(\varepsilon)M(\varepsilon)=\Phi(\varepsilon)=1$.

Conversely, every indecomposable integral homogeneous function $\Phi$ of $x_E,x_A,x_B,\ldots$ that satisfies condition (9.) must be a factor of the group determinant $\Theta(x)$. For, if in this equation one puts for $(y)$ the matrix adjugate to $(x)$, then $z_R=\varepsilon_R\Theta(x)$, hence
\[
\Phi(z)=\Theta(x)^f=\Phi(x)\Phi(y),
\]
where $y_A$ is an integral function of the $h$ variables $x_R$. Hence the function $\Phi(x)$, because it is indecomposable, must be a factor of $\Theta(x)$.

With the help of relation (9.) all properties of determinants that flow from the multiplication theorem can be transferred to the prime factors of the group determinant, especially the properties that I developed in my paper, cited below as \emph{V.}, \emph{On Commuting Matrices} (p. 601 of this volume).

\section*{§ 2.}

Every prime factor $\Phi(x)$ of the group determinant satisfies the condition
\frobtagc{1.}{\Phi(z)=\Phi(x)\Phi(y),}
provided
\frobtagc{2.}{z_C=\sum x_Ay_B\qquad (AB=C)}
is set. With the help of this relation the linear factors
\frobtagc{3.}{\Phi(x)=\sum\chi(A)x_A}
can be completely determined. For from the equation
\[
\left(\sum\chi(A)x_A\right)\left(\sum\chi(B)y_B\right)=\left(\sum\chi(C)z_C\right)=\sum\chi(AB)x_Ay_B
\]
one obtains for the coefficients $\chi(A)$ of these functions the relation
\frobtagc{4.}{\chi(AB)=\chi(A)\chi(B).}
Thus $\chi(E)=1$, $\chi(A)\chi(A^{-1})=1$, and more generally
\[
\chi(ABCD\cdots)=\chi(A)\chi(B)\chi(C)\chi(D)\cdots,
\]
and consequently
\[
\begin{aligned}
\chi(B^{-1}A^{-1}BA)
&=\chi(B^{-1})\chi(A^{-1})\chi(B)\chi(A)\\
&=1.
\end{aligned}
\]
The element $F$ that is obtained from $A$ and $B$ by means of the equation
\frobtagc{5.}{BA=ABF}
I call, following Dedekind, the \emph{commutator} of $A$ and $B$. Accordingly $\chi(F)=1$ for every commutator $F$ of any two elements of the group $\Hgrp$. If $T$ is an arbitrary element of $\Hgrp$, and if
\[
T^{-1}AT=A',\qquad T^{-1}BT=B',\qquad T^{-1}FT=F',
\]
then also $B'A'=A'B'F'$. Thus, if $F$ is a commutator, every element $F'$ conjugate to $F$ is also one. If the elements of $\Hgrp$ are divided into classes of conjugate elements, then the commutators are formed from all the elements of some of these classes. The group $\Ggrp$ generated by them is therefore an invariant subgroup of $\Hgrp$. (It may also be equal to $\Hgrp$, or also equal to the principal group $\Egrp$, the latter always and only when $\Hgrp$ is a commutative group.) If $G$ is an element of $\Ggrp$, then there are commutators $F,F',F'',\ldots$ (which need not be distinct) such that $G=FF'F''\cdots$. Hence $\chi(G)=\chi(F)\chi(F')\chi(F'')\cdots=1$.

Now let
\[
\Hgrp=\Ggrp A+\Ggrp B+\Ggrp C+\cdots,
\]
so that $A,B,C,\ldots$ are the elements of $\Hgrp$ distinct modulo $\Ggrp$. Their number is $\frac{h}{g}$, if $g$ is the order of $\Ggrp$. The $\frac{h}{g}$ complexes $\Ggrp A=A\Ggrp$, $\Ggrp B=B\Ggrp$, $\ldots$ form a group, which is denoted by $\frac{\Hgrp}{\Ggrp}$. If $\Ggrp$ is the commutator group, then $\frac{\Hgrp}{\Ggrp}$ is a commutative (Abelian) group; and, in order that $\frac{\Hgrp}{\Ggrp}$ be a commutative group, it is necessary and sufficient that $\Ggrp$ be divisible by the commutator group. For if $\Ggrp A$ and $\Ggrp B$ are two elements of $\frac{\Hgrp}{\Ggrp}$, then there is in $\Ggrp$ an element $F$ such that $BA=ABF$, and hence also $\Ggrp BA=\Ggrp ABF$. Now $\Ggrp ABF=(\Ggrp A)(\Ggrp B)(\Ggrp F)$ and $\Ggrp F=\Ggrp$, so
\[
(\Ggrp B)(\Ggrp A)=(\Ggrp A)(\Ggrp B).
\]
Dedekind found these properties of the commutator group in the year 1880. They were first published, however, by Miller, \emph{The regular substitution groups whose order is less than 48}. Quarterly Journal of Math. 1896, vol. 28, p. 266.

If $G$ is any element of $\Ggrp$, then $\chi(GA)=\chi(G)\chi(A)=\chi(A)$. Hence $\chi(R)$ has the same value for all elements $R$ of the complex $\Ggrp A$. Thus one may also assign the number $\chi(A)$ to the complex $\Ggrp A$. Since these complexes form a commutative group, and the numbers $\chi(A)$ assigned to them have property (4.), the numbers $\chi(A),\chi(B),\chi(C),\ldots$ form a character of the commutative group $\frac{\Hgrp}{\Ggrp}$. It is well known that for such a group there are always $\frac{h}{g}$ distinct characters, all of whose values are roots of unity. If $\psi$ is one of them, and for every element $R$ contained in the complex $\Ggrp A$ one sets $\chi(R)=\psi(\Ggrp A)$, then for every element $G$ of the group $\Ggrp$ one has $\chi(G)=1$, and equation (4.) holds for any two elements of $\Hgrp$. Furthermore the function (3.), whose coefficients are these values $\chi(A)$, is then a linear factor of the group determinant $\Theta$. For if one sets $y_R=\chi(R)x_R$, then
\[
y_{PQ^{-1}}=\chi(PQ^{-1})x_{PQ^{-1}}=\chi(P)\chi(Q^{-1})x_{PQ^{-1}},
\]
and therefore $|y_{PQ^{-1}}|=|x_{PQ^{-1}}|$. This determinant, however, contains the factor $\sum y_R=\sum\chi(R)x_R=\Phi$, and indeed only in the first power. For if one adds the elements of all rows to those of the first row, then they all become equal to $\sum y_R=\Phi$, and after the factor $\Phi$ is then removed, all equal to 1. If now one subtracts the elements of the first column from those of the following columns, one recognizes that $\Theta:\Phi$ depends only on the differences $y_R-y_A$. Thus this quotient cannot be divisible once more by the sum $\sum y_R$.

Consequently the number of linear factors of the group determinant is equal to the quotient of the order of the group by the order of its commutator group, and every linear factor is contained in $\Theta$ only in the first power.

Dedekind found this theorem by induction. There is always one linear factor, namely $\sum x_R$. The corresponding character, $\chi(R)=1$ for every element $R$, is called the principal character. If $\Ggrp=\Hgrp$, then there is no other linear factor. This must always occur when $\Hgrp$ is a simple group whose order is a composite number.

\section*{§ 3.}

Now choose an arbitrary integer $f\le h$ and try to form an integral function $\Phi$ of degree $f$ in the $h$ variables $x_E,x_A,x_B,\ldots$, satisfying condition (9.) of § 1. In it the coefficient of $x_E^f$ must be equal to 1. For if one sets $y_R=\varepsilon_R$, then $z_R=x_R$, hence $\Phi(x)=\Phi(x)\Phi(\varepsilon)$, and therefore $\Phi(\varepsilon)=1$. Now, when $R$ is different from $E$, I denote the coefficient of $x_E^{f-1}x_R$ in $\Phi(x)$ by $\chi(R)$, but set
\frobtagc{1.}{\chi(E)=f.}
For every element $R$, $\chi(R)$ is then the coefficient of $x_E^{f-1}$ in $\frac{\partial\Phi}{\partial x_R}$. If $u$ is a variable, then for abbreviation I set
\[
\Phi(x+u\varepsilon)=\Phi(x_E+u,x_A,x_B,x_C,\ldots).
\]
The symbol $u$, which denotes a variable quantity, is here to be sharply distinguished from the empty symbols $x$ and $\varepsilon$, which acquire meaning only by appending the indices $E,A,B,C,\ldots$.

If now
\frobtagc{2.}{\Phi(x+u\varepsilon)=u^f+\Phi_1u^{f-1}+\Phi_2u^{f-2}+\cdots+\Phi_f,}
then $\Phi_n$ is an integral homogeneous function of degree $n$ in the $h$ variables $x_E,x_A,x_B,x_C,\ldots$, and in particular $\Phi_f=\Phi$, and
\frobtagc{3.}{\Phi_n=\frac{1}{(f-n)!}\frac{\partial^{f-n}\Phi}{\partial x_E^{f-n}},\qquad \frac{\partial\Phi_{n+1}}{\partial x_E}=(f-n)\Phi_n.}
Finally
\frobtagc{4.}{\Phi_1=\sum\chi(R)x_R.}
I regard the $h$ constants $\chi(R)$ as the values of a function $\chi$. If $\Phi$ is indecomposable, I call $\chi$ the character of degree $f$ of the group $\Hgrp$ corresponding to the prime function $\Phi$ of degree $f$.

If, in particular, $\sum\psi(R)x_R$ is a linear factor of $\Theta$, then $\psi(R)$ is called a character of first degree of $\Hgrp$. For this, as shown above, the relation
\[
\psi(A)\psi(B)=\psi(AB)
\]
is the necessary and sufficient condition. If one sets $\psi(R)x_R=y_R$, then $\Theta(y)=\Theta(x)$. Thus, if $\Phi(x)$ is a prime factor of degree $f$ of $\Theta(x)$, then $\Phi(y)$ is also such a factor. In this latter function the coefficient of $x_E^{f-1}x_R$ is equal to $\chi(R)\psi(R)$. Hence the following theorem holds:

\emph{If $\chi(R)$ is a character of degree $f$ and $\psi(R)$ a character of first degree of $\Hgrp$, then $\chi(R)\psi(R)$ is also a character of degree $f$ of $\Hgrp$.}

This new character is equal to $\chi(R)$ when $\psi(R)$ is the principal character. But in other cases also it need not be different from $\chi(R)$, namely when $\chi(R)=0$ for every element $R$ for which $\psi(R)$ is different from 1.

Let $u_1,u_2,\ldots,u_f$ be the $f$ values of $-u$ for which
\frobtagc{5.}{\Phi(x+u\varepsilon)=(u+u_1)(u+u_2)\cdots(u+u_f)}
vanishes. Let $a,v_1,v_2,\ldots,v_n$ be constants and
\[
g(u)=a(u+v_1)(u+v_2)\cdots(u+v_n)
\]
an integral function of $u$. Then, if the variable $u$ is replaced by the matrix $(x)$, one has also
\[
(y)=g((x))=a((x)+v_1(\varepsilon))((x)+v_2(\varepsilon))\cdots((x)+v_n(\varepsilon)).
\]
Thus, by repeated application of relation (9.) of § 1, one finds
\[
\Phi(y)=a^f\Phi(x+v_1\varepsilon)\Phi(x+v_2\varepsilon)\cdots\Phi(x+v_n\varepsilon).
\]
If here one sets
\[
\Phi(x+v\varepsilon)=(u_1+v)(u_2+v)\cdots(u_f+v),
\]
then the theorem results: if the matrix $(y)=g((x))$ is an integral function of the matrix $(x)$, then
\[
\Phi(y)=g(u_1)g(u_2)\cdots g(u_f).
\]
If here $g(u)$ is replaced by $g(u)+v$, where $v$ is a parameter, while the abbreviation $(y)$ for $g((x))$ is retained, one obtains the equation
\[
\Phi(y+v\varepsilon)=(v+g(u_1))(v+g(u_2))\cdots(v+g(u_f)).
\]
For example, if $g(u)=u^n$, then
\[
\Phi(x^{(n)}+v\varepsilon)=(v+u_1^n)(v+u_2^n)\cdots(v+u_f^n).
\]
By comparing the coefficients of $v^{f-1}$, one obtains from this, by (4.),
\frobtagc{6.}{S_n=\sum_R \chi(R)x_R^{(n)},}
where $S_n$ is the sum of the $n$th powers of the $f$ quantities $u_1,u_2,\ldots,u_f$. By formula (6.) of § 1 one may also write for this
\frobtagc{7.}{S_n=\sum_{R_1,R_2,\ldots,R_n}\chi(R_1R_2\cdots R_n)x_{R_1}x_{R_2}\cdots x_{R_n},}
where each of the $n$ summation letters $R_1,R_2,\ldots,R_n$ runs independently of the others through the $h$ elements of $\Hgrp$.

From the power sums $S_n$, however, one can calculate the coefficients $\Phi_n$ of the function (2.) by the formula
\frobtag{8.}{(-1)^n\Phi_n=\sum\frac{(-1)^{a+b+c+\cdots}S_1^aS_2^bS_3^c\cdots}{1^a2^b3^c\cdots a!b!c!\cdots},\qquad (a+2b+3c+\cdots=n)}
where $a,b,c,\ldots$ run through all integers $(\ge0)$ that satisfy the condition $a+2b+3c+\cdots=n$. This formula also holds when $n>f$, provided that then $\Phi_n=0$ is set. With its help we shall represent the functions $\Phi_2,\Phi_3,\ldots$ and especially $\Phi_f=\Phi$. The coefficients of $\Phi$ are integral functions of the $h$ constants $\chi(R)$ whose coefficients are rational numbers. If one chooses $n>f$, then relations result from that formula which the $h$ quantities $\chi(R)$ must satisfy. Before I turn to these calculations, however, I must first state an important property of the function $\chi(R)$.

If in equation (9.) of § 1 one sets the variables $y_R=0$ except one, $y_A$, then one obtains
\frobtagc{9.}{\Phi(x_{RA^{-1}})=\zchar(A)\Phi(x_R),}
where $\zchar(A)$ denotes the coefficient of $x_A^f$ in $\Phi(x_R)$. If, however, one sets the variables $x_R=0$ except one, $x_A$ (and then replaces the letter $y$ by $x$), one obtains
\frobtagc{10.}{\Phi(x_{A^{-1}R})=\zchar(A)\Phi(x_R).}
If here, when $B$ is a fixed element, the variable $x_R$ is replaced by $x_{B^{-1}R}$ for every index $R$, one finds
\[
\Phi(x_{B^{-1}A^{-1}R})=\zchar(A)\Phi(x_{B^{-1}R})=\zchar(A)\zchar(B)\Phi(x_R).
\]
If, on the other hand, $A$ is replaced by $AB$ in equation (10.), then
\[
\Phi(x_{B^{-1}A^{-1}R})=\zchar(AB)\Phi(x_R).
\]
Thus $\zchar(AB)=\zchar(A)\zchar(B)$, and consequently $\zchar(R)$ is a character of first degree of the group $\Hgrp$, a root of unity. If $A$ is a fixed element, and for every $R$ one sets $y_R=x_{RA^{-1}}$, then
\[
y_{P,Q}=y_{PQ^{-1}}=x_{PQ^{-1}A^{-1}}=x_{P,AQ}.
\]
But if, in the determinant $|x_{P,Q}|$, one replaces $Q$ by $AQ$, then only the order of the $h$ columns is changed thereby. If $m$ is the order of the element $A$, then that permutation of the $h$ columns consists entirely of cycles of order $m$. The number of these cycles is therefore $\frac{h}{m}$, and hence the permutation is even or odd according as $h-\frac{h}{m}$ is even or odd. Hence
\[
\Theta(x_{RA^{-1}})=(-1)^{h-\frac{h}{m}}\Theta(x_R),
\]
and therefore
\frobtagc{11.}{\Pi(\zchar A)^e=(-1)^{h-\frac{h}{m}}.}
If in equation (10.) one replaces $x_R$ by $x_{RA}$, one obtains
\[
\Phi(x_{A^{-1}RA})=\zchar(A)\Phi(x_{RA})=\zchar(A)\zchar(A^{-1})\Phi(x_R),
\]
and hence
\frobtagc{12.}{\Phi(x_{A^{-1}RA})=\Phi(x_R).}
The function $\Phi(x_R)$ therefore remains unchanged if, for every index $R$, the variable $x_R$ is replaced by $x_{A^{-1}RA}$, where $A$ is a fixed element of $\Hgrp$. In this process the variable $x_E$ remains unchanged. By comparing the coefficients of $x_E^{f-1}x_R$, however, one obtains from (12.) $\chi(ARA^{-1})=\chi(R)$, or, if one sets $R=BA$,
\frobtagc{13.}{\chi(AB)=\chi(BA).}
Thus, if the $h$ elements of $\Hgrp$ are divided into classes of conjugate elements, $\chi(R)$ has the same value for all elements $R$ of one class.

I now turn to the calculation of the functions $\Phi_n$ with the help of formula (8.). First, by (6.),
\[
\begin{aligned}
\Phi_1&=S_1=\sum\chi(A)x_A,\\
2\Phi_2&=S_1^2-S_2=\sum(\chi(A)\chi(B)-\chi(AB))x_Ax_B,\\
6\Phi_3&=S_1^3-3S_1S_2+2S_3\\
&=\sum\bigl(\chi(A)\chi(B)\chi(C)-\chi(A)\chi(BC)-\chi(B)\chi(AC)\\
&\qquad-\chi(C)\chi(AB)+\chi(ABC)+\chi(ACB)\bigr)x_Ax_Bx_C.
\end{aligned}
\]
I therefore set
\frobtag{14.}{\begin{aligned}
\chi(A,B)&=\chi(A)\chi(B)-\chi(AB),\\
\chi(A,B,C)&=\chi(A)\chi(B)\chi(C)-\chi(A)\chi(BC)-\chi(B)\chi(AC)\\
&\qquad-\chi(C)\chi(AB)+\chi(ABC)+\chi(ACB).
\end{aligned}}
This expression is symmetric in $A,B,C$, because $\chi(ABC)$ remains unchanged under cyclic interchange of the elements $A,B,C$ by (13.). The general law of formation for the coefficients of the function
\frobtagc{15.}{n!\Phi_n(x)=\sum_{R_1,R_2,\ldots,R_n}\chi(R_1,R_2,\ldots,R_n)x_{R_1}x_{R_2}\cdots x_{R_n}}
is somewhat complicated. Let $A,B,C,D,F,G,H,\ldots,L,M$ be any $n$ elements, distinct or equal, among the elements of $\Hgrp$. Form the $n!$ permutations of $n$ symbols and decompose each of them into cyclic factors. If for the $n$ symbols one then puts the $n$ elements $A,B,C,\ldots,L,M$, let, say,
\frobtagc{16.}{(ABCD)(FGH)\cdots(LM)}
be one of these $n!$ permutations. To it one assigns the product
\frobtagc{17.}{\pm\chi(ABCD)\chi(FGH)\cdots\chi(LM),}
where the sign $+$ or $-$ is to be chosen according as the permutation (16.) is even or odd. In the permutation (16.) the symbol $(FGH)$ means that the three symbols $F,G,H$ are to be cyclically interchanged; in the expression (17.), however, $FGH$ means the product of the three elements $F,G,H$.

A given permutation can be represented as a product of cyclic factors in only one way. Yet the individual cycles may be arranged arbitrarily and, within each cycle, the symbols may be cyclically interchanged without changing its meaning. Other transpositions, however, are not allowed. Thus the permutation (16.) is equal to
\[
(GHF)(ML)\cdots(CDAB).
\]
In the corresponding product (17.), however, the same alterations are permitted, for $\chi(ABCD)$ remains unchanged when one cyclically interchanges the elements $A,B,C,D$. Let the sum of the $n!$ products corresponding to the $n!$ permutations be
\frobtagc{18.}{\chi(A,B,C,\cdots L,M)=\sum^{(n)}(\pm\Pi\chi).}
This function remains unchanged when the $n$ elements
\[
A,B,C,\cdots L,M
\]
are arbitrarily interchanged among themselves. If the permutation (16.) consists of $l$ cycles, then it is even or odd according as $n-l$ is even or odd. Hence one may also write
\frobtagc{19.}{(-1)^n\chi(A,B,C,\cdots L,M)=\sum^{(n)}(\Pi-\chi).}
For example, regardless of whether the elements denoted by $A,B,C,D$ are distinct or not,
\frobtag{20.}{\begin{aligned}
\chi(A,B,C,D)={}&\chi(A)\chi(B)\chi(C)\chi(D)-\chi(B)\chi(C)\chi(AD)\nobreak\\
&{}-\chi(A)\chi(C)\chi(BD)-\chi(A)\chi(B)\chi(CD)\nobreak\\
&{}-\chi(A)\chi(D)\chi(BC)-\chi(B)\chi(D)\chi(AC)\nobreak\\
&{}-\chi(C)\chi(D)\chi(AB)+\chi(BC)\chi(AD)+\chi(AC)\chi(BD)\nobreak\\
&{}+\chi(AB)\chi(CD)+\chi(A)\chi(BCD)+\chi(B)\chi(ACD)\nobreak\\
&{}+\chi(C)\chi(ABD)+\chi(D)\chi(ABC)+\chi(A)\chi(BDC)\nobreak\\
&{}+\chi(B)\chi(ADC)+\chi(C)\chi(ADB)+\chi(D)\chi(ACB)\nobreak\\
&{}-\chi(ABCD)-\chi(ACBD)-\chi(BACD)\nobreak\\
&{}-\chi(BCAD)-\chi(CABD)-\chi(CBAD).
\end{aligned}}
I now form the sum
\[
(-1)^n\sum\chi(A,B,C,\cdots L,M)x_Ax_Bx_C\cdots x_Lx_M,
\]
in which each of the $n$ summation letters $A,B,C,\cdots L,M$ runs independently of the others through the $h$ elements of $\Hgrp$. The expression $(-1)^n\chi(A,B,C,\cdots L,M)$ is a sum of $n!$ products $\Pi(-\chi)$. Suppose that the permutation of the $n$ symbols from which one of these products is formed consists of $a$ cycles of one symbol, $b$ of two symbols, $c$ of three symbols, and so on, so that $a+2b+3c+\cdots=n$. If one then multiplies this product $\Pi(-\chi)$ by $x_Ax_Bx_C\cdots x_Lx_M$ and sums, one obtains by (7.)
\[
(-1)^{a+b+c+\cdots}S_1^aS_2^bS_3^c\cdots .
\]
This term occurs as many times as there are permutations which can be represented in the specified way as a product of cyclic factors, namely (\textsc{Cauchy}, \emph{Comptes rendus} tom. 21, p. 604)
\[
\frac{n!}{1^a2^b3^c\cdots a!b!c!\cdots}
\]
times. Therefore the sum under consideration is equal to
\[
n!\sum\frac{(-1)^{a+b+c+\cdots}S_1^aS_2^bS_3^c\cdots}{1^a2^b3^c\cdots a!b!c!\cdots}=n!(-1)^n\Phi_n.
\]
Thus formula (15.) is proved. Thus, if $n>f$, then
\frobtagc{21.}{\chi(R_1,R_2,\cdots R_n)=0\qquad(n>f).}
In each of the $(n+1)!$ products of the sum $\chi(R,R_1,R_2,\ldots R_n)$ place the factor which contains the element denoted by $R$ in the first position, and in this factor itself place $R$ in the first position with the aid of a cyclic interchange. Then take first the products which contain the factor $\chi(R)$, then those in which the element $R_1$ follows $R$, then those in which the element $R_2$ follows $R$, and so on. In this way one obtains the recursion formula
\frobtag{22.}{\begin{aligned}
\chi(R,R_1,R_2,R_3,\cdots R_n)={}&\chi(R)\chi(R_1,R_2,R_3,\cdots R_n)\nobreak\\
&{}-\chi(RR_1,R_2,R_3,\cdots R_n)-\chi(R_1,RR_2,R_3,\cdots R_n)\nobreak\\
&{}-\chi(R_1,R_2,RR_3,\cdots R_n)-\cdots-\chi(R_1,R_2,R_3,\cdots RR_n).
\end{aligned}}
It follows from this that, if for one value of $n$ the quantities $\chi(R_1,R_2,\cdots R_n)$ all vanish, the same must also occur for every larger value of $n$. In particular,
\frobtagc{23.}{\chi(E,R_1,R_2,\cdots R_n)=(f-n)\chi(R_1,R_2,\cdots R_n).}

\section*{§ 4.}

Differentiating the equation $\Phi(z)=\Phi(x)\Phi(y)$ with respect to $y_A$, one obtains
\[
\sum_R\frac{\partial\Phi(z)}{\partial z_R}x_{RA^{-1}}=\Phi(x)\frac{\partial\Phi(y)}{\partial y_A},
\]
and if one sets $y_R=\varepsilon_R$,
\frobtagc{1.}{\sum_R\frac{\partial\Phi(x)}{\partial x_R}x_{RA^{-1}}=\chi(A)\Phi(x).}
If, however, one differentiates that equation with respect to $x_A$, one finds in the same way
\frobtagc{2.}{\sum_R\frac{\partial\Phi(x)}{\partial x_R}x_{A^{-1}R}=\chi(A)\Phi(x).}
In these equations I replace $x_E$ by $x_E-u$. Then there results, since
\[
\frac{\partial\log\Phi(x-u\varepsilon)}{\partial u}=\frac{1}{u-u_1}+\cdots+\frac{1}{u-u_f}=\sum_{n=0}^{\infty}\frac{S_n}{u^{n+1}}
\]
is, the recursion formula
\frobtag{3.}{\frac{1}{n+1}\frac{\partial S_{n+1}}{\partial x_A}=\frac{1}{n}\sum_R\frac{\partial S_n}{\partial x_R}x_{RA^{-1}}=\frac{1}{n}\sum_R\frac{\partial S_n}{\partial x_R}x_{A^{-1}R},\qquad(n>0)}
from which a new proof of formula (7.) § 3 can be derived. If one differentiates equation (1.) with respect to $x_S$, then multiplies it by $x_{SB^{-1}}$ and sums with respect to $S$, one obtains the equation
\[
\sum_{R,S}\frac{\partial^2\Phi}{\partial x_R\partial x_S}x_{RA^{-1}}x_{SB^{-1}}=(\chi(A)\chi(B)-\chi(BA))\Phi,
\]
from which formula (13.) § 3 follows directly. Likewise, in general, if $A,B,\cdots M$ are any $n$ elements of $\Hgrp$,
\[
\sum_{R,S,\ldots,V}\frac{\partial^n\Phi}{\partial x_R\partial x_S\cdots\partial x_V}x_{RA^{-1}}x_{SB^{-1}}\cdots x_{VM^{-1}}=\chi(A,B,\cdots M)\Phi.
\]
Here I make a different use of those relations. If one sets
\[
y_{R^{-1}}=\frac{\partial\Phi(x-u\varepsilon)}{\partial x_R},
\]
then they read
\[
\sum_R y_{PR^{-1}}(x_{RQ^{-1}}-u\varepsilon_{RQ^{-1}})=\sum_R(x_{PR^{-1}}-u\varepsilon_{PR^{-1}})y_{RQ^{-1}}=\chi(QP^{-1})\Phi(x-u\varepsilon).
\]
I further set $\chi(R^{-1})=\chi'(R)$ and denote the matrix $(\chi'(PQ^{-1}))$ briefly by $(\chi')$. Then this formula expresses the following relation between matrices:
\frobtagc{4.}{(y)((x)-u(\varepsilon))=((x)-u(\varepsilon))(y)=(\chi')\Phi(x-u\varepsilon).}
Hence $(y)(x)=(x)(y)$. Thus, if, expanded according to powers of $u$, $(y)=(p)+(q)u+(r)u^2+\cdots$, then $(x)$ is commutable with each of the matrices $(p),(q),(r),\cdots$. If one now also expands $\Phi(x-u\varepsilon)$ in relation (4.) according to powers of $u$, the coefficients of the separate powers of $u$ on both sides must agree. The equations so obtained are then recombined after being multiplied, not by the powers of the variable $u$, but by the corresponding powers of the matrix $(x)$. This procedure leads to the same result as if in equation (4.) one directly replaces the variable $u$ by the matrix $(x)$. (This argument is developed more fully in V., p. 605.) Then one obtains $(\chi')\Phi(x-(x)\varepsilon)=0$, or more explicitly $(\chi')\Phi(x_E-(x),x_A,x_B,x_C,\cdots)=0$, and still more fully, by (2.) § 3,
\frobtag{5.}{(\chi')\bigl((x)^f-(x)^{f-1}\Phi_1+(x)^{f-2}\Phi_2-\cdots+(-1)^f(x)^0\Phi_f\bigr)=0.}
Multiplying also by $(x)^{n-f}$, one obtains
\frobtag{6.}{\sum_R\chi(AR)\bigl(x_R^{(n)}-x_R^{(n-1)}\Phi_1+x_R^{(n-2)}\Phi_2-\cdots+(-1)^f x_R^{(n-f)}\Phi_f\bigr)=0.}
Setting $n=f$ and determining the coefficient of $x_S^f$, one finds
\frobtag{7.}{\chi(AS^f)-\mathfrak{z}_1(S)\chi(AS^{f-1})+\mathfrak{z}_2(S)\chi(AS^{f-2})-\cdots+(-1)^f\mathfrak{z}_f(S)\chi(A)=0,}
where
\frobtagc{8.}{\mathfrak{z}_n(S)=\frac{1}{n!}\chi(S,S,\cdots S)}
is the coefficient of $x_S^n$ in $\Phi_n$, and so is independent of $A$. In particular, $\mathfrak{z}_1(S)=\chi(S)$ and $\mathfrak{z}_f(S)=\mathfrak{z}(S)$. If in the function $\Phi(x_E+u,x_A,x_B,\cdots,x_S,\cdots)$ all variables are set equal to zero except $u$ and $x_S$, it becomes
\frobtag{9.}{\Phi(u,0,0,\cdots x_S,0,\cdots)=u^f+\mathfrak{z}_1(S)u^{f-1}x_S+\mathfrak{z}_2(S)u^{f-2}x_S^2+\cdots+\mathfrak{z}_f(S)x_S^f.}
Therefore $\mathfrak{z}_n(E)=\binom{f}{n}$.

\section*{§ 5.}

The preceding results have been derived by me solely from relation (9.) § 1, without using the irreducibility of $\Phi$ and the exponent $e$ of the power of $\Phi$ by which the group determinant $\Theta$ is divisible. Each of the $h$ variables $x_R$ occurs once in each row and in each column of $\Theta$, altogether therefore in $h$ positions. At each of these $h$ positions, however, the same complementary minor corresponds to it, as I showed in \emph{Ch.} § 6. Thus, if $\Theta_{P,Q}$ is the minor complementary to the element $x_{P,Q}$ in the determinant $\Theta=|x_{P,Q}|$, then
\frobtagc{1.}{\Theta_{P,Q}=\frac{1}{h}\frac{\partial\Theta}{\partial x_{PQ^{-1}}}=\Theta_{PQ^{-1}},}
provided one sets $h\Theta_R=\frac{\partial\Theta}{\partial x_R}$. The minors of $\Theta$ therefore form a matrix which has the same symmetry properties as the matrix $(x)$. By the known relations between the elements of a determinant and their complementary minors,
\[
\sum_R x_{P,R}\Theta_{Q,R}=\sum_R x_{R,P}\Theta_{R,Q}=\varepsilon_{P,Q}\Theta
\]
or
\frobtagc{2.}{\sum_R x_{AR}\frac{\partial\Theta}{\partial x_R}=\sum_R x_{RA}\frac{\partial\Theta}{\partial x_R}=\varepsilon_A h\Theta.}
By (7.) § 1, $\Theta=\Phi^e\Psi$, where $\Psi$ is relatively prime to $\Phi$. Therefore
\[
\sum_R x_{AR}\left(e\frac{\partial\log\Phi}{\partial x_R}+\frac{\partial\log\Psi}{\partial x_R}\right)=\varepsilon_A h
\]
or
\[
\sum_R x_{AR}\left(e\Psi\frac{\partial\Phi}{\partial x_R}+\Phi\frac{\partial\Psi}{\partial x_R}\right)=\varepsilon_A h\Phi\Psi.
\]
Since $\Psi$ is relatively prime to $\Phi$, it follows that $\sum_Rx_{AR}\frac{\partial\Phi}{\partial x_R}$ is divisible by $\Phi$; and since the two functions have the same degree, they can differ only by a constant factor. This is obtained by comparing the coefficients of $x_E^f$, and hence the formula holds:
\frobtagc{3.}{\sum_R x_{AR}\frac{\partial\Phi}{\partial x_R}=\sum_R x_{RA}\frac{\partial\Phi}{\partial x_R}=\chi(A^{-1})\Phi,}
which was already derived in another way in § 4. Multiplying the equation
\[
\sum_S x_{RS}\frac{\partial\Phi}{\partial x_S}=\chi(R^{-1})\Phi
\]
by $\Theta_{RA}$ and summing with respect to $R$, one obtains
\[
h\Theta\frac{\partial\Phi}{\partial x_A}=\Phi\sum_R\chi(R^{-1})\frac{\partial\Theta}{\partial x_{RA}}=\Phi\sum_R\chi(AR^{-1})\frac{\partial\Theta}{\partial x_R}.
\]
If one sets here $\Theta=\Phi^e\Psi$, then one obtains
\[
h\Psi\frac{\partial\Phi}{\partial x_A}=\sum_R\chi(AR^{-1})\left(e\Psi\frac{\partial\Phi}{\partial x_R}+\Phi\frac{\partial\Psi}{\partial x_R}\right),
\]
and consequently
\frobtagc{4.}{\frac{h}{e}\frac{\partial\Phi}{\partial x_A}=\sum_R\chi(AR^{-1})\frac{\partial\Phi}{\partial x_R},}
because the difference of these two functions is divisible by $\Phi$ and is only of degree $f-1$. If, however, one sets $\Theta=\Phi^{e'}\Psi'$, where $\Phi'$ is a prime factor of $\Theta$ different from $\Phi$, then one obtains
\[
h\Phi'\Psi'\frac{\partial\Phi}{\partial x_A}=\Phi\sum_R\chi(AR^{-1})\left(e'\Psi'\frac{\partial\Phi'}{\partial x_R}+\Phi'\frac{\partial\Psi'}{\partial x_R}\right),
\]
and consequently
\frobtagc{5.}{\sum_R\chi(AR^{-1})\frac{\partial\Phi'}{\partial x_R}=0,}
because this function is divisible by $\Phi'$. Comparing in these relations the coefficients of $x_E^{f-1}$, respectively $x_E^{f'-1}$, gives
\[
\sum_R\chi(AR^{-1})\chi(R)=\frac{h}{e}\chi(A),\qquad \sum_R\chi(AR^{-1})\psi(R)=0,
\]
where $\psi(R)$ is the character corresponding to the prime function $\Phi'$. One can also write these equations as
\frobtag{6.}{\sum\chi(R)\chi(S)=\frac{h}{e}\chi(A),\qquad \sum\chi(R)\psi(S)=0\qquad(RS=A)}
or
\frobtag{7.}{\sum_R\chi(PR^{-1})\chi(RQ^{-1})=\frac{h}{e}\chi(PQ^{-1}),\qquad \sum_R\chi(PR^{-1})\psi(RQ^{-1})=0.}
For $A=E$,
\frobtag{8.}{\sum_R\chi(R)\chi(R^{-1})=\frac{hf}{e},\qquad \sum_R\chi(R)\psi(R^{-1})=0.}
From this flows the conclusion that the values $\chi(E),\chi(A),\chi(B),\cdots$ of the character $\chi$ cannot be proportional to the corresponding values $\psi(E),\psi(A),\psi(B),\cdots$ of the character $\psi$. If the matrix $(\chi(PQ^{-1}))$ is denoted briefly by $(\chi)$, then the relations (7.) can also be put into the form
\frobtagc{9.}{\left(\frac{e}{h}\chi\right)^2=\left(\frac{e}{h}\chi\right),\qquad (\chi)(\psi)=0.}

\section*{§ 6.}

Simultaneously with $P$, $P^{-1}$ also runs through the $h$ elements of $\Hgrp$, only in a different order. Therefore
\[
|x_{P,Q}|=|x_{P^{-1},Q^{-1}}|=|x_{Q^{-1},P^{-1}}|,
\]
hence
\frobtagc{1.}{|x_{PQ^{-1}}|=|x_{Q^{-1}P}|.}
In each of these two determinants one obtains the rows by putting for $P$, and the columns by putting for $Q$, the elements $G_1,G_2,\cdots G_h$ of $\Hgrp$. Thus, if $x_R$ and $y_R$ are two systems of $h$ variables each, then
\[
|x_{PQ^{-1}}|=\Pi(\Phi(x)^e),\qquad |y_{Q^{-1}P}|=\Pi(\Phi(y)^e).
\]
The two matrices $(x_{PQ^{-1}})$ and $(y_{Q^{-1}P})$, however, are commutable with one another; one has
\[
\sum_R x_{PR^{-1}}y_{Q^{-1}R}=\sum_S y_{S^{-1}P}x_{SQ^{-1}}.
\]
For if one sets $SQ^{-1}=PR^{-1}$, thus $S=PR^{-1}Q$, then $S$ runs simultaneously with $R$ through the $h$ elements of $\Hgrp$, and also $S^{-1}P=Q^{-1}R$. Let $a_1,a_2,a_3,\cdots$ be the $h$ roots of the characteristic equation of the matrix $(x_{PQ^{-1}})$, thus the roots of the equation $|x_{PQ^{-1}}-u\varepsilon_{PQ^{-1}}|=0$, and let $b_1,b_2,b_3,\cdots$ be the $h$ roots of the equation $|y_{Q^{-1}P}-u\varepsilon_{PQ^{-1}}|=0$, hence also of the equation $|y_{PQ^{-1}}-u\varepsilon_{PQ^{-1}}|=0$. Then (V., p. 602, III) these two sequences of $h$ roots each can be assigned to one another so that $a_1+b_1,a_2+b_2,a_3+b_3,\cdots$ become the roots of the matrix $(x_{PQ^{-1}}+y_{Q^{-1}P})$.

I shall return to this general theorem later (§ 10). Here I now make the assumption that, for every two elements of $\Hgrp$,
\frobtagc{2.}{y_{BA}=y_{AB}}
holds. Thus, if one divides the $h$ elements of $\Hgrp$ into classes of conjugate elements, then $y_R$ has the same value for all elements of a class, for instance the value $y_R=y_\rho$ for the elements $R$ of the $\rho^{\mathrm{th}}$ class. If $k$ is the number of classes, let the $k$ variables $y_0,y_1,\cdots,y_{k-1}$ which correspond to the $k$ classes $(0),(1),\ldots,(k-1)$ be independent of one another. Since now $y_{Q^{-1}P}=y_{PQ^{-1}}$, the matrix $(y_{PQ^{-1}})$ is commutable with the matrix $(x_{PQ^{-1}})$. Every matrix $(x)$ which possesses the symmetry properties defined in § 1 is commutable with every other matrix $(y)$ whose elements in addition satisfy conditions (2.). If therefore $u,v,w$ are variables, then the determinant
\[
|ux_{PQ^{-1}}+vy_{PQ^{-1}}+w\varepsilon_{PQ^{-1}}|=\Pi\bigl(\Phi(ux+vy+w\varepsilon)^e\bigr)
\]
is a product of linear functions of $u,v,w$, and hence also
\frobtagc{3.}{\Phi(ux+vy+w\varepsilon)=(u_1u+v_1v+w)(u_2u+v_2v+w)\cdots(u_fu+v_fv+w),}
where $u_1,v_1,\cdots,u_f,v_f$ are independent of $u,v,w$. The coefficient of $w$ in each of these $f$ linear functions may be assumed equal to 1, because the left side becomes $w^f$ when $u=v=0$. Setting $v=0$ and $u=1$ gives
\frobtagc{4.}{\Phi(x+w\varepsilon)=(u_1+w)(u_2+w)\cdots(u_f+w),}
and setting $u=0$ and $v=1$ gives
\frobtagc{5.}{\Phi(y+w\varepsilon)=(v_1+w)(v_2+w)\cdots(v_f+w).}
Therefore $u_1,u_2,\cdots,u_f$ depend only on the $h$ variables $x_R$ and have the same meaning as in § 3, while $v_1,v_2,\cdots,v_f$ depend only on the $k$ variables $y_\rho$.

Since $\Phi(x)$ is irreducible, $\Phi(x+w\varepsilon)$ is also irreducible as a function of $w$, that is, this expression cannot be represented as the product of two integral functions of $w$ whose coefficients are rational functions of the $h$ independent variables $x_R$. If one regards the $k$ quantities $y_\rho$ as constant, then $v_1$ is also a constant, and consequently
\[
\Phi(x_E+v_1+w,x_A,x_B,\cdots)=(u_1+v_1+w)(u_2+v_1+w)\cdots(u_f+v_1+w)
\]
is also irreducible as a function of $w$. But this function has the linear factor $u_1+v_1+w$ in common with the function
\[
\Phi(x+y+w\varepsilon)=(u_1+v_1+w)(u_2+v_2+w)\cdots(u_f+v_f+w).
\]
Consequently both functions must be identical. Setting $w=0$ and $v_1=\eta$, therefore,
\frobtagc{6.}{\Phi(x_E+y_E,x_A+y_A,x_B+y_B,\cdots)=\Phi(x_E+\eta,x_A,x_B,\cdots),}
where $\eta$ depends only on the variables $y_R$. If one sets all variables $x_R$ equal to zero except $x_E=u$, one obtains
\frobtagc{7.}{\Phi(y_E+u,y_A,y_B,\cdots)=(u+\eta)^f,\qquad(y_{BA}=y_{AB})}
and from this, by comparison of the coefficients of $u^{f-1}$,
\frobtagc{8.}{f\eta=\sum_R\chi(R)y_R.}
Thus the theorem follows:

\emph{If, for every two elements $A$ and $B$ of the group $\Hgrp$, $x_{BA}=x_{AB}$, then the prime function of $f^{\mathrm{th}}$ degree $\Phi(x)$ becomes equal to the $f^{\mathrm{th}}$ power of a linear function $\xi=\frac{1}{f}\sum_R\chi(R)x_R$.}

For all elements $R$ of the $\rho^{\mathrm{th}}$ class, whose number is $h_R=h_\rho$, $\chi(R)$ has the same value $\chi_\rho$. If $x_{BA}=x_{AB}$, then $x_R$ also has the same value $x_\rho$ for these $h_\rho$ elements. Therefore $\sum\chi(R)x_R=\sum h_\rho\chi_\rho x_\rho$. If $\Phi'$ is a prime factor of $\Theta$ distinct from $\Phi$, and $\psi(R)$ is the character corresponding to it, then by § 5 the $k$ quantities $\psi_\rho$ cannot be proportional to the $k$ quantities $\chi_\rho$, and consequently the two linear functions $\sum h_\rho\psi_\rho x_\rho$ and $\sum h_\rho\chi_\rho x_\rho$ cannot differ merely by a constant factor. Applying the theorem above to every prime factor of $\Theta$, one obtains
\frobtagc{9.}{|x_{PQ^{-1}}|=\Pi\left(\frac{1}{f}\sum_R\chi(R)x_R\right)^f=\Pi(\xi^f).\qquad(x_{BA}=x_{AB})}
To each prime factor $\Phi$ of the determinant (7.) § 1, in which the $h$ variables $x_R$ are independent of one another, there corresponds a linear factor $\xi$ of the determinant (9.), in which $x_{BA}=x_{AB}$ is set. If the $k$ variables $x_\rho$ are independent, then to two different prime factors of that determinant there correspond two essentially different linear factors of determinant (9.). If $f$ is the degree of the prime function $\Phi$, and $e$ the exponent of the power of $\Phi$ occurring in $\Theta$, then $g=ef$ is the exponent of the power of the corresponding linear factor $\xi$ by which determinant (9.) is divisible.

Comparing the result obtained with formula (22.) \emph{Ch.} § 5, one sees that the quantities $\chi_\rho$, which there were introduced in a quite different way as the characters of the group $\Hgrp$, agree completely with the quantities $\chi(R)$ defined here; likewise the numbers $g=ef$ agree. Only there the two factors $e$ and $f$ of $g$ were left arbitrary, whereas here each separately has a definite meaning. One could therefore now make use of all the properties of the characters developed there. Nevertheless, I prefer to derive the results of that work once more from the starting point chosen here.

\section*{§ 7.}

If one differentiates equation (6.) § 6 with respect to $y_\rho$ and then sets all $k$ variables $y_\rho$ equal to zero, one finds
\[
\sum_{(\rho)}'\frac{\partial\Phi(x)}{\partial x_R}=\frac{h_\rho\chi_\rho}{f}\frac{\partial\Phi(x)}{\partial x_E},
\]
where the summation letter $R$ runs through the $h_\rho$ elements of the $\rho^{\mathrm{th}}$ class. If $A$ is a fixed element, and for every $R$ one replaces the variable $x_R$ by $x_{AR}$, then by (10.) § 3, $\Phi$ changes only by a constant factor different from zero. Consequently also
\frobtagc{1.}{\sum_{(\rho)}'\frac{\partial\Phi(x)}{\partial x_{AR}}=\frac{h_\rho\chi_\rho}{f}\frac{\partial\Phi(x)}{\partial x_A},}
where again $R$ runs through the $h_\rho$ elements of class $(\rho)$. Comparing the coefficients of $x_E^{f-1}$ gives
\frobtagc{2.}{\sum'\chi(AS)=\frac{h_B}{f}\chi(A)\chi(B),}
where $S$ runs through the $h_B$ elements conjugate to $B$. One obtains this relation directly by computing, for the function (7.) § 6, the expression $S_2(y)=\sum\chi(PQ)y_Py_Q=f\eta^2$. If one sets $S=R^{-1}BR$ and lets $R$ run through all $h$ elements of $\Hgrp$, then $S$ becomes equal to each element of the class of $B$, and indeed each one $\frac{h}{h_B}$ times. Therefore (\emph{Ch.} § 5, (5.))
\frobtagc{3.}{h\chi(A)\chi(B)=f\sum_R\chi(AR^{-1}BR).}
One reaches the same result directly from formula (9.) § 6 by the path indicated in \emph{Ch.} p. 1001, thus by means of the same conclusions that led in § 1 to formula (9.).

Let the number of different prime factors of the group determinant $\Theta$ be $l$. These $l$ functions $\Phi$ and the characters $\chi$ corresponding to them may be distinguished from one another by upper indices $\lambda=0,1,\cdots,l-1$. Now in the expansion of the determinant $\Theta(x+u\varepsilon)$ by powers of $u$, the coefficient of $u^{h-1}$ is equal to $hx_E$. If, therefore, in the product
\frobtagc{4.}{\Theta=\Pi_\lambda\bigl(\Phi^{(\lambda)e^{(\lambda)}}\bigr)}
one replaces $x_E$ by $x_E+u$ and then compares the coefficients of $x_E^{h-1}$, one obtains
\frobtagc{5.}{\sum_\lambda e^{(\lambda)}\chi^{(\lambda)}(R)=h\varepsilon_R.}
Now by (8.) § 5 and (3.)
\[
\sum_R\chi(S^{-1}R^{-1}SR)=\frac{h}{f}\chi(S)\chi(S^{-1}),\qquad \sum_S\chi(S)\chi(S^{-1})=\frac{hf}{e},
\]
and hence, for every value of $\lambda$,
\[
\sum_{R,S}\frac{e^{(\lambda)}}{h}\chi^{(\lambda)}(S^{-1}R^{-1}SR)=h,
\]
and consequently, since this value is independent of $\lambda$,
\[
\sum_\lambda\sum_{R,S}\frac{e^{(\lambda)}}{h}\chi^{(\lambda)}(S^{-1}R^{-1}SR)=hl.
\]
I now reverse the order of the two summations. Then by (5.)
\[
\sum_\lambda\frac{e^{(\lambda)}}{h}\chi^{(\lambda)}(S^{-1}R^{-1}SR)=0,
\]
except when $S^{-1}R^{-1}SR=E$, that is, when $SR=RS$; in that case the sum is equal to 1. Therefore $hl$ is equal to the number of solutions of the equation $SR=RS$. But this, as I showed in \emph{Ch.} p. 987 by the simplest considerations, is equal to $hk$, and consequently $l=k$.

\emph{The number of different prime factors of the group determinant is equal to the number of classes of conjugate elements into which the elements of the group are divided.}

As $R$ runs through the class $(\alpha)$, $R^{-1}$ runs through the inverse class, which I denote by $(\alpha')$. Therefore equations (8.) § 5 are identical with
\frobtag{6.}{\sum_\alpha h_\alpha\chi_\alpha^{(\kappa)}\chi_{\alpha'}^{(\kappa)}=\frac{h f^{(\kappa)}}{e^{(\kappa)}},\qquad \sum_\alpha h_\alpha\chi_\alpha^{(\kappa)}\chi_{\alpha'}^{(\lambda)}=0.}
Accordingly the two matrices of $k^{\mathrm{th}}$ degree
\frobtagc{7.}{\left(\frac{h_\alpha}{h}\chi_{\alpha'}^{(\kappa)}\right),\qquad \left(\frac{e^{(\kappa)}}{f^{(\kappa)}}\chi_\alpha^{(\kappa)}\right),}
are complementary, and consequently the equations
\[
\sum_\kappa\frac{e^{(\kappa)}}{f^{(\kappa)}}\chi_\alpha^{(\kappa)}\chi_{\alpha'}^{(\kappa)}=\frac{h}{h_\alpha},\qquad \sum_\kappa\frac{e^{(\kappa)}}{f^{(\kappa)}}\chi_\alpha^{(\kappa)}\chi_\beta^{(\kappa)}=0
\]
also hold. If therefore $h_{\alpha\beta}$ is set equal to $h_\alpha$ or to 0 according as $(\beta)=(\alpha')$ or not, then
\frobtagc{8.}{\sum_\kappa\frac{e^{(\kappa)}}{f^{(\kappa)}}\chi_\alpha^{(\kappa)}\chi_\beta^{(\kappa)}=\frac{h\,h_{\alpha\beta}}{h_\alpha h_\beta}.}
This equation is obtained immediately from relation (2.). Namely, if $\chi(A)=\chi_\alpha$ and $\chi(B)=\chi_\beta$, then it reads
\[
\sum_S'\chi^{(\kappa)}(AS)=\frac{h_\beta}{f^{(\kappa)}}\chi_\alpha^{(\kappa)}\chi_\beta^{(\kappa)},
\]
where $S$ runs through the $h_\beta$ elements of the $\beta^{\mathrm{th}}$ class. Therefore
\[
\sum_\kappa\sum_S'\frac{e^{(\kappa)}}{h}\chi^{(\kappa)}(AS)=\frac{h_\beta}{h}\sum_\kappa\frac{e^{(\kappa)}}{f^{(\kappa)}}\chi_\alpha^{(\kappa)}\chi_\beta^{(\kappa)}.
\]
But now $\sum_\kappa\frac{e^{(\kappa)}}{h}\chi^{(\kappa)}(AS)=0$ except when $AS=E$, hence $S=A^{-1}$. Then the sum is equal to 1. The latter occurs always and only when $(\beta)=(\alpha')$, and so equation (8.) results. Both from this and from (6.) it follows that the determinant of $k^{\mathrm{th}}$ degree
\frobtagc{9.}{|\chi_\alpha^{(\kappa)}|}
is different from zero.

More generally, let $A$ be a fixed element of the $\alpha^{\mathrm{th}}$ class, and let $R$ be a variable element running through the $h_\alpha$ elements conjugate with $A$. The signs $B$ and $S$ shall have the same meaning for the $\beta^{\mathrm{th}}$ class, and the signs $C$ and $T$ for the $\gamma^{\mathrm{th}}$ class. Then by (2.)
\[
\sum_R'\chi(RSC)=\frac{h_\alpha}{f}\chi(A)\chi(SC),\qquad \sum_S'\chi(SC)=\frac{h_\beta}{f}\chi(B)\chi(C),
\]
therefore
\[
\sum_{R,S}'\chi(RSC)=\frac{h_\alpha h_\beta}{f^2}\chi_\alpha\chi_\beta\chi_\gamma
\]
and hence
\[
\sum_\kappa\sum_{R,S}'\frac{e^{(\kappa)}}{h}\chi^{(\kappa)}(RSC)=\frac{h_\alpha h_\beta}{h}\sum_\kappa\frac{e^{(\kappa)}}{f^{(\kappa)2}}\chi_\alpha^{(\kappa)}\chi_\beta^{(\kappa)}\chi_\gamma^{(\kappa)}.
\]
The left side is equal to the number of solutions of the equation $RSC=E$. The right side shows that this number does not depend on the element $C$, but only on the class $(\gamma)$ to which $C$ belongs. If it is denoted by $\frac{h_{\alpha\beta\gamma}}{h_\gamma}$, then accordingly
\frobtagc{10.}{\frac{h\,h_{\alpha\beta\gamma}}{h_\alpha h_\beta h_\gamma}=\sum_\kappa\frac{e^{(\kappa)}}{f^{(\kappa)2}}\chi_\alpha^{(\kappa)}\chi_\beta^{(\kappa)}\chi_\gamma^{(\kappa)}.}
If for $C$ one puts, one after another, the $h_\gamma$ elements $T$ of class $(\gamma)$, it follows that $h_{\alpha\beta\gamma}$ is the number of solutions of the equation $RST=E$, when $R$ runs through the $h_\alpha$ elements of class $(\alpha)$, $S$ through the $h_\beta$ elements of class $(\beta)$, and $T$ through those of class $(\gamma)$. As the right side of this equation shows, the number $h_{\alpha\beta\gamma}$ remains unchanged under every interchange of its three indices with one another. Therefore $\frac{h_{\alpha\beta\gamma}}{h_\beta}$ is also the number of solutions of the equation $RBT=E$, in which in place of $S$ a fixed element $B$ of the $\beta^{\mathrm{th}}$ class has entered.

From (6.) and (10.) the relations follow:
\[
h_\beta h_\gamma\chi_\beta^{(\kappa)}\chi_\gamma^{(\kappa)}=f^{(\kappa)}\sum_\alpha h_{\alpha\beta\gamma}\chi_\alpha^{(\kappa)}.
\]
These show that the quantities $\chi_\alpha$ satisfy the equations
\frobtagc{11.}{h_\beta h_\gamma\chi_\beta\chi_\gamma=f\sum_\alpha h_{\alpha'\beta\gamma}\chi_\alpha}
Thus the latter have $k$ systems of solutions
\[
\chi_\alpha=\chi_\alpha^{(\kappa)}\qquad(\kappa=0,1,\cdots k-1)
\]
but no more. For if $x_0,x_1,\cdots,x_{k-1}$ are variables and one sets
\frobtagc{12.}{\sum_\alpha h_\alpha\chi_\alpha^{(\kappa)}x_\alpha=f^{(\kappa)}\xi^{(\kappa)},}
then from (11.) it follows that
\[
h_\alpha\chi_\alpha^{(\kappa)}\xi^{(\kappa)}=\sum_\beta\left(\sum_\gamma h_{\alpha\beta\gamma}x_\gamma\right)\chi_\beta^{(\kappa)}
\]
or
\[
\sum_\beta\left(\left(\sum_\gamma h_{\alpha\beta\gamma}x_\gamma\right)-h_{\alpha\beta}\xi^{(\kappa)}\right)\chi_\beta^{(\kappa)}=0.
\]
Consequently the determinant
\frobtagc{13.}{\left|\left(\sum_\gamma h_{\alpha\beta\gamma}x_\gamma\right)-h_{\alpha\beta}r\right|,\qquad(\alpha,\beta=0,1,\cdots k-1)}
which is an integral function of $r$ of $k^{\mathrm{th}}$ degree, vanishes for the $k$ values $r=\xi^{(\kappa)}$, which are mutually distinct (cf. \textsc{Dedekind}, \emph{Zur Theorie der aus $n$ Haupteinheiten gebildeten complexen Grössen}. Göttinger Nachrichten 1885, p. 146).

Equations (11.) determine the quantities $\frac{\chi_\alpha}{f}$. Equation (6.) then supplies
\[
\frac{h}{ef}=\sum_\alpha h_\alpha\frac{\chi_\alpha}{f}\frac{\chi_{\alpha'}}{f}
\]
for each of the $k$ systems of values $\frac{\chi_\alpha}{f}$ the corresponding value of $ef=g$.

By § 3, therefore, for the complete calculation of the $k$ prime functions $\Phi$ nothing further is required than to determine separately the positive integers $e$ and $f$, whose product $g$ is already known. This problem, the most difficult among all questions concerning the group determinant, is solved in § 9 by the theorem that always $e=f=\sqrt{g}$. If one adds to this the result from § 12 that $\chi_\alpha$ and $\chi_{\alpha'}$ are conjugate complex quantities, then it follows from (7.) that the matrix
\frobtagc{14.}{\left(\sqrt{\frac{h_\alpha}{h}}\,\chi_\alpha^{(\kappa)}\right)}
is complementary to the conjugate complex matrix.

\section*{§ 8.}

In equation (7.) § 6 I assign to the variables $y_R$ the values $\chi(R^{-1})=\chi'(R)$, which satisfy the condition $y_{BA}=y_{AB}$ by (13.) § 3. Then by (8.) § 5, $\eta=\frac{h}{e}$, and consequently
\frobtagc{1.}{\Phi\left(\frac{e}{h}\chi'\right)=1.}

By virtue of property (13.) § 3, the matrix $(\chi')$ is interchangeable with every matrix $(x)$. I set
\frobtag{2.}{\frac{h}{e}\xi_A=\sum\chi(R^{-1})x_S=\sum\chi(S^{-1})x_R,\qquad (RS=A)}
and form from the quantities $\xi_{PQ^{-1}}$ the matrix $(\xi)$. Then
\[
((x)+u(\varepsilon))\left(\frac{e}{h}\chi'\right)=(\xi)+u\left(\frac{e}{h}\chi'\right)
\]
and therefore, by (9.) § 1,
\[
\Phi(x+u\varepsilon)=\Phi\left(\xi+u\frac{e}{h}\chi'\right).
\]
But by equation (6.) § 6,
\[
\Phi\left(\xi+u\frac{e}{h}\chi'\right)=\Phi(\xi+u\varepsilon),
\]
because the value of $\eta$ corresponding to the quantities $y_R=\frac{e}{h}\chi'(R)$ is equal to 1 by (8.) § 5. Consequently
\frobtagc{3.}{\Phi(x+u\varepsilon)=\Phi(\xi+u\varepsilon).}
If, on the other hand, $\Phi'$ is a prime factor of the group determinant different from $\Phi$, if $f'$ is its degree and $\Psi$ the corresponding character, then by (7.) § 6
\[
\Phi'(y_R+u\varepsilon_R)=(u+\eta')^{f'},\qquad (y_{BA}=y_{AB})
\]
where $f'\eta'=\sum\Psi(R)y_R$. If one therefore sets $y_R=\frac{e}{h}\chi'(R)$, then by (8.) § 5 $\eta'=0$. Accordingly, by (6.) § 6, for $h$ independent variables $z_R$,
\[
\Phi'\left(uz+\frac{e}{h}\chi'\right)=\Phi'(uz)=u^{f'}\Phi'(z).
\]
Let $(z)=(x)^{-1}$, hence $(z)(x)=(\varepsilon)$ and $\Phi(z)\Phi(x)=1$. Then
\[
\Phi'(x)\,\Phi'\left(uz+\frac{e}{h}\chi'\right)=u^{f'},
\]
and therefore, since $(x)(u(z)+(\frac{e}{h}\chi'))=u(\varepsilon)+(\xi)$,
\frobtagc{4.}{\Phi'(\xi+u\varepsilon)=u^{f'}.}
From these two equations, which are only special cases of a more general formula, it follows by (7.) § 1 that
\frobtagc{5.}{\left|\xi_{PQ^{-1}}+u\varepsilon_{PQ^{-1}}\right|=\Phi(x+u\varepsilon)^e\,u^{h-ef}}
and, for $x_R=\varepsilon_R$,
\frobtagc{6.}{\left|\frac{e}{h}\chi(PQ^{-1})+u\varepsilon_{PQ^{-1}}\right|=(u+1)^f\,u^{h-ef}.}

Therefore the characteristic determinant of the matrix $(\xi)$ vanishes for the $f+1$ values $u=0,u_1,u_2,\cdots u_f$. For since the group determinant $\Theta(x)$ always contains the linear factor $\sum x_R$, and indeed only in the first power, it is necessary that $h-ef>0$, if the trivial case $h=1$ is excluded.

Now let $G((\xi))=0$ be the equation of lowest degree satisfied by the matrix $(\xi)$ (V. § 1, VI). Then the integral function $G(u)$, which is a divisor of the characteristic determinant $\left|\xi_{PQ^{-1}}-u\varepsilon_{PQ^{-1}}\right|$, must vanish for every root of the matrix $(\xi)$, hence for each of the $f+1$ values $u=0,u_1,u_2,\cdots u_f$, and consequently must be divisible by $u\Phi(x-u\varepsilon)$. Since, furthermore, the matrix $(\chi')$ is interchangeable with every matrix $(x)$, by (9.) § 5 one has
\frobtag{7.}{(\xi)^n=\left(\left(\frac{e}{h}\chi'\right)(x)\right)^n=\left(\frac{e}{h}\chi'\right)^n(x)^n=\left(\frac{e}{h}\chi'\right)(x)^n.}
Multiplying relation (5.) § 4 by $(x)$, one therefore obtains
\frobtag{8.}{(\xi)^{f+1}-(\xi)^f\Phi_1+(\xi)^{f-1}\Phi_2-\cdots+(-1)^f(\xi)\Phi_f=0,}
where by (3.)
\frobtagc{9.}{\Phi_n=\Phi_n(x)=\Phi_n(\xi)}
is true, or, more briefly,
\[
(\xi)\Phi(x-(\xi)\varepsilon)=0.
\]
Thus $u\Phi(x-u\varepsilon)$ is divisible by $G(u)$, hence equal to $G(u)$; consequently (8.) is the equation of lowest degree satisfied by the matrix $(\xi)$.

Without using relation (5.) § 4, one can also see this theorem as follows. The function $G(u)$ is obtained by dividing the determinant of degree $h$, $(\xi_{PQ^{-1}}-u\varepsilon_{PQ^{-1}})$, by the greatest common divisor of its minors of degree $h-1$. If the $h$ variables $x_R$ are independent, it follows from equation (1.) § 5 that all minors of degree $h-1$ of $\Theta(x)$ are divisible by $\Phi(x)^{e-1}$. Therefore all minors of $\Theta(\xi-u\varepsilon)$ are divisible by $\Phi(\xi-u\varepsilon)^{e-1}=\Phi(x-u\varepsilon)^{e-1}$, and hence also by $(u-u_1)^{e-1}$, but not by a higher power of $u-u_1$, because otherwise $\Theta(\xi-u\varepsilon)$ would have to be divisible by a higher than the $e^{\mathrm{th}}$ power of $u-u_1$. By means of the same theorems it follows from equations (9.) § 5 and (6.) that the rank of the matrix $(\chi)$ is equal to $ef$, as I have shown more fully in Ch. § 5. Therefore the rank of the matrix
\frobtagc{10.}{(\xi)=\left(\frac{e}{h}\chi'\right)(x)=(x)\left(\frac{e}{h}\chi'\right)}
is also equal to $ef$. Consequently $h-ef$ elementary divisors of the characteristic determinant of $(\xi)$ vanish for $u=0$; and because their product is equal to $u^{h-ef}$ by (6.), each of them must be linear. Thus the greatest common divisor of the minors of degree $h-1$ of that determinant contains the factor $u$ exactly in the $(h-ef-1)^{\mathrm{th}}$ power, and accordingly
\[
(-1)^fG(u)=u\Phi(x-u\varepsilon)
\]
must hold.

\section*{§ 9.}

After these preparations I now turn to the proof of the fundamental theorem of the theory of group determinants:

\emph{The exponent of the power in which the group determinant contains a prime factor is equal to the degree of the factor.}

I have already disposed of the case $f=1$ in § 2. Because of the difficulty of the general proof, I first put forward the special cases $f=2$ and 3.

If $f=2$, then
\frobtagc{1.}{2\Phi(x)=S_1^2-S_2,}
and the corresponding character $\chi$ satisfies the relations $\chi(A,B,C)=0$, or
\frobtagc{2.}{\begin{aligned}
&\chi(A)\chi(B)\chi(C)-\chi(A)\chi(BC)-\chi(B)\chi(AC)\nobreak\\
&\qquad{}-\chi(C)\chi(AB)+\chi(ABC)+\chi(ACB)=0.
\end{aligned}}
In this equation I replace $B$ by $BC^{-1}$ and then sum over $C$ through the $h$ elements of $\Hgrp$ (or one sums in (2.) over all elements $B,C$ satisfying the condition $BC=B'$, where $B'$ is a fixed element). With the aid of formulas (7.) § 5 and (3.) § 7 one finds
\[
\frac{h}{e}\chi(A)\chi(B)-h\chi(A)\chi(B)-\frac{h}{e}\chi(AB)-\frac{h}{e}\chi(AB)+h\chi(AB)+\frac{h}{f}\chi(A)\chi(B)=0,
\]
and hence, since $f=2$,
\[
\left(\frac{h}{e}-\frac{h}{f}\right)(\chi(A)\chi(B)-2\chi(AB))=0.
\]
Therefore $e=f$, since it cannot be true for every two elements that
\[
\chi(A)\chi(B)-2\chi(AB)=0.
\]
For otherwise, multiplying this equation by $x_Ax_B$ and summing over $A$ and $B$, one would obtain $S_1^2-2S_2=0$. By (1.), therefore, $4\Phi=S_1^2$, whereas $\Phi$ is irreducible.

If $f=3$, then
\frobtagc{3.}{6\Phi=S_1^3-3S_1S_2+2S_3,}
and the corresponding character satisfies the relations $\chi(A,B,C,D)=0$ ((20.) § 3). If one replaces $C$ there by $CD^{-1}$ and then sums over $D$, one obtains
\[
\begin{aligned}
&\frac{h}{e}\bigl(\chi(A)\chi(B)\chi(C)-\chi(B)\chi(AC)-\chi(A)\chi(BC)-\chi(A)\chi(BC)\nobreak\\
&\qquad{}-\chi(B)\chi(AC)-\chi(C)\chi(AB)+\chi(ABC)+\chi(ACB)\nobreak\\
&\qquad{}+\chi(ABC)+\chi(ABC)+\chi(ACB)+\chi(ACB)\bigr)\nobreak\\
&{}+h\bigl(-\chi(A)\chi(B)\chi(C)+\chi(C)\chi(AB)+\chi(A)\chi(BC)+\chi(B)\chi(AC)\nobreak\\
&\qquad{}-\chi(ABC)-\chi(ACB)\bigr)\nobreak\\
&{}+\frac{h}{f}\bigl(\chi(A)\chi(B)\chi(C)+\chi(A)\chi(B)\chi(C)-\chi(B)\chi(AC)-\chi(A)\chi(BC)\nobreak\\
&\qquad{}-\chi(C)\chi(AB)-\chi(C)\chi(AB)\bigr)=0,
\end{aligned}
\]
therefore, if one replaces the factor $h$ by $3\frac{h}{f}$,
\frobtag{4.}{\begin{aligned}
&\left(\frac{h}{e}-\frac{h}{f}\right)\bigl(\chi(A)\chi(B)\chi(C)-2\chi(A)\chi(BC)-2\chi(B)\chi(AC)\nobreak\\
&\qquad{}-\chi(C)\chi(AB)+3\chi(ABC)+3\chi(ACB)\bigr)=0.
\end{aligned}}
If now the second factor were always zero, then, multiplying by $x_Ax_Bx_C$ and summing, one would obtain $S_1^3-5S_1S_2+6S_3=0$; and if, by means of this equation, one eliminates $S_3$ from (3.), one obtains $9\Phi=S_1(S_1^2-2S_2)$, whereas $\Phi$ is irreducible. Therefore $e=f=3$.

In the general case the character $\chi$ satisfies the relations
\[
\chi(A,B,C,\cdots Q,R,S)=0,
\]
where $A,B,\cdots R,S$ are any $f+1$ elements, or, briefly,
\[
\sum^{(f+1)}(\Pi-\chi)=0.
\]
In each of the $(f+1)!$ terms of this sum I replace $R$ by $RS^{-1}$ and then sum also over $S$. Each term corresponds to a certain permutation of $f+1$ symbols, decomposed into cyclic factors. With respect to this permutation I distinguish three cases:

1. $R$ and $S$ occur in two different cycles of the permutation, for example
\[
(-\chi(ABCD\cdots FR))(-\chi(S))(-\chi(G\cdots K))\cdots .
\]
Replacing $R$ by $RS^{-1}$ and then summing over $S$, one obtains by (7.) § 5
\[
-\frac{h}{e}(-\chi(ABCD\cdots FR))(-\chi(G\cdots K))\cdots .
\]
The same result is obtained in the same way from the term
\[
(-\chi(BCD\cdots FR))(-\chi(SA))(-\chi(G\cdots K))\cdots
\]
and from
\[
(-\chi(CD\cdots FR))(-\chi(SAB))(-\chi(G\cdots K))\cdots
\]
and so on, and finally from
\[
(-\chi(R))(-\chi(SABCD\cdots F))(-\chi(G\cdots K))\cdots,
\]
but from no other term. Thus if $r$ is the number of the elements $ABCD\cdots FR$, one obtains in this way
\frobtagc{5.}{-\frac{h}{e}\sum^{(f)}r\,(\Pi-\chi).}
The $f!$ terms of this sum are formed in a way analogous to those of the sum (19.) § 3 from the $f!$ permutations of the $f$ elements $A,B,C,\cdots Q,R$. Only in this sum, in which the element $R$ is distinguished, each term $(\Pi-\chi)$ receives in addition a numerical factor $r$. This is equal to the number of elements of the cycle in which $R$ occurs (cf. (4.)).

2. $R$ and $S$ both occur in the same cycle of the permutation, and $S$ immediately follows $R$ in the cycle (so $S$ may also be the first and $R$ the last element of the cycle), for example
\[
(-\chi(AB\cdots FRS))(-\chi(G\cdots K))\cdots .
\]
Replacing $R$ by $RS^{-1}$ and then summing over $S$, one obtains
\[
h(-\chi(AB\cdots FR))(-\chi(G\cdots K))\cdots,
\]
and each term only once, hence altogether
\frobtagc{6.}{h\sum^{(f)}(\Pi-\chi).}
3. $R$ and $S$ both occur in the same cycle, without $S$ immediately following $R$, for example
\[
(-\chi(A\cdots RBCD\cdots FS))(-\chi(G\cdots K))((-\chi(L\cdots N)))\cdots .
\]
Replacing $R$ by $RS^{-1}$ and then summing over $S$, one obtains by (3.) § 7
\[
-\frac{h}{f}(-\chi(A\cdots R))(-\chi(BCD\cdots F))(-\chi(G\cdots K))(-\chi(L\cdots N))\cdots .
\]
The same result is obtained in the same way from the term
\[
(-\chi(A\cdots RCD\cdots FBS))(-\chi(G\cdots K))(-\chi(L\cdots N))\cdots,
\]
which arises from the preceding one by cyclic interchange of the elements $BCD\cdots F$ standing between $R$ and $S$. The number of cyclic interchanges that can be so performed is equal to the number of elements $BCD\cdots F$. The same result is further obtained from the term
\[
(-\chi(A\cdots RG\cdots KS))(-\chi(BCD\cdots F))(-\chi(L\cdots N))\cdots .
\]
Here too one can cyclically interchange the elements $G\cdots K$ standing between $R$ and $S$, which is possible in as many ways as the number of elements $G\cdots K$ amounts to.

The same result is obtained from the term
\[
(-\chi(A\cdots RL\cdots NS))(-\chi(BCD\cdots F))(-\chi(G\cdots K))\cdots
\]
and so on, altogether therefore in as many ways as the number of the elements $BCD\cdots\allowbreak FG\cdots\allowbreak KL\cdots\allowbreak N\cdots$ amounts to, and in no more ways; hence in $f-r$ ways, if, as above, the number of elements $A\cdots R$ is denoted by $r$. Accordingly one obtains
\frobtagc{7.}{-\frac{h}{f}\sum^{(f)}(f-r)(\Pi-\chi).}
Combining the three expressions (5.), (6.), and (7.), one obtains the equation
\frobtagc{8.}{\left(\frac{h}{f}-\frac{h}{e}\right)\left(\sum^{(f)}r\,(\Pi-\chi)\right)=0.}
Therefore $e=f$ if one can show that, for not every $f$ elements $AB\cdots QR$, one has
\[
\sum^{(f)}r\,(\Pi-\chi)=0.
\]
Multiplying this equation by $x_Ax_B\cdots x_Qx_R$ and summing over each of the $f$ elements $A,B,\cdots Q,R$, one obtains a relation between $S_1,S_2,\cdots S_f$. This relation is not satisfied identically (without regard to the meaning of $S_1,S_2,\cdots S_f$), since with respect to $S_f$ it is linear and the coefficient of $S_f$ is a non-vanishing integer. I shall show, however, that $S_1,S_2,\cdots S_f$ are $f$ independent functions of the $h$ variables $x_R$, so that there is no relation among them whose coefficients are independent of the variables $x_R$. It then follows that $\Phi_1,\Phi_2,\cdots\Phi_f$ are also independent, and likewise $u_1,u_2,\cdots u_f$.

If there were an equation among the $f$ functions $S_1,S_2,\cdots S_f$ of the $h$ independent variables $x_R$, then, differentiating it with respect to $x_R$, one would obtain a relation of the form
\[
\Psi_1\frac{\partial S_1}{\partial x_R}+\Psi_2\frac{\partial S_2}{\partial x_R}+\cdots+\Psi_f\frac{\partial S_f}{\partial x_R}=0,
\]
where $\Psi_1,\Psi_2,\cdots\Psi_f$ are integral functions of the $h$ variables, independent of $R$. But now
\[
S_n=\sum_{R_1,R_2,R_3,\cdots R_n}\chi(R_1R_2R_3\cdots R_n)x_{R_1}x_{R_2}x_{R_3}\cdots x_{R_n}
\]
and consequently
\[
\begin{aligned}
\frac{\partial S_n}{\partial x_R}={}&\sum_{R_2,R_3,\cdots R_n}\chi(RR_2R_3\cdots R_n)x_{R_2}x_{R_3}\cdots x_{R_n}\nobreak\\
&{}+\sum_{R_1,R_3,\cdots R_n}\chi(R_1RR_3\cdots R_n)x_{R_1}x_{R_3}\cdots x_{R_n}+\cdots\nobreak\\
&{}+\sum_{R_1,R_2,\cdots R_{n-1}}\chi(R_1R_2\cdots R_{n-1}R)x_{R_1}x_{R_2}\cdots x_{R_{n-1}},
\end{aligned}
\]
and since $\chi(ABC\cdots F)$ remains unchanged under cyclic interchange of the elements,
\frobtagc{9.}{\frac{\partial S_n}{\partial x_R}=n\sum_{R_1,R_2,\cdots R_{n-1}}\chi(RR_1R_2\cdots R_{n-1})x_{R_1}x_{R_2}\cdots x_{R_{n-1}}}
or, by (6.) § 1 and (7.) § 8,
\frobtagc{10.}{\frac{\partial S_n}{\partial x_R}=n\sum_S\chi(RS)x_S^{(n-1)}=\frac{nh}{e}\xi_{R^{-1}}^{(n-1)},}
where $\xi_R^{(n)}$ is formed from the quantities $\xi_R$ in the same way as $x_R^{(n)}$ from the quantities $x_R$. In particular,
\frobtagc{11.}{\frac{\partial S_n}{\partial x_E}=nS_{n-1}.}
Substituting these expressions into the above relation, one obtains, if one also replaces $R$ by $R^{-1}$,
\[
\Psi_1\frac{e}{h}\chi(R^{-1})+2\Psi_2\xi_R+3\Psi_3\xi_R^{(2)}+\cdots+f\Psi_f\xi_R^{(f-1)}=0.
\]
Setting $R=PQ^{-1}$, this becomes an equation between matrices, which, multiplied by $(x)$, is
\[
(\xi)\Psi_1+2(\xi)^2\Psi_2+3(\xi)^3\Psi_3+\cdots+f(\xi)^f\Psi_f=0.
\]
But I have shown in § 8 that the equation of lowest degree satisfied by the matrix $(\xi)$ is of degree $f+1$. The second factor of expression (8.) therefore cannot vanish for every system of $f$ elements, and hence
\frobtagc{12.}{e=f}
must hold.

\section*{§ 10.}

If $x_R$ and $y_R$ are two systems each of $h$ variables, then by § 6 the matrix $(x_{PQ^{-1}})$ is interchangeable with the matrix $(y_{Q^{-1}P})$, and consequently the determinant
\frobtagc{1.}{\left|ux_{PQ^{-1}}+vy_{Q^{-1}P}+w\varepsilon_{PQ^{-1}}\right|}
is a product of $h$ linear functions of the three variables $u,v,w$ of the form $ua_\alpha+vb_\alpha+w$. Here $a_1,a_2,a_3\cdots$ are the $h$ roots of the matrix $(x_{PQ^{-1}})$, and $b_1,b_2,b_3,\cdots$ those of the matrix $(y_{Q^{-1}P})$ or, what is the same by § 6, of the matrix $(y_{PQ^{-1}})$. The question now is in what way the roots of these two matrices must be assigned to one another so that $ua_\alpha+vb_\alpha+w$ is a linear factor of determinant (1.). I set
\[
\sum_R x_R\chi(RS^{-1})=\frac{h}{e}\xi_S,
\qquad
\sum_R y_R\chi(RS^{-1})=\frac{h}{e}\eta_S,
\]
and further
\[
(x_{PQ^{-1}})=(x),\qquad (\xi_{PQ^{-1}})=(\xi),\qquad (y_{Q^{-1}P})=(\overline{y}),\qquad (\eta_{Q^{-1}P})=(\overline{\eta}),
\]
where throughout $P$ denotes the rows and $Q$ the columns of the matrix.

\clearpage
Then
\[
\left(\frac{e}{h}\chi'\right)(x)=(x)\left(\frac{e}{h}\chi'\right)=(\xi),
\qquad
\left(\frac{e}{h}\chi'\right)(\overline{y})=(\overline{y})\left(\frac{e}{h}\chi'\right)=(\overline{\eta}).
\]
I distinguish the $k$ different characters $\chi$ by upper indices $\chi^{(\varkappa)}$ $(\varkappa=0,1,\cdots k-1)$. I denote the matrices $(\xi)$ and $(\overline{\eta})$ corresponding to the character $\chi^{(\varkappa)}$ by $(\xi^{(\varkappa)})$ and $(\overline{\eta}^{(\varkappa)})$. If $\varkappa$ and $\lambda$ are different, then from (9.) § 5 it follows that
\frobtagc{2.}{(\xi^{(\varkappa)})(\xi^{(\lambda)})=0,\quad (\overline{\eta}^{(\varkappa)})(\overline{\eta}^{(\lambda)})=0,\quad (\xi^{(\varkappa)})(\overline{\eta}^{(\lambda)})=0,\quad (\overline{\eta}^{(\varkappa)})(\xi^{(\lambda)})=0.}
By equation (5.) § 7,
\[
\sum_\varkappa\left(\frac{e^{(\varkappa)}}{h}\chi'\right)=(\varepsilon)
\]
and hence
\[
\sum_\varkappa(\xi^{(\varkappa)})=(x),\qquad \sum_\varkappa(\overline{\eta}^{(\varkappa)})=(\overline{y}).
\]
If one therefore expands the product of the $k$ matrices
\[
\Pi_\varkappa\left(u(\xi^{(\varkappa)})+v(\overline{\eta}^{(\varkappa)})+w(\varepsilon)\right)
\]
according to powers of $w$, one obtains
\[
w^k(\varepsilon)+w^{k-1}\bigl(u(x)+v(\overline{y})\bigr),
\]
while the remaining terms vanish by (2.). Between the determinants of these matrices there consequently results the relation
\frobtag{3.}{\Pi_\varkappa\left|u\xi_{PQ^{-1}}^{(\varkappa)}+v\eta_{Q^{-1}P}^{(\varkappa)}+w\varepsilon_{PQ^{-1}}\right|=w^{h(k-1)}\left|ux_{PQ^{-1}}+vy_{Q^{-1}P}+w\varepsilon_{PQ^{-1}}\right|}
(compare the analogous development V. p. 610). Let any one of the $k$ factors on the left-hand side be
\frobtagc{4.}{\left|u\xi_{PQ^{-1}}+v\eta_{Q^{-1}P}+w\varepsilon_{PQ^{-1}}\right|.}
As the right-hand side shows, this determinant is equal to a power of $w$, multiplied by a number of the linear factors $ua_\alpha+vb_\alpha+w$ of the determinant (1.). On the other hand, one can regard the determinant (4.) as a special case of the determinant (1.): by (5.) § 8, the roots of the matrix $(\xi)$ are the $f$ roots $u_1,u_2,\cdots u_f$ of the equation $\Phi(x-u\varepsilon)=0$, each counted $e$ times, and in addition the number $0$ counted $(h-ef)$ times. Likewise the roots of the matrix $(\overline{\eta})$ are the number $0$ and the roots $v_1,v_2,\cdots v_f$ of the equation $\Phi(y-v\varepsilon)=0$. Hence the determinant (4.) is a product of linear factors $au+bv+w$, where $a$ is one of the $f+1$ quantities $0,u_1,u_2,\cdots u_f$ and $b$ is one of the $f+1$ quantities $0,v_1,v_2,\cdots v_f$. A combination such as $a=u_1,b=0$, however, cannot occur, as the right-hand side of equation (3.) shows. Apart from a power of $w$, the determinant (4.) therefore contains only linear factors of the form $uu_\alpha+vv_\beta+w$, in which $u_\alpha$ is one of the $f$ quantities $u_1,u_2,\cdots u_f$ and $v_\beta$ is one of the $f$ quantities $v_1,v_2,\cdots v_f$.

Now
\[
u(\xi)+v(\overline{\eta})=\left(\frac{e}{h}\chi'\right)\bigl(u(x)+v(\overline{y})\bigr),
\]
and therefore this matrix has rank $ef$. Hence the determinant contains the factor $w$ at least to the power $h-ef$, but also to no higher power, because by (5.) § 8 this is not the case even for $v=0$. The remaining $ef$ linear factors are accordingly all of the form $uu_\alpha+vv_\beta+w$. Let $uu_1+vv_1+w$ be one of them. If one regards the $h$ quantities $y_R$, and hence the $f$ quantities $v_\beta$, as constant, then the determinant (4.), regarded as a function of $w$, has in common with the irreducible function $\Phi(ux_E+vv_1+w,ux_A,ux_B,\cdots)$ the linear factor $uu_1+vv_1+w$. Consequently it has all factors $uu_\alpha+vv_1+w$ $(\alpha=1,2,\cdots f)$ in common with it. Likewise one sees that the determinant contains all the $f^2$ linear functions
\[
uu_\alpha+vv_\beta+w\qquad (\alpha,\beta=1,2,\cdots f)
\]
and each equally often. If each factor occurs $m$ times, then $ef=mf^2$, hence $e=fm$. In this way one can therefore prove, without using the result of § 9, that $e$ is divisible by $f$. By that result, however, $e=f$, and hence $m=1$, so that
\frobtag{5.}{\left|u\xi_{PQ^{-1}}+v\eta_{Q^{-1}P}+w\varepsilon_{PQ^{-1}}\right|=w^{h-ef}\prod_{\alpha,\beta}'(uu_\alpha+vv_\beta+w).}
By this consideration the manner is now determined in which one must assign the roots of the two matrices $(x)$ and $(y)$ to one another in order to obtain the linear factors of the determinant (1.). If $\Phi$ and $\Phi'$ are two different prime factors of $\Theta$, then each root of the equation $\Phi(x-w\varepsilon)=0$ is to be combined with each root of the equation $\Phi(y-w\varepsilon)=0$, but with no root of the equation $\Phi'(y-w\varepsilon)=0$. The generality of the formula obtained is not diminished if one sets $u=v=1$ and $w=0$. If
\frobtagc{6.}{\Psi(x,y)=\prod_{\alpha,\beta}'(u_\alpha+v_\beta)}
is the resultant of the two functions $\Phi(x-\varepsilon w)$ and $\Phi(y+\varepsilon w)$ in the variable $w$, then
\frobtagc{7.}{\left|x_{PQ^{-1}}+y_{Q^{-1}P}\right|=\Pi\Psi(x,y).}
If one could prove directly that this determinant, considered as a function of the $2h$ independent variables $x_R,y_R$, has no multiple factor, then a new proof of the equation $e=f$ would thereby be furnished. If one sets the $h$ quantities $y_R=0$, then $\Psi(x,y)=\Phi(x)^f$. By this route one obtains a deeper insight into the reason for the remarkable phenomenon that the group determinant contains each prime factor in a power whose exponent is equal to the degree of the factor.

Of special interest is the case where $y_R=-x_R$. Then the determinant is
\frobtagc{8.}{\left|x_{PQ^{-1}}-x_{Q^{-1}P}+w\varepsilon_{PQ^{-1}}\right|=w^s\Pi\Psi,}
where
\frobtagc{9.}{\Psi=\prod_{\alpha>\beta}\left(w-(u_\beta-u_\alpha)^2\right)}
and
\frobtagc{10.}{s=\sum f}
Thus the equation $\Psi=0$ is satisfied by the squares of the differences of the roots of the equation $\Phi(x-w\varepsilon)=0$. I shall now show that the elementary divisors of that determinant which vanish for $w=0$ are all linear, or, what is the same thing, that the rank $r$ of the matrix
\frobtagc{11.}{(x_{PQ^{-1}}-x_{Q^{-1}P})}
is equal to
\frobtagc{12.}{r=h-s}
By formula (8.) one has $r\ge h-s$. Since the two matrices $(x)$ and $(x)'$ are commutative with one another, one has
\[
x_A^{(n+1)}=\sum x_Rx_S^{(n)}=\sum x_R^{(n)}x_S\qquad (RS=A)
\]
or
\[
x_A^{(n+1)}=\sum_R x_{AR^{-1}}x_R^{(n)}=\sum_R x_R^{(n)}x_{R^{-1}A},
\]
and hence
\[
\sum_R\left(x_{AR^{-1}}-x_{R^{-1}A}\right)x_R^{(n)}=0.
\]
Putting for $A$, successively, all $h$ elements of $\Hgrp$, one obtains
\frobtagc{13.}{\sum_R\left(x_{AR^{-1}}-x_{R^{-1}A}\right)y_R=0}
a system of $h$ linear equations between the $h$ unknowns $y_R$. The rank of the matrix formed from their coefficients is $r$. Hence the complete system of their solutions forms a matrix of rank $h-r$, and the rank of any system of their solutions, for example the system
\[
y_R=x_R^{(n)}\qquad (n=0,1,2,\cdots)
\]
is $\le h-r$. If the characteristic determinant $\left|x_{PQ^{-1}}-u\varepsilon_{PQ^{-1}}\right|$ of the matrix $(x)$ contains any linear factor $u-u_i$ in the $e$th power, then by § 8 the greatest common divisor of its minors of degree $(h-1)$ contains that factor in the $(e-1)$st power. Consequently the equation of lowest degree satisfied by the matrix $(x)$ is of degree $\sum f=s$, namely $G((x))=0$, where
\[
(-1)^sG(u)=\Pi\Phi(x-u\varepsilon)
\]
Thus the matrices $(x)^0,(x)^1,\cdots(x)^{s-1}$ are linearly independent; consequently the $s$ solutions
\frobtagc{14.}{y_R=x_R^{(n)}\qquad (n=0,1,2,\cdots s-1)}
are also independent, and hence $s\le h-r$. Therefore $r=h-s$, and the quantities (14.) form a complete system of independent solutions of the $h$ linear equations (13.), among which $r$ are independent.

\section*{§ 11.}

By equations (9.) § 8,
\frobtagc{1.}{\Phi(x)=\Phi(\xi),\quad \Phi_n(x)=\Phi_n(\xi),\quad S_n(x)=S_n(\xi).}
hence
\frobtag{2.}{n!\Phi_n(x)=\sum_{R_1,R_2,\cdots R_n}\chi(R_1,R_2,\cdots R_n)\xi_{R_1}\xi_{R_2}\cdots\xi_{R_n}}
and
\frobtag{3.}{S_n(x)=\sum_{R_1,R_2,\cdots R_n}\chi(R_1R_2\cdots R_n)\xi_{R_1}\xi_{R_2}\cdots\xi_{R_n}.}
Another representation follows from the formulae (10.) § 9, namely
\frobtag{4.}{\frac{e}{h}S_n(x)=\xi_E^{(n)}=\sum_{R_1,R_2,\cdots R_n}\xi_{R_1}\xi_{R_2}\cdots\xi_{R_n}\qquad (R_1R_2\cdots R_n=E)}
Accordingly the functions $S_n$ and $\Phi_n$, and in particular $\Phi$ itself, can be expressed through the $h$ variables
\frobtagc{5.}{\xi_S=\frac{e}{h}\sum_R\chi(RS^{-1})x_R}
among which only $ef$ are independent, because by § 8 the rank of the matrix formed from the coefficients of these $h$ linear functions of the $h$ variables $x_R$ is equal to $ef$. If one carries out this transformation for every prime factor of $\Theta$, then the group determinant is expressed through
\frobtagc{6.}{\sum ef=h}
new variables.

Transform each of the $k$ prime functions $\Phi,\Phi',\cdots$ separately by a linear substitution into a function of as few new variables as possible. If their numbers for $\Phi,\Phi',\cdots$ are $g,g',\cdots$, then $g\le ef$, $g'\le e'f'$, $\cdots$. It could then happen that the functions $\Phi,\Phi',\cdots$ together could be represented by still fewer than $g+g'+\cdots$ new variables, linear combinations of the $h$ independent variables $x_R$. If this were the case, or if $g<ef$, or $g'<e'f'$, $\cdots$, then $\Theta$ could be transformed by a linear substitution into a function of fewer than $h$ variables. But now
\[
\sum_R\frac{\partial l\Theta}{\partial x_R}x_{RA^{-1}}=\varepsilon_A h.
\]
Differentiating this equation with respect to $x_{B^{-1}}$, one obtains
\[
\sum_R\frac{\partial^2 l\Theta}{\partial x_R\partial x_{B^{-1}}}x_{RA^{-1}}=-\frac{\partial l\Theta}{\partial x_{B^{-1}A}}=-\frac{h}{\Theta}\Theta_{B^{-1},A^{-1}}
\]
and hence
\[
\left|\frac{\partial^2 l\Theta}{\partial x_P\partial x_{Q^{-1}}}\right|\left|x_{PQ^{-1}}\right|=\left(-\frac{h}{\Theta}\right)^h\left|\Theta_{P,Q}\right|,
\]
so that
\frobtagc{7.}{\left|\frac{\partial^2 l\Theta}{\partial x_P\partial x_{Q^{-1}}}\right|=\frac{(-h)^h}{\Theta^2}}
If, however, one could transform $\Theta$ by a linear substitution into a function of fewer than $h$ variables, this determinant would have to vanish. Consequently $\Phi$ can be expressed through $ef$, but not through fewer than $ef$, variables which are linear combinations of the $h$ variables $x_R$; and the $ef$ variables of $\Phi$, the $e'f'$ of $\Phi'$, $\cdots$, are all independent of one another.

In a particularly simple way, $\Phi^e$ can be represented through the variables $\xi_R$. For this I use the following determinant theorem (compare my paper \emph{On the Pfaffian Problem}, Crelle's Journal vol. 82; § 4, I):

\emph{If $r$ is the rank of the matrix}
\[
a_{\alpha\beta},\qquad (\alpha,\beta=1,2,\cdots n),
\]
\emph{then the determinants of degree $r$ that can be formed from the elements of $r$ columns of this matrix stand in the same ratios as the corresponding determinants of degree $r$ that can be formed from the elements of any other $r$ columns of this matrix.}

Here two determinants are called corresponding if the same rows, and in the same order, are used in their formation. The same theorem holds when the rows and the columns are interchanged. It can be generalized as follows:

\emph{If the matrix $(c_{\alpha\beta})$ is composed from the two matrices $(a_{\alpha\beta})$ and $(b_{\alpha\beta})$, and if $r$ is the rank of the matrix $(a_{\alpha\beta})$, then the determinants of degree $r$ that can be formed from the elements of $r$ columns of the matrix $(a_{\alpha\beta})$ stand in the same ratios as the corresponding determinants of degree $r$ that can be formed from the elements of any $r$ columns of the matrix $(c_{\alpha\beta})$.}

The rank of the matrix
\[
c_{\alpha\gamma}=a_{\alpha1}b_{1\gamma}+a_{\alpha2}b_{2\gamma}+\cdots+a_{\alpha n}b_{n\gamma}
\]
is then at most equal to $r$. But this theorem has significance only when the rank of $(c_{\alpha\gamma})$ is equal to $r$. This must be the case if the determinant of degree $n$, $(b_{\alpha\beta})$, is different from zero.

With the sign
\[
\left|\,a\begin{array}{cccc}
\alpha_1&\alpha_2&\cdots&\alpha_r\\
\beta_1&\beta_2&\cdots&\beta_r
\end{array}\,\right|
\]
I denote the determinant of degree $r$ formed from the elements of the rows $\alpha_1,\alpha_2,\cdots\alpha_r$ and the columns $\beta_1,\beta_2,\cdots,\beta_r$ of the matrix $(a_{\alpha\beta})$, in the indicated order. Thus the ratio
\[
\left|\,a\begin{array}{cccc}
\rho_1&\rho_2&\cdots&\rho_r\\
\alpha_1&\alpha_2&\cdots&\alpha_r
\end{array}\,\right|:
\left|\,c\begin{array}{cccc}
\rho_1&\rho_2&\cdots&\rho_r\\
\beta_1&\beta_2&\cdots&\beta_r
\end{array}\,\right|
\]
is independent of the choice of the indices $\rho_1,\rho_2,\cdots\rho_r$.

If again $(c_{\alpha\beta})=(a_{\alpha\beta})(b_{\alpha\beta})$, and if $r$ is the rank of the matrix $(b_{\alpha\beta})$, then the determinants of degree $r$ that can be formed from the elements of $r$ rows of the matrix $b_{\alpha\beta}$ stand in the same ratios as the corresponding determinants of degree $r$ that can be formed from the elements of any $r$ rows of the matrix $(c_{\alpha\beta})$.

I apply this theorem to the matrix
\frobtagc{8.}{(\xi)=\left(\frac{e}{h}\chi'\right)(x)=(x)\left(\frac{e}{h}\chi'\right)}
The rank of the matrix $(\chi')$ is $g=ef$, and likewise that of the matrix $(\xi)$. Hence the theorem above holds both for the rows and for the columns, and one has
\[
\left|\,\xi\begin{array}{cccc}
A_1&A_2&\cdots&A_g\\
B_1&B_2&\cdots&B_g
\end{array}\,\right|
=\left(\frac{e}{h}\right)^g
\left|\,\chi'\begin{array}{cccc}
A_1&A_2&\cdots&A_g\\
B_1&B_2&\cdots&B_g
\end{array}\,\right|\Psi,
\]
where $\Psi$ is independent of the choice of the elements $A_1,A_2,\cdots A_g,B_1,B_2,\cdots B_g$. Comparing, in the relations
\[
\left|\xi_{PQ^{-1}}+u\varepsilon_{PQ^{-1}}\right|=\Phi(x+u\varepsilon)^e u^{h-g},\qquad
\left|\frac{e}{h}\chi'(PQ^{-1})+u\varepsilon_{PQ^{-1}}\right|=(u+1)^g u^{h-g}
\]
the coefficients of $u^{h-g}$, one obtains: the sum of all principal minors of degree $r$ is, for the matrix $(\xi)$, equal to $\Phi(x)^e$, and, for the matrix $\left(\frac{e}{h}\chi'\right)$, equal to 1. Consequently $\Psi=\Phi^e$, hence
\frobtag{9.}{\left|\,\xi\begin{array}{cccc}
A_1&A_2&\cdots&A_g\\
B_1&B_2&\cdots&B_g
\end{array}\,\right|
=\left(\frac{e}{h}\right)^g
\left|\,\chi\begin{array}{cccc}
B_1&B_2&\cdots&B_g\\
A_1&A_2&\cdots&A_g
\end{array}\,\right|\Phi(x)^e.}
If one chooses the $2g$ elements $A_1,B_1,\cdots A_g,B_g$ so that the determinant of degree $g$ of the matrix $(\chi)$ occurring here is different from zero, then the corresponding determinant of the matrix $(\xi)$ is, up to a constant factor, equal to $\Phi^e$. Among the non-zero determinants of degree $r$ of these two matrices there are also principal minors (in which $B_1=A_1,\cdots B_g=A_g$), because the sum of all principal minors of degree $g$ does not vanish.

In the same way one obtains the more general formula
\frobtag{10.}{\left|\xi_{PQ^{-1}}+\eta_{Q^{-1}P}\right|=\left(\frac{e}{h}\right)^g\left|\chi(QP^{-1})\right|\Psi(x,y),\qquad \left(\begin{array}{l}P=A_1,A_2,\cdots A_g\\ Q=B_1,B_2,\cdots B_g\end{array}\right)}
where $\Psi(x,y)$ has the same meaning as in equation (6.) § 10.

\section*{§ 12.}

The determination of the $k$ prime factors into which the group determinant decomposes has been reduced to the determination of the constants $\chi_\alpha^{(\varkappa)}$, which depends on the solution of an equation of degree $k$. I shall now investigate more closely the algebraic and arithmetic nature of these quantities. First I determine the algebraic field to which these numbers belong.

Let $\Ggrp$ be a subgroup of $\Hgrp$, let $g$ be its order, let $h=gm$, and let $H$ be the group determinant belonging to $\Ggrp$. Let $E,A,B,\cdots$ be the $g$ elements of $\Ggrp$, and let $L,M,\cdots$ be the elements of $\Hgrp$ not contained in $\Ggrp$. If one then sets in $\Theta$ all the variables $x_L,x_M,\cdots$ equal to zero, then, as I have shown in \emph{Ch.} § 7, (10.),
\frobtagc{1.}{\Theta=H^m.}
Hence the coefficients of those terms of $\Theta$ that depend only on $x_E,x_A,x_B,\cdots$ are equal to the coefficients of the corresponding terms of $H^m$.

If $\Ggrp$ is a commutative group, then $H$ is a product of $g$ linear factors,
\[
H=\left(\sum\psi_1(R)x_R\right)\left(\sum\psi_2(R)x_R\right)\cdots,
\]
and the characters $\psi_1(R),\psi_2(R),\cdots$ are all of the first degree, hence roots of unity. In particular $\psi_1(E)=\psi_2(E)=\cdots=1$. Accordingly, if $\Phi$ is a prime factor of degree $f$ of $\Theta$, then this function, when one sets in it $x_L=x_M=\cdots=0$, becomes equal to the product of $f$ of these linear factors, say
\frobtagc{2.}{\Phi=\left(\sum\psi_1(R)x_R\right)\left(\sum\psi_2(R)x_R\right)\cdots\left(\sum\psi_f(R)x_R\right).}
If one replaces $x_E$ by $x_E-u$, one sees that the $f$ roots of the equation $\Phi(x-u\varepsilon)=0$, if one sets $x_L=x_M=\cdots=0$, become integral linear functions of the variables $x_E,x_A,x_B,\cdots$, say
\frobtagc{3.}{u_\lambda=\sum_R\psi_\lambda(R)x_R,\qquad (\lambda=1,2,\cdots f)}
whose coefficients are roots of unity. Thus, if $A$ is one of the elements of $\Ggrp$, the coefficient of $x_E^{f-1}x_A$ in $\Phi(x)$ is equal to
\frobtagc{4.}{\chi(A)=\psi_1(A)+\psi_2(A)+\cdots+\psi_f(A),}
and this equation also holds for $A=E$. If furthermore $A$ and $B$ are two elements of $\Ggrp$, then, because $\psi_\lambda$ is a character of the first degree, $\psi_\lambda(AB)=\psi_\lambda(A)\psi_\lambda(B)$, and hence
\frobtagc{5.}{\chi(AB)=\psi_1(A)\psi_1(B)+\psi_2(A)\psi_2(B)+\cdots+\psi_f(A)\psi_f(B).}
Therefore the matrix $(\chi(PQ^{-1}))$, which for the group $\Hgrp$ has rank $ef$, has at most rank $f$ if $P$ and $Q$ run only through the elements of $\Ggrp$.

Let $A$ be any element of $\Hgrp$, and let $m$ be its order. The powers of $A$ then form a commutative group of order $m$. If one chooses this for $\Ggrp$, the $m$ characters $\psi_1(A)=\rho_1$, $\psi_2(A)=\rho_2,\cdots$ become all the $m$th roots of unity. Hence
\frobtagc{6.}{\chi(A)=\rho_1+\rho_2+\cdots+\rho_f,\qquad \chi(A^n)=\rho_1^n+\rho_2^n+\cdots+\rho_f^n}
for every value of $n$. For $n=m-1$ it follows that $\chi(A)$ and $\chi(A^{-1})$ are conjugate complex quantities (\emph{Ch.} § 3).

\emph{If $A$ is an element of order $m$ and $\chi$ is a character of degree $f$, then $\chi(A)$ can be represented as a sum of $f$ $m$th roots of unity.}

These are determined individually by the fact that $\chi(A^n)$ is equal to the sum of their $n$th powers. If in $\Phi$ one sets all variables equal to zero except $x_E,x_A,x_{A^2},\cdots x_{A^{m-1}}$, then by (2.)
\[
\Phi=(x_E+\rho_1x_A+\rho_1^2x_{A^2}+\cdots+\rho_1^{m-1}x_{A^{m-1}})\cdots(x_E+\rho_fx_A+\rho_f^2x_{A^2}+\cdots+\rho_f^{m-1}x_{A^{m-1}}).
\]
If one therefore sets in $\Phi$ all variables equal to zero except $x_E$ and $x_A$, then
\frobtagc{7.}{\Phi(x_E,x_A,0,0,\cdots)=(x_E+\rho_1x_A)(x_E+\rho_2x_A)\cdots(x_E+\rho_fx_A)}
and consequently, by (9.) § 4,
\frobtagc{8.}{u^f+\zchar_1(A)u^{f-1}+\zchar_2(A)u^{f-2}+\cdots+\zchar_f(A)=(u+\rho_1)(u+\rho_2)\cdots(u+\rho_f).}
In particular, since $\zchar_f(A)=\zchar(A)$,
\frobtagc{9.}{\zchar(A)=\rho_1\rho_2\cdots\rho_f.}
Applying formula (1.) to the group $\Ggrp$ formed by the powers of $A$, one obtains
\[
\Theta=\left(\Pi_\rho(x_E+\rho x_A+\rho^2x_{A^2}+\cdots+\rho^{m-1}x_{A^{m-1}})\right)^{\frac{h}{m}},
\]
where $\rho$ runs through all $m$th roots of unity. Thus, if one sets in $\Theta$ all variables equal to zero except $x_E$ and $x_A$, then
\frobtagc{10.}{\Theta(x_E,x_A,0,0,\cdots)=(x_E^m+(-x_A)^m)^{\frac{h}{m}}.}
Consequently, as I showed in another way in § 3, the coefficient of $x_A^h$ in $\Theta$ is equal to
\frobtagc{11.}{\Pi\zchar(A)^e=(-1)^{h-\frac{h}{m}}.}
Let $B$ be a second element of $\Hgrp$, and let $\sigma_1,\sigma_2,\cdots\sigma_f$ be the $f$ roots of unity corresponding to the character $\chi(B)$, so that
\frobtagc{12.}{\chi(B)=\sigma_1+\sigma_2+\cdots+\sigma_f,\qquad \chi(B^n)=\sigma_1^n+\sigma_2^n+\cdots+\sigma_f^n.}
If now $A$ and $B$ are commutative with one another, then they generate a group $\Ggrp$. Therefore one can apply formula (5.) and sees: the roots of unity of the two systems can be assigned to one another in such a way that
\frobtagc{13.}{\chi(AB)=\rho_1\sigma_1+\cdots+\rho_f\sigma_f,\\ \chi(A^rB^s)=\rho_1^r\sigma_1^s+\cdots+\rho_f^r\sigma_f^s}
results. If in the prime function $\Phi$ that corresponds to the character $\chi$ one sets all variables equal to zero except $x_E,x_A$ and $x_B$, then by (2.) one obtains
\frobtagc{14.}{\Phi(x_E,x_A,x_B,0,0,\cdots)=(x_E+\rho_1x_A+\sigma_1x_B)\cdots(x_E+\rho_fx_A+\sigma_fx_B),}
which at the same time determines the assignment of the roots of unity.

To give an example for the theorems developed, I consider an invariant element $B$ of the group $\Hgrp$, that is, one which is commutative with every element $R$ of $\Hgrp$. In formula (3.) § 7 one then has $R^{-1}BR=B$, and consequently, if $A$ is any other element of $\Hgrp$,
\[
\chi(A)\chi(B)=f\chi(AB).
\]
All invariant elements of $\Hgrp$ form a commutative group $\Ggrp$. If for each element $G$ of it one sets $\chi(G)=f\psi(G)$, then for any two elements $A$ and $B$ of $\Ggrp$ one accordingly has $\psi(A)\psi(B)=\psi(AB)$. Hence $\psi(G)$ is a character of $\Ggrp$, and therefore a root of unity $\rho$. Further, $\psi(G^n)=\psi(G)^n$, so $\chi(G)=f\rho$ and $\chi(G^n)=f\rho^n$. For an invariant element $G$ of $\Hgrp$, consequently, the $f$ roots of unity occurring in formula (6.) are all equal to one another.

The same result is obtained from the observation that an invariant element $A$ by itself alone forms a class of conjugate elements. If one therefore sets in $\Phi$ all variables equal to zero except $x_E$ and $x_A$, then by formula (7.) § 6, $\Phi$ becomes the $f$th power of a linear function of $x_E$ and $x_A$, and consequently by formula (7.) one has $\rho_1=\rho_2=\cdots=\rho_f$.

Now let $A$ again be any element of $\Hgrp$, and let $m$ be its order. If $\rho$ is a primitive $m$th root of unity, then the quantities $\rho_1,\rho_2,\cdots\rho_f$ in formula (6.) are all powers of $\rho$, and therefore $\chi(A)$ is a number of the field $K(\rho)$, which is formed from all rational functions of $\rho$. Among the elements of the group $\Hgrp$ conjugate with $A$ there may also be powers of $A$, namely $A^r,A^s,A^t,\cdots$. Their exponents are relatively prime to $m$ and form a group, that is, one among them is $\equiv rs\pmod m$ when $r$ and $s$ are any two of them. Since $A$ and $A^r$ are conjugate, $\chi(A)=\chi(A^r)$, hence $\rho_1+\cdots+\rho_f=\rho_1^r+\cdots+\rho_f^r$. If one expresses $\chi(A)$ through $\rho$, this number therefore remains unchanged when one replaces $\rho$ by $\rho^r$. Those numbers of the field $K(\rho)$ that remain unchanged when one replaces $\rho$ by $\rho^r$ or $\rho^s$, or by $\rho^t,\cdots$, form a field $\Lambda(\rho)$, a divisor of $K(\rho)$. The number $\chi(A)$ consequently belongs to this field $\Lambda(\rho)$.

If, for example, $\Hgrp$ is the symmetric group of degree $n$, so that $h=n!$, then $A$ is conjugate (similar) with every power $A^r$ whose exponent $r$ is relatively prime to $m$. Hence the characters of the symmetric group are all integral rational numbers (compare the examples $n=4$ and 5, \emph{Ch.} § 8).

In the field $\Lambda(\rho)$, $\chi(A)$, being a sum of roots of unity, is an algebraic integer. The same is true, however, of every coefficient of $\Phi$, and therefore also of $\Phi_n$. For if, in a product $\Theta=\Phi\Psi$ of two integral functions of arbitrarily many variables whose coefficients are algebraic numbers, all coefficients are algebraic integers, then the product of each coefficient of $\Phi$ and each coefficient of $\Psi$ is also an algebraic integer (compare Dedekind, \emph{On an Arithmetic Theorem of Gauss}. Prague Math. Soc. 1892). If $A,B,C,\cdots$ are distinct elements of $\Hgrp$ and $r+s+t+\cdots=n$, then in $\Phi_n$ the coefficient of $x_A^rx_B^sx_C^t\cdots$ is
\frobtagc{15.}{\frac{1}{r!s!t!\cdots}\chi(A,\cdots A,B,\cdots B,C,\cdots C,\cdots).}
Consequently this expression is an algebraic integer. Applying the same theorem to the factors of the product (9.) § 6, one sees that also
\frobtagc{16.}{\frac{h_\alpha\chi_\alpha}{f}}
is an algebraic integer. Consequently, by (6.) § 7, also
\[
\sum\frac{h_\alpha\chi_\alpha}{f}\chi_{\alpha'}=\frac{h}{e}
\]
is an integer. Therefore the number $e=f$ is a divisor of the order $h$.

\begin{center}
\vspace{1.0em}
\rule{0.32\linewidth}{0.4pt}\\[0.35em]
Issued on 10 December.\\[1.8em]
{\scriptsize Berlin, printed in the Reichsdruckerei.}
\end{center}
% ===== END WORK 054 =====

\clearpage

% ===== BEGIN WORK 056: 056_ber_die_Darstellung_der_endlichen_Gruppen =====
\begin{center}
{\Large On the representation of finite groups\\by linear substitutions.}\par
\vspace{1.2em}
By G. \textsc{Frobenius}.
\end{center}
\vspace{1em}

In my paper \emph{On the prime factors of the group determinant} (Proceedings 1896; cited in the following as \emph{Pr.}) I have assigned to each finite group $\Hgrp$ of order $h$ a matrix of degree $h$, whose elements depend on $h$ variables. Its importance rests on the connection in which it itself and the prime factors of its determinant stand with the linear substitutions by which the group $\Hgrp$ and the groups isomorphic to it can be represented. From each such representation one can derive a matrix belonging to the group $\Hgrp$, whose determinant is contained in a power of the group determinant (§ 2). If the determinant is irreducible, hence equal to a prime factor of the group determinant, I call the representation a primitive one. Conversely, to each prime factor of degree $f$ of the group determinant there corresponds one, and, apart from the choice of variables, only one primitive representation of the group by substitutions of $f$ variables (§ 4).

The group matrix can be transformed into a similar matrix which decomposes into partial matrices. If one uses for this purpose only the $k$ group characters, one can decompose it into $k$ matrices, each of which has as determinant the $f^{\ten}$ power of a prime function $\Phi$ of degree $f$ (§ 3). If, however, one uses higher irrationalities, then one can decompose it into $\sum f$ matrices, each of which has a prime function $\Phi$ itself as determinant (§ 5). With the aid of some remarkable theorems on determinants of degree $n$, whose elements are $n^2$ independent variables (§ 7), I then show: one can arrange the transformation in such a way that, for each set of $f$ partial matrices whose determinants are equal to the same prime factor $\Phi$ of degree $f$, the matrices are identically equal to one another. Then the elements of all partial matrices together are $\sum f^2=h$ variables independent of one another. From such a partial matrix, whose determinant is $\Phi$, there result $h$ linear substitutions, which form a group isomorphic to $\Hgrp$. The isomorphism can also be a merohedric one. This depends on a special relation in which the character $\chi$ of the group $\Hgrp$, which corresponds to the prime function $\Phi$, can stand to an invariant subgroup of $\Hgrp$ (§ 1). The primitive representations of a group by linear substitutions cast new light on the significance of the relations from which the characters of the group and therefore the coefficients of the prime factors of the group determinant are computed (§ 6). I communicated the most important results of this paper in April of this year to \textsc{Dedekind}, to whom I owe the stimulus for these investigations.

\section*{§ 1.}

Among the characters of a group $\Hgrp$ the principal character, whose values are all equal to $1$, occupies a privileged position. If $\Hgrp$ is composite, and if $\Ggrp$ is an invariant subgroup of $\Hgrp$, then there are further certain characters of $\Hgrp$ which show a special behavior with respect to $\Ggrp$. Such a character has, namely, the same value for any two elements of $\Hgrp$ that are equivalent modulo $\Ggrp$. I therefore say that it belongs to the group $\frac{\Hgrp}{\Ggrp}$, on the basis of the theorem:

\begin{quote}\itshape
If $\Ggrp$ is an invariant subgroup of the group $\Hgrp$, then every character of $\frac{\Hgrp}{\Ggrp}$ is also a character of $\Hgrp$.
\end{quote}

This is to be understood as follows. If $A,B,C,\cdots$ form a complete residue system of $\Hgrp$ (mod. $\Ggrp$), then
\[
\Hgrp=A\Ggrp+B\Ggrp+C\Ggrp+\cdots.
\]
The complexes $A\Ggrp=\Ggrp A$, $B\Ggrp=\Ggrp B$, $\ldots$ can then be regarded as the elements of the group $\frac{\Hgrp}{\Ggrp}$. Let $\psi$ be a character of this group, and let $\psi(A\Ggrp)$ be its value for the element $A\Ggrp$. Each element $R$ of the group $\Hgrp$ belongs to one and only one of these complexes. For all elements $R$ of the complex $A\Ggrp$ let $\chi(R)=\psi(A\Ggrp)$. The $h$ quantities $\chi(R)$ so defined form a character of $\Hgrp$. Conversely, let $\chi(R)$ be a character of $\Hgrp$ which has the same value for every two elements of $\Hgrp$ that are equivalent modulo $\Ggrp$. If one then puts $\psi(A\Ggrp)=\chi(R)$ for each element $R$ of $\Hgrp$ that is equivalent to $A$ (mod. $\Ggrp$), and so belongs to the complex $A\Ggrp$, then $\psi$ is a character of $\frac{\Hgrp}{\Ggrp}$. The two corresponding characters have the same degree,
\[
f=\chi(E)=\psi(\Ggrp).
\]
If the character $\chi$ belongs to the group $\frac{\Hgrp}{\Ggrp}$, then it also belongs to $\frac{\Hgrp}{\Ggrp'}$, where $\Ggrp'$ is an invariant subgroup of $\Hgrp$ contained in $\Ggrp$.

For $h$ quantities $\chi(R)$ to form a character of $\Hgrp$, the following conditions are necessary and sufficient:
\begin{equation}\tag{1.}
\chi(E)=f
\end{equation}
is a positive (whole) number. For every two elements $A$ and $B$ of $\Hgrp$ one has
\begin{equation}\tag{2.}
\chi(AB)=\chi(BA)
\end{equation}
and
\begin{equation}\tag{3.}
h\chi(A)\chi(B)=f\sum_R\chi(AR^{-1}BR).
\end{equation}
Finally,
\begin{equation}\tag{4.}
h=\sum_R\chi(R)\chi(R^{-1}).
\end{equation}
By the second condition, $\chi(R)$ has the same value for all elements of a class of conjugate elements, say for the elements of the $\rho^{\ten}$ class the value $\chi_\rho$ $(\rho=0,1,\cdots k-1)$. Therefore equations (3.) and (4.) are identical with (compare \emph{Pr.} § 7)
\begin{equation}\tag{3.}
h_\alpha h_\beta\chi_\alpha\chi_\beta=f\sum_\gamma h_\gamma c_{\alpha\beta\gamma}\chi_\gamma,
\end{equation}
\begin{equation}\tag{4.}
h=\sum_\rho h_\rho\chi_\rho\chi_{\rho'}.
\end{equation}
Every solution of equations (3.) gives the ratios of the values of a character $\chi$. Equation (4.) then supplies the square of the still undetermined factor of proportionality, and finally equation (1.) supplies its sign.

Now $\chi(E)=\psi(\Ggrp)=f$. Further, $A\Ggrp B\Ggrp=AB\Ggrp$. Thus, if $R$ is an element of the complex $A\Ggrp$, and $S$ an element of $B\Ggrp$, then $RS$ is an element of $AB\Ggrp$. Hence $\chi(RS)=\psi(AB\Ggrp)=\psi(A\Ggrp B\Ggrp)$ and $\chi(SR)=\psi(B\Ggrp A\Ggrp)$, and consequently, since $\psi$ possesses property (2.), $\chi(RS)=\chi(SR)$. Further,
\[
\frac{h}{g}\psi(A\Ggrp)\psi(B\Ggrp)=f\sum_S\psi(AS^{-1}BS\Ggrp),
\]
where $S$ runs through a complete residue system of $\Hgrp$ (mod. $\Ggrp$). Now $AS^{-1}BS\Ggrp=A\Ggrp S^{-1}BS\Ggrp$ and therefore remains unchanged when $S$ is replaced by $SG$, where $G$ is any element of $\Ggrp$. Hence also
\[
h\psi(A\Ggrp)\psi(B\Ggrp)=f\sum_R\psi(AR^{-1}BR\Ggrp),
\]
where $R$ runs through all elements of $\Hgrp$, and therefore also
\[
h\chi(A)\chi(B)=f\sum_R\chi(AR^{-1}BR).
\]
Likewise, from the relation
\[
\frac{h}{g}=\sum_S\psi(S\Ggrp)\psi(S^{-1}\Ggrp)
\]
one obtains the equation
\[
h=\sum_R\chi(R)\chi(R^{-1}).
\]
Thus $\chi(R)$ is a character of $\Hgrp$.

Conversely, let $\chi(R)$ be a character of $\Hgrp$, and let $\chi(R)=\chi(S)$ always hold when $R\sim S$ (mod. $\Ggrp$). Then one recognizes in the same way that $\psi(R\Ggrp)=\chi(R)$ is a character of $\frac{\Hgrp}{\Ggrp}$.

The principal character belongs to the group $\frac{\Hgrp}{\Hgrp}$; the characters of the first degree belong to the commutative group $\frac{\Hgrp}{\Ggrp}$ if $\Ggrp$ is the commutator group of $\Hgrp$.

If $G$ runs through the elements of the group $\Ggrp$, let
\begin{equation}\tag{5.}
\sum_G x_{RG}=y_{R\Ggrp}.
\end{equation}
Then under the above hypotheses (compare \emph{Pr.} § 8)
\[
\frac{h}{f}\xi_A=\sum_R\chi(R)x_{RA}=\sum_S\psi(S\Ggrp)y_{SA\Ggrp}=\frac{h:g}{f}\eta_{A\Ggrp}.
\]
The rank of the matrix $(\xi_{PQ^{-1}})$ is equal to $f^2$. If $\Phi(x)$ is the prime factor of the group determinant $\Theta(x)$ of $\Hgrp$ which corresponds to the character $\chi$, then every subdeterminant of degree $f^2$ of that matrix is, up to a constant factor, equal to $\Phi(x)^f$. If $\Psi(y)$ is the prime factor of the group determinant $H(y)$ of $\frac{\Hgrp}{\Ggrp}$ which corresponds to the character $\psi$, then every subdeterminant of degree $f^2$ of the matrix $\eta_{PQ^{-1}\Ggrp}$ is equal to $\Psi(y)^f$. Now $g\xi_A=\eta_{A\Ggrp}$, and hence, under assumption (5.), $\Phi(x)=\Psi(y)$. Since $\Theta(x)$ and $H(y)$ contain the factors $\Phi(x)$ and $\Psi(y)$ exactly in the $f^{\ten}$ power, $\Theta(x)$ is divisible by $H(y)$, and the quotient is relatively prime to $H(y)$ (compare \emph{Pr.} § 2).

\section*{§ 2.}

A finite number of linear substitutions
\begin{align*}
(A)\qquad x_\alpha&=a_{\alpha1}y_1+a_{\alpha2}y_2+\cdots+a_{\alpha n}y_n,\qquad (\alpha=1,2,\cdots n)\\
(B)\qquad x_\alpha&=b_{\alpha1}y_1+b_{\alpha2}y_2+\cdots+b_{\alpha n}y_n,
\end{align*}
\[
\cdot\quad\cdot\quad\cdot\quad\cdot\quad\cdot\quad\cdot\quad\cdot\quad\cdot\quad\cdot\quad\cdot,
\]
whose determinants are different from zero form a group $\Hgrp'$ if the substitution $(C)=(A)(B)$ composed from any two of them, $(A)$ and $(B)$, is likewise contained among them. The coefficients of $(C)$ are
\[
c_{\alpha\beta}=a_{\alpha1}b_{1\beta}+a_{\alpha2}b_{2\beta}+\cdots+a_{\alpha n}b_{n\beta}.
\]
Following the procedure of \textsc{Gauss}, one may also entirely disregard the notation of the variables and understand by $(A),(B),(C),\cdots$ the matrices of degree $n$ formed from the systems of coefficients $a_{\alpha\beta},b_{\alpha\beta},c_{\alpha\beta},\cdots$.

Let $\Hgrp$ be an abstract group, $A,B,C,\cdots$ its elements. If one assigns to the element $A$ the matrix $(A)$, to the element $B$ the matrix $(B)$, and so on, let the group $\Hgrp'$ be isomorphic to the group $\Hgrp$, that is, let $(A)(B)=(AB)$. I then say that the substitutions, or the matrices, $(A),(B),(C),\cdots$ represent the group $\Hgrp$. The isomorphism can also be a merohedric one. Then $\Hgrp'$ is holohedrically isomorphic to a group $\frac{\Hgrp}{\Ggrp}$, where $\Ggrp$ is a certain invariant subgroup of $\Hgrp$, formed from the elements of $\Hgrp$ to which the principal element $(E)$ of $\Hgrp'$ corresponds. This substitution $(E)$ corresponds to the principal element $E$ of $\Hgrp$. Since $E^2=E$, one also has $(E)^2=(E)$. Hence $(E)$ is the identical substitution, since no other substitution of non-vanishing determinant satisfies this condition.

If $(P)$ is any matrix of degree $n$ with non-vanishing determinant, then the matrices $(P)^{-1}(A)(P),(P)^{-1}(B)(P),\cdots$ also represent the group $\Hgrp$. The corresponding substitutions arise from the original ones by subjecting each system of $n$ variables $x_1,\cdots x_n;y_1,\cdots y_n;\cdots$ to the substitution $(P)$. I call two such representations of $\Hgrp$ equivalent. All representations that are equivalent to a given one form a class of equivalent representations of $\Hgrp$.

To the $h$ elements $A,B,C,\cdots$ of $\Hgrp$ I assign $h$ independent variables $x_A,x_B,x_C,\cdots$, and from the matrices $(A),(B),(C),\cdots$ I then form the matrix
\[
(A)x_A+(B)x_B+(C)x_C+\cdots=\sum_R(R)x_R,
\]
whose elements are linear functions of the independent variables. Every matrix which corresponds in this way to a representation of $\Hgrp$ I call a matrix belonging to the group $\Hgrp$.

Let $y_A,y_B,y_C,\cdots$ be a second system of $h$ independent variables, and let (\emph{Pr.} § 1)
\begin{equation}\tag{1.}
z_T=\sum x_Ry_S,\qquad (RS=T).
\end{equation}
Then the product of the two matrices is
\begin{equation}\tag{2.}
\left(\sum(R)x_R\right)\left(\sum(S)y_S\right)=\sum(RS)x_Ry_S=\sum(T)z_T.
\end{equation}
Conversely, this property characterizes the matrix $\sum(R)x_R$ as a matrix belonging to the group $\Hgrp$. Such a matrix is therefore, for example, the group matrix $(x_{PQ^{-1}})$. If now the determinant $\left|\sum(R)x_R\right|=F(x)$, then $F(x)F(y)=F(z)$. By \emph{Pr.} § 1 it follows from this that $F(x)$ is absorbed in a power of the group determinant
\[
\Theta(x)=\left|x_{PQ^{-1}}\right|=\Phi^f{\Phi'}^{f'}{\Phi''}^{f''}\cdots,
\]
that is, it is a product of powers of the prime factors $\Phi,\Phi',\Phi'',\cdots$ of $\Theta$:
\begin{equation}\tag{3.}
\left|\sum(R)x_R\right|=\Phi^s{\Phi'}^{s'}{\Phi''}^{s''}\cdots,
\end{equation}
where the exponents $s,s',s'',\cdots$ may in part also be zero.

If $\Phi$ is a prime factor of degree $f$ of $\Theta$, then $\Theta$ is divisible exactly by the $f^{\ten}$ power of $\Phi$. By \emph{Pr.} § 5, the subdeterminant of degree $(h-1)$ of $\Theta$ complementary to the element $x_{PQ^{-1}}$ is equal to
\[
\frac1h\frac{\partial\Theta}{\partial x_{PQ^{-1}}}.
\]
Therefore the greatest common divisor of all subdeterminants of degree $(h-1)$ of $\Theta$ contains the prime factor $\Phi$ exactly in the $(f-1)^{\ten}$ power. By the known properties of the elementary divisors of a determinant, the greatest common divisor of all subdeterminants of degree $(h-m)$ of $\Theta$ therefore contains the factor $\Phi$ exactly in the $(f-m)^{\ten}$ power $(m<f)$, and the subdeterminants of degree $(h-f)$ of $\Theta$ are not all divisible by $\Phi$. The characteristic determinant $\Theta(x-u\eps)$ of the group matrix consequently has, as a function of $u$, only linear elementary divisors, and if one sets $u$ equal to a root of the equation $\Phi(x-u\eps)=0$, that determinant has rank $h-f$.

If, in formula (1.), one puts $y_R=\dfrac{1}{h\Phi^{f-1}}\dfrac{\partial\Theta}{\partial x_R}$, then $z_R=\dfrac{\Theta}{\Phi^{f-1}}\eps_R$, hence, since $\sum(R)\eps_R=(E)$,
\begin{equation}\tag{4.}
\left(\sum(R)x_R\right)\left(\sum(S)y_S\right)=\frac{\Theta}{\Phi^{f-1}}(E).
\end{equation}
Now the elementary divisors of the determinant of the latter matrix which are equal to powers of $\Phi$ are all equal to the first power of $\Phi$. On the basis of (4.), however, these elementary divisors are divisible by the corresponding elementary divisors of $\left|\sum(R)x_R\right|$ (\emph{On the elementary divisors of determinants}, Theorem IX, Proceedings 1894). Hence also the elementary divisors of the determinant $\left|\sum x_R(R)\right|$ which are powers of $\Phi$ are all equal to the first power of $\Phi$, and their number is equal to $s$. Replacing $x_R$ by $x_R-u\eps_R$, one sees that the characteristic determinant of that matrix $\left|\left(\sum x_R(R)\right)-u(E)\right|$ has only linear elementary divisors.

\section*{§ 3.}

If a matrix belonging to the group $\Hgrp$ decomposes into partial matrices, then by formula (2.) of § 2 each single one of them is also a matrix belonging to the group. The product of the determinants of the partial matrices is equal to the determinant of the whole matrix. As I shall show, one can find a matrix equivalent to the group matrix which decomposes into $\sum f$ partial matrices. Their determinants are equal to the $\sum f$ prime factors of $\Theta$. The decomposition of the group matrix cannot be carried further. Before I pass to the presentation of this transformation, however, I shall first carry out another, simpler one, in which the determinant of each partial matrix is equal to the $f^{\ten}$ power of a prime factor $\Phi$ of degree $f$. Thus the individual partial matrices have only different determinants, of which any two are relatively prime. To carry out this decomposition, knowledge of the characters is sufficient.

Let $\Phi$ be a prime factor of degree $f$ of $\Theta$, and let $\chi$ be the corresponding character. Then the matrix of degree $h$, $(\chi(PQ^{-1}))$, has rank $r=f^2$, and in it there is a non-zero principal minor of degree $r$ (\emph{Pr.} § 11). Let such a minor be obtained by putting the $r$ elements $A_1,A_2,\ldots,A_r$ for $P$ and $Q$. Let $\psi$ be a character different from $\chi$, let $s$ be the rank of the matrix $(\psi(PQ^{-1}))$, and let a non-zero principal minor of degree $s$ be obtained in it by putting $P$ and $Q$ equal to $B_1,B_2,\ldots,B_s$. These $s$ elements may also agree wholly or partly with $A_1,A_2,\ldots,A_r$. If one makes such a determination for each of the $k$ characters of $\Hgrp$, then the sum of the $k$ numbers is
\[
r+s+\cdots=\sum r=\sum f^2=h.
\]
Now let $M$ be the matrix of degree $h$ whose $h$ rows are obtained from the row
\[
\chi(PA_1^{-1}),\ldots,\chi(PA_r^{-1}),\psi(PB_1^{-1}),\ldots,\psi(PB_s^{-1}),\ldots
\]
by putting for $P$ the $h$ elements of $\Hgrp$. Further, let $L'$ be the matrix whose rows arise, in an analogous way, from
\[
\chi(A_1Q^{-1}),\ldots,\chi(A_rQ^{-1}),\psi(B_1Q^{-1}),\ldots,\psi(B_sQ^{-1}),\ldots.
\]
If in $L'$ one interchanges rows and columns, one obtains the conjugate matrix $L$. I now form the composite matrix $LM$. If $\alpha$ and $\beta$ are two of the indices $1,2,\ldots,r$, then in it the $\beta^{\te}$ element of the $\alpha^{\ten}$ row is
\[
\sum_R\chi(A_\alpha R^{-1})\chi(RA_\beta^{-1})=\frac{h}{f}\chi(A_\alpha A_\beta^{-1}).
\]
The determinant of degree $r$ formed from these $r^2$ elements is, by hypothesis, different from zero. Further, if $\alpha$ is one of the indices $1,2,\ldots,r$ and $\beta$ one of the indices $1,2,\ldots,s$, then the $(r+\beta)^{\te}$ element of the $\alpha^{\ten}$ row is
\[
\sum_R\chi(A_\alpha R^{-1})\psi(RB_\beta^{-1})=0.
\]
Consequently the matrix $LM$ decomposes into $k$ partial matrices, and its determinant is
\[
\prod\left|\frac{h}{f}\chi(A_\alpha A_\beta^{-1})\right|,
\]
and therefore different from zero. Hence the determinants $|L|$ and $|M|$ are both different from zero.

Now let $X$ be the group matrix $(x_{PQ^{-1}})$ and let $Y=LXM$. If $\alpha$ and $\beta$ are two of the numbers $1,2,\ldots,r$, then in it the $\beta^{\te}$ element of the $\alpha^{\ten}$ row is
\[
\sum_{R,S}\chi(A_\alpha R^{-1})x_{RS^{-1}}\chi(SA_\beta^{-1}).
\]
But since $\chi(PQ)=\chi(QP)$, the two matrices $(x_{PQ^{-1}})$ and $(\chi(PQ^{-1}))$ are interchangeable with one another (\emph{Pr.} § 6), and that sum is equal to
\[
\sum_{R,S}\chi(A_\alpha R^{-1})\chi(RS^{-1})x_{SA_\beta^{-1}}=\frac{h}{f}\sum_S\chi(A_\alpha S^{-1})x_{SA_\beta^{-1}}.
\]
If $\alpha$ is one of the numbers $1,2,\ldots,r$, and $\beta$ one of the numbers $1,2,\ldots,s$, then the $(r+\beta)^{\te}$ element of the $\alpha^{\ten}$ row is
\[
\sum_{R,S}\chi(A_\alpha R^{-1})x_{RS^{-1}}\psi(SB_\beta^{-1})
=\sum_{R,S}\chi(A_\alpha R^{-1})\psi(RS^{-1})x_{SB_\beta^{-1}}=0.
\]
Consequently $Y$ decomposes into $k$ partial matrices, the first of which is formed from the $r^2$ elements
\[
\frac{h}{f}\sum_R\chi(A_\alpha R^{-1})x_{RA_\beta^{-1}}
=\sum_R\frac{h}{f}\chi(A_\alpha A_\beta^{-1}R^{-1})x_R.
\]
In this sum, $x_E$ is multiplied by the matrix
\[
N_1=\left(\frac{h}{f}\chi(A_\alpha A_\beta^{-1})\right)\qquad(\alpha,\beta=1,2,\ldots,r),
\]
whose determinant is different from zero. Similarly let
\[
N_2=\left(\frac{h}{f'}\psi(B_\alpha B_\beta^{-1})\right)\qquad(\alpha,\beta=1,2,\ldots,s)
\]
and so on. From these matrices of degrees $r,s,\ldots$ form the matrix
\[
\begin{pmatrix}
N_1&0&0&\cdots\\
0&N_2&0&\cdots\\
0&0&N_3&\cdots\\
\cdot&\cdot&\cdot&\cdots
\end{pmatrix}=N
\]
of degree $h$; then $YN^{-1}=Z$ also decomposes into $k$ partial matrices, the first of which is
\begin{equation}\tag{1.}
\left(\sum_R\frac{h}{f}\chi(A_\alpha A_\beta^{-1}R^{-1})x_R\right)
\left(\frac{h}{f}\chi(A_\alpha A_\beta^{-1})\right)^{-1}\qquad(\alpha,\beta=1,2,\ldots,f^2).
\end{equation}
In it, $x_E$ is multiplied by the principal matrix. Therefore, comparing in the equation $LXMN^{-1}=Z$ the matrices multiplied by $x_E$, one obtains $LMN^{-1}=E$, and hence
\begin{equation}\tag{2.}
LX L^{-1}=Z
\end{equation}
is a matrix equivalent to $X$. Thus $Z$ and every partial matrix of $Z$, such as (1.), belongs to the group $\Hgrp$. By \emph{Pr.} § 11, the determinant of (1.) is equal to $\Psi^f$, where $\Psi$ is the prime factor of $\Theta$ conjugate to $\Phi$.

If $\Ggrp$ is an invariant subgroup of $\Hgrp$, if the character $\chi$ belongs to the group $\frac{\Hgrp}{\Ggrp}$, and if the matrix (1.) is equal to $\sum(R)x_R$, then $(R)=(S)$ whenever $R\sim S$ (mod. $\Ggrp$). Therefore the matrices $(A),(B),(C),\cdots$ represent the group $\frac{\Hgrp}{\Ggrp}$.

The above transformation can be carried out in the same way, without the use of new auxiliary means, for any matrix belonging to the group. Such a matrix too, if its determinant (3.) § 2 is divisible by different prime factors of the group determinant, can, by use of the characters alone, be transformed into an equivalent decomposed matrix whose individual parts have determinants $\Phi^s,{\Phi'}^{s'},\ldots$.

\section*{§ 4.}

I now turn to the complete decomposition of the group matrix discussed at the beginning of the preceding paragraph. Its characteristic determinant $\Theta(x-u\eps)$ has only simple elementary divisors and therefore, for a root $u$ of the equation $\Phi(x-u\eps)=0$, has rank $h-f$. I now put, for the variables $x_R$, such constants $k_R$ that the equation $\Phi(k-u\eps)=0$ has a simple root $\rho$ for which none of the functions $\Phi'(k-u\eps),\Phi''(k-u\eps),\ldots$ vanishes. Then the matrix $(k_{PQ^{-1}}-\rho\eps_{PQ^{-1}})$ also has rank $h-f$. Consequently the $h$ linear equations
\begin{equation}\tag{1.}
\sum_Q k_{PQ^{-1}}a_Q=\rho a_P
\end{equation}
have $f$ independent solutions $a'_Q,a''_Q,\ldots,a^{(f)}_Q$. From them every other solution can be composed by multiplying them by certain factors and then adding. If $P$ and $Q$ are replaced by $PR^{-1}$ and $QR^{-1}$, then
\[
\sum_Q k_{PQ^{-1}}a_{QR^{-1}}=\rho a_{PR^{-1}},
\]
and consequently, if $x_A,x_B,x_C,\ldots$ are independent variables,
\[
\sum_{Q,R} k_{PQ^{-1}}a_{QR^{-1}}x_R=\rho\sum_R a_{PR^{-1}}x_R.
\]
Thus $a_Q=\sum_R a^{(\kappa)}_{QR^{-1}}x_R$ is also a solution of equations (1.). Hence there are factors $x_{\kappa\lambda}$ such that
\begin{equation}\tag{2.}
\sum_R a^{(\kappa)}_{QR^{-1}}x_R=\sum_\lambda x_{\kappa\lambda}a^{(\lambda)}_Q,\qquad(\kappa,\lambda=1,2,\ldots,f).
\end{equation}
The factors $x_{\kappa\lambda}$ are completely determined by these conditions and are consequently linear functions of the variables $x_R$. The $f$ new solutions of equations (1.) so obtained are linearly independent or not according as the determinant of degree $f$, $|x_{\kappa\lambda}|$, is different from zero or not. Since $x_{\kappa\lambda}=\eps_{\kappa\lambda}$ for $x_R=\eps_R$, that determinant cannot vanish identically.

If $x_{\kappa\lambda}$ passes into $y_{\kappa\lambda}$ or into $z_{\kappa\lambda}$ when the variables $x_R$ are replaced by other variables $y_R$ or $z_R$, then also
\[
\sum_R a^{(\kappa)}_{QR^{-1}}y_R=\sum_\lambda y_{\kappa\lambda}a^{(\lambda)}_Q,
\qquad
\sum_R a^{(\kappa)}_{QR^{-1}}z_R=\sum_\lambda z_{\kappa\lambda}a^{(\lambda)}_Q,
\]
and hence
\[
\sum_{Q,R} a^{(\kappa)}_{PQ^{-1}}x_{QR^{-1}}y_R
=\sum_\lambda x_{\kappa\lambda}\sum_R a^{(\lambda)}_{PR^{-1}}y_R
=\sum_{\lambda,\mu}x_{\kappa\lambda}y_{\lambda\mu}a^{(\mu)}_Q,
\]
therefore, under the hypothesis (1.), § 2,
\[
\sum_\mu z_{\kappa\mu}a^{(\mu)}_Q
=\sum_Q a^{(\kappa)}_{PQ^{-1}}z_Q
=\sum_{\lambda,\mu}x_{\kappa\lambda}y_{\lambda\mu}a^{(\mu)}_Q,
\]
and consequently
\[
z_{\kappa\mu}=\sum_\lambda x_{\kappa\lambda}y_{\lambda\mu}.
\]
Therefore $(x_{\kappa\lambda})$ is a matrix of degree $f$ belonging to $\Hgrp$, and hence its determinant is
\[
|x_{\kappa\lambda}|=\Phi^s{\Phi'}^{s'}\cdots.
\]
If $\chi$ is the character corresponding to the prime function $\Phi$, and one puts $\frac{h}{f}\chi(R^{-1})=c_R$, then by \emph{Pr.} § 8
\[
\Phi(c)=1,
\quad \Phi'(c)=0,\ldots,
\quad \Phi(c-u\eps)=(1-u)^f,
\quad \Phi'(c-u\eps)=(-u)^{f'},\ldots,
\]
where $\Phi'$ is any prime function of degree $f'$ different from $\Phi$. If $\rho,\rho',\rho'',\ldots$ are the $f$ roots of the equation $\Phi(k-u\eps)=0$, then by \emph{Pr.} § 6, $\rho-\lambda,\rho'-\lambda,\rho''-\lambda,\ldots$ are those of the equation $\Phi(k-\lambda c-u\eps)=0$. Thus if $\lambda$ is a constant different from $0,\rho,\rho',\rho'',\ldots$, then $\Phi(k-\lambda c-\rho\eps)$ is different from zero. Further, $\Phi'(k-\lambda c-\rho\eps)=\Phi'(k-\rho\eps)$ is different from zero, and therefore the determinant $\Theta(k-\lambda c-\rho\eps)$ cannot vanish.

For $x_R=c_R$ the determinant $|x_{\kappa\lambda}|$ does not vanish. For otherwise the $f$ solutions
\[
\sum_Q a^{(\kappa)}_{PQ^{-1}}c_Q\qquad(\kappa=1,2,\ldots,f)
\]
would not be linearly independent. One could then compose from them a solution $\sum_Q a_{PQ^{-1}}c_Q=0$ in which the quantities $a_P$ satisfy equations (1.) and are not all zero. Since $c_{PQ}=c_{QP}$, the matrices $a_{PQ^{-1}}$ and $c_{PQ^{-1}}$ are interchangeable, and hence one also has $\sum_Q c_{PQ^{-1}}a_Q=0$. Thus also
\[
\sum_Q(k_{PQ^{-1}}-\lambda c_{PQ^{-1}}-\rho\eps_{PQ^{-1}})a_Q=0.
\]
But since the determinant of these $h$ equations is different from zero, only the values $a_Q=0$ can satisfy them.

The product $\Phi(c)^s{\Phi'}(c)^{s'}{\Phi''}(c)^{s''}\cdots$ can be different from zero only if $s'=s''=\cdots=0$. Since $|x_{\kappa\lambda}|$ is of degree $f$, it follows that
\begin{equation}\tag{3.}
|x_{\kappa\lambda}|=\Phi(x),\qquad |x_{\kappa\lambda}-u\eps_{\kappa\lambda}|=\Phi(x-u\eps).
\end{equation}
Now I have shown in \emph{Pr.} § 11 that $\Phi(x)$ can be represented by $f^2$, but not by fewer, linear combinations of the variables $x_R$. Hence the $f^2$ linear functions $x_{\kappa\lambda}$ of the $h$ variables $x_R$ are independent of one another.

This remarkable theorem, that there is a matrix belonging to the group $\Hgrp$ whose $f^2$ elements are independent variables, was also found by \textsc{Molien} in his excellent work \emph{Über Systeme höherer complexer Zahlen} (Math. Ann. vol. 41, p. 124), to which \textsc{Study} recently drew my attention. In a further work, \emph{Eine Bemerkung zur Theorie der homogenen Substitutionsgruppen}, Sitzungsberichte der Naturforscher-Gesellschaft zu Dorpat 1897, year 18, p. 259, \textsc{Molien} applied the general results found there specifically to the group determinant.

If $X=(x_{\kappa\lambda})=\sum(R)x_R$, then comparison of the coefficients of $u^{f-1}$ in (3.) gives
\begin{equation}\tag{4.}
\sum_R\chi(R)x_R=\sum_\lambda x_{\lambda\lambda}.
\end{equation}
Thus if $(R)=(r_{\kappa\lambda})$, then
\begin{equation}\tag{5.}
\chi(R)=\sum_\lambda r_{\lambda\lambda}.
\end{equation}
If now the group formed from the matrices $(A),(B),(C)\cdots$ is isomorphic to the group $\frac{\Hgrp}{\Ggrp}$, then $(R)=(S)$ whenever $R\sim S$ (mod. $\Ggrp$). Hence also $\chi(R)=\chi(S)$. The character $\chi$ therefore likewise belongs to the group $\frac{\Hgrp}{\Ggrp}$.

Let $L=(l_{\kappa\lambda})$ be any constant matrix of degree $f$ with non-vanishing determinant. If one then chooses for the $f$ solutions $a'_Q,a''_Q,\ldots$ of equations (1.) any other $f$ independent solutions, then instead of $X$ one always obtains a matrix $L^{-1}XL$ similar to $X$, and, with a suitable choice of the $f$ solutions, every such matrix.

But one may also choose for $\rho$ another root of the equation $\Phi(k-u\eps)=0$, and further, for $k_R$ every system of $h$ numbers is admissible which satisfies certain inequalities. Then other quantities take the place of $a'_Q,a''_Q,\ldots$, and instead of the matrix $X$ one obtains a matrix $U=(u_{\kappa\lambda})$, whose elements $u_{\kappa\lambda}$ are linear functions of the variables $x_R$. For every choice of the arbitrary quantities, however, $U$ is a matrix of degree $f$ belonging to $\Hgrp$, and its characteristic determinant is $\Phi(x-u\eps)$. Since $\Phi(x)$ contains the $h$ quantities $x_R$ only in the $f^2$ independent linear combinations $x_{\kappa\lambda}$, the $f^2$ quantities $u_{\kappa\lambda}$ must be linear combinations of the variables $x_{\kappa\lambda}$. By a theorem which I shall develop in § 7, one can therefore determine a constant matrix $L$ so that either $U=L^{-1}XL$ or $U=L^{-1}X'L$, where $X'$ is the matrix conjugate to $X$; and in fact, if $f>1$, only one of these two cases is possible.

For the matrices considered here, however, when $f>1$, the case $U=L^{-1}X'L$ cannot occur. For if the variables $x_R$ are replaced by $y_R$ or $z_R$, let $X$ pass into $Y$ or $Z$, and $U$ into $V$ or $W$. Then $XY=Z$ and $UV=W$. Now if $U=L^{-1}X'L$, $V=L^{-1}Y'L$, $W=L^{-1}Z'L$, then also $X'Y'=Z'$, and consequently, since $X'Y'=(YX)'$, one has $Z=YX=XZ$. Thus if $X=\sum(R)x_R$, any two of the matrices $(A),(B),(C)\cdots$ are interchangeable. The group formed from them and the group $\frac{\Hgrp}{\Ggrp}$ holohedrically isomorphic to it is therefore commutative, and the character $\chi$ belonging to it has degree $f=1$ (\emph{Pr.} § 2). Accordingly $U=L^{-1}XL$, and for every disposition of the arbitrary quantities one can choose the $f$ solutions $a'_Q,a''_Q,\ldots$ so that $U=X$.

If the character $\chi$ belongs to the group $\frac{\Hgrp}{\Ggrp}$, then there is a prime function $\Psi(y)$ of degree $f$ of the determinant of this group which passes into $\Phi(x)$ by the substitution (5.) § 1. This can be represented as the determinant of a matrix of degree $f$, $(y_{\kappa\lambda})$, whose elements are linear functions of the variables $y_{R\Ggrp}$ and which belongs to the group $\frac{\Hgrp}{\Ggrp}$. If in it one makes the substitution (5.) § 1, it passes into a matrix $X=(x_{\kappa\lambda})$ which belongs to the group $\Hgrp$ and whose determinant is equal to $\Phi(x)$. Since every other matrix having the same properties is equal to $L^{-1}XL$, the converse of the theorem obtained above is thereby proved, namely that, if $\chi$ belongs to the group $\frac{\Hgrp}{\Ggrp}$, then one always also has $(R)=(S)$ whenever $R\sim S$ (mod. $\Ggrp$).

\section*{§ 5.}

Now let the quantities $k_R$ be so chosen that the equation
\[
\Phi(k-u\eps)\Phi'(k-u\eps)\Phi''(k-u\eps)\cdots=0
\]
has no multiple roots. Then the determinant of the family of bilinear forms
\begin{equation}\tag{1.}
\sum_{P,Q}(k_{PQ^{-1}}-u\eps_{PQ^{-1}})u_Pv_Q
\end{equation}
has only linear elementary divisors. The form $\sum k_{PQ^{-1}}u_Pv_Q$ can therefore, by a theorem of \textsc{Weierstrass}, be transformed by a linear substitution of non-vanishing determinant
\begin{equation}\tag{2.}
u'_\nu=\sum_R a^{(\nu)}_R u_R,\qquad v_R=\sum_\nu a^{(\nu)}_R v'_\nu\qquad(\nu=1,2,\ldots,h),
\end{equation}
which carries $\sum u_Rv_R$ into $\sum u'_\nu v'_\nu$, into
\begin{equation}\tag{3.}
\rho u'_1v'_1+\rho u'_2v'_2+\cdots+\rho u'_fv'_f+
ho'u'_{f+1}v'_{f+1}+\cdots+
ho'u'_{f+f'}v'_{f+f'}+\cdots
\end{equation}
if $\rho$ is an $f$-fold root, $\rho'$ an $f'$-fold root, $\ldots$ of the characteristic equation $\Theta(k-u\eps)=0$. Then
\[
\sum_{P,Q}k_{PQ^{-1}}u_P\sum_\nu a^{(\nu)}_Qv'_\nu
=\rho\left(\sum_P a^{(1)}_P u_P\right)v'_1+
ho\left(\sum_P a^{(2)}_P u_P\right)v'_2+\cdots,
\]
and hence
\[
\sum_Q k_{PQ^{-1}}a^{(\kappa)}_Q=\rho a^{(\kappa)}_P\qquad(\kappa=1,2,\ldots,f).
\]
Since the determinant of degree $h$, $|a^{(\nu)}_R|$, is different from zero, not all determinants of degree $f$ vanish in the system of the $fh$ quantities $a^{(\kappa)}_Q$, and hence they form $f$ independent solutions of equations (1.) § 4. If they are replaced by any other $f$ independent solutions, substitution (2.) retains the double property that its determinant is different from zero and that it transforms the family of forms (1.) into the normal form (3.). Therefore the $h^2$ quantities $a^{(\nu)}_R$ can be defined by the linear equations (1.) § 4 and by the analogous equations which correspond to the roots $\rho',\rho'',\ldots$, and the essential content of the theorem of \textsc{Weierstrass} used above consists in the fact that the determinant of degree $h$, $|a^{(\nu)}_R|$, is then different from zero. Now let, as above,
\[
\sum_Q a^{(\kappa)}_{PQ^{-1}}x_Q=\sum_\lambda x_{\kappa\lambda}a^{(\lambda)}_P
\]
or
\[
\sum_P a^{(\kappa)}_{QP^{-1}}x_{PQ^{-1}}=\sum_\lambda x_{\kappa\lambda}a^{(\lambda)}_{Q^{-1}}\qquad(\kappa,\lambda=1,2,\ldots,f).
\]
In the same way let, corresponding to the $f'$-fold root $\rho'$,
\[
\sum_P a^{(\kappa)}_{QP^{-1}}x_{PQ^{-1}}=\sum_\lambda x_{\kappa\lambda}a^{(\lambda)}_{Q^{-1}}\qquad(\kappa,\lambda=f+1,f+2,\ldots,f+f').
\]
In the special case where $\rho'$ is a root of the same equation $\Phi(x-u\eps)=0$ as $\rho$ (thus $f'=f$), one can, by suitable choice of the solution of the equations analogous to (1.) § 4, achieve that $x_{f+\kappa,f+\lambda}=x_{\kappa\lambda}$.

Then the bilinear form
\begin{equation}\tag{4.}
\sum_{P,Q}x_{PQ^{-1}}u_Pv_Q
\end{equation}
passes, by the substitution
\begin{equation}\tag{5.}
u_R=\sum_\nu a^{(\nu)}_{R^{-1}}u'_\nu,
\qquad
v'_\nu=\sum_R a^{(\nu)}_{R^{-1}}v_R
\end{equation}
into
\begin{equation}\tag{6.}
\sum_{\kappa,\lambda}^{1\ldots f}x_{\kappa\lambda}u'_\kappa v'_\lambda
+\sum_{\kappa,\lambda}^{f+1\ldots f+f'}x_{\kappa\lambda}u'_\kappa v'_\lambda+\cdots,
\end{equation}
which decomposes into $f$ forms in $f^2$ variables, $f'$ forms in ${f'}^2$ variables, and so on. By the same substitution $\sum u_Rv_R$ passes into $\sum u'_\nu v'_\nu$. Thus if $X$ is the group matrix $(x_{PQ^{-1}})$, there is a constant matrix $L$ of degree $h$ with non-vanishing determinant such that the equivalent matrix $L^{-1}XL$ decomposes into $f$ mutually equal partial matrices of degree $f$, whose characteristic determinants are equal to $\Phi(x-u\eps)$, into $f'$ mutually equal partial matrices of degree $f'$, whose characteristic determinants are equal to $\Phi'(x-u\eps)$, and so on; and the $f^2+{f'}^2+\cdots=h$ elements of these partial matrices are $h$ mutually independent variables, because $\Theta(x)$ cannot be expressed by fewer than $h$ linear combinations of the variables $x_R$.

A representation of a group by linear substitutions for which the corresponding determinant (3.) § 2 is irreducible I call a \emph{primitive representation}. Then the number of classes of primitive representations for the group $\Hgrp$ and for the groups isomorphic to it is equal to the number $k$ of classes of conjugate elements into which the elements of $\Hgrp$ decompose. If $f$ is the number of variables transformed by one of the substitutions, then the $k$ numbers $g=f^2$ are found by solving an equation of degree $k$ which I developed in my work \emph{Über Gruppencharaktere} (Sitzungsberichte 1896, § 4, (12.)).

To illustrate this transformation of the group matrix, I choose the example that \textsc{Dedekind} found in the year 1886 and communicated to me in April 1896. Let
\[
\begin{array}{cccccccccccc}
1.&abc&2.&bca&3.&cab&4.&acb&5.&cba&6.&bac
\end{array}
\]
be the 6 permutations of 3 symbols. The substitutions which carry $abc$ into these 6 permutations shall be denoted, instead of by $A,B,C$, by the numerals $1,2,\ldots,6$. Let $\rho$ be a third root of unity and
\[
\begin{array}{lll}
u=x_1+x_2+x_3, &\qquad& v=x_4+x_5+x_6,\\
u_1=x_1+\rho x_2+\rho^2x_3, && v_1=x_4+\rho x_5+\rho^2x_6,\\
u_2=x_1+\rho^2x_2+\rho x_3, && v_2=x_4+\rho^2x_5+\rho x_6.
\end{array}
\]
Further, let $X$, $L$, and $U$ be the three matrices
\[
\begin{pmatrix}
x_1&x_3&x_2&x_4&x_5&x_6\\
x_2&x_1&x_3&x_5&x_6&x_4\\
x_3&x_2&x_1&x_6&x_4&x_5\\
x_4&x_5&x_6&x_1&x_3&x_2\\
x_5&x_6&x_4&x_2&x_1&x_3\\
x_6&x_4&x_5&x_3&x_2&x_1
\end{pmatrix},
\qquad
\begin{pmatrix}
1&-1&1&0&0&1\\
1&-1&\rho^2&0&0&\rho\\
1&-1&\rho&0&0&\rho^2\\
1&1&0&1&1&0\\
1&1&0&\rho&\rho^2&0\\
1&1&0&\rho^2&\rho&0
\end{pmatrix},
\]
\[
\begin{pmatrix}
u+v&0&0&0&0&0\\
0&u-v&0&0&0&0\\
0&0&u_1&v_1&0&0\\
0&0&v_2&u_2&0&0\\
0&0&0&0&u_1&v_1\\
0&0&0&0&v_2&u_2
\end{pmatrix}.
\]
Then $X$ is the group matrix, and
\[
XL=LU,
\qquad
L^{-1}XL=U,
\]
and by forming $L'L$ one recognizes that the determinant of $L$ does not vanish.

\section*{§ 6.}

Formula (5.) § 4 leads to a deeper insight into the meaning of the equations which I developed for the calculation of the characters of a group $\Hgrp$. If the determinant of the matrix belonging to the group
\[
X=(x_{\kappa\lambda})=\sum_R(R)x_R\qquad(\kappa,\lambda=1,2,\ldots,f)
\]
is a prime factor $\Phi$ of degree $f$ of the group determinant, then the $f^2$ quantities $x_{\kappa\lambda}$ are independent variables. Therefore one can give the $h$ variables $x_R$ such values that $X$ becomes equal to an arbitrary matrix of degree $f$.

The two similar matrices $(A)$ and $(R)^{-1}(A)(R)=(B)$ have the same characteristic function $|(A)-u(E)|=|(B)-u(E)|$. Comparing on both sides of this equation the coefficients of $u^{f-1}$, one obtains, by (5.) § 4, $\chi(A)=\chi(B)$. Hence $\chi(R^{-1}AR)=\chi(A)$, or, what is the same, $\chi(QP)=\chi(PQ)$.

Divide the elements of $\Hgrp$ into the $k$ classes of conjugate elements, and let the $\alpha^{\te}$ class consist of $h_\alpha$ elements $(\alpha=0,1,\ldots,k-1)$. If $R$ runs through all $h$ elements of $\Hgrp$, and $A$ is a fixed element of the $\alpha^{\ten}$ class, then $R^{-1}AR$ represents every element of the class and represents each $\frac{h}{h_\alpha}$ times. Hence
\[
\sum_R(R^{-1}AR)=\frac{h}{h_\alpha}\sum_{(\alpha)}'(R),
\]
where the sum on the right is to be extended over the $h_\alpha$ elements of the $\alpha^{\ten}$ class. But now
\[
(S)\sum_R(R^{-1}AR)=\sum_R(SR^{-1}AR)=\sum_R(R^{-1}ARS)=\left(\sum_R(R^{-1}AR)\right)(S).
\]
The third sum is obtained from the second by replacing $R$ by $RS$. Hence the matrix $\sum_{(\alpha)}'(R)$ is interchangeable with $(S)$, and therefore also with $\sum(S)x_S$, consequently with every matrix of degree $f$. But only the principal matrix, multiplied by a scalar factor, has this property. Consequently
\begin{equation}\tag{1.}
f\sum_{(\alpha)}'(R)=(E)h_\alpha\chi_\alpha,
\qquad
f\sum_R(R^{-1}AR)=(E)h\chi_\alpha.
\end{equation}
To determine the still unknown factor $\chi_\alpha$, form on both sides the sum of the diagonal elements. For each of the $h_\alpha$ matrices $(R)$ similar to $(A)$, this sum is equal to $\chi(R)=\chi(A)$, and for $(E)$ it is equal to $f$. Hence
\[
\chi_\alpha=\chi(A).
\]
Composing the matrix $(A)$ with the matrix
\[
f\sum_{(\beta)}'(S)=(E)h_\beta\chi_\beta,
\]
one obtains
\[
f\sum_{(\beta)}'(AS)=(A)h_\beta\chi_\beta,
\]
and, if one again forms on both sides the sum of the diagonal elements,
\begin{equation}\tag{2.}
f\sum_{(\beta)}'\chi(AS)=h_\beta\chi(A)\chi(B),
\qquad
f\sum_R\chi(AR^{-1}BR)=h\chi(A)\chi(B).
\end{equation}
These are the equations which determine the ratios of the $k$ values of each of the $k$ characters.

One can, however, give these formulas another interpretation. Let $e_A,e_B,e_C,\cdots$ be the $h$ independent units of a system of hypercomplex numbers, for which the ordinary rules of addition and the distributive and associative law of multiplication, but not the commutative law, are assumed. The multiplication law
\begin{equation}\tag{3.}
e_Ae_B=e_{AB}
\end{equation}
satisfies these conditions. Accordingly $e_E$ represents the number 1. If one now sets
\begin{equation}\tag{4.}
\sum_{(\alpha)}'e_R=\frac{h_\alpha}{f}e_\alpha,
\end{equation}
then the $k$ complex numbers $e_0,e_1,\ldots,e_{k-1}$ are interchangeable with every number of the system, and hence also with one another. The product of two of these commutative numbers is
\[
\frac{h_\alpha h_\beta}{f^2}e_\alpha e_\beta
=\left(\sum_{(\alpha)}'e_R\right)\left(\sum_{(\beta)}'e_S\right)
=\sum_\tau' h_{\alpha,\beta,\tau} e_T,
\]
where $h_{\alpha,\beta,\tau}$ indicates how many of the $h_\alpha h_\beta$ products $RS$ are equal to the fixed element $T$. By \emph{Pr.} § 7 this number, if $T$ belongs to the $\gamma^{\ten}$ class, is equal to $\frac{1}{h_\gamma}h_{\alpha\beta\gamma}$, and thus has the same value for conjugate elements $T$. Hence
\begin{equation}\tag{5.}
\frac{h_\alpha h_\beta}{f^2}e_\alpha e_\beta
=\sum_\gamma^{k-1}\frac{1}{h_\gamma}h_{\alpha\beta\gamma}e_\gamma.
\end{equation}
Consequently the $k$ independent numbers $e_0,e_1,\ldots,e_{k-1}$ by themselves form the basis of a system of complex numbers for which the commutative law of multiplication also holds. As I have shown in \emph{Über Gruppencharaktere}, p. 991, the restrictions under which \textsc{Weierstrass} and \textsc{Dedekind} investigated such number systems also hold for them. Consequently there are actual numbers $\chi_\alpha$ which, substituted for $e_\alpha$, satisfy equations (5.):
\begin{equation}\tag{6.}
h_\alpha h_\beta\chi_\alpha\chi_\beta=f\sum_\gamma h_{\alpha\beta\gamma}\chi_\gamma.
\end{equation}
Therefore one cannot arrive at contradictions if the former independence of the $h$ complex units $e_R$ is restricted by the linear equations $e_\alpha=\chi_\alpha e_0$, or
\begin{equation}\tag{7.}
f\sum_{(\alpha)}'e_R=h_\alpha\chi_\alpha e_E.
\end{equation}
By multiplication with $e_B$, the further linear equations
\begin{equation}\tag{8.}
f\sum_{(\alpha)}'e_{BR}=h_\alpha\chi_\alpha e_B
\end{equation}
flow from these by (3.). In consequence of these relations the number of linearly independent units among the $e_R$ is reduced to $f^2$, and these can be so chosen that the formulas for their multiplication agree with those for the composition of all matrices of degree $f$.

The calculation indicated here was carried out by \textsc{Dedekind} for the example of a group of order $h=6$ communicated in § 5 and for some other groups of the smallest orders, especially for the group $\Omega$ of order 8, which he called the quaternion group because the system of hypercomplex numbers derived from it agrees with the system of \textsc{Hamilton}'s quaternions. In his work \emph{Über Gruppen, deren sämmtliche Theiler Normaltheiler sind}, Math. Ann. vol. 48, p. 551, he alludes to these relations with the words: ``But there is, as I already recognized in February 1886, a still deeper relation between the group $\Omega$ and \textsc{Hamilton}'s quaternions.'' Calculation with quaternions is indeed, if ordinary complex numbers are admitted as scalar coefficients, completely equivalent to calculation with matrices of the second degree, because the two quaternary quadratic forms $x^2+y^2+z^2+t^2$ and $xt-yz$ can be transformed into one another. Only if one restricts oneself to real coefficients can the quaternions claim an independent significance, the essence of which consists in the fact that, apart from them and the ordinary complex numbers, there is no number system in which a product cannot vanish without one of the factors being zero, as I first set forth in my work \emph{Über lineare Substitutionen und bilineare Formen}, Crelle's Journal, vol. 84, § 14.

\section*{§ 7.}

Let the elements $x_{\alpha\beta}$ $(\alpha,\beta=1,2,\ldots,n)$ of a matrix $X$ of degree $n$ be $n^2$ mutually independent variables. If one interchanges in $X$ the rows with the columns, one obtains the matrix $X'$ conjugate to $X$. Let the elements of the two matrices $A$ and $B$ be constant quantities. Let their determinants $|A|$ and $|B|$ be non-zero quantities whose product is equal to $k$. Then the two matrices $AXB$ and $AX'B$ have determinant $k|X|$, and in each of them the elements are linear functions of the $n^2$ variables $x_{\alpha\beta}$. Conversely, the following theorem holds:

\begin{quote}\itshape
I. If the elements of the matrix $X$ are independent variables and those of the matrix $Y$ are linear functions of these variables, and if the determinant of the matrix $Y$ differs from that of the matrix $X$ only by a non-zero constant factor, then either $Y=AXB$ or $Y=AX'B$, where $A$ and $B$ are constant matrices; and in fact, when the degree of $X$ is greater than 1, only one of these two cases occurs, and the matrices $A$ and $B$ are completely determined up to a scalar factor.
\end{quote}

The second part of this theorem is easy to prove. For let $Y=AXB=CXD$. If one puts therein $x_{\alpha\beta}=0$ or 1 according as $\alpha$ and $\beta$ are different or equal, it follows that $AB=CD$, and if one puts $BD^{-1}=A^{-1}C=F$, then $XF=FX$. Accordingly $F$ is interchangeable with every matrix, and hence $F=hE$, where $h$ is a scalar factor and $E$ is the principal matrix. Consequently $C=hA$ and $D=\frac{1}{h}B$. Further, if $n>1$, neither can $AXB=CX'D$ hold. For from it there follows in the same way $XF=FX'$, thus, for example, for $n=2$,
\[
\begin{pmatrix}x&y\\z&t\end{pmatrix}
\begin{pmatrix}a&b\\c&d\end{pmatrix}
=\begin{pmatrix}a&b\\c&d\end{pmatrix}
\begin{pmatrix}x&z\\y&t\end{pmatrix}.
\]
These four linear equations can be satisfied, if the unknowns $a,b,c,d$ are to be independent of the variables $x,y,z,t$, only by vanishing values of the constants.

Between the minors of degree $m$ $(0<m<n)$ of the matrix $X$ there is no linear relation with constant coefficients. For let $u,v,w,\ldots$ be these minors, $a,b,c,\ldots$ constants, and let $au+bv+cw+\cdots=0$. If one sets therein all variables $x_{\alpha\beta}=0$ except the $m^2$ which occur in $u$, then $v,w,\ldots$ vanish, and hence $a$ must be zero. The minors of degree $(n-1)$ are the first derivatives of $|X|$ with respect to the variables $x_{\alpha\beta}$. Since there is no linear relation with constant coefficients between them, $|X|$ cannot be represented as a function of fewer than $n^2$ variables which are linear functions of the $n^2$ variables $x_{\alpha\beta}$. Thus if $|Y|=k|X|$, the $n^2$ variables $y_{\alpha\beta}$ are independent of one another. Therefore there is also no linear relation with constant coefficients between the minors of degree $m$ of the matrix $Y$.

Let the coefficient of $x_{\kappa\kappa}$ in $y_{\alpha\beta}$ be $c_{\alpha\beta}^{(\kappa)}$, or, since $\kappa$ is at first a fixed index, simply $c_{\alpha\beta}$. Let $r$ be a new variable, and replace $x_{\kappa\kappa}$ by $x_{\kappa\kappa}+r$; then $y_{\alpha\beta}$ passes into $y_{\alpha\beta}+rc_{\alpha\beta}$. Hence $|y_{\alpha\beta}+rc_{\alpha\beta}|$ is equal to the determinant which arises from $k|x_{\alpha\beta}|$ by replacing $x_{\kappa\kappa}$ by $x_{\kappa\kappa}+r$. This, however, is a function of the first degree in $r$, and consequently the coefficient of $r^2$ in that determinant also vanishes. That coefficient is equal to the sum of the products of each determinant of the second degree of the matrix $c_{\alpha\beta}$ and the complementary determinant of degree $(n-2)$ of the matrix $y_{\alpha\beta}$. But between the latter, if $n>2$, there is no linear relation with constant coefficients. Consequently all determinants of the second degree of the matrix $c_{\alpha\beta}$ must vanish. For $n=2$ this is evident without further argument. Therefore one can determine $2n$ quantities $p_\alpha,q_\alpha$ so that $c_{\alpha\beta}=p_\alpha q_\beta$.

Let $c_{\alpha\beta}^{(\kappa)}=p_{\alpha\kappa}q_{\kappa\beta}$. If one sets the variables $x_{\alpha\beta}$ whose indices $\alpha$ and $\beta$ are different equal to zero, then let $X$ and $Y$ pass into $X_0$ and $Y_0$. The elements of the matrix $Y_0$ are then
\[
\sum_\kappa c_{\alpha\beta}^{(\kappa)}x_{\kappa\kappa}=\sum_\kappa p_{\alpha\kappa}x_{\kappa\kappa}q_{\kappa\beta}.
\]
Thus if $P$ is the matrix formed from the quantities $p_{\alpha\beta}$, and $Q$ the matrix formed from the quantities $q_{\alpha\beta}$, then
\[
Y_0=PX_0Q,
\]
and because $|Y_0|=k|X_0|=kx_{11}x_{22}\cdots x_{nn}$, one has $|P|\,|Q|=k$. Accordingly $|P|$ and $|Q|$ are different from zero.

The elements $z_{\alpha\beta}$ of the matrix $Z=P^{-1}YQ^{-1}$ therefore become equal to those of $X_0$ if one sets equal to zero the variables $x_{\alpha\beta}$ whose indices are different; that is, they become equal to $x_{\alpha\alpha}$ or 0 according as $\beta=\alpha$ or not. Or the quantities $z_{\alpha\beta}$ $(\alpha\ne\beta)$ and $z_{\alpha\alpha}-x_{\alpha\alpha}=v_\alpha$ depend only on the variables $x_{\alpha\beta}$ with different indices. Expanding the identical equation $|X|=|Z|$ in powers of $x_{11},x_{22},\ldots,x_{nn}$, comparison of the coefficients of the product $x_{22}x_{33}\cdots x_{nn}$ gives $v_1=0$. Similarly $v_\alpha=0$. Comparing then the terms multiplied by $x_{33}x_{44}\cdots x_{nn}$, one finds $x_{12}x_{21}=z_{12}z_{21}$ and in general
\begin{equation}\tag{1.}
x_{\alpha\beta}x_{\beta\alpha}=z_{\alpha\beta}z_{\beta\alpha}.
\end{equation}
Finally, comparing the coefficients of $x_{44}x_{55}\cdots x_{nn}$, one obtains
\[
\left|\begin{array}{ccc}
0&x_{12}&x_{13}\\
x_{21}&0&x_{23}\\
x_{31}&x_{32}&0
\end{array}\right|=
\left|\begin{array}{ccc}
0&z_{12}&z_{13}\\
z_{21}&0&z_{23}\\
z_{31}&z_{32}&0
\end{array}\right|
\]
and likewise in general
\[
x_{\beta\gamma}x_{\gamma\alpha}x_{\alpha\beta}+x_{\gamma\beta}x_{\alpha\gamma}x_{\beta\alpha}
=z_{\beta\gamma}z_{\gamma\alpha}z_{\alpha\beta}+z_{\gamma\beta}z_{\alpha\gamma}z_{\beta\alpha}.
\]
By (1.) the product of the two summands on the left is also equal to that of the summands on the right. Therefore either
\[
x_{\beta\gamma}x_{\gamma\alpha}x_{\alpha\beta}=z_{\beta\gamma}z_{\gamma\alpha}z_{\alpha\beta},
\qquad
x_{\gamma\beta}x_{\alpha\gamma}x_{\beta\alpha}=z_{\gamma\beta}z_{\alpha\gamma}z_{\beta\alpha},
\]
or
\[
x_{\beta\gamma}x_{\gamma\alpha}x_{\alpha\beta}=z_{\gamma\beta}z_{\alpha\gamma}z_{\beta\alpha},
\qquad
x_{\gamma\beta}x_{\alpha\gamma}x_{\beta\alpha}=z_{\beta\gamma}z_{\gamma\alpha}z_{\alpha\beta}.
\]
Since the quantities $z_{\alpha\beta}$ are linear functions of the $n^2$ independent variables $x_{\alpha\beta}$, it follows from (1.), up to a constant factor which we shall at first disregard, that either $z_{\alpha\beta}=x_{\alpha\beta}$, $z_{\beta\alpha}=x_{\beta\alpha}$, or $z_{\alpha\beta}=x_{\beta\alpha}$, $z_{\beta\alpha}=x_{\alpha\beta}$.

Let $z_{12}=x_{12}$, $z_{21}=x_{21}$. Then $x_{12}x_{2\alpha}x_{\alpha1}=z_{21}z_{\alpha2}z_{1\alpha}$ $(\alpha>2)$ cannot hold, because the right-hand side is not divisible by $x_{12}$. Therefore $x_{12}x_{2\alpha}x_{\alpha1}=z_{12}z_{2\alpha}z_{\alpha1}$, hence $z_{\alpha1}=x_{\alpha1}$, $z_{2\alpha}=x_{2\alpha}$, and consequently $z_{1\alpha}=x_{1\alpha}$, $z_{\alpha2}=x_{\alpha2}$. Further, one must then have $x_{1\alpha}x_{\alpha\beta}x_{\beta1}=z_{1\alpha}z_{\alpha\beta}z_{\beta1}$, and hence $z_{\alpha\beta}=x_{\alpha\beta}$. If, on the other hand, $z_{12}=x_{21}$, then one recognizes in the same way that in general $z_{\alpha\beta}=x_{\beta\alpha}$.

These equations, however, are exact only up to constant factors. If $z_{\alpha1}=\frac{k_\alpha}{k_1}x_{\alpha1}$, then by (1.) $z_{1\alpha}=\frac{k_1}{k_\alpha}x_{1\alpha}$ and, by the equation $z_{1\alpha}z_{\alpha\beta}z_{\beta1}=x_{1\alpha}x_{\alpha\beta}x_{\beta1}$, in general $z_{\alpha\beta}=\frac{k_\alpha}{k_\beta}x_{\alpha\beta}$. Thus if one sets
\[
R=\begin{pmatrix}
k_1&0&0&\cdots\\
0&k_2&0&\cdots\\
0&0&k_3&\cdots\\
\cdot&\cdot&\cdot&\cdots
\end{pmatrix},
\]
then $Z=RXR^{-1}$ and $Y=PRXP^{-1}Q=AXB$. Likewise in the other case $z_{\alpha\beta}=\frac{k_\alpha}{k_\beta}x_{\beta\alpha}$, $Y=PRX'R^{-1}Q=AX'B$.

\begin{quote}\itshape
II. If the hypotheses of Theorem I are fulfilled, and if the characteristic functions of the matrices $X$ and $Y$ are also equal to one another, then either $Y=AXA^{-1}$ or $Y=AX'A^{-1}$.
\end{quote}

Let $e_{\alpha\beta}=1$ or 0 according as $\alpha=\beta$ or not. Then $|X-rE|=|Y-rE|$ are the characteristic functions of $X$ and $Y$. If one puts $x_{\alpha\beta}=e_{\alpha\beta}$, let $y_{\alpha\beta}=c_{\alpha\beta}$ result. If one then replaces every $x_{\alpha\beta}$ in the equation $|Y|=|X|$ by $x_{\alpha\beta}-re_{\alpha\beta}$, one obtains
\[
|Y-rC|=|X-rE|=|Y-rE|.
\]
Comparing on both sides the coefficients of the first powers of $r$, one finds
\[
\sum c_{\alpha\beta}Y_{\alpha\beta}=\sum e_{\alpha\beta}Y_{\alpha\beta},
\]
where $Y_{\alpha\beta}$ is the minor of degree $(n-1)$ of $Y$ complementary to the element $y_{\alpha\beta}$. Since there is no linear relation between these minors, $c_{\alpha\beta}=e_{\alpha\beta}$. Now, by Theorem I, the matrices $A$ and $B$ can be so chosen that $Y=AXB$ or $Y=AX'B$. If here one sets $X=E$, then, as just shown, $Y=E$ also, and hence $AB=E$, so that $B=A^{-1}$.

The theorems just developed remain valid if the variability of the $n^2$ variables $x_{\alpha\beta}$ is restricted by the relations $x_{\alpha\beta}=x_{\beta\alpha}$:

\begin{quote}\itshape
III. If in a symmetric matrix $X$ the elements $x_{\alpha\beta}$ $(\beta\geqq\alpha)$ are independent variables, and if the elements of the symmetric matrix $Y$ are linear functions of these variables, and if the determinant $|Y|$ differs from $|X|$ only by a non-zero constant factor, then $Y=AXA'$, where $A$ is a constant matrix completely determined up to sign. If, in addition, the characteristic functions of $X$ and $Y$ are equal to one another, then $A$ is an orthogonal matrix.
\end{quote}

First suppose $AXA'=BXB'$. By setting $X=E$ one obtains $AA'=BB'$. Putting $A^{-1}B=A'B'^{-1}=F$, one has $XF=FX$ and $FF'=E$. If one sets the variables $x_{\alpha\beta}$ whose indices are different equal to zero, then $x_{\alpha\alpha}f_{\alpha\beta}=f_{\beta\alpha}x_{\beta\beta}$, hence $f_{\alpha\beta}=0$ when $\alpha$ is different from $\beta$. From $FF'=E$ it follows that $f_{\alpha\alpha}=\pm1$. The general equation $XF=FX$ therefore gives $x_{\alpha\beta}f_{\beta\beta}=f_{\alpha\alpha}x_{\alpha\beta}$. Hence $F=\pm E$, and $B=\pm A$, and consequently $A$ is completely determined up to sign.

Between the first derivatives of $|X|$ with respect to the variables $x_{\alpha\beta}$ $(\beta\geqq\alpha)$ there is no relation, because conversely the elements $x_{\alpha\beta}$ can be represented as functions of these derivatives. Between the minors of degree $m$, $u,v,w,\ldots$, there are indeed linear relations (\textsc{Kronecker}, \emph{Über die Subdeterminante symmetrischer Systeme}, Sitzungsberichte 1882). If $au+bv+cw+\cdots=0$ is such a relation, and if $u$ is a principal minor, then $a$ must be zero. For if one sets all variables $x_{\alpha\beta}=0$ except those which occur in $u$, then $v=w=\cdots=0$.

If now $c_{\alpha\beta}^{(\kappa)}$ is the coefficient of $x_{\kappa\kappa}$ in $y_{\alpha\beta}$, it follows from this, in the same way as above, that in the symmetric matrix $c_{\alpha\beta}^{(\kappa)}$ all principal minors of the second and third degrees vanish. As I have shown in my work \emph{Über das Trägheitsgesetz der quadratischen Formen}, § 2, Theorem 2 (Sitzungsberichte 1894), all minors of the second degree consequently vanish, and hence $c_{\alpha\beta}^{(\kappa)}=p_{\alpha\kappa}p_{\beta\kappa}$. The same conclusions as above now follow. Only, since $z_{\alpha\beta}^2=x_{\alpha\beta}^2$, one may put the quantities $k_\alpha=\pm1$. Then $R^{-1}=R=R'$ and $Y=PRXRP'=AXA'$.

If the characteristic functions of $X$ and $Y$ are equal to one another, then, as above, one shows that for $X=E$ one also has $Y=E$, and consequently $AA'=E$, so that $A$ is an orthogonal matrix.

\begin{center}
\rule{40mm}{0.4pt}

Issued on 25 November.

\rule{40mm}{0.4pt}
\end{center}

\begin{flushright}
Berlin, printed in the Reichsdruckerei.
\end{flushright}
% ===== END WORK 056 =====

\clearpage

% ===== BEGIN WORK 057: 057_ber_Relationen_zwischen_den_Charakteren_e =====
\begin{center}
{\Large On Relations between the Characters of a Group and Those of Its Subgroups}\\[0.75em]
By G. \textsc{Frobenius}.
\end{center}

In my paper \emph{On group characters} (Sitzungsberichte 1896) I developed, for the computation of the characters of a finite group of known constitution, a general method, and I demonstrated its practical usability on a series of simple examples. Since, however, its application to more complicated groups is bound up with considerable difficulties, I have sought other ways of obtaining the characters of a group and thereby its primitive representations by linear substitutions, and I have found two quite different methods which in special cases can lead more easily to this goal than that general method.

The first, which I shall set out here, is based on the consideration of the groups $\Ggr$ contained in the given group $\Hgr$ and on the relations that exist between the characters of $\Ggr$ and $\Hgr$. These relations arise in two different ways. One (§ 1) leads from the prime factors of the determinant of the group $\Hgr$ to those of the determinant of the group $\Ggr$; the other (§ 3) leads conversely from the latter to the former. Particularly simple results are obtained from these considerations in the case where $\Ggr$ is an invariant subgroup of $\Hgr$ (§§ 2, 4). The formulae obtained stand in close relation to the decomposition of the group $\Hgr$ into complexes of elements which are equivalent modulo a double modulus formed from two groups $\Ggr$ and $\Ggr'$ contained in $\Hgr$. The investigation of the special case where $\Ggr'=\Ggr$ leads directly to the determination of a character of every doubly transitive group of permutations (§ 5).

A second method for computing the characters of a group arises from the theory of the composition of characters, which I shall develop on another occasion.

\section*{§ 1.}

Let $\Hgr$ be a group of order $h$, and let
\[
  \Theta=\prod_\lambda \Phi_\lambda^{f_\lambda}
  \qquad(\lambda=0,1,\ldots,l-1)
\]
be its group determinant decomposed into prime factors. Let $\Ggr$ be a group contained in $\Hgr$ of order $g=\frac{h}{n}$, and let
\[
  \mathrm H=\prod_\kappa \Psi_\kappa^{e_\kappa}
  \qquad(\kappa=0,1,\ldots,k-1)
\]
be its group determinant. If in $\Theta$ all variables $x_R=0$ are set except those whose indices are the elements of $\Ggr$, then one obtains (\emph{Group characters} § 7)
\[
\tag{1.}
  \Theta=\mathrm H^n,
\]
and consequently each prime factor $\Phi_\lambda$ of $\Theta$ also becomes a product of prime factors of $\mathrm H$,
\[
\tag{2.}
  \Phi_\lambda=\prod_\kappa \Psi_\kappa^{r_{\kappa\lambda}}.
\]
Here $r_{\kappa\lambda}=0$ is to be set when $\Psi_\kappa$ does not divide $\Phi_\lambda$. If in this equation $x_E$ is replaced by $x_E+u$, comparison of the coefficients of $u^{f_\lambda-1}$ gives
\[
\tag{3.}
  \sum_\kappa r_{\kappa\lambda}\psi^{(\kappa)}(P)=\chi^{(\lambda)}(P),
\]
where $\psi^{(\kappa)}$ is the character of $\Psi_\kappa$ and $\chi^{(\lambda)}$ that of $\Phi_\lambda$, and where $P$ denotes an element of $\Ggr$. With the aid of the equations
\[
  \sum_P \psi^{(\kappa)}(P^{-1})\psi^{(\kappa)}(P)=g,
  \qquad
  \sum_P \psi^{(\kappa)}(P^{-1})\psi^{(\mu)}(P)=0
\]
it follows that
\[
\tag{4.}
  \sum_P \psi^{(\kappa)}(P^{-1})\chi^{(\lambda)}(P)=g r_{\kappa\lambda}.
\]
If therefore $R$ is an element of $\Hgr$, then
\[
  g\sum_\lambda r_{\kappa\lambda}\chi^{(\lambda)}(R)
  =
  \sum_P\psi^{(\kappa)}(P)
  \left(\sum_\lambda \chi^{(\lambda)}(P^{-1})\chi^{(\lambda)}(R)\right).
\]
But by \emph{Group characters} § 3, (4.) one has
\[
  \sum_\lambda \chi^{(\lambda)}(P^{-1})\chi^{(\lambda)}(R)=0,
\]
except when $P$ is conjugate to $R$ with respect to $\Hgr$. In that case, however, the sum is equal to $\frac{h}{h_R}=\frac{h}{h_\rho}$, if $R$ and $P$ are elements of the $\rho$-th class of $\Hgr$. Consequently
\[
\tag{5.}
  \sum_\lambda r_{\kappa\lambda}\chi^{(\lambda)}(R)
  =
  \frac{h}{g h_\rho}\sum_{(\rho)}\psi^{(\kappa)}(P),
\]
where $P$ runs through the elements of the $\rho$-th class that are contained in $\Ggr$. If one sets $\psi(R)=0$ when $R$ is not contained in $\Ggr$, then for $P$ one may also take all elements of $\Hgr$ conjugate to $R$. Under this convention special care must be taken that equation (3.) holds only for the elements $P$ that belong to the group $\Ggr$.

If the $k$ characters $\psi^{(\kappa)}$ of a subgroup $\Ggr$ are known, then the new property of the $l$ characters $\chi^{(\lambda)}$, expressed by equations (3.), (4.) and (5.) in three different but equivalent forms, consists in the fact that
\[
\tag{4a.}
  \frac{1}{g}\sum_P \psi(P^{-1})\chi(P)=r
\]
is a positive integer.

For a fixed value of $\lambda$ the $k$ numbers $r_{\kappa\lambda}$ are not all zero, and for a fixed value of $\kappa$ the $l$ numbers $r_{\kappa\lambda}$ are not all zero. This follows from equation (2.) together with relation (1.).

If $(0)$ is the principal class, then by putting $P=R=E$ in (3.) and (5.) one obtains the equations
\[
\tag{6.}
  \sum_\kappa r_{\kappa\lambda}e_\kappa=f_\lambda,
  \qquad
  \sum_\lambda r_{\kappa\lambda}f_\lambda=\frac{h}{g}e_\kappa,
\]
which also follow immediately from (2.) and (1.). Hence
\[
\tag{7.}
  e_\kappa r_{\kappa\lambda}\le f_\lambda,
  \qquad
  r_{\kappa\lambda}\le f_\lambda.
\]
If one chooses the index $0$ for the principal character, then
\[
\tag{8.}
  r_{00}=1,\qquad r_{\kappa 0}=0,\qquad(\kappa>0),
\]
and, if one puts $r_{0\lambda}=r_\lambda$,
\[
\tag{9.}
  \sum_\rho g_\rho\chi_\rho^{(\lambda)}=g r_\lambda,
  \qquad
  \sum_\lambda r_\lambda\chi_\rho^{(\lambda)}
  =
  \frac{h g_\rho}{g h_\rho},
\]
where $g_\rho$ denotes the number of elements in the $\rho$-th class of $\Hgr$ which belong to the group $\Ggr$.

Since the values of the characters are algebraic integers, it follows from the last formula that $h g_\rho$ is divisible by $g h_\rho$. This theorem can easily be proved directly. Let $H$ run through the $h$ elements of $\Hgr$, and let $R$ be a fixed element of the $\rho$-th class. Then the $h$ elements $H^{-1}RH$ are the $h_\rho$ different elements of the $\rho$-th class, each counted $\frac{h}{h_\rho}$ times. In $\Hgr$ there are therefore $g_\rho\frac{h}{h_\rho}$ different elements $H$ such that $H^{-1}RH$ is contained in $\Ggr$. Let $\mathfrak R$ be the complex of these elements $H$. If $A^{-1}RA$ is contained in $\Ggr$, and if $G$ is an element of $\Ggr$, then also
\[
  G^{-1}(A^{-1}RA)G=(AG)^{-1}R(AG)
\]
is contained in $\Ggr$. Thus if $A$ is an element of $\mathfrak R$, all elements of the complex $A\Ggr$ also belong to the complex $\mathfrak R$. Consequently $\mathfrak R$ decomposes into a number of complexes $A\Ggr+B\Ggr+C\Ggr+\cdots$, any two of which are disjoint. Hence the order $\frac{h g_\rho}{h_\rho}$ of $\mathfrak R$ is divisible by $g$. The number $\frac{h g_\rho}{g h_\rho}$ indicates how many among the $n$ complexes
\[
\tag{10.}
  A_0\Ggr+A_1\Ggr+\cdots+A_{n-1}\Ggr=\Hgr
\]
satisfy the condition $A_\nu\Ggr=R A_\nu\Ggr$.

By formula (3.) one has
\[
  \sum_P \chi^{(\lambda)}(P^{-1})\chi^{(\mu)}(P)
  =
  \sum_{\alpha,\beta} r_{\alpha\lambda}r_{\beta\mu}
  \left(\sum_P\psi^{(\alpha)}(P^{-1})\psi^{(\beta)}(P)\right),
\]
and consequently
\[
\tag{11.}
  \sum_P \chi^{(\lambda)}(P^{-1})\chi^{(\mu)}(P)
  =
  g\sum_\kappa r_{\kappa\lambda}r_{\kappa\mu},
\]
where $P$ runs only through the elements of $\Ggr$, or
\[
\tag{12.}
  \sum_\rho g_\rho\chi_{\rho'}^{(\lambda)}\chi_{\rho}^{(\mu)}
  =
  g\sum_\kappa r_{\kappa\lambda}r_{\kappa\mu}.
\]
But if $R$ runs through all elements of $\Hgr$, then
\[
  \sum_R \chi^{(\lambda)}(R^{-1})\chi^{(\lambda)}(R)=h.
\]
Since $\chi(R)$ and $\chi(R^{-1})$ are conjugate complex quantities, their product is a real positive quantity. Hence
\[
\tag{13.}
  \sum_\kappa r_{\kappa\lambda}^{2}\le \frac{h}{g}.
\]
If, however, in (11.) one puts $\mu=\lambda$ and then sums with respect to $\lambda$, one obtains
\[
\tag{14.}
  \sum_{\kappa,\lambda}r_{\kappa\lambda}^{2}
  =
  \sum_\lambda \frac{h g_\lambda}{g h_\lambda},
\]
where on the right the summation is over the $l$ classes of $\Hgr$. Multiplied by $g$, this number is equal to the number of solutions of the equation $QR=RQ$, when $Q$ runs through the $g$ elements of $\Ggr$ and $R$ through the $h$ elements of $\Hgr$.

Let $\Ggr$ and $\Ggr'$ be two subgroups of $\Hgr$, and let the numbers denoted for the group $\Ggr$ by $g,g_\rho,r_\lambda$ be denoted for the group $\Ggr'$ by $g',g'_\rho,r'_\lambda$. Then by (2.)
\[
  gg'\sum_\lambda r_\lambda r'_\lambda
  =
  \sum_{\rho,\rho'}g_\rho g'_{\rho'}
  \left(\sum_\lambda \chi_\rho^{(\lambda)}\chi_{\rho'}^{(\lambda)}\right).
\]
The last sum, taken with respect to $\lambda$, is zero unless $(\rho)=(\rho')$, but in that case it is equal to $\frac{h}{h_\rho}$. Hence
\[
  gg'\sum_\lambda r_\lambda r'_\lambda
  =
  h\sum_\rho\frac{g_\rho g'_\rho}{h_\rho}.
\]
In my paper \emph{On the congruence modulo a double modulus formed from two finite groups}, \textsc{Crelle}'s Journal vol. 101, § 2, (8.), I established for the number $(\Hgr:\Ggr,\Ggr')$ of classes into which the elements of $\Hgr$ decompose modulo the double modulus $(\Ggr,\Ggr')$ the relation
\[
\tag{15.}
  \frac{gg'}{h}(\Hgr:\Ggr,\Ggr')
  =
  \sum_\rho\frac{g_\rho g'_\rho}{h_\rho}.
\]
From this follows the formula
\[
\tag{16.}
  \sum_\lambda r_\lambda r'_\lambda=(\Hgr:\Ggr,\Ggr'),
  \qquad
  \sum_\lambda r_\lambda^{2}=(\Hgr:\Ggr,\Ggr).
\]

If all variables $x_R=0$ whose index $R$ is not contained in $\Ggr$ are set equal to zero, then $\Phi_\lambda$ contains the factor $\Psi_0^{r_\lambda}$. If $\Ggr'$ is a subgroup of $\Ggr$, one can insert those values into the expression (2.) for $\Phi_\lambda$; then $\Psi_0^{r_\lambda}$ passes into $\Psi_0^{\prime\,r'_\lambda}$. Consequently
\[
\tag{17.}
  r'_\lambda\ge r_\lambda,\quad\hbox{when }\Ggr'<\Ggr
\]
that is, when $\Ggr'$ is contained in $\Ggr$.

If $g<h$, then there are always classes in $\Hgr$ of which no element is contained in $\Ggr$. For under the assumption (10.) the groups $A_\nu\Ggr A_\nu^{-1}=\Ggr_\nu$ are the $n$ groups conjugate with $\Ggr$, which need not be distinct. Let $t_\mu$ be the number of elements of $\Hgr$ that are contained in exactly $\mu$ of these $n$ groups. Then
\[
  t_0+t_1+t_2+\cdots+t_n=h.
\]
On the other hand, the complex $\Ggr_0+\Ggr_1+\cdots+\Ggr_{n-1}$ contains $gn=h$ elements if elements occurring more than once are also counted more than once; thus each of the $t_\mu$ elements which occurs in exactly $\mu$ of those $n$ groups is counted $\mu$ times. Hence
\[
  t_1+2t_2+\cdots+nt_n=h
\]
and consequently
\[
  t_0=t_2+2t_3+\cdots+(n-1)t_n.
\]
Since the principal element $E$ occurs in all $n$ groups, $t_n>0$. Thus if $n>1$, then $t_0>0$. If $R$ is one of these $t_0$ elements and $H$ runs through the $h$ elements of $\Hgr$, then $R$ occurs in none of the groups $H\Ggr H^{-1}$, and hence none of the elements $H^{-1}RH$ conjugate with $R$ is contained in $\Ggr$. By equation (5.), therefore,
\[
  \sum_\lambda r_{\kappa\lambda}\chi^{(\lambda)}(R)=0.
\]
Thus if $k\ge l$, then in the matrix
\[
\tag{18.}
  r_{\kappa\lambda}\qquad(\kappa=0,1,\ldots,k-1;\ \lambda=0,1,\ldots,l-1)
\]
all determinants of degree $l$ vanish.

\section*{§ 2.}

The $kl$ numbers $r_{\kappa\lambda}$ can be determined more closely when $\Ggr$ is an invariant subgroup of $\Hgr$. For this case I developed the following theorems in \emph{Group characters} § 7. Let $S$ be a fixed element of $\Hgr$, and let $P$ be a variable element of $\Ggr$. Since $S^{-1}\Ggr S=\Ggr$, $S^{-1}PS$ runs at the same time with $P$ through the $g$ elements of $\Ggr$. If one puts $\psi^{(\kappa)}(S^{-1}PS)=\psi^{(\kappa')}(P)$, then $\psi^{(\kappa')}(P)$ is a character of $\Ggr$. If it is different from $\psi^{(\kappa)}(P)$, I call it conjugate to $\psi^{(\kappa)}(P)$ with respect to $\Hgr$. Conjugate characters have the same degree $e_\kappa=e_{\kappa'}$ and differ only by the ordering of their values. Let $s_\kappa=s_{\kappa'}=\cdots$ be the number of different characters of $\Ggr$ which are conjugate with $\psi^{(\kappa)}$ and therefore also with one another. To each character $\psi^{(\kappa)}$ of $\Ggr$ there correspond one or more characters $\chi^{(\lambda)},\chi^{(\lambda')},\ldots$ of $\Hgr$ such that
\[
\tag{1.}
  \psi^{(\kappa)}(P)+\psi^{(\kappa')}(P)+\cdots
  =
  \frac{s_\kappa e_\kappa}{f_\lambda}\chi^{(\lambda)}(P)
  =
  \frac{s_\kappa e_\kappa}{f_{\lambda'}}\chi^{(\lambda')}(P)
  =\cdots
\]
for all $g$ elements $P$ of $\Ggr$. Conversely, to each character $\chi^{(\lambda)}$ of $\Hgr$ there corresponds at least one character $\psi^{(\kappa)}$ of $\Ggr$ such that equation (1.) is fulfilled; and if several correspond to it, then any two of them are conjugate. If $\chi^{(\lambda)}$ is given, then $\chi^{(\lambda')},\chi^{(\lambda'')},\ldots$ are determined by the condition that the $g$ values $\chi^{(\lambda)}(P)$ are proportional to the $g$ values $\chi^{(\lambda')}(P)$. If $R$ runs through the $h$ elements of $\Hgr$, and if $\chi^{(\lambda)}$ and $\chi^{(\mu)}$ are different characters, then by \emph{Group characters} § 3 the $h$ values $\chi^{(\lambda)}(R)$ cannot be proportional to the $h$ values $\chi^{(\mu)}(R)$.

If now one inserts the expression (1.) for $\chi^{(\lambda)}(P)$ in formula (3.), § 1, one sees that always $r_{\kappa\lambda}=0$ unless $\psi^{(\kappa)}$ and $\chi^{(\lambda)}$ correspond. In that case, however,
\[
\tag{2.}
  r_{\kappa\lambda}=r_{\kappa'\lambda}=\cdots=\frac{f_\lambda}{s_\kappa e_\kappa}
\]
is different from zero, where $\psi^{(\kappa)},\psi^{(\kappa')},\ldots$ are the $s_\kappa$ mutually conjugate characters of $\Ggr$ corresponding to $\chi^{(\lambda)}$. Thus equation (2.), § 1 reduces to
\[
\tag{3.}
  \Phi_\lambda=(\Psi_\kappa\Psi_{\kappa'}\cdots)^{\frac{f_\lambda}{s_\kappa e_\kappa}},
\]
and from (6.), § 1 it follows that
\[
\tag{4.}
  f_\lambda^2+f_{\lambda'}^2+\cdots=ns_\kappa e_\kappa^2,
  \qquad
  r_{\kappa\lambda}^2+r_{\kappa\lambda'}^2+\cdots=\frac{n}{s_\kappa}.
\]
Consequently $f_\lambda$ is divisible by $s_\kappa e_\kappa$, and the quotient $\frac{f_\lambda}{s_\kappa e_\kappa}$ is a common divisor of the $g$ algebraic integers $\chi^{(\lambda)}(P)$. That $n$ is divisible by $s_\kappa$ I have already shown in \emph{Group characters} § 7.

If $R$ is not contained in $\Ggr$, then no element $P$ conjugate with $R$ belongs to the group $\Ggr$. Thus from (5.), § 1 it follows that
\[
\tag{5.}
  f_\lambda\chi^{(\lambda)}(R)+f_{\lambda'}\chi^{(\lambda')}(R)+\cdots=0,
\]
when $R$ is not contained in $\Ggr$.

\section*{§ 3.}

We have passed from the $l$ prime factors $\Phi$ of the determinant $\Theta$ of the group $\Hgr$ to the $k$ prime factors $\Psi$ of the determinant $\mathrm H$ of the subgroup $\Ggr$ by setting all variables $x_R=0$ whose index $R$ is an element of $\Hgr$ not contained in $\Ggr$. I shall now show how, by another construction, one can ascend from the $k$ functions $\Psi$ to the $l$ functions $\Phi$.

Let $X$ be a matrix of degree $e$ belonging to the group $\Ggr$. Its elements are linear functions of the $g$ variables $x_P$, and its determinant does not vanish identically. It is characterized by the following property. If $x_P$ is replaced by $y_P$ or $z_P$, let $X$ pass into $Y$ or $Z$. If then
\[
  z_Q=\sum_P x_{P^{-1}}y_{PQ},
\]
then $Z=XY$. If in $X$ each of the $g$ variables $x_P$ is replaced by $x_{APB^{-1}}$, where $A$ and $B$ are two elements of $\Hgr$, one obtains a matrix that I denote by $X_{A,B}$. Let the $n$ elements $A_0,A_1,\ldots,A_{n-1}$ form a complete residue system of $\Hgr(\mathrm{mod.}\,\Ggr)$, so that
\[
\tag{1.}
  \Hgr=A_0\Ggr+A_1\Ggr+\cdots+A_{n-1}\Ggr
  =
  \Ggr A_0^{-1}+\Ggr A_1^{-1}+\cdots+\Ggr A_{n-1}^{-1}.
\]
I then consider the $n^2$ matrices of degree $e$ which are obtained from $X_{A,B}$ by putting for $A$ and $B$ each of the $n$ elements $A_0,A_1,\ldots,A_{n-1}$, and from them I form a matrix $(X_{A,B})$ of degree $ne$. If $x_R$ is replaced by $y_R$ or $z_R$, let $X_{A,B}$ pass into $Y_{A,B}$ or $Z_{A,B}$. If $A,B$ and $N$ are elements of $\Hgr$, then the matrix $X_{A,N}Y_{N,B}$ arises from $X$ by replacing $x_Q$ by
\[
  \sum_P x_{AP^{-1}N^{-1}}y_{NPQB^{-1}}.
\]
Here $P$ and $Q$, as above, are elements of $\Ggr$. Now the elements of $X$ are linear functions of the $g$ variables $x_Q$. Therefore the elements of the matrix
\[
  \sum_N X_{A,N}Y_{N,B}
  \qquad(N=A_0,A_1,\ldots,A_{n-1})
\]
arise from $X$ by replacing $x_Q$ by
\[
  \sum_{N,P}x_{AP^{-1}N^{-1}}y_{NPQB^{-1}}
  =
  \sum_R x_{AR^{-1}}y_{RQB^{-1}}.
\]
As $P$ runs through the $g$ elements of $\Ggr$ and $N$ through the $n$ elements $A_0,A_1,\ldots,A_{n-1}$, $R=NP$ runs, by (1.), through the $gn=h$ different elements of $\Hgr$. Thus if now, when $R$ and $S$ are elements of $\Hgr$, one puts
\[
  z_S=\sum_R x_{R^{-1}}y_{RS},
\]
then the last sum is equal to $z_{AQB^{-1}}$, and hence
\[
  \sum_N X_{A,N}Y_{N,B}=Z_{A,B}
\]
or
\[
  (X_{A,B})(Y_{A,B})=(Z_{A,B}).
\]
Consequently $(X_{A,B})$ is a matrix of degree $ne$ belonging to the group $\Hgr$, and its determinant is therefore a product of prime factors of the group determinant $\Theta$,
\[
\tag{2.}
  |(X_{A,B})|=\prod_\lambda \Phi_\lambda^{r_\lambda}.
\]

In every matrix belonging to a group, $x_E$ is multiplied by the identity matrix. Thus if $x_E$ is replaced there by $x_E+u$, only the term $u$ is added to each element of the diagonal. Therefore the coefficient of $u^{ne-1}$ in the determinant of degree $ne$ in (2.) is equal to the sum of the diagonal elements.

If $\Psi(x)$ is a prime factor of degree $e$ of $\mathrm H$, then one can find a matrix $X$ belonging to $\Ggr$ whose determinant is equal to $\Psi(x)$ (\emph{On the representation of finite groups by linear substitutions}, Sitzungsberichte 1897). The coefficient of $u^{e-1}$ in $\Psi(x+u\varepsilon)$ is $\sum_P\psi(P)x_P$. For the matrix $X_{N,N}$, the sum of the diagonal elements is therefore equal to
\[
  \sum_P\psi(P)x_{NPN^{-1}}
  =
  \sum_R\psi(N^{-1}RN)x_R.
\]
Here $P$ and $R=NPN^{-1}$ may run through all elements of $\Hgr$, if as above one sets $\psi(R)=0$ when $R$ does not belong to the group $\Ggr$. Comparing the coefficients of $u^{ne-1}$ in equation (2.) gives, if $r_\lambda=r_{\kappa\lambda}$ is put for $\Psi=\Psi_\kappa$,
\[
\tag{3.}
  \sum_N\psi^{(\kappa)}(N^{-1}RN)
  =
  \sum_\lambda r_{\kappa\lambda}\chi^{(\lambda)}(R).
\]
If $P$ is an element of $\Ggr$, then $\psi(P^{-1}SP)=\psi(S)$. For if $S$ is an element of $\Ggr$, this is the equation \emph{Group characters} § 5, (2.). But if $S$ is not contained in $\Ggr$, then $P^{-1}SP$ also does not belong to the group $\Ggr$, and both sides of the equation are zero. Therefore
\[
  g\sum_N\psi(N^{-1}RN)
  =
  \sum_{N,P}\psi(P^{-1}N^{-1}RNP)
  =
  \sum_S\psi(S^{-1}RS),
\]
where $S=NP$ runs through the $h$ elements of $\Hgr$, so that
\[
\tag{4.}
  g\sum_\lambda r_{\kappa\lambda}\chi^{(\lambda)}(R)
  =
  \sum_S\psi^{(\kappa)}(S^{-1}RS).
\]
If $R$ is an element of the $\rho$-th class in $\Hgr$, then $S^{-1}RS$ represents each of the $h_\rho$ different elements of this class $\frac{h}{h_\rho}$ times. Hence
\[
\tag{5.}
  \sum_\lambda r_{\kappa\lambda}\chi^{(\lambda)}(R)
  =
  \frac{h}{g h_\rho}\sum_{(\rho)}\psi^{(\kappa)}(P),
\]
where $P$ runs through the $h_\rho$ elements of the $\rho$-th class, or also only through those among them that are contained in $\Ggr$. Since the $kl$ numbers $r_{\kappa\lambda}$ are completely determined by these equations, they are identical with the exponents occurring in the equations
\[
\tag{6.}
  \Phi_\lambda=\prod_\kappa\Psi_\kappa^{r_{\kappa\lambda}}.
\]
In consequence of the inequalities (7.) § 1, the $e_\kappa$-th power of the expression
\[
\tag{7.}
  |(X^{(\kappa)}_{A,B})|=\prod_\lambda\Phi_\lambda^{r_{\kappa\lambda}}
\]
is a divisor of the group determinant $\Theta$.

If $\Ggr$ is an invariant subgroup of $\Hgr$, then, with the above notation,
\[
\tag{8.}
  |(X^{(\kappa)}_{A,B})|^{s_\kappa e_\kappa}
  =
  \Phi_\lambda^{f_\lambda}\Phi_{\lambda'}^{f_{\lambda'}}\cdots,
\]
and therefore it is a divisor of $\Theta$ relatively prime to the complementary divisor.

\section*{§ 4.}

In the particularly noteworthy case $\kappa=0$, for the simplification of the notation, if
\[
  \mathfrak A=P+Q+R+\cdots
\]
is a complex of elements, I put
\[
\tag{1.}
  x_{\mathfrak A}=x_P+x_Q+x_R+\cdots.
\]
Then the matrix of degree $n$
\[
\tag{2.}
  (x_{A\Ggr B^{-1}})\qquad(A,B=A_0,A_1,\ldots,A_{n-1})
\]
is a matrix belonging to $\Hgr$, and its determinant is
\[
\tag{3.}
  |x_{A\Ggr B^{-1}}|=\prod_\lambda \Phi_\lambda^{r_\lambda}.
\]
To every matrix belonging to $\Hgr$ there corresponds a representation of the group $\Hgr$, or of a group meroedrically isomorphic to $\Hgr$, by linear substitutions. To the matrix (2.) there corresponds the representation of a group isomorphic to $\Hgr$ by permutations of $n$ symbols, which I developed in my paper \emph{On finite groups} § 4 (Sitzungsberichte 1895).

If $\Ggr$ is an invariant subgroup of $\Hgr$, then, since $e_0=s_0=1$,
\[
\tag{4.}
  |x_{AB^{-1}\Ggr}|=\prod_\lambda\Phi_\lambda^{r_\lambda}
  =
  \Phi_\nu^{f_\nu}\Phi_{\nu'}^{f_{\nu'}}\cdots,
\]
where $\chi^{(\nu)},\chi^{(\nu')},\ldots$ are the characters of $\Hgr$ corresponding to the character $\psi^{(0)}$ of $\Ggr$. The left hand side is the determinant of the group $\frac{\Hgr}{\Ggr}$, which I considered in \emph{Representation} § 1. There too I showed that the non-zero exponents $r_\lambda$ are equal to $f_\lambda$. From that equation there follows a further noteworthy consequence: the indices $\lambda=\nu,\nu',\ldots$ are characterized by the fact that $\chi^{(\lambda)}(P)$ has the same value for all $g$ elements $P$ of the group $\Ggr$. But since the left hand side contains the $h$ quantities $x_R$ only in the $n$ linear combinations $x_{\mathfrak A}$, the same holds for each of the prime functions $\Phi_\nu,\Phi_{\nu'},\ldots$; and consequently, if $N$ is a fixed element of $\Hgr$ and $P$ is a variable element of $\Ggr$, $\chi^{(\lambda)}(NP)$ also has the same value for all elements $P$ of $\Ggr$. Thus one obtains the theorem:

\emph{I. In order that a character of $\Hgr$ belong to the group $\frac{\Hgr}{\Ggr}$, it is necessary and sufficient that it have the same value for all elements of $\Ggr$. Then it also has equal values for every two elements of $\Hgr$ which are equivalent modulo $\Ggr$.}

This theorem can also be derived from the formula (\emph{Group characters} § 5)
\[
\tag{5.}
  h_\beta\chi(A)\chi(B)=f\sum_{(S)}\chi(AS),
\]
where $S$ runs through the $h_\beta$ elements conjugate with $B$. Since therefore $\chi(S)=\chi(B)$, one can also write this equation in the form
\[
  \chi(R)\sum_{(\Ggr)}\chi(S)=f\sum_{(\Ggr)}\chi(RS).
\]
If $\Ggr$ is an invariant subgroup of $\Hgr$, it contains either no element of the $\rho$-th class or all of them. If for $(\Ggr)$ one successively takes all classes whose elements are contained in $\Ggr$, then by summing the corresponding equations one finds
\[
  \chi(R)\sum\chi(P)=f\sum\chi(P'),
\]
where $P$ runs through the elements of the group $\Ggr$, and $P'$ through those of the complex $R\Ggr$. If therefore $\sum\chi(P)=0$, then also for every element $R$ of $\Hgr$ one has $\sum\chi(P')=0$. If further $S\sim R(\mathrm{mod.}\,\Ggr)$, then $R\Ggr=S\Ggr$, and consequently also
\[
  \chi(S)\sum\chi(P)=f\sum\chi(P'),
\]
where $P'$ runs through the same elements as above. Thus if $\sum\chi(P)$ is different from zero, it must be that
\[
\tag{6.}
  \chi(R)=\chi(S),\qquad \hbox{if } R\sim S(\mathrm{mod.}\,\Ggr).
\]
The last case certainly occurs when the $g$ values $\chi(P)$ are all equal to one another, hence all equal to $\chi(E)=f$. Then the character $\chi=\chi^{(\nu)}$ belongs to the group $\frac{\Hgr}{\Ggr}$; the function $\Phi_\nu$ contains the $h$ variables $x_R$ only in the $n$ linear combinations $x_{\mathfrak A}$ and, up to a numerical factor, becomes equal to the corresponding prime factor of the group $\frac{\Hgr}{\Ggr}$ if one sets in it
\[
\tag{7.}
  x_R=x_S,\qquad \hbox{if } R\sim S(\mathrm{mod.}\,\Ggr).
\]
If, however, the character $\chi=\chi^{(\lambda)}$ does not belong to the group $\frac{\Hgr}{\Ggr}$, then equation (6.) is not fulfilled for every two equivalent elements; therefore $\sum_P\chi(P)=0$, and consequently also, for every element $N$ of $\Hgr$,
\[
\tag{8.}
  \sum_P\chi(NP)=0.
\]

Theorem I can be partly reversed. From formula (5.), if $\chi(A)=\chi(B)=f$, one obtains the equation
\[
  h_\beta f=\sum\chi(AS).
\]
Since $\chi(R)$ is a sum of $f$ roots of unity, the right hand side is a sum of $h_\beta f$ roots of unity. Such a sum can be equal to $h_\beta f$ only when each term is equal to 1. Hence $\chi(AS)=f$. Thus if $\chi(A)=f$ and $\chi(B)=f$, then also $\chi(AB)=f$. Consequently all elements $R$ of $\Hgr$ for which $\chi(R)=f$ form a group $\Ggr$. If this group contains the element $B$, then it also contains all elements $S$ conjugate with $B$, since $\chi(S)=\chi(B)$. Therefore $\Ggr$ is an invariant subgroup of $\Hgr$.

\emph{II. If $\chi(R)$ is a character of degree $f$ of the group $\Hgr$, then all elements $R$ of $\Hgr$ for which $\chi(R)=f$ form an invariant subgroup $\Ggr$ of $\Hgr$, and the character $\chi$ belongs to the group $\frac{\Hgr}{\Ggr}$.}

With the aid of these two theorems one can, when the characters of a group are known, specify all its invariant subgroups. In the special case in which $\Hgr$ is a commutative group, this procedure was already given by \textsc{Weber}.

Among the $h$ linear functions
\[
  \frac{h}{f}\xi_R=\sum_S\chi(S)x_{RS}
\]
there are $f^2$ mutually independent ones, and through these $f^2$ combinations of the $h$ variables $x_R$ the function $\Phi$ can be expressed. If $P$ runs through the $g$ elements of $\Ggr$ and $N$ through a complete residue system of $\Hgr(\mathrm{mod.}\,\Ggr)$, then $S=NP$ runs through the $h$ elements of $\Hgr$. Therefore
\[
  \frac{h}{f}\xi_R=\sum_{N,P}\chi(NP)x_{RNP},
\]
and consequently $\xi_R=0$ if assumption (7.) is made. If $\Phi$ is represented in any way as a function of $f^2$ independent variables, then these are linear combinations of the variables $\xi_R$ and therefore have the same property. Thus the following theorem holds:

\emph{III. Let $\Phi$ be a prime factor of degree $f$ of the determinant of the group $\Hgr$, let $\Ggr$ be an invariant subgroup of $\Hgr$, and suppose that in $\Phi$ one always sets $x_R=x_S$ if $R\sim S(\mathrm{mod.}\,\Ggr)$. Then the function $\Phi$, if its character $\chi$ belongs to $\frac{\Hgr}{\Ggr}$, becomes equal to a prime factor of this group; but if $\chi$ does not belong to $\frac{\Hgr}{\Ggr}$, then each of the $f^2$ independent variables through which $\Phi$ can be represented vanishes.}

If now one brings the group matrix $X$ to the reduced form $L^{-1}XL$ (\emph{Representation}, § 5), then under assumption (7.) this passes into the reduced form of the matrix of the group $\frac{\Hgr}{\Ggr}$; the rank of $L^{-1}XL$ becomes $\frac{h}{g}=n$, and of its minors of degree $n$ only one is different from zero. Consequently the rank of the group matrix $X$ also becomes equal to $n$, and every minor of degree $n$ of $X$ becomes, up to a constant factor, equal to the determinant of the group $\frac{\Hgr}{\Ggr}$.

\section*{§ 5.}

Formula (5.), § 1 and the formula (9.), § 1 contained in it,
\[
\tag{1.}
  \sum_\lambda r_{\kappa\lambda}\chi^{(\lambda)}(R)
  =
  \frac{h}{g h_\rho}\sum_{(\rho)}\psi^{(\kappa)}(P),
  \qquad
  \sum_\lambda r_\lambda\chi_\rho^{(\lambda)}
  =
  \frac{h g_\rho}{g h_\rho},
\]
are especially suited to deriving characters of the group $\Hgr$ from the characters of the subgroup $\Ggr$. With the aid of the latter formula I have succeeded, as I shall explain on another occasion, in determining generally the characters of the symmetric group of degree $n$. A particularly simple application of this equation is given by the following theorem:

\emph{Let the group $\Hgr$ of order $h$ contain the group $\Ggr$ of order $g$; let the $\rho$-th class of conjugate elements in $\Hgr$ consist of $h_\rho$ elements, and let $g_\rho$ of them belong to the group $\Ggr$. Then the necessary and sufficient condition that the quantities}
\[
\tag{1.}
  \chi_\rho=\frac{h g_\rho}{g h_\rho}-1
\]
\emph{form a character of $\Hgr$ is that the number of classes into which the elements of $\Hgr$ decompose modulo the double modulus $(\Ggr,\Ggr)$ be equal to two.}

If $\chi_\rho$ is a character, then
\[
\tag{2.}
  \sum_\rho h_\rho\chi_\rho\chi_{\rho'}=h,
\]
and hence
\[
  \sum_\rho h_\rho\left(\frac{h g_\rho}{g h_\rho}-1\right)^2=h,
  \qquad
  \sum_\rho\frac{g_\rho^2}{h_\rho}=\frac{2g^2}{h}.
\]
By formula (15.), § 1, consequently
\[
\tag{3.}
  (\Hgr:\Ggr,\Ggr)=2.
\]
The principal element $E$ represents, modulo $\Ggr,\Ggr$, the complex $\Ggr E\Ggr=\Ggr$. If $L$ is an element of $\Hgr$ not contained in $\Ggr$, then the elements of $\Hgr$ equivalent to $L$ form the complex $\Ggr L\Ggr$. Since $(\Hgr:\Ggr,\Ggr)=2$, it follows that
\[
\tag{4.}
  \Hgr=\Ggr+\Ggr L\Ggr.
\]
If therefore one decomposes $\Hgr(\mathrm{mod.}\,\Ggr)$ into $n$ different complexes
\[
  \Hgr=\Ggr+P\Ggr+Q\Ggr+R\Ggr+\cdots,
\]
then there is in $\Ggr$ an element $G$ such that $G(P\Ggr)=Q\Ggr$.

Finally one can also express the condition found as follows. Let $\Dgr$ be the greatest common divisor of all subgroups of $\Hgr$ conjugate with $\Ggr$. Then $\frac{\Hgr}{\Dgr}$ can always be represented as a transitive group of permutations of $n$ symbols in such a way that the subgroup $\frac{\Ggr}{\Dgr}$ is formed by all permutations which leave a specified symbol fixed. The above condition now consists in the fact that this group of permutations is doubly transitive. If the permutation corresponding to the element $R$ is decomposed into its cyclic factors, then by (10.), § 1, $1+\chi(R)$ is equal to the number of cycles of first degree, or equal to the number of symbols which that permutation leaves unchanged.

That condition (3.) is also sufficient follows from formula (16.), § 1, according to which $\sum r_\lambda^2=2$. Since $r_0=1$, there is consequently one and only one number $r_\nu=1$, and each of the other $l-2$ numbers $r_\lambda=0$. If then $\chi_\rho^{(\nu)}=\chi_\rho$, it follows from (9.), § 1 that
\[
  \chi_\rho^{(0)}+\chi_\rho^{(\nu)}
  =
  \frac{h g_\rho}{g h_\rho},
  \qquad
  \chi_\rho=\frac{h g_\rho}{g h_\rho}-1.
\]

One can also prove directly that under condition (4.) the quantities $\chi_\rho$ satisfy the equations that serve for the computation of characters. Of the elements of the $\rho$-th class, $h_\rho-g_\rho$ are not contained in $\Ggr$. Of these $h_\rho-g_\rho$ elements, equally many occur in each of the $n-1$ complexes $P\Ggr,Q\Ggr,R\Ggr,\ldots$, namely $\frac{h_\rho-g_\rho}{n-1}$. For let
\[
  PA,\qquad PB,\qquad PC,\ldots
\]
be the elements of the $\rho$-th class which belong to the complex $P\Ggr$, so that $A,B,C,\ldots$ are elements of $\Ggr$. Now there is in $\Ggr$ an element $G$ such that $GP\Ggr=Q\Ggr$. Then the elements
\[
  GPAG^{-1},\ GPBG^{-1},\ GPCG^{-1},\ldots
\]
are mutually distinct and conjugate to the above elements, and hence are also contained in the $\rho$-th class. Finally they belong to the complex $Q\Ggr$. This complex therefore contains no fewer elements of the $\rho$-th class than does $P\Ggr$, and since the same holds conversely, each of the $n-1$ complexes $P\Ggr,Q\Ggr,R\Ggr,\ldots$ contains equally many elements of the $\rho$-th class.

Now let $(\alpha')$ and $(\beta)$ be any two classes, equal or different, and let $A$ be an element of $(\alpha')$, so that $A^{-1}$ is an element of $(\alpha)$, and let $B$ be an element of $(\beta)$. If $A$ and $B$ are both contained in $\Ggr$, then $A^{-1}B$ is also contained in $\Ggr$. If of these two elements one is contained in $\Ggr$ and the other is not, then $A^{-1}B$ is not contained in $\Ggr$. If both are not contained in $\Ggr$, then $A^{-1}B$ is contained in $\Ggr$ or not according as $A$ and $B$ both belong to the same complex $P\Ggr$ (or $Q\Ggr$, or $R\Ggr,\ldots$), or not. Therefore, as $A$ runs through the $h_\alpha$ elements of $(\alpha')$ and $B$ through the $h_\beta$ elements of $(\beta)$, of the $h_\alpha h_\beta$ elements $A^{-1}B$ the number
\[
  g_\alpha g_\beta+(n-1)\frac{h_\alpha-g_\alpha}{n-1}\frac{h_\beta-g_\beta}{n-1}
\]
is contained in $\Ggr$.

If $C$ is an element of the class $(\gamma)$, then $\frac{h_{\alpha\beta\gamma'}}{h_\gamma}$ of the $h_\alpha h_\beta$ elements $A^{-1}B$ are equal to $C$. Therefore of them $\frac{h_{\alpha\beta\gamma'}}{h_\gamma}g_\gamma$ are contained in the class $(\gamma)$ and at the same time in the group $\Ggr$, and consequently, of those $h_\alpha h_\beta$ elements,
\[
  \sum_\gamma \frac{h_{\alpha\beta\gamma'}}{h_\gamma}g_\gamma
  =
  g_\alpha g_\beta+\frac{1}{n-1}(h_\alpha-g_\alpha)(h_\beta-g_\beta)
\]
belong to the group $\Ggr$. With the aid of the relation
\[
  \sum_\gamma h_{\alpha\beta\gamma'}=h_\alpha h_\beta
\]
this equation is easily transformed into
\[
\tag{5.}
  h_\alpha h_\beta\chi_\alpha\chi_\beta
  =
  f\sum_\gamma h_{\alpha\beta\gamma'}\chi_\gamma,
\]
where $\chi_\rho$ is defined by equation (1.). From these relations, however, together with (2.), it follows that the quantities $\chi_\rho$ form a character of $\Hgr$.

For the group of degree $\frac12 p(p^2-1)$, which I investigated in \emph{Group characters} §§ 9, 10, formula (1.) gives the character of degree $p$. For it can be represented as a doubly transitive group of degree $p+1$.

\begin{center}
\small Issued on 21 July.\\[1.5em]
Berlin, printed at the Reichsdruckerei.
\end{center}
% ===== END WORK 057 =====

\clearpage

% ===== BEGIN WORK 058: 058_ber_die_Composition_der_Charaktere_einer =====
\begin{center}
{\Large On the Composition of the Characters of a Group}\\[0.75em]
By G. \textsc{Frobenius}.
\end{center}

In order to make easier the computation of the characters of a group, I have in my last paper (Sitzungsberichte 1898) derived relations which exist between the characters of a group and those of its subgroups. Another method, serving the same purpose, results from the theorem which I shall develop in the present paper. According to it, the product of two characters of a group can be represented as a linear combination of its characters, whose coefficients are positive integers. These coefficients, which I denote by $f_{\varkappa\lambda\mu}$, have properties similar to the numbers $h_{\alpha\beta\gamma}$ which I introduced in my paper \emph{On group characters} (Sitzungsberichte 1896). I have admittedly not succeeded in investigating the meaning of the numbers $f_{\varkappa\lambda\mu}$ for a given group. But already the certainty that relations of the indicated kind exist between the characters of a group permits, in many cases, new characters to be derived from one or more known characters.

\section*{§ 1.}

From two linear substitutions
\begin{equation*}\tag{a}
u_\alpha=\sum_\beta a_{\alpha\beta}v_\beta\qquad(\alpha,\beta=1,2,\cdots f)
\end{equation*}
and
\begin{equation*}\tag{a'}
u'_\gamma=\sum_\delta a'_{\gamma\delta}v'_\delta\qquad(\gamma,\delta=1,2,\cdots f')
\end{equation*}
one can derive a third,
\begin{equation*}\tag{A}
u_\alpha u'_\gamma=\sum_{\beta,\delta}a_{\alpha\beta}a'_{\gamma\delta}v_\beta v'_\delta,
\end{equation*}
by denoting the $ff'$ products $u_\alpha u'_\gamma$, in some order, by $U_\lambda$ $(\lambda=1,2,\cdots ff')$, and the products $v_\alpha v'_\gamma$, in the same order, by $V_\lambda$. If, following \textsc{Dedekind}, one calls the sum of the diagonal elements of a substitution or matrix its \emph{trace}, then the trace of (A) is equal to the product of the traces of (a) and (a'). In the same way, from the matrices (b) and (b') of degrees $f$ and $f'$ form the matrix (B) of degree $ff'$, and from (c) and (c') the matrix (C). If then $(c)=(a)(b)$ and $(c')=(a')(b')$, then also $(C)=(A)(B)$, as is immediately evident from the construction just given of (A) from (a) and (a').

Let $\Phi_\nu$ $(\nu=0,1,\cdots k-1)$ be the $k$ prime factors of the determinant of the group $\Hgr$, and let $\chi^{(\nu)}$ be their characters. To the elements $A,B,C,\cdots$ of $\Hgr$ let the matrices $(a),(b),(c),\cdots$ correspond in the primitive representation of $\Hgr$ by linear substitutions which belongs to $\Phi_\varkappa$; in the representation belonging to $\Phi_\lambda$, let the matrices $(a'),(b'),(c'),\cdots$ correspond. If then $AB=C$, then $(a)(b)=(c)$ and $(a')(b')=(c')$, and consequently also $(A)(B)=(C)$. It follows that $(A)x_A+(B)x_B+(C)x_C+\cdots$ is a matrix belonging to the group $\Hgr$, and its determinant is
\begin{equation*}\tag{1.}
\prod_\mu \Phi_{\mu'}^{\,f_{\varkappa\lambda\mu}},
\end{equation*}
where $f_{\varkappa\lambda\mu}$ is a positive integer or zero, and where by $\Phi_{\mu'}$ the complex prime function conjugate to $\Phi_\mu$ is to be denoted.

Since the trace of (a) is equal to $\chi^{(\varkappa)}(A)$, and that of (a') is equal to $\chi^{(\lambda)}(A)$, the trace of (A) is equal to $\chi^{(\varkappa)}(A)\chi^{(\lambda)}(A)$. From formula (1.), however, one obtains for this trace the expression
\[
\chi^{(\varkappa)}(A)\chi^{(\lambda)}(A)
 =\sum_\mu f_{\varkappa\lambda\mu}\chi^{(\mu')}(A)
 =\sum_\mu f_{\varkappa\lambda\mu}\chi^{(\mu)}(A^{-1}),
\]
or also, because $\mu'$ runs through the values $0,1,\cdots k-1$ simultaneously with $\mu$,
\begin{equation*}\tag{2.}
\chi^{(\varkappa)}(R)\chi^{(\lambda)}(R)
=\sum_\mu f_{\varkappa\lambda\mu'}\chi^{(\mu)}(R),
\qquad
\chi_\rho^{(\varkappa)}\chi_\rho^{(\lambda)}
=\sum_\mu f_{\varkappa\lambda\mu'}\chi_\rho^{(\mu)}.
\end{equation*}
This formula contains the rules according to which the \emph{composition} of characters is effected. If one sets
\begin{equation*}\tag{3.}
h f_{\varkappa\varkappa'}=\sum_R \chi^{(\varkappa)}(R)\chi^{(\varkappa')}(R)=h,
\qquad
h f_{\varkappa\lambda}=\sum_R \chi^{(\varkappa)}(R)\chi^{(\lambda)}(R)=0
\end{equation*}
(where $\lambda$ is different from $\varkappa'$), then with the aid of these relations one obtains
\begin{equation*}\tag{4.}
h f_{\varkappa\lambda\mu}=\sum_R \chi^{(\varkappa)}(R)\chi^{(\lambda)}(R)\chi^{(\mu)}(R),
\end{equation*}
or, if $\rho$ runs through the $k$ classes of conjugate elements,
\begin{equation*}\tag{5.}
h f_{\varkappa\lambda\mu}
=\sum_\rho h_\rho\chi_\rho^{(\varkappa)}\chi_\rho^{(\lambda)}\chi_\rho^{(\mu)}.
\end{equation*}
Thus $f_{\varkappa\lambda\mu}$ remains unchanged under all permutations of the indices, and since $R^{-1}$ runs through the $h$ elements of $\Hgr$ simultaneously with $R$, one has
\begin{equation*}\tag{6.}
f_{\varkappa'\lambda'\mu'}=f_{\varkappa\lambda\mu}.
\end{equation*}

That the right-hand side of equation (4.) is an integer can be seen directly without difficulty. For this sum is an integral rational function of a primitive $h$th root of unity $\vartheta$, and it remains unchanged if $\vartheta$ is replaced by any conjugate quantity $\vartheta^n$, where $n$ is prime to $h$. For, in order to carry out the substitution of $\vartheta$ by $\vartheta^n$, one needs only to replace $R$ by $R^n$ in each term of the sum (\emph{Prime factors}, § 12). In doing so the sum remains unchanged, because $R^n$ runs through the $h$ elements of $\Hgr$ simultaneously with $R$. In this way, however, one does not see that that integer is positive, nor that it is divisible by $h$.

Since $\chi^{(\varkappa)}(E)=f_\varkappa=f_{\varkappa'}$ is the degree of $\Phi_\varkappa$, it follows from (2.) that
\begin{equation*}\tag{7.}
f_\varkappa f_\lambda=\sum_\mu f_{\varkappa\lambda\mu}f_\mu.
\end{equation*}
Therefore
\begin{equation*}\tag{8.}
f_{\varkappa\lambda\mu}\leq \frac{f_\varkappa f_\lambda}{f_\mu},
\end{equation*}
where for $f_\mu$ one may take the greatest of the three numbers $f_\varkappa,f_\lambda,f_\mu$, and $f_{\varkappa\lambda\mu}\leq f_\varkappa$.

If one sets
\begin{equation*}\tag{9.}
\sum_\rho f_{\varkappa\lambda\rho'}f_{\rho\mu\nu}=f_{\varkappa\lambda\mu\nu},
\end{equation*}
then
\[
h^2 f_{\varkappa\lambda\mu\nu}
=\sum_{R,S}\chi^{(\varkappa)}(R)\chi^{(\lambda)}(R)\chi^{(\mu)}(S)\chi^{(\nu)}(S)
\left(\sum_\rho\chi^{(\rho')}(R)\chi^{(\rho)}(S)\right).
\]
But now
\[
\sum_\rho\chi^{(\rho')}(R)\chi^{(\rho)}(S)=0,
\]
except when $S$ is conjugate with $R$, and then it is equal to $\frac{h}{h_R}$. If $R$ is given, the latter case occurs for $h_R$ different values of $S$, for which $\chi(S)=\chi(R)$. Consequently
\begin{equation*}\tag{10.}
h f_{\varkappa\lambda\mu\nu}
=\sum_R \chi^{(\varkappa)}(R)\chi^{(\lambda)}(R)\chi^{(\mu)}(R)\chi^{(\nu)}(R).
\end{equation*}
This number also therefore remains unchanged under every permutation of the indices, something which is not apparent from formula (9.). In a similar way the symbol $f_{\varkappa\lambda\mu\ldots}$ is to be explained for arbitrarily many indices, but not for only one index, where $f_\varkappa$ denotes the degree of $\Phi_\varkappa$. As already remarked above, $f_{\varkappa\varkappa'}=1$ and $f_{\varkappa\lambda}=0$ if $\chi_\varkappa$ and $\chi_\lambda$ are not inverse characters. Further, $f_{\varkappa\lambda 0}=f_{\varkappa\lambda}$ and $f_{\varkappa\lambda\mu0}=f_{\varkappa\lambda\mu}$.

If $\chi_\varkappa$ is a character of the first degree, hence $f_\varkappa=1$, then by (8.) $f_{\varkappa\lambda\mu}$ is not greater than $\frac{f_\lambda}{f_\mu}$ or than $\frac{f_\mu}{f_\lambda}$, and hence is equal to zero. Only when $f_\lambda=f_\mu$ can $f_{\varkappa\lambda\mu}=1$ occur. By (7.) one has $f_\lambda=\sum_\mu f_{\varkappa\lambda\mu}f_\mu$, and consequently the latter case must occur for a given $\lambda$ for one, but also for only one, value of $\mu$. Thus, if $f_\varkappa=1$, then for every value of $\lambda$, of the $k$ numbers $f_{\varkappa\lambda\mu}$ one is equal to 1 and the other $k-1$ are zero; and if $f_{\varkappa\lambda\mu}=1$, then $f_\lambda=f_\mu$. The product of a character of degree $f$, $\chi^{(\lambda)}(R)$, and a character of the first degree, $\chi^{(\varkappa)}(R)$, is therefore a character of degree $f$, $\chi^{(\mu)}(R)=\chi^{(\varkappa)}(R)\chi^{(\lambda)}(R)$, as I have already shown by another route, \emph{Prime factors}, § 3. Formula (2.) is to be regarded as a generalization of this theorem.

With the aid of the formula
\begin{equation*}\tag{11.}
\sum_\varkappa \chi^{(\varkappa')}(R)\chi^{(\varkappa)}(R)=\frac{h}{h_R}
\end{equation*}
one obtains easily
\begin{equation*}\tag{12.}
\sum_\lambda f_{\varkappa\lambda\lambda'}=\sum_\rho \chi_\rho^{(\varkappa)},\qquad
\sum_\mu f_{\varkappa\lambda\mu\mu'}=\sum_\rho \chi_\rho^{(\varkappa)}\chi_\rho^{(\lambda)},\qquad
\sum_\nu f_{\varkappa\lambda\mu\nu\nu'}=\sum_\rho \chi_\rho^{(\varkappa)}\chi_\rho^{(\lambda)}\chi_\rho^{(\mu)}.
\end{equation*}
Therefore
\[
\sum_{\lambda,\mu}f_{\lambda\lambda'\mu\mu'}
=\sum_\rho\left(\sum_\lambda \chi_\rho^{(\lambda')}\chi_\rho^{(\lambda)}\right)
=\sum_\rho\frac{h}{h_\rho}
\]
and hence, by (9.) and (6.),
\begin{equation*}\tag{13.}
\sum_{\varkappa,\lambda,\mu} f_{\varkappa\lambda\mu}^{\,2}
=\sum_\rho\frac{h}{h_\rho}.
\end{equation*}
Accordingly the numbers $f_{\varkappa\lambda\mu}$ have relatively small values.

That the sums (12.) are rational integers can be seen directly in the same way as was shown above for the sum (5.). As formulas (12.) show, these numbers are positive, and they can be expressed through the numbers $f_{\varkappa\lambda\mu}$.

The numbers $f_{\varkappa\lambda\mu}$ introduced here have several properties in common with the numbers $h_{\alpha\beta\gamma}$ used earlier. For example, the determinant of degree $k$ is
\begin{equation*}\tag{14.}
\left|\left(\sum_\mu f_{\varkappa\lambda'\mu}x_\mu\right)-f_{\varkappa\lambda'}u\right|
=\prod_\lambda\left(\left(\sum_\varkappa \chi_\lambda^{(\varkappa)}x_\varkappa\right)-u\right)
\qquad(\varkappa,\lambda=0,1,\cdots k-1).
\end{equation*}
Its indices $\varkappa,\lambda,\mu$ refer to the $k$ prime factors of the group determinant, while the indices $\alpha,\beta,\gamma$ refer to the $k$ classes of conjugate elements into which the $h$ elements of $\Hgr$ decompose. The formulas developed here also can be generalized by adjoining, as in my last paper, to the characters of $\Hgr$ the characters of a subgroup.

\section*{§ 2.}

From two linear substitutions (a) and (a') in $f$ and $f'$ variables we have formed a linear substitution (A) in $ff'$ variables. The roots of its characteristic equation are the $ff'$ products obtained by multiplying each of the $f$ roots of the characteristic equation of (a) by each of the $f'$ roots of the characteristic equation of (a'). This follows easily from the observation that, if $(c)^{-1}(a)(c)=(b)$ and $(c')^{-1}(a')(c')=(b')$, then also $(C)^{-1}(A)(C)=(B)$ must hold.

If the matrices $(a),(b),(c)\cdots$ form the primitive representation of $\Hgr$ which corresponds to the prime function $\Phi(x)$, then one obtains the characteristic function of (a) by setting $x_A=-1$ in $\Phi(x+\varepsilon u)$ and setting the other variables $x_R=0$. If it is equal to
\[
F_A(u)=(u-\alpha_1)(u-\alpha_2)\cdots(u-\alpha_f),
\]
then $\alpha_1,\alpha_2,\cdots,\alpha_f$ are $m$th roots of unity, where $m$ is the order of the element $A$ (\emph{Prime factors}, § 12). Their sum is
\begin{equation*}\tag{1.}
\chi(A)=\alpha_1+\alpha_2+\cdots+\alpha_f.
\end{equation*}
If $\Phi=\Phi_\varkappa$, let $F_A(u)$ be denoted by $F_A^{(\varkappa)}(u)$, or by $F_\alpha^{(\varkappa)}(u)$ if $A$ belongs to the $\alpha$th class. From formula (2.) of § 1, in conjunction with equation (1.), there follows the theorem:
\begin{quote}\itshape
The roots of the equation
\begin{equation*}\tag{2.}
\prod_\mu\left(F_\alpha^{(\mu)}(u)\right)^{f_{\varkappa\lambda\mu'}}
\end{equation*}
of degree $f_\varkappa f_\lambda$ are obtained by multiplying each of the $f_\varkappa$ roots of the equation $F_\alpha^{(\varkappa)}(u)=0$ by each of the $f_\lambda$ roots of the equation $F_\alpha^{(\lambda)}(u)=0$.
\end{quote}
If one expands the logarithmic derivative of expression (2.) according to descending powers of $u$, then the coefficient of $u^{-n-1}$ is
\[
\sum_\mu f_{\varkappa\lambda\mu'}\chi^{(\mu)}(A^n)
=\chi^{(\varkappa)}(A^n)\chi^{(\lambda)}(A^n).
\]
From this it is apparent that the relation stated in that theorem is in no way more general than formula (2.), § 1.

In the function
\begin{equation*}\tag{3.}
(-1)^f F_\alpha^{(\varkappa)}(-u)
=(u+\alpha_1)(u+\alpha_2)\cdots(u+\alpha_f)
\end{equation*}
the coefficient of $u^{f-1}$ is equal to $\chi_\alpha^{(\varkappa)}$. The constant term is a character of the first degree (\emph{Prime factors}, § 12). More generally, the coefficient of $u^{f-n}$, for which I choose the notation introduced in \emph{Prime factors}, § 4 (8.), is equal to an expression of the form
\begin{equation*}\tag{4.}
\FrobP_n(A)=\frac{1}{n!}\chi(A,A,\cdots A)=\sum_\lambda s_{n\lambda}\chi^{(\lambda)}(A),
\end{equation*}
where the quantities $s_{n\lambda}$ are positive integers independent of $\alpha$.

To prove this, form from the substitution (a) $n$ different substitutions
\begin{equation*}\tag{5.}
u_\alpha^{(\nu)}=\sum_\beta a_{\alpha\beta}v_\beta^{(\nu)}
\qquad(\alpha,\beta=1,2,\cdots f),
\end{equation*}
where
\[
u_1^{(\nu)},u_2^{(\nu)},\cdots,u_f^{(\nu)}
\qquad(\nu=1,2,\cdots n)
\]
are $nf$ independent variables. The $\binom{f}{n}=g$ determinants of degree $n$ which can be formed from them are to be denoted in some order by $U_1,U_2,\cdots,U_g$, and the determinants formed analogously from the variables
\[
v_1^{(\nu)},v_2^{(\nu)},\cdots,v_f^{(\nu)}
\qquad(\nu=1,2,\cdots n)
\]
by $V_1,V_2,\cdots,V_g$. According to equations (5.), $U_1,U_2,\cdots,U_g$ are then linear functions of $V_1,V_2,\cdots,V_g$, whose coefficients are the $g^2$ minors of degree $n$ of the matrix $a_{\alpha\beta}$. The linear substitution so obtained is now to be denoted by (A). The roots of its characteristic equation are the products of $n$ at a time of the roots $\alpha_1,\alpha_2,\cdots,\alpha_f$; its trace is the sum of the principal minors of degree $n$ of the matrix (a), and hence is equal to the expression $\FrobP_n(A)$. If in the same way to the matrices $(b),(c),\cdots$ of degree $f$ there correspond the matrices $(B),(C),\cdots$ of degree $g$, then, if $(c)=(a)(b)$, one also has $(C)=(A)(B)$. Consequently $(A)x_A+(B)x_B+(C)x_C+\cdots$ is a matrix belonging to $\Hgr$. This was also considered by \textsc{Moliên}, \emph{On the invariants of linear substitution groups}, § 5 (Sitzungsberichte 1897). Its determinant is an expression of the form $\prod_\lambda\Phi_\lambda^{s_{\varkappa\lambda}^{n}}$, whence relation (4.) results. Further, the roots of the equation
\[
\prod_\lambda\left(F_\alpha^{(\lambda)}(u)\right)^{s_{\varkappa\lambda}^{n}}=0
\]
of degree $\binom{f_\varkappa}{n}$ are the products of $n$ at a time of the $f_\varkappa$ roots $\alpha_1,\alpha_2,\cdots$ of the equation $F_\alpha^{(\varkappa)}(u)=0$.

The expressions (4.) have the form
\begin{equation*}\tag{6.}
2\FrobP_2(R)=\chi(R)^2-\chi(R^2),\qquad
6\FrobP_3(R)=\chi(R)^3-3\chi(R)\chi(R^2)+2\chi(R^3),
\end{equation*}
in general
\begin{equation*}\tag{7.}
\FrobP_n(R)=\sum (-1)^{\beta+\delta+\cdots}
\frac{\chi(R)^\alpha}{1^\alpha\alpha!}
\frac{\chi(R^2)^\beta}{2^\beta\beta!}
\frac{\chi(R^3)^\gamma}{3^\gamma\gamma!}
\frac{\chi(R^4)^\delta}{4^\delta\delta!}\cdots,
\end{equation*}
where the sum is to be extended over all positive solutions of the equation
\begin{equation*}\tag{8.}
\alpha+2\beta+3\gamma+4\delta+\cdots=n.
\end{equation*}

As \textsc{Moliên}, loc. cit., § 2, has shown, a formula analogous to (4.) holds for the coefficient $\zeta_n(R)$ of $u^{-f-n}$ in the expansion of the function $\frac{1}{F_\alpha^{(\varkappa)}(u)}$. This coefficient is equal to the sum
\begin{equation*}\tag{9.}
\zeta_n(R)=\sum
\frac{\chi(R)^\alpha}{1^\alpha\alpha!}
\frac{\chi(R^2)^\beta}{2^\beta\beta!}
\frac{\chi(R^3)^\gamma}{3^\gamma\gamma!}
\frac{\chi(R^4)^\delta}{4^\delta\delta!}\cdots,
\end{equation*}
extended over the positive solutions of equation (8.).

\section*{§ 3.}

In conclusion I shall add a few examples of groups whose characters I have computed with the aid of the methods which I have developed here and in my last paper.

The elements of the symmetric group of degree $n=6$ are the 720 permutations of six symbols $a,b,c,d,e,f$; those of the alternating group are the 360 even permutations. In the first column of the following table, each class of conjugate elements is represented by a permutation decomposed into its cycles. By composition with the character of the first degree $\chi^{(1)}$, the characters $\chi^{(3)},\chi^{(5)},\chi^{(7)},\chi^{(9)}$ arise from the characters $\chi^{(2)},\chi^{(4)},\chi^{(6)},\chi^{(8)}$, while $\chi^{(10)}$ remains unchanged. Especially noteworthy is the character $\chi=\chi^{(2)}$. If $\lambda$ symbols are not moved by the permutation $R$ (if $R$ contains $\lambda$ cycles of the first degree), then $\chi(R)=\lambda-1$.

According to the developments of my last paper (§ 5), one can also express the rule by which $\chi_2$ is to be computed as follows. The permutations of $\Hgr$ which leave a specified symbol unchanged form a group $\Ggr$ of degree $g=\frac{h}{n}$. If the $\rho$th class contains $h_\rho$ elements, and if $g_\rho$ of them are contained in $\Ggr$, then
\begin{equation*}\tag{1.}
\chi_\rho=\frac{h g_\rho}{g h_\rho}-1.
\end{equation*}
From the same formula one finds the character $\chi^{(6)}$ by taking for $\Ggr$ the group of order $g=72$ contained in $\Hgr$. This imprimitive group is obtained by dividing the $n=6$ symbols into $s=2$ systems of $r=3$ symbols each, $abc,def$, and then permuting the $s$ systems and the $r$ symbols of each system in all possible ways. The order of this group is $g=(r!)^s s!=(3!)^2 2!=72$.

Finally, by means of formula (1.) one can also compute the character $\chi^{(4)}$ by taking for $\Ggr$ the triply transitive group of order 120 which is isomorphic to the symmetric group of degree 5. Instead one can derive $\chi^{(4)}$ from $\chi^{(2)}$ by the known isomorphism of $\Hgr$ into itself, by which the classes 1 and 7, 4 and 10, 5 and 9 are interchanged with one another, while the others remain unchanged. Through this isomorphism $\chi^{(8)}$ also goes over into $\chi^{(9)}=\chi^{(8)}\chi^{(1)}$.

\begin{center}
\emph{Symmetric group of degree 6.}\\[0.4em]
$h=720$.\\[0.4em]
{\scriptsize
\resizeMath{
\begin{array}{ll|rrrrrrrrrrr|r}
 & & \chi^{(0)}&\chi^{(1)}&\chi^{(2)}&\chi^{(3)}&\chi^{(4)}&\chi^{(5)}&\chi^{(6)}&\chi^{(7)}&\chi^{(8)}&\chi^{(9)}&\chi^{(10)}&h_\alpha\\ \hline
 &\chi_0&1&1&5&5&5&5&9&9&10&10&16&1\\
(ab)&\chi_1&1&-1&3&-3&-1&1&3&-3&2&-2&0&15\\
(ab)(cd)&\chi_2&1&1&1&1&1&1&1&1&-2&-2&0&45\\
(abcd)&\chi_3&1&-1&1&-1&1&-1&-1&1&0&0&0&90\\
(abc)&\chi_4&1&1&2&2&-1&-1&0&0&1&1&-2&40\\
(abc)(de)&\chi_5&1&-1&0&0&-1&1&0&0&-1&1&0&120\\
(abcde)&\chi_6&1&1&0&0&0&0&-1&-1&0&0&1&144\\
(ab)(cd)(ef)&\chi_7&1&-1&-1&1&3&-3&3&-3&-2&2&0&15\\
(abcd)(ef)&\chi_8&1&1&-1&-1&-1&-1&1&1&0&0&0&90\\
(abcdef)&\chi_9&1&-1&-1&1&0&0&0&0&1&-1&0&120\\
(abc)(def)&\chi_{10}&1&1&-1&-1&2&2&0&0&1&1&-2&40
\end{array}}}
\end{center}

\begin{center}
\emph{Alternating group of degree 6.}\\[0.4em]
$h=360$.\\[0.4em]
{\scriptsize
\resizeMath{
\begin{array}{ll|rrrrrrr|r}
 & & \chi^{(0)}&\chi^{(2)}&\chi^{(4)}&\chi^{(6)}&\chi^{(8)}&\chi^{(10)}&\chi^{(\overline{10})}&h_\alpha\\ \hline
 &\chi_0&1&5&5&9&10&8&8&1\\
(ab)(cd)&\chi_2&1&1&1&1&-2&0&0&45\\
(abc)&\chi_4&1&2&-1&0&1&-1&-1&40\\
(abcde)&\chi_6&1&0&0&-1&0&\half(1+\sqrt5)&\half(1-\sqrt5)&72\\
(acebd)&\chi_{\overline6}&1&0&0&-1&0&\half(1-\sqrt5)&\half(1+\sqrt5)&72\\
(abcd)(ef)&\chi_8&1&-1&-1&1&0&0&0&90\\
(abc)(def)&\chi_{10}&1&-1&2&0&1&-1&-1&40
\end{array}}}
\end{center}

\section*{§ 4.}

The binary groups of the tetrahedron, octahedron, and icosahedron have the orders $h=24,48$, and 120. The classes are for the most part already determined by the orders of their elements, which are found in the first column of each table. Only where the order does not suffice for the definition of the class is an element of the class found in that column; specifically, in the group of the tetrahedron (octahedron, icosahedron) the symbol $R$ ($S,T$) denotes an arbitrary element of order 6 (8, 10). The class (4) of the octahedral group contains every element of order 4 which cannot be represented as the square of an element of order 8; the remaining elements of order 4 form the class (5) characterized by the symbol $S^2$.

Each of the three groups contains an invariant subgroup $\Fgr=E+F$ of order 2. The element $F$ of order 2 is interchangeable with every element of the group and forms, by itself, the class (1). The group $\frac{\Hgr}{\Fgr}$ is the group of order 12 (24, 60), whose characters I computed in \emph{Group characters}, § 8. From these tables one can directly read off the values of the characters for which $\chi_0=\chi_1$, and which belong to the group $\frac{\Hgr}{\Fgr}$. For these characters one has in general $\chi(RF)=\chi(R)$, but for all the others $\chi(RF)=-\chi(R)$.

The most interesting property of these groups consists in the fact that they possess proper characters of the second degree. (In the table for the octahedral group, \emph{Group characters}, § 8, the character $\chi^{(3)}$ of degree $f=2$ belongs to the group $\frac{\Hgr}{\Ggr}$, where $\Ggr$ is the invariant subgroup of order 4 of $\Hgr$.) Their values result at once from the known representations of these groups by binary linear substitutions with the aid of the formula in \emph{Representation}, § 4 (5). For the octahedral group there are exactly two non-equivalent representations of this kind, but they differ only by the sign of $\sqrt2$ (for the character $\chi^{(4)}$ belongs to the group $\frac{\Hgr}{\Fgr}$).

The same holds for the icosahedral group. For the tetrahedral group, however, there are three non-equivalent representations. Their characters $\chi^{(4)},\chi^{(5)}$, and $\chi^{(6)}$ are obtained from one of them by composition with the three characters of the first degree. In this way one finds all the characters of these three groups without using the numbers $h_{\alpha\beta\gamma}$. For the still missing characters $\chi^{(7)}$ of the octahedral group and $\chi^{(7)}$ and $\chi^{(8)}$ of the icosahedral group are obtained from the quadratic relations existing between the characters.

\clearpage

\begin{center}
\emph{Tetrahedron.}\\[0.25em]
$h=24$.\\[0.25em]
{\scriptsize
\resizeMath{
\begin{array}{ll|rrrrrrr|r}
 & & \chi^{(0)}&\chi^{(1)}&\chi^{(2)}&\chi^{(3)}&\chi^{(4)}&\chi^{(5)}&\chi^{(6)}&h_\alpha\\ \hline
1&\chi_0&1&1&1&3&2&2&2&1\\
2&\chi_1&1&1&1&3&-2&-2&-2&1\\
4&\chi_2&1&1&1&-1&0&0&0&6\\
R^2&\chi_3&1&\rho&\rho^2&0&-1&-\rho&-\rho^2&4\\
R^4&\chi_4&1&\rho^2&\rho&0&-1&-\rho^2&-\rho&4\\
R^5&\chi_5&1&\rho&\rho^2&0&1&\rho&\rho^2&4\\
R&\chi_6&1&\rho^2&\rho&0&1&\rho^2&\rho&4
\end{array}}}
\end{center}

\begin{center}
\emph{Octahedron.}\\[0.25em]
$h=48$.\\[0.25em]
{\scriptsize
\resizeMath{
\begin{array}{ll|rrrrrrrr|r}
 & & \chi^{(0)}&\chi^{(1)}&\chi^{(2)}&\chi^{(3)}&\chi^{(4)}&\chi^{(5)}&\chi^{(6)}&\chi^{(7)}&h_\alpha\\ \hline
1&\chi_0&1&1&3&3&2&2&2&4&1\\
2&\chi_1&1&1&3&3&2&-2&-2&-4&1\\
3&\chi_2&1&1&0&0&-1&-1&-1&1&8\\
6&\chi_3&1&1&0&0&-1&1&1&-1&8\\
4&\chi_4&1&-1&1&-1&0&0&0&0&12\\
S^2&\chi_5&1&1&-1&-1&2&0&0&0&6\\
S^3&\chi_6&1&-1&-1&1&0&\sqrt2&-\sqrt2&0&6\\
S&\chi_7&1&-1&-1&1&0&-\sqrt2&\sqrt2&0&6
\end{array}}}
\end{center}

\begin{center}
\emph{Icosahedron.}\\[0.25em]
$h=120$.\\[0.25em]
{\scriptsize
\resizeMath{
\begin{array}{ll|rrrrrrrrr|r}
 & & \chi^{(0)}&\chi^{(1)}&\chi^{(2)}&\chi^{(3)}&\chi^{(4)}&\chi^{(5)}&\chi^{(6)}&\chi^{(7)}&\chi^{(8)}&h_\alpha\\ \hline
1&\chi_0&1&5&4&3&3&2&2&4&6&1\\
2&\chi_1&1&5&4&3&3&-2&-2&-4&-6&1\\
4&\chi_2&1&1&0&-1&-1&0&0&0&0&30\\
3&\chi_3&1&-1&1&0&0&-1&-1&1&0&20\\
6&\chi_4&1&-1&1&0&0&1&1&-1&0&20\\
T^4&\chi_5&1&0&-1&\half(1+\sqrt5)&\half(1-\sqrt5)&\half(-1+\sqrt5)&\half(-1-\sqrt5)&-1&1&12\\
T^2&\chi_6&1&0&-1&\half(1-\sqrt5)&\half(1+\sqrt5)&\half(-1-\sqrt5)&\half(-1+\sqrt5)&-1&1&12\\
T^3&\chi_7&1&0&-1&\half(1-\sqrt5)&\half(1+\sqrt5)&\half(1+\sqrt5)&\half(1-\sqrt5)&1&-1&12\\
T&\chi_8&1&0&-1&\half(1+\sqrt5)&\half(1-\sqrt5)&\half(1-\sqrt5)&\half(1+\sqrt5)&1&-1&12
\end{array}}}
\end{center}
% ===== END WORK 058 =====

\clearpage

% ===== BEGIN WORK 059: 059_ber_die_Darstellung_der_endlichen_Gruppen =====
\frobtitle{On the representation of finite groups\\by linear substitutions. II.}

To the $h$ distinct elements $A,B,C,\cdots$ of a finite group $\Hgr$ let $h$ homogeneous linear substitutions $a,b,c,\cdots$ in $n$ variables be so assigned that always, when $AB=C$, also $ab=c$. Then these substitutions form a group which is holohedrically or merohedrically isomorphic with $\Hgr$, or, expressed otherwise, the substitutions form a \emph{representation} of the group $\Hgr$. By the signs $a,b,c,\cdots$ one may also understand the matrices of the substitutions, and by $ab$ the matrix composed from $a$ and $b$.

From the given representation one can derive a new one by introducing other variables into the substitutions. This procedure amounts to replacing the matrices $a,b,c,\cdots$ by $p^{-1}ap,p^{-1}bp,p^{-1}cp,\cdots$, where $p$ is an arbitrary matrix with non-vanishing determinant. In the first part of this paper (Sitzungsberichte 1897, cited in what follows as $D.$, § 2) I called two such representations \emph{equivalent}.

If for the same group $\Hgr$ one knows a second representation by the matrices $a',b',c',\cdots$ of degree $n'$, then the matrices
\[
\begin{pmatrix}a&0\\0&a'\end{pmatrix},\quad
\begin{pmatrix}b&0\\0&b'\end{pmatrix},\quad
\begin{pmatrix}c&0\\0&c'\end{pmatrix},\cdots
\]
form a new representation of degree $n+n'$. The case is not thereby excluded in which the second representation is identical with, or equivalent to, the first. In the same way, from several known representations, whether equal or different, one can derive a new representation. This group of substitutions can under certain circumstances be holohedrically isomorphic to the group $\Hgr$, although the given representations were only merohedrically isomorphic to it; and this is the reason why, in the investigation of all representations of a given group $\Hgr$, those representations too must not be excluded which in reality represent a group merohedrically isomorphic with $\Hgr$.

Every representation which is obtained in the above-discussed way from several representations I call a \emph{decomposed} or \emph{decomposable} one, and every representation which is equivalent to a decomposed one an \emph{imprimitive} or \emph{reducible} one. If, however, a representation is equivalent to no decomposable representation, then I call it \emph{primitive} or \emph{irreducible} (compare the earlier provisional definition, $D.$ § 5).

If equivalent representations are not regarded as distinct, then there is only a finite number of different primitive representations of a group $\Hgr$ by linear substitutions or their matrices. This number $k$ is equal to the number of the classes of conjugate elements into which the elements of $\Hgr$ decompose. Every imprimitive representation is equivalent to another which decomposes into purely primitive representations, but each individual one of the $k$ primitive representations can occur several times. Moreover such a decomposition is possible in only one way.

Let $x_A,x_B,x_C,\cdots$ be $h$ independent variables, and let
\[
\Theta=|x_{PQ^{-1}}|=\Phi^f{\Phi'}^{f'}{\Phi''}^{f''}\cdots
\]
be the determinant of the group $\Hgr$, and let $\Phi,\Phi',\Phi'',\cdots$ be its different prime factors. If the matrices $a,b,c,\cdots$ of degree $n$ form a representation of $\Hgr$, then I call ($D.$ § 2) the matrix of degree $n$, $ax_A+bx_B+cx_C+\cdots$, the matrix corresponding to the representation of $\Hgr$, or a matrix belonging to the group $\Hgr$. Its $n^2$ elements are linear functions of the $h$ variables $x_A,x_B,x_C,\cdots$, and its determinant
\[
|ax_A+bx_B+cx_C+\cdots|=\Phi^s{\Phi'}^{s'}{\Phi''}^{s''}\cdots
\]
is a product of prime factors of the group determinant. Conversely, to each such product there corresponds one and only one representation of $\Hgr$, that is, two representations whose corresponding matrices have equal determinants are equivalent. To each one of the $k$ prime factors $\Phi$ of the group determinant $\Theta$ there corresponds one of the $k$ primitive representations, which I shall denote by $[\Phi]$. To a product $\Phi^s{\Phi'}^{s'}\cdots$ there corresponds a representation which decomposes into $s$ representations $[\Phi]$, $s'$ representations $[\Phi']$, $\cdots$. By investigating the determinant of the matrix which corresponds to a given representation, one can therefore recognize whether the representation is primitive or not, and in the latter case into which primitive representations it can be decomposed.

The results developed here, distinguished by their simplicity, form the conclusion of my general investigations on the group determinant. \textsc{Molien} obtained them by another route in the work cited in $D.$ § 4.

In conclusion I consider more closely the construction of the primitive representations. To obtain such a representation, one needs only to calculate a solution of a certain system of linear and quadratic equations. Every such solution I therefore call a \emph{system of values determining the representation}.

\section*{§ 1.}

In the matrix $u$ of degree $f$ let the elements
\[
u_{\alpha\beta}\qquad(\alpha,\beta=1,2,\cdots f)
\]
be $f^2$ variables independent of one another. Let $v_{\alpha\beta}$ be $f^2$ other variables, and let
\[
w_{\alpha\beta}=u_{\alpha1}v_{1\beta}+u_{\alpha2}v_{2\beta}+\cdots+u_{\alpha f}v_{f\beta}.
\]
If then the matrix $u$ passes into $v$ or $w$ when $u_{\alpha\beta}$ is replaced by $v_{\alpha\beta}$ or $w_{\alpha\beta}$, then $w=uv$.

Let the $n^2$ elements $x_{\kappa\lambda}$ of the matrix $X$ of degree $n$ be linear functions of the $f^2$ variables $u_{\alpha\beta}$. Let it pass into $Y$ or $Z$ when $u_{\alpha\beta}$ is replaced by $v_{\alpha\beta}$ or $w_{\alpha\beta}$. We shall investigate how those linear functions must be constituted in order that $Z=XY$ when $w=uv$. I restrict myself here to the case in which the determinant $|X|$ is different from zero.

If $X$ has this property, then so also has $PXP^{-1}$, where $P$ is a matrix of $n^2$ constant elements and $|P|$ is different from zero. Further, if $0$ is the matrix of degree $f$ all of whose elements vanish, and if $n=fg$ is a multiple of $f$, then
\[
U=\begin{pmatrix}
u&0&0&\cdots\\
0&u&0&\cdots\\
0&0&u&\cdots\\
\cdot&\cdot&\cdot&\cdots
\end{pmatrix}
\]
has the required property, and so also does $P^{-1}UP=X$. I now wish conversely to show: if $X$ is any matrix of the kind under consideration, then $n=fg$ must be a multiple of $f$, and one can determine a constant matrix $P$ so that $PXP^{-1}=U$.

If $B$ is the identity matrix of degree $g$, then by a certain permutation of the rows and the same permutation of the columns one can bring the matrix $U$ of degree $fg$ into the form
\[
V=\begin{pmatrix}
u_{11}B&u_{12}B&\cdots&u_{1f}B\\
u_{21}B&u_{22}B&\cdots&u_{2f}B\\
\cdot&\cdot&\cdots&\cdot\\
u_{f1}B&u_{f2}B&\cdots&u_{ff}B
\end{pmatrix}.
\]
The permutation of the columns is effected by composing $U$ with a matrix $Q$ in which in each row and in each column one element is equal to $1$ and the others are equal to $0$. If one then performs on $UQ$ the same permutation of the rows, one obtains $Q'UQ$, where $Q'$ is the matrix conjugate to $Q$. But since $Q$ is an orthogonal matrix, $Q'=Q^{-1}$, and hence $V=Q^{-1}UQ$.

Since the elements of the matrix $X$ are linear functions of $u_{11},u_{12},\cdots,u_{ff}$, one can set $X=\sum A_{\alpha\beta}u_{\alpha\beta}$, where $A_{\alpha\beta}$ is a constant matrix of degree $n$. Now $XY=Z$ is required, hence
\[
\Bigl(\sum A_{\alpha\beta}u_{\alpha\beta}\Bigr)
\Bigl(\sum A_{\gamma\delta}v_{\gamma\delta}\Bigr)
 =\sum A_{\alpha\gamma}w_{\alpha\gamma}
 =\sum A_{\alpha\gamma}u_{\alpha\beta}v_{\beta\gamma}.
\]
It follows that $A_{\alpha\beta}A_{\gamma\delta}=0$ when $\beta$ is different from $\gamma$, whereas $A_{\alpha\beta}A_{\beta\gamma}=A_{\alpha\gamma}$. Since $A_{11}^2=A_{11}$, one has
\[
|A_{11}+rE|=r^{n-g_1}(1+r)^{g_1}.
\]
If $g_1=0$, then $A_{11}=0$, and hence also $A_{\alpha\beta}=A_{\alpha1}A_{11}A_{1\beta}=0$. Therefore $g_1>0$, and one can determine $P$ so that $P^{-1}A_{11}P$ becomes a matrix in which $c_{11}=\cdots=c_{g_1,g_1}=1$ and all other elements are equal to zero. We suppose $X$ to have been replaced by $P^{-1}XP$, and therefore assume that $A_{11}$ itself is that matrix. Now $A_{11}A_{22}=A_{22}A_{11}=0$, and consequently in $A_{22}$ the elements of the first $g_1$ rows and columns are all zero, so that $A_{22}$ can also be regarded as a matrix of degree $n-g_1$ and transformed as such. Consequently one can determine a matrix $P$ of degree $n$, in which the elements of the first $g_1$ rows and columns agree with the corresponding elements of $A_{11}$, in such a way that in $P^{-1}A_{22}P$ one has $c_{g_1+1,g_1+1}=\cdots=c_{g_1+g_2,g_1+g_2}=1$, while all other elements vanish. Then in $A_{12}=A_{11}A_{12}A_{22}$ only those elements need not be zero which the rows $1,2,\cdots,g_1$ have in common with the columns $g_1+1,g_1+2,\cdots,g_1+g_2$. After carrying out these transformations one has
\[
A_{11}=\begin{pmatrix}
B_{11}&N_{12}&N_{13}&\cdots\\
N_{21}&N_{22}&N_{23}&\cdots\\
N_{31}&N_{32}&N_{33}&\cdots\\
\cdot&\cdot&\cdot&\cdots
\end{pmatrix},\quad
A_{12}=\begin{pmatrix}
N_{11}&B_{12}&N_{13}&\cdots\\
N_{21}&N_{22}&N_{23}&\cdots\\
N_{31}&N_{32}&N_{33}&\cdots\\
\cdot&\cdot&\cdot&\cdots
\end{pmatrix},\quad\cdots,
\]
where $B_{\alpha\beta}$ is a matrix with $g_\alpha$ rows and $g_\beta$ columns, and $N_{\alpha\beta}$ is a matrix of the same kind with all elements vanishing. $B_{\alpha\alpha}$ is the identity matrix of degree $g_\alpha$. From $A_{\alpha\beta}A_{\beta\gamma}=A_{\alpha\gamma}$ it follows that $B_{\alpha\beta}B_{\beta\gamma}=B_{\alpha\gamma}$. Thus $B_{12}B_{21}=B_{11}$. If therefore $g_2<g_1$, the determinant of $B_{11}$ would have to vanish. Likewise $B_{21}B_{12}=B_{22}$ implies that $g_2\geq g_1$. Hence $g_1=g_2=g_3=\cdots=g$, and $n=fg$. For if $n>fg$, the last rows and columns of $X$ would consist entirely of zeros. If $B$ is the identity matrix of degree $g$, then $B_{\alpha\alpha}=B$ and $B_{\alpha\beta}B_{\beta\alpha}=B$. Consequently $B_{\alpha\beta}$ and $B_{\beta\alpha}$ are reciprocal matrices of degree $g$, and their determinants are non-zero. Let $N$ be a matrix of $g^2$ vanishing elements. Then the determinant of the matrix of degree $n$
\[
L=\begin{pmatrix}
B_{11}&N&N&\cdots\\
N&B_{12}&N&\cdots\\
N&N&B_{13}&\cdots\\
\cdot&\cdot&\cdot&\cdots
\end{pmatrix}
\]
is non-zero, and
\[
L^{-1}=\begin{pmatrix}
B_{11}&N&N&\cdots\\
N&B_{21}&N&\cdots\\
N&N&B_{31}&\cdots\\
\cdot&\cdot&\cdot&\cdots
\end{pmatrix}.
\]
The matrix $LA_{\alpha\beta}L^{-1}$ differs from $A_{\alpha\beta}$ only in that $B_{1\alpha}B_{\alpha\beta}B_{\beta1}=B_{11}=B$ occurs in place of $B_{\alpha\beta}$. Hence $LXL^{-1}=V$, and the asserted proposition is thereby proved.

\section*{§ 2.}

Let $(x_{PQ^{-1}})$ be the matrix of the group $\Hgr$ of order $h$, and let
\begin{equation}\tag{1.}
z_R=\sum x_Py_Q\qquad(PQ=R).
\end{equation}
Then $(z_{PQ^{-1}})=(x_{PQ^{-1}})(y_{PQ^{-1}})$. Let
\[
X=(x_{\kappa\lambda})\qquad(\kappa,\lambda=1,2,\cdots,n)
\]
be a matrix of degree $n$ whose determinant does not vanish, and whose $n^2$ elements $x_{\kappa\lambda}$ are linear functions of the $h$ variables $x_R$. Let it pass into $Y$ or $Z$ when $x_R$ is replaced by $y_R$ or $z_R$. If then $Z=XY$ under the condition (1.), $X$ is called a matrix belonging to the group $\Hgr$.

Let $\Phi,\Phi',\cdots$ be different prime factors of the determinant $\Theta$ of the group $\Hgr$, let $f,f',\cdots$ be their degrees, let $\chi(R),\psi(R),\cdots$ be their characters, and let $\chi'(R)=\chi(R^{-1}),\psi'(R)=\psi(R^{-1}),\cdots$ be the conjugate characters. Then, by \emph{Prime Factors}, § 8,
\begin{equation}\tag{2.}
\Phi\!\left(\eps+u\frac{f}{h}\chi'\right)=(1+u)^f,
\qquad
\Phi\!\left(\eps+u\frac{f'}{h}\psi'\right)=1,\cdots .
\end{equation}
If one sets $x_R=\frac{f}{h}\chi'(R)$, let $X=A$; if one sets $x_R=\frac{f'}{h}\psi'(R)$, let $X=B$; and so on. Then
\[
\sum \frac{f}{h}\chi'(P)\frac{f}{h}\chi'(Q)=\frac{f}{h}\chi'(R),
\qquad
\sum \frac{f}{h}\chi'(P)\frac{f'}{h}\psi'(Q)=0\qquad(PQ=R),
\]
and hence $A^2=A$, $B^2=B$, $AB=BA=0$. Since
\[
\sum \frac{f}{h}\chi'(P)x_Q=\sum x_P\frac{f}{h}\chi'(Q)\qquad(PQ=R)
\]
$A$ (and likewise $B,C,\cdots$) is interchangeable with every matrix $X$. Since finally
\[
\frac{f}{h}\chi'(R)+\frac{f'}{h}\psi'(R)+\cdots=\eps_R
\]
when the sum is extended over all $k$ characters of $\Hgr$, one has
\begin{equation}\tag{3.}
A+B+C+\cdots=E.
\end{equation}
Therefore
\[
(E+uAX)(E+uBX)(E+uCX)\cdots=E+u(A+B+C+\cdots)X=E+uX.
\]
All other terms in the expansion of the symbolic product by powers of $u$ vanish, because $AXBX=ABXX=0$. Consequently also
\begin{equation}\tag{4.}
|E+uAX|\,|E+uBX|\cdots=|E+uX|.
\end{equation}
But this determinant, since $X$ is a matrix belonging to the group $\Hgr$, is a product of powers of the prime functions $\Phi(\eps+ux),\Phi'(\eps+ux),\cdots$. Therefore also
\[
|E+uAX|=\Phi(\eps+ux)^s\Phi'(\eps+ux)^t\cdots.
\]
But one must have $t=0$. For, if one sets $x_R=\frac{f'}{h}\psi'(R)$, hence $X=B$, then $AX=AB=0$, so the determinant on the left is equal to $1$. On the right, however, $\Phi'(\eps+u\frac{f'}{h}\psi')=(1+u)^{f'}$. Consequently
\begin{equation}\tag{5.}
|E+uAX|=\Phi(\eps+ux)^s,
\qquad
|E+uBX|=\Phi'(\eps+ux)^{s'},\cdots
\end{equation}
and therefore by (4.)
\begin{equation}\tag{6.}
|X|=\Phi(x)^s\Phi'(x)^{s'}\cdots,
\end{equation}
and
\[
n=fs+f's'+\cdots.
\]
If one sets in equation (5.) $x_R=\frac{f}{h}\chi'(R)$, so that $X=A$, one obtains
\begin{equation}\tag{7.}
|E+uA|=(1+u)^{fs},
\qquad |uE-A|=(u-1)^{fs}u^{n-fs}.
\end{equation}
Since $A^2=A$, the elementary divisors of this determinant are all linear. Hence the rank of $A$ is equal to $r=fs$, and the principal minors of degree $r$ of $A$ are not all zero, since their sum is equal to $1$. If $s=0$, then $r=0$ and consequently $A=0$. In the sum (3.), therefore, only those matrices $A,B,C,\cdots$ actually occur which are formed from the characters $\chi$ of those prime functions $\Phi$ which divide $|X|$. For example, if $|X|$ is a power of $\Phi$, then $A=E$. Let
\begin{equation}\tag{8.}
|a_{\kappa\lambda}|\qquad(\kappa,\lambda=\alpha_1,\alpha_2,\cdots,\alpha_r)
\end{equation}
be a non-zero principal minor of degree $r$ of $A$. Let $r'=f's'$ be the rank of $B$, and let
\[
|b_{\kappa\lambda}|\qquad(\kappa,\lambda=\beta_1,\beta_2,\cdots,\beta_{r'})
\]
be a non-zero principal minor of degree $r'$ of $B$, and so on.

Let $M$ be the matrix of degree $n=r+r'+\cdots$ whose $n$ rows are obtained from the row
\[
a_{\lambda,\alpha_1},\cdots,a_{\lambda,\alpha_r},
 b_{\lambda,\beta_1},\cdots,b_{\lambda,\beta_{r'}},\cdots
\]
by putting $\lambda=1,2,\cdots,n$. Similarly let $L'$ be the matrix
\[
a_{\alpha_1,\kappa},\cdots,a_{\alpha_r,\kappa},
 b_{\beta_1,\kappa},\cdots,b_{\beta_{r'},\kappa},\cdots,
\]
and let $L$ be the conjugate matrix. I then form the matrix $LXM$. If $\rho$ and $\sigma$ are two of the indices $1,2,\cdots,r$, then the $\sigma$-th element of the $\rho$-th row is
\[
\sum_{\kappa,\lambda}a_{\alpha\kappa}x_{\kappa\lambda}a_{\lambda\beta},
\]
where $\alpha=\alpha_\rho$ and $\beta=\alpha_\sigma$. This is an element of the matrix $AXA=XAA=XA=AX$, and is therefore equal to
\begin{equation}\tag{9.}
\xi_{\alpha\beta}=\sum_\lambda a_{\alpha\lambda}x_{\lambda\beta}
=\sum_\lambda x_{\alpha\lambda}a_{\lambda\beta}
\end{equation}
and it is obtained from $x_{\alpha\beta}$ by replacing $x_R$ in it by
\begin{equation}\tag{10.}
\xi_R=\sum_S \frac{f}{h}\chi'(RS^{-1})x_S.
\end{equation}
If $\rho$ is one of the numbers $1,2,\cdots,r$, and $\sigma$ is one of the numbers $1,2,\cdots,r'$, then the $(r+\sigma)$-th element of the $\rho$-th row is
\[
\sum_{\kappa,\lambda}a_{\alpha\kappa}x_{\kappa\lambda}b_{\lambda\beta},
\]
where $\alpha=\alpha_\rho$, $\beta=\beta_\sigma$. This is an element of the matrix $AXB=XAB=0$. Consequently $X$ decomposes into submatrices of degrees $r,r',\cdots$, the first of which is formed from the $r^2$ elements
\[
(\xi_{\kappa\lambda})\qquad(\kappa,\lambda=\alpha_1,\alpha_2,\cdots,\alpha_r).
\]
In it $x_E$ is multiplied by the matrix
\[
N_1=(a_{\kappa\lambda})\qquad(\kappa,\lambda=\alpha_1,\alpha_2,\cdots,
\alpha_r),
\]
whose determinant does not vanish. Similarly let
\[
N_2=(b_{\kappa\lambda})\qquad(\kappa,\lambda=\beta_1,\beta_2,\cdots,
\beta_{r'})
\]
and
\[
N=\begin{pmatrix}
N_1&0&0&\cdots\\
0&N_2&0&\cdots\\
0&0&N_3&\cdots
\end{pmatrix}.
\]
Then $LXMN^{-1}=Z$ also decomposes into submatrices of degrees $r,r',\cdots$, the first of which is
\begin{equation}\tag{11.}
(\xi_{\kappa\lambda})(a_{\kappa\lambda})^{-1}
\qquad(\kappa,\lambda=\alpha_1,\alpha_2,\cdots,\alpha_r).
\end{equation}
In the matrix $Z$, therefore, $x_E$ is multiplied by the identity matrix $E$. If now one puts $x_R=\eps_R$, so that $X=E$, then also $Z=E$, and hence $LMN^{-1}=E$. Therefore $|L|$ and $|M|$ are non-zero, and $MN^{-1}=L^{-1}$. Consequently $Z=LXL^{-1}$ is a matrix belonging to the group $\Hgr$, and likewise each of the submatrices, such as (11.), into which it decomposes.

The matrix $A$ has rank $r$. By \emph{Prime Factors}, § 11, the determinants of degree $r$ of the matrix $AX=XA$ therefore stand in the same ratios as the corresponding determinants of the matrix $A$, and hence differ from one another only by constant factors. Now by (5.) and (7.) the sum of the principal minors of degree $r$ of $AX$ is equal to $\Phi(x)^s$, and that of $A$ is equal to $1$. Consequently the determinant of the matrix (11.) is equal to $\Phi(x)^s$.

\section*{§ 3.}

In \emph{Representation}, § 5, I transformed the matrix $(x_{PQ^{-1}})$ of the group $\Hgr$ into an equivalent one which decomposes into $f$ mutually equal submatrices
\begin{equation}\tag{1.}
(u_{\alpha\beta})\qquad(\alpha,\beta=1,2,\cdots,f)
\end{equation}
of degree $f$, into $f'$ mutually equal submatrices
\begin{equation}\tag{2.}
(u'_{\alpha\beta})\qquad(\alpha,\beta=1,2,\cdots,f')
\end{equation}
of degree $f'$, and so on. The $f^2+f'^2+\cdots=h$ elements $u_{\alpha\beta},u'_{\alpha\beta},\cdots$ of these matrices are $h$ linear functions, independent of one another, of the $h$ variables $x_R$. Their determinants are
\[
|u_{\alpha\beta}|=\Phi(x),\qquad |u'_{\alpha\beta}|=\Phi'(x),\cdots.
\]
In the same way the matrix $X$ belonging to the group $\Hgr$ and considered in the preceding paragraph can be transformed into an equivalent one which decomposes into $s$ submatrices (1.), into $s'$ submatrices (2.), and so on.

If $|X|$ is divisible by several different prime functions $\Phi,\Phi',\cdots$, then by § 2 the matrix $X$ can be transformed into a decomposed matrix in which each submatrix has a power of a prime function for its determinant. It therefore suffices to prove the assertion for the case in which $|X|=\Phi^s$ is a power of one prime function $\Phi$.

The $h$ variables $x_R$ are linear functions of the $h$ variables $u_{\alpha\beta},u'_{\alpha\beta},\cdots$. If the elements $x_{\kappa\lambda}$ of $X$ are expressed by means of these, then $\Phi$, and therefore also $|X|$, depends only on the $f^2$ variables $u_{\alpha\beta}$. But that each of the $n^2$ elements $x_{\kappa\lambda}$ also depends only on these quantities follows from the developments of § 2. According to them, in the present case $A=E$, hence $X=XA$, and therefore $x_{\kappa\lambda}=\xi_{\kappa\lambda}$; consequently $x_{\kappa\lambda}$ can be represented as a function of the $h$ variables $\xi_R$. Among these there are exactly $f^2$ linearly independent ones. On the other hand $\Phi$ can be represented by the $f^2$ mutually independent variables $u_{\alpha\beta}$, and cannot be transformed by a linear substitution into a function of fewer than $f^2$ variables. Consequently the $h$ quantities $\xi_R$, and hence also the $n^2$ quantities $x_{\kappa\lambda}$, are linear functions of the $f^2$ variables $u_{\alpha\beta}$.

If $x_R$ is replaced by $y_R$ or $z_R$, let $u=(u_{\alpha\beta})$ pass into $(v)$ or $(w)$, and let $X$ pass into $Y$ or $Z$. Then $w=uv$ and $Z=XY$, if $z_R$ is defined by equation (1.) of § 2. Thus the elements $x_{\kappa\lambda}$ of $X$ are linear functions of the elements $u_{\alpha\beta}$ of $u$ of such a kind that $Z=XY$ when $w=uv$. Therefore the possibility of the asserted transformation of $X$ follows from the developments of § 1.

In this proof one could also have dispensed with the theorems of § 2, and would only have had to give the theorems of § 1 a somewhat more general form. Instead of starting from one matrix $u$, one can start from several matrices $u,u',\cdots$ of degrees $f,f',\cdots$, whose $f^2+f'^2+\cdots$ elements are all variables independent of one another. The elements of the matrix $X$ are then linear functions of all these variables, and $Z=XY$ when simultaneously $w=uv,w'=u'v',\cdots$. Nevertheless the transformation carried out in § 2 is also of interest in itself, because it requires only the knowledge of the characters.

\section*{§ 4.}

Every prime factor $\Phi$ of degree $f$ of the group determinant $\Theta$ can be represented as a determinant of degree $f$ whose elements are $f^2$ mutually independent linear functions of the $h$ variables $x_R$. This representation can be chosen so that the elements of the determinant form a matrix belonging to the group $\Hgr$. In \emph{Representation}, § 3, I expressed the $f$-th power of $\Phi$ by a determinant of degree $f^2$. This representation is distinguished by the fact that its elements are all obtained from the linear function $\sum_R\chi(R)x_R$ by permuting the $h$ variables $x_R$ in a special way. Every determinant of degree $f^2$ formed from the elements of the matrix
\begin{equation}\tag{1.}
\sum_R\chi(R)x_{RPQ^{-1}}
\end{equation}
is, up to a constant factor, equal to $\Phi^f$.

I shall now show that the determinant of degree $f$ which is equal to $\Phi$ can also be brought into a similar form. For the determination of this determinant it is enough to calculate a single linear function $\sum a_{R^{-1}}x_R$. Then every determinant of degree $f$ of the matrix
\begin{equation}\tag{2.}
x_{P,Q}=\sum_R a_{R^{-1}}x_{PRQ^{-1}}
\end{equation}
is, up to a constant factor, equal to $\Phi$. Whereas the character $\chi(R)$ is completely determined by the prime function $\Phi$, the $h$ quantities $a_R$ can be chosen in infinitely many ways, apart from the case $f=1$, where the two representations coincide. This is because $\chi(PQ)=\chi(QP)$, whereas the quantities $a_R$ need not satisfy this condition.

To arrive at this representation of $\Phi$, I need some auxiliary propositions about matrices which, although known, shall be developed briefly here. For convenience of presentation I denote by the signs $A,B,C,\cdots$ not only matrices of degree $h$, but at the same time bilinear forms whose coefficients are the elements of those matrices, as in my paper \emph{On linear substitutions and bilinear forms}, \textsc{Crelle}'s Journal, vol. 84. The following propositions hold for any matrix $K$ (loc. cit. § 13), but here they shall be derived only for one whose characteristic determinant $|uE-K|$ has only linear elementary divisors. One can then determine a substitution $P$ (that is, a matrix $P$ with non-vanishing determinant) so that the bilinear form $P^{-1}KP$ assumes the normal form
\[
\rho(u_1v_1+\cdots+u_fv_f)+\rho'(u_{f+1}v_{f+1}+\cdots+u_{f+f'}v_{f+f'})+\cdots.
\]
Here $\rho,\rho',\cdots$ are the different roots of that characteristic equation, with $\rho$ an $f$-fold root, $\rho'$ an $f'$-fold root, and so on. Since the elementary divisors of $|uE-K|$ which vanish for $u=\rho$ are all linear, the expansion
\begin{equation}\tag{3.}
(uE-K)^{-1}=\frac{A}{u-\rho}+\cdots
\end{equation}
begins with the $(-1)$-st power of $u-\rho$. From the normal form it is clear that $P^{-1}AP=u_1v_1+\cdots+u_fv_f$. Hence $A$ has rank $f$, satisfies the equation
\begin{equation}\tag{4.}
A^2=A,
\end{equation}
and
\begin{equation}\tag{5.}
|E+\mu A|=(\mu+1)^f.
\end{equation}
Consequently the sum of the principal minors of degree $f$ of $A$ is equal to $1$. In the matrix $A$, in which the minors of degree $f+1$ all vanish, the principal minors of degree $f$ therefore cannot all be zero. If further
\[
(uE-K)^{-1}=\frac{B}{u-\rho'}+\cdots,
\]
then $P^{-1}BP=u_{f+1}v_{f+1}+\cdots+u_{f+f'}v_{f+f'}$, and consequently
\begin{equation}\tag{6.}
AB=BA=0.
\end{equation}
Without using the normal form one can arrive at these propositions by the following considerations. Squaring equation (3.) once, and differentiating it with respect to $u$ the other time, one obtains the two relations
\[
(uE-K)^{-2}=\frac{A^2}{(u-\rho)^2}+\cdots,
\qquad
-(uE-K)^{-2}=-\frac{A}{(u-\rho)^2}+\cdots,
\]
from which property (4.) follows at once. (Compare \textsc{Weierstrass}, Monatsberichte 1858, p. 215.)

From the identity
\[
(uE-K)\bigl((uE-K)^{-1}-(vE-K)^{-1}\bigr)(vE-K)=(vE-K)-(uE-K)=(v-u)E
\]
there follows the relation
\[
-\frac{(uE-K)^{-1}-(vE-K)^{-1}}{u-v}=(uE-K)^{-1}(vE-K)^{-1}.
\]
If these expressions are expanded by ascending powers of $u-\rho$ and $v-\rho'$, one obtains
\[
-\frac{1}{(u-\rho)-(v-\rho')+(\rho-\rho')}
\left(\frac{A}{u-\rho}+\cdots-\frac{B}{v-\rho'}-\cdots\right)
=\frac{AB}{(u-\rho)(v-\rho')}+\cdots.
\]
If $u-\rho$ and $v-\rho'$ are sufficiently small, then in the expansion of the first factor on the left-hand side only positive powers of $u-\rho$ and $v-\rho'$ occur, and consequently on the left no negative power of $u-\rho$ multiplied by a negative power of $v-\rho'$ appears. Therefore $AB=0$.

Since $(uE-K)^{-1}$ is a proper rational function of $u$, one has
\begin{equation}\tag{7.}
(uE-K)^{-1}=\frac{A}{u-\rho}+\frac{B}{u-\rho'}+\cdots.
\end{equation}
Expanding both sides of this equation by descending powers of $u$, comparison of the coefficients gives
\begin{equation}\tag{8.}
E=A+B+C+\cdots,
\qquad
K=\rho A+\rho'B+\rho''C+\cdots.
\end{equation}
Finally, that the rank of $A$ is equal to $f$ can be seen as follows. If $\psi(u)=(u-\rho)(u-\rho')(u-\rho'')\cdots$, then $\psi(K)=0$ is the equation of lowest degree satisfied by $K$. If then
\[
\frac{\psi(u)-\psi(v)}{u-v}=\psi(u,v),
\]
then also
\[
\frac{\psi(u)E-\psi(K)}{uE-K}=\psi(u,K),
\]
and hence
\[
(uE-K)^{-1}=\frac{\psi(u,K)}{\psi(u)}.
\]
Therefore
\[
A=\frac{\psi(\rho,K)}{\psi'(\rho)}=g(K),
\]
where
\[
\psi'(\rho)g(v)=\psi(\rho,v)=\frac{\psi(v)}{v-\rho}=(v-\rho')(v-\rho'')\cdots.
\]
Thus $g(\rho)=1$, $g(\rho')=0$, $g(\rho'')=0,\cdots$.

Now from (3.) follows
\[
\bigl((\rho E-K)+(u-\rho)E\bigr)\left(\frac{A}{u-\rho}+\cdots\right)=E,
\]
and hence
\begin{equation}\tag{9.}
(\rho E-K)A=0.
\end{equation}
Since the determinant $|uE-K|$ has $f$ linear elementary divisors $u-\rho$, the rank of the matrix $\rho E-K$ is equal to $h-f$. By equation (9.) the rank of the matrix $A$ can therefore be at most $f$. It must be equal to $f$ if among the different columns of $A$ there are $f$ independent solutions of the linear equations
\[
(\rho e_{\alpha1}-k_{\alpha1})v_1+(\rho e_{\alpha2}-k_{\alpha2})v_2+\cdots+(\rho e_{\alpha h}-k_{\alpha h})v_h=0,
\]
that is, if it can be shown that: if $X$ is any solution of the equation $(\rho E-K)X=0$, then one can determine $Y$ so that $X=AY$. Now $KX=\rho X$ implies $K^2X=\rho KX=\rho^2X,\cdots,K^nX=\rho^nX$, and hence also $g(K)X=g(\rho)X$, if $g(v)$ is an integral function of $v$. If, as above, $(v-\rho)\psi'(\rho)g(v)=\psi(v)$, then $g(K)=A$ and $g(\rho)=1$, hence $AX=X$, which proves the assertion.

\section*{§ 5.}

In the group matrix $X=(x_{PQ^{-1}})$ I give the $h$ independent variables $x_R$ such constant values $x_R=k_R$ that the $f+f'+\cdots$ roots $\rho,\rho',\cdots$ of the equation $\Phi(k-u\eps)\Phi'(k-u\eps)\cdots=0$ are all distinct from one another. If then $\rho$ is a root of the equation of degree $f$, $\Phi(k-u\eps)=0$, then the characteristic determinant of the matrix $K=(k_{PQ^{-1}})$ has $f$ linear elementary divisors $u-\rho$. If now, expanded by powers of $u-\rho$ (or $u-\rho',\cdots$),
\begin{equation}\tag{1.}
(uE-K)^{-1}=\frac{A}{u-\rho}+\cdots=\frac{B}{u-\rho'}+\cdots,
\end{equation}
then the matrices $A,B,\cdots$ exhibit the same symmetry relations as $X$, that is, one can set $A=(a_{PQ^{-1}})$, $B=(b_{PQ^{-1}})\cdots$. Then $A^2=A$, $AB=0$, hence
\begin{equation}\tag{2.}
\sum a_Pa_Q=a_R\qquad(PQ=R)
\end{equation}
and
\begin{equation}\tag{3.}
\sum a_Pb_Q=0\qquad(PQ=R).
\end{equation}
The matrix $A$ has rank $f$ and is an integral function of $K$, $A=g(K)$; moreover $g(\rho)=1$, while for every other root $g(\rho')=0$. Thus, if $\rho,\rho_1,\cdots,\rho_{f-1}$ are the roots of the equation $\Phi(k-u\eps)=0$, then $g(\rho)=1,g(\rho_1)=0,\cdots,g(\rho_{f-1})=0$ are the roots of the equation $\Phi(a-u\eps)=0$, as I showed in \emph{Prime Factors}, § 3.

Consequently
\begin{equation}\tag{4.}
\Phi(a+u\eps)=u^{f-1}(u+1),
\qquad
\Phi'(a+u\eps)=u^{f'},
\qquad
\Phi''(a+u\eps)=u^{f''},\cdots .
\end{equation}
If $x_{PQ}=x_{QP}$, then $\Phi(x)$ is equal to the $f$-th power of a linear function of the variables $x_R$. By the first equation, therefore, the quantities $a_R$ cannot have that property if $f>1$.

Now the sum of the roots of the equation $\Phi(x-u\eps)=0$ is equal to $\sum_R\chi(R)x_R$. Hence, if $\psi$ denotes any one of the $k-1$ characters distinct from $\chi$,
\begin{equation}\tag{5.}
\sum_R\chi(R)a_R=1,
\qquad
\sum_R\psi(R)a_R=0.
\end{equation}
A system of $h$ quantities $a_R$ satisfying equations (2.) and (5.) I call a system of values determining the prime function $\Phi$, or the corresponding primitive representation of the group $\Hgr$.

If $R$ runs through the $h_\rho$ elements of the $\rho$-th class, let
\[
\sum_{(\rho)} a_R=h_\rho a_\rho.
\]
Since $\chi(R)$ has the same value $\chi_\rho$ for these $h_\rho$ elements, (5.) gives
\[
\sum_\rho h_\rho\chi_\rho a_\rho=1,
\qquad
\sum_\rho h_\rho\psi_\rho a_\rho=0.
\]
By these $k$ equations the $k$ quantities $a_\rho$ are completely determined. By \emph{Group Characters}, § 3, they are satisfied by the values
\[
a_\rho=\frac1h\chi_{\rho'}.
\]
Therefore
\begin{equation}\tag{6.}
\sum_{(\rho)}a_R=\frac{h_\rho}{h}\chi_{\rho'},
\qquad
\sum_S a_{S^{-1}RS}=\chi(R^{-1}).
\end{equation}
For example,
\begin{equation}\tag{7.}
a_E=\frac{f}{h}.
\end{equation}
The $k$ linear equations (5.) are completely equivalent to the $k$ linear equations (6.). The equations (2.) and (5.), by which we defined a system of values $a_R$ determining the prime function $\Phi$, are also equivalent to the equations (2.) and (4.). For above we obtained (5.) from (4.). But if the notation of \emph{Prime Factors}, § 1, (6.), is used, then by (2.) $a_R^{(2)}=a_R$, and hence also $a_R^{(n)}=a_R$. Now the sum of the $n$-th powers of the roots of the equation $\Phi(a-u\eps)=0$ is equal to
\[
\sum_R\chi(R)a_R^{(n)}=\sum_R\chi(R)a_R=1,
\]
and since the roots are completely determined by the power sums, one of them is $1$ and the other $f-1$ are zero. Similarly one sees that the $f'$ roots of the equation $\Phi'(a-u\eps)=0$ all vanish. Therefore the determinant of degree $h$ is
\[
|a_{PQ^{-1}}+u\eps_{PQ^{-1}}|=(u+1)^f u^{h-f}.
\]
In consequence of the equation $A^2=A$, all its elementary divisors are linear. Hence the matrix $A$ has rank $f$ and the matrix $A-E$ has rank $h-f$.

Let $H$ be a fixed element and $R$ a variable element of $\Hgr$, and set, independently of the notation used in equation (3.), $b_R=a_{H^{-1}RH}$. Then the $h$ quantities $b_R$ satisfy equations (2.) and (5.), and therefore form a system of values determining the function $\Phi$. More generally, if $v_R$ are $h$ independent variables, if $V=(v_{PQ^{-1}})$ and $V^{-1}AV=B=(b_{PQ^{-1}})$, then the $h$ quantities $b_R$ form such a system of values (\emph{Prime Factors}, § 1, (9.)). As I shall now show, every such system of values is obtained in this way, that is, by taking for $V$ any group matrix with non-vanishing determinant.

If in the group matrix $X=(x_{PQ^{-1}})$ the $h$ quantities $x_R$ are independent variables, then one can ($D.$ § 5) determine a constant matrix $L$ (whose $h^2$ elements are independent of the variables $x_R$, and whose determinant does not vanish) so that $L^{-1}XL=Z$ decomposes into $f$ matrices $Z_1,Z_2,\cdots,Z_f$ of degree $f$, agreeing in their elements, into $f'$ matrices $Z_{f+1},\cdots,Z_{f+f'}$ of degree $f'$, and so on. The $f^2+f'^2+\cdots=h$ elements of the matrices $Z_1,Z_{f+1},\cdots$ are $h$ linearly independent functions of the $h$ variables $x_R$. Therefore, if $C$ is any matrix which decomposes in the same way as $Z$, and for which $C_1=C_2=\cdots=C_f$, $C_{f+1}=\cdots=C_{f+f'}$, $\cdots$, then $LCL^{-1}$ can be obtained from $X$ by giving the variables $x_R$ certain values. The characteristic determinant of $Z_1$ is $\Phi(x-u\eps)$, that of $Z_{f+1}$ is $\Phi'(x-u\eps)$, and so on. The signs $L,Z,C$ here do not denote group matrices.

If one sets $x_R=a_R$, let $Z$ and $Z_\lambda$ pass into $C$ and $C_\lambda$. Since $A^2=A$, the elementary divisors of $|uE-A|$ are all linear, and consequently so are those of $|uE-C|$. But if a determinant decomposes into parts, its elementary divisors are the elementary divisors of the individual parts taken together. Consequently the elementary divisors of
\[
|uE_1-C_1|=(u-1)u^{f-1}
\]
are also linear. Now let $B=(b_{PQ^{-1}})$ be a second matrix whose elements satisfy conditions (2.) and (6.). Then $L^{-1}BL=D$ decomposes correspondingly into the parts $D_1,D_2,\cdots$, and $uE_1-C_1$ and $uE_1-D_1$ agree in their elementary divisors. Hence one can determine a matrix $M_1$ with non-vanishing determinant so that $M_1^{-1}C_1M_1=D_1$. If $M_2,\cdots,M_f$ agree in their elements with $M_1$, then also $M_2^{-1}C_2M_2=D_2$, and so on. In the same way one determines $M_{f+1},M_{f+2},\cdots$, and assembles the matrix $M$ from these parts. Then $LML^{-1}=K$ is a group matrix, and from $M^{-1}CM=D$ it follows that
\[
(LML^{-1})^{-1}(LCL^{-1})(LML^{-1})=LDL^{-1}
\]
or $K^{-1}AK=B$.

It is now also easy to specify the most general group matrix $V$ which transforms $A$ into $B$. Namely, if $u_R$ are $h$ independent variables and $U=(u_{PQ^{-1}})$, then
\begin{equation}\tag{8.}
V=AUB+(E-A)U(E-B).
\end{equation}
For from the equations $A^2=A$ and $B^2=B$ follows $AV=AUB$ and $VB=AUB$, hence $AV=VB$. Further, if $K$ is any given group matrix satisfying $K^{-1}AK=B$, then one can give the variables $u_R$ such values that $V=K$. For then $AK=KB$ and $(E-A)K=K(E-B)$. Thus, putting $U=K$, one obtains
\[
\begin{aligned}
V&=A(KB)+(E-A)(K(E-B))\\
&=A(AK)+(E-A)((E-A)K)\\
&=AK+(E-A)K=K.
\end{aligned}
\]
Above I proved the existence of a group matrix $K$ whose determinant does not vanish and which satisfies the condition $AK=KB$. For $U=K$, therefore, $|V|$ is non-zero, and consequently $|V|$ cannot vanish for indeterminate $u_R$.

With the help of the proposition just developed, one can derive from the quadratic and linear equations (2.) and (5.) more general linear relations among the quantities $a_R$. The group matrix
\[
C=\left(\frac{f}{h}\chi'(PQ^{-1})\right)
\]
is interchangeable with every group matrix and satisfies $C^2=C$. If therefore $AC=L=(l_{PQ^{-1}})$, then $L^2=L$. The quantities $l_R$ also satisfy equations (5.), as follows easily from \emph{Prime Factors}, § 5, (6.), or § 8. Hence one can determine a group matrix $K$ so that $A=K^{-1}ACK=K^{-1}AKC$, or, if $K^{-1}AK=M$, so that $MC=A$.

Thus $AC=MC^2=MC=A$. If $\Psi$ is a prime factor of $\Theta$ different from $\Phi$, if $\psi$ is its character and $D=\left(\frac{f'}{h}\psi'(PQ^{-1})\right)$, then $CD=0$, and therefore also $AD=(AC)D=A(CD)=0$. If $b_R$ is any system of values determining the prime function $\Psi$ and $B=(b_{PQ^{-1}})$, then $B=BD$, and hence $AB=(AC)(BD)=(AB)(CD)=0$. If $X$ is an arbitrary group matrix, one can replace $B$ by $XBX^{-1}$. Therefore also $A(XBX^{-1})=0$, and consequently
\begin{equation}\tag{9.}
AXB=0.
\end{equation}
The $h$ quantities $a_R$ therefore satisfy the equations $AC=A$, or
\begin{equation}\tag{10.}
\sum_R\chi(RS^{-1})a_R=\frac{h}{f}a_S,
\end{equation}
and $AD=0$, or
\begin{equation}\tag{11.}
\sum_R\psi(RS^{-1})a_R=0.
\end{equation}
The latter equations, as the above derivation of the relation $AD=0$ from $AC=0$ shows, are a consequence of the former. But these equations (10.), or $(E-C)A=0$, contain, since the rank of the matrix $E-C$ is $h-f^2$, exactly $h-f^2$ independent homogeneous linear relations among the $h$ quantities $a_R$. The $k$ equations (5.), however, are contained among equations (10.) and (11.) if one adds to them the non-homogeneous equation (7.). Indeed, the number of independent linear equations set up among the $h$ quantities $a_R$ is $h-f^2+1\geq k$, because $h-k=(f^2-1)+(f'^2-1)+(f''^2-1)+\cdots$.

The equations (5.) can be brought into the form (6.). In the same way the equations (10.) and (11.) can be transformed into
\begin{equation}\tag{12.}
f\sum_{(\rho)}a_{PRQ}=h_\rho\chi_{\rho'}a_{PQ},
\end{equation}
where $R$ runs through the $h_\rho$ elements of the $\rho$-th class.

\section*{§ 6.}

If $X=(x_{PQ^{-1}})$ is any group matrix, I put $\overline{X}=(x_{Q^{-1}P})$. If then $Z=XY$, then $\overline{Z}=\overline{Y}\,\overline{X}$. The two matrices $X$ and $\overline{Y}$ are interchangeable with one another. Let $u_1,\cdots,u_f$ be the roots of the equation $\Phi(x-u\eps)=0$, and let $v_1,\cdots,v_f$ be those of $\Phi(y-u\eps)=0$. If then $g(u,v)$ is an integral function of $u$ and $v$, the roots of the characteristic equation of the matrix $g(X,\overline{Y})$ are the $f^2$ quantities
\[
g(u_1,v_1),g(u_1,v_2),\cdots,g(u_1,v_f),g(u_2,v_1),\cdots,g(u_f,v_f),
\]
and the $f'^2+f''^2+\cdots$ quantities which are formed in the analogous way from the other prime factors $\Phi',\Phi'',\cdots$ of $\Theta$ (\emph{Prime Factors}, § 10).

Therefore the roots of the characteristic equation of $X\overline{A}$ are the roots $u_1,\cdots,u_f$ followed by zeros. Consequently
\begin{equation}\tag{1.}
|E+uX\overline{A}|=(1+u_1u)\cdots(1+u_fu)=\Phi(\eps+ux).
\end{equation}
The elements of the matrix $X\overline{A}=\overline{A}X$ are
\begin{equation}\tag{2.}
\begin{aligned}
x_{P,Q}&=\sum_R x_{PR^{-1}}a_{Q^{-1}R}
       =\sum_R a_{R^{-1}P}x_{RQ^{-1}}\\
      &=\sum_R a_{Q^{-1}R^{-1}P}x_R
       =\sum_R a_{R^{-1}}x_{PRQ^{-1}}.
\end{aligned}
\end{equation}
If in the matrix $X=(x_{PQ^{-1}})$ one replaces $P$ by $P^{-1}$ and $Q$ by $Q^{-1}$, then one obtains the matrix $(x_{P^{-1}Q})$. This matrix is therefore obtained from the former by permuting the rows and the columns in the same way. Hence the two matrices are similar. The same result follows from the easily proved identity
\[
(y_{PQ})(x_{PQ^{-1}})=(x_{P^{-1}Q})(y_{PQ}).
\]
If one then interchanges the rows and columns in $(x_{P^{-1}Q})$, one obtains the matrix $\overline{X}=(x_{Q^{-1}P})$. For the matrix $\overline{X}$, therefore, the totality of all minors of degree $f$ is the same as for $X$. Hence the rank of $\overline{A}$ is also equal to $f$. By \emph{Prime Factors}, § 11, the minors of degree $f$ of the matrix $X\overline{A}=\overline{A}X$ therefore stand in the same ratios as the corresponding minors of degree $f$ of the matrix $\overline{A}$. Now by (1.) the sum of all minors of degree $f$ of $X\overline{A}$ is equal to $\Phi(x)$. Hence every minor of degree $f$ of $X\overline{A}$ is equal to the corresponding minor of $\overline{A}$ multiplied by $\Phi(x)$,
\begin{equation}\tag{3.}
|x_{P,Q}|=\Phi(x)|a_{Q^{-1}P}|
\qquad(P=A_1,\cdots,A_f;\\ Q=B_1,\cdots,B_f).
\end{equation}
In the matrix of degree $f$
\[
(x_{P,Q})_f\qquad(P=A_1,\cdots,A_f;\ Q=B_1,\cdots,B_f)
\]
the variable $x_E$ is multiplied by the matrix $(a_{Q^{-1}P})_f$. If its determinant is non-zero, then, as I shall now show,
\begin{equation}\tag{4.}
(x_{P,Q})_f(a_{Q^{-1}P})_f^{-1}
\qquad(P=A_1,\cdots,A_f;\ Q=B_1,\cdots,B_f)
\end{equation}
is a matrix belonging to the group $\Hgr$.

Putting $E-A=B$, the rank of $B$ is equal to $h-f=g$, and
\[
\sum a_Pa_Q=a_R,
\qquad
\sum b_Pb_Q=b_R,
\qquad
\sum a_Pb_Q=\sum b_Pa_Q=0\qquad(PQ=R).
\]
Let the two determinants of degrees $f$ and $g$
\begin{equation}\tag{5.}
|a_{Q^{-1}P}|\\[-0.5em]
\qquad(P=A_1,\cdots,A_f;\ Q=B_1,\cdots,B_f)
\end{equation}
and
\[
|b_{Q^{-1}P}|
\qquad(P=C_1,\cdots,C_g;\ Q=D_1,\cdots,D_g)
\]
be non-zero. Let $M$ be the matrix of degree $h$ whose rows are obtained from the row
\[
a_{B_1^{-1}R},\cdots,a_{B_f^{-1}R},b_{D_1^{-1}R},\cdots,b_{D_g^{-1}R}
\]
by letting $R$ run through the $h$ elements of $\Hgr$. Let $L'$ be the matrix whose rows are obtained in the same way from
\[
a_{R^{-1}A_1},\cdots,a_{R^{-1}A_f},b_{R^{-1}C_1},\cdots,b_{R^{-1}C_g},
\]
and let $L$ be the matrix conjugate to $L'$.

I now form the matrix $LXM$. Let $\alpha$ and $\beta$ be two of the indices $1,2,\cdots,f$, and $\gamma$ and $\delta$ two of the indices $1,2,\cdots,g$. Then in that matrix the $\beta$-th element of the $\alpha$-th row is
\[
\sum_{R,S}a_{R^{-1}A_\alpha}x_{RS^{-1}}a_{B_\beta^{-1}S}
=\sum_{R,S}a_{R^{-1}P}x_{RS^{-1}}a_{Q^{-1}S},
\]
if one sets $A_\alpha=P$, $B_\beta=Q$. This is an element of the matrix $\overline{A}X\overline{A}=X\overline{A}\,\overline{A}=X\overline{A}$, hence equal to $x_{P,Q}$. The $(f+\delta)$-th element of the $\alpha$-th row, if $A_\alpha=P$, $D_\delta=Q$, is
\[
\sum_{R,S}a_{R^{-1}P}x_{RS^{-1}}b_{Q^{-1}S}.
\]
This is an element of the matrix $\overline{A}X\overline{B}=X\overline{A}\,\overline{B}=0$, and therefore vanishes. Similarly the $\beta$-th element of the $(f+\gamma)$-th row vanishes. Finally the $(f+\delta)$-th element of the $(f+\gamma)$-th row, if $C_\gamma=P$, $D_\delta=Q$, is
\[
\sum_{R,S}b_{R^{-1}P}x_{RS^{-1}}b_{Q^{-1}S}.
\]
This is an element of the matrix $\overline{B}X\overline{B}=X\overline{B}\,\overline{B}=X\overline{B}=X-X\overline{A}$, hence equal to $x_{PQ^{-1}}-x_{P,Q}$.

The elements of the matrix $LXM$ are linear functions of the $h$ variables $x_R$. It can therefore be regarded as a linear combination of $h$ constant matrices, each of which is multiplied by one of the $h$ variables $x_R$. In particular, $x_E$ is multiplied by the matrix
\[
\begin{pmatrix}
a_{B_\beta^{-1}A_\alpha}&0\\
0&b_{D_\delta^{-1}C_\gamma}
\end{pmatrix}=N,
\]
whose determinant does not vanish. The matrix $LXMN^{-1}$ decomposes into a matrix of degree $f$
\[
(x_{P,Q})(a_{Q^{-1}P})^{-1}
\qquad(P=A_1,\cdots,A_f;\ Q=B_1,\cdots,B_f)
\]
and a matrix of degree $g$
\[
(x_{PQ^{-1}}-x_{P,Q})(\eps_{PQ^{-1}}-a_{Q^{-1}P})^{-1}
\qquad(P=C_1,\cdots,C_g;\ Q=D_1,\cdots,D_g).
\]
If in this representation one gives the variables $x_R$ the values $\eps_R$, and so sets $X=E$, then one obtains $LMN^{-1}=E$. Therefore the determinants of $L$ and $M$ are non-zero, and $MN^{-1}=L^{-1}$.

But now $LXL^{-1}$ is a matrix belonging to the group $\Hgr$. Since it decomposes into two matrices of degrees $f$ and $g$, each of the two parts has the same property. Hence (4.) is a matrix belonging to the group $\Hgr$, whose determinant is equal to $\Phi(x)$.

Thus if $\Phi$ is one of the $k$ prime factors of the determinant of the group $\Hgr$, one only needs to know any system of values $a_R$ determining that function in order to be able to give the primitive representation of the group $\Hgr$ by linear substitutions corresponding to the prime function $\Phi$.
% ===== END WORK 059 =====

\clearpage

% ===== BEGIN WORK 060: 060_ber_die_Charaktere_der_symmetrischen_Grup =====
\frobtitle{On the characters of the symmetric group}

All permutations of $n$ symbols $1,2,\cdots,n$ are the elements of a group $\Hgr$ of order $h=n!$, which is called the \emph{symmetric} group. A permutation can always, and in only one way, be represented as a product of cyclic interchanges of which no two have a symbol in common. In order that two permutations be conjugate with respect to $\Hgr$, it is necessary and sufficient that they be composed of equally many cycles of the same orders. If the permutations of the $\rho^{\ten}$ class $(\rho=0,1,\cdots,k-1)$ consist of $\alpha$ cycles of order $1$, $\beta$ cycles of order $2$, $\gamma$ cycles of order $3$, $\cdots$, then their number, according to \textsc{Cauchy}, is
\begin{equation}\tag{1.}
 h_\rho=\frac{n!}{1^\alpha\alpha!\,2^\beta\beta!\,3^\gamma\gamma!\cdots}.
\end{equation}
The number $k$ of classes of conjugate elements into which the $h$ permutations decompose is equal to the number of solutions of the equation
\begin{equation}\tag{2.}
 n=\alpha+2\beta+3\gamma+\cdots
\end{equation}
in positive integers $\alpha,\beta,\gamma,\cdots$, which may also be zero, or also equal to the number of the different decompositions of the number
\begin{equation}\tag{3.}
 n=n_1+n_2+\cdots
\end{equation}
into an indeterminate number of positive parts.

To compute the $k$ characters of the group $\Hgr$
\[
 \chi_\rho^{(\kappa)}\qquad(\rho=0,1,\cdots,k-1;\ \kappa=0,1,\cdots,k-1)
\]
I use $k$ definite subgroups of $\Hgr$. After the number $n$ has been decomposed into the summands (3.), one forms all permutations which interchange only the first $n_1$ symbols among themselves, likewise only the following $n_2$ symbols, and so on. These permutations form a group $\Ggr$ of order
\begin{equation}\tag{4.}
 g=n_1!\,n_2!\cdots.
\end{equation}
Every permutation $R$ of $\Ggr$ decomposes into a permutation $R_1$ among the first $n_1$ symbols, a permutation $R_2$ among the following $n_2$ symbols, and so on. If $R_\nu$ consists of $\alpha_\nu$ cycles of order $1$, $\beta_\nu$ cycles of order $2$, $\cdots$, then
\begin{equation}\tag{5.}
 n_1=\alpha_1+2\beta_1+\cdots,\\qquad n_2=\alpha_2+2\beta_2+\cdots,\quad\cdots,
\end{equation}
and, if $R$ belongs in $\Hgr$ to the $\rho^{\ten}$ class,
\begin{equation}\tag{6.}
 \alpha=\alpha_1+\alpha_2+\cdots,\qquad \beta=\beta_1+\beta_2+\cdots.
\end{equation}
The number of permutations of this class contained in $\Ggr$ is therefore
\[
 g_\rho=\sum \frac{n_1!}{1^{\alpha_1}\alpha_1!\,2^{\beta_1}\beta_1!\cdots}
           \frac{n_2!}{1^{\alpha_2}\alpha_2!\,2^{\beta_2}\beta_2!\cdots}\cdots,
\]
where the sum extends over all solutions of equations (5.) and (6.). If these equations have no solution, then $g_\rho=0$. Otherwise, since
\begin{equation}\tag{7.}
 \frac{h}{h_\rho}=1^\alpha\alpha!\,2^\beta\beta!\,3^\gamma\gamma!\cdots
\end{equation}
one has
\[
 \frac{h g_\rho}{g h_\rho}=\sum
 \frac{\alpha!}{\alpha_1!\alpha_2!\cdots}
 \frac{\beta!}{\beta_1!\beta_2!\cdots}\cdots.
\]
Let $x_1,x_2,\cdots,x_m$ be independent variables, whose number $m$ is not smaller than the number of the parts $n_1,n_2,\cdots$ of $n$. Then this sum, extended over all solutions of equations (5.) and (6.), is equal to the coefficient of $x_1^{n_1}x_2^{n_2}\cdots$ in the expansion of the product
\begin{equation}\tag{8.}
 (x_1+x_2+\cdots+x_m)^\alpha
 (x_1^2+x_2^2+\cdots+x_m^2)^\beta
 (x_1^3+x_2^3+\cdots+x_m^3)^\gamma\cdots
\end{equation}
according to powers of the variables.

This result, which is fundamental for the following investigation, can also be seen without any calculation. Let $R=C_1C_2\cdots C_s$ be a permutation of the $\rho^{\ten}$ class which consists of $s$ cycles on respectively $c_1,c_2,\cdots,c_s$ symbols, so that
\begin{equation}\tag{9.}
 n=c_1+c_2+\cdots+c_s
\end{equation}
It follows from the developments in § 1 of my paper \emph{On relations between the characters of a group and those of its subgroups}, Sitzungsberichte 1898, which I shall cite in what follows as $U.$, that the number $\dfrac{h g_\rho}{g h_\rho}$ indicates how many of the $\dfrac{h}{g}$ subgroups of $\Hgr$ which are conjugate with $\Ggr$ contain the permutation $R$. One finds these groups by dividing the $n$ symbols, in all possible ways, into systems of respectively $n_1,n_2,\cdots$ symbols, and by permuting the symbols of each of these systems among themselves in every way. The number of groups so obtained is
\[
 \frac{h}{g}=\frac{n!}{n_1!n_2!\cdots},
\]
they need not all be distinct from one another. In order that $R$ occur in one of these groups, it must be possible to form, from some of the cycles $C_1,C_2,\cdots,C_s$, a permutation of $n_1$ symbols, from some of the others a permutation of $n_2$ symbols, and so on. This can be done in as many ways as the numbers $n_1,n_2,\cdots$ can be represented as sums of the numbers $c_1,c_2,\cdots,c_s$. But this number is equal to the coefficient of $x_1^{n_1}x_2^{n_2}\cdots$ in the expansion of the product
\begin{equation}\tag{10.}
 (x_1^{c_1}+x_2^{c_1}+\cdots+x_m^{c_1})(x_1^{c_2}+x_2^{c_2}+\cdots+x_m^{c_2})\cdots
 (x_1^{c_s}+x_2^{c_s}+\cdots+x_m^{c_s}),
\end{equation}
which, apart from the chosen notation, agrees with the expression (8.). If among the numbers $n_1,n_2,\cdots$, or among the numbers $c_1,c_2,\cdots$, several are equal to one another, then two decompositions are also to be regarded as distinct when they differ only in the designation of the numbers occurring in them. For even if $c_1=c_2$, $C_1$ is nevertheless distinct from $C_2$.

By the formula $U.$ § 1, (9.), $\dfrac{h g_\rho}{g h_\rho}$ can be represented as a linear combination of the characters of $\Hgr$ with integral coefficients; or, expressed otherwise, the system of the $k$ numbers $\dfrac{h g_\rho}{g h_\rho}$ is divisible by the module formed by the $k^2$ values of the $k$ characters. I shall call every such combination a composite character of $\Hgr$. In the expansion of the integral symmetric function of degree $n$ in (8.) according to powers of $x_1,x_2,\cdots,x_m$, every coefficient is therefore a composite character of $\Hgr$, that is, the value of such a character for the class $(\rho)$ determined by the exponents $\alpha,\beta,\gamma,\cdots$.

\section*{§ 2.}

Let the product of all differences of the $m$ variables $x_1,x_2,\cdots,x_m$ be
\[
 \Delta(x_1,x_2,\cdots,x_m)=(x_2-x_1)(x_3-x_1)(x_3-x_2)\cdots(x_m-x_{m-1}).
\]
In the expansion of this function of degree $\frac12m(m-1)$, each term has the form
\[
 \pm x_1^{k_1}x_2^{k_2}\cdots x_m^{k_m},
\]
where $k_1,k_2,\cdots,k_m$ are the numbers $0,1,\cdots,m-1$ in some order. According as this permutation is positive or negative, the coefficient of the term is equal to $+1$ or to $-1$. Even when $k_1,k_2,\cdots,k_m$ are any $m$ positive or negative numbers, I shall use in what follows the symbol
\begin{equation}\tag{1.}
 [k_1,k_2,\cdots,k_m]=\sign.\,\Delta(k_1,k_2,\cdots,k_m).
\end{equation}
It has the value $0$ when two of the $m$ numbers are equal; otherwise it has the value $+1$ or $-1$ according as, among the differences $k_2-k_1,k_3-k_1,k_3-k_2,\cdots$, an even or an odd number are negative, that is, according as the sequence $k_1,k_2,\cdots,k_m$ contains an even or an odd number of inversions.

Then in the product
\begin{equation}\tag{2.}
 (x_1+x_2+\cdots+x_m)^\alpha(x_1^2+x_2^2+\cdots+x_m^2)^\beta(x_1^3+x_2^3+\cdots+x_m^3)^\gamma\cdots\Delta(x_1,x_2,\cdots,x_m)
\end{equation}
every coefficient is an integral linear combination of certain coefficients of the expression (8.), § 1, and hence is likewise a composite character of $\Hgr$. How I arrived at transforming the symmetric function (8.), § 1 into an alternating one by multiplication with the product of differences will become clear from the following proof. In this alternating function I denote the coefficient of
\[
 x_1^{\lambda_1}x_2^{\lambda_2}\cdots x_m^{\lambda_m}\quad\text{with}\quad [\lambda_1,\lambda_2,\cdots,\lambda_m]\chi_\rho^{(\lambda)}.
\]
The symbol $(\lambda)$ denotes the system of the $m$ exponents $\lambda_1,\lambda_2,\cdots,\lambda_m$, whose sum is
\begin{equation}\tag{3.}
 \lambda_1+\lambda_2+\cdots+\lambda_m=n+\frac12m(m-1)
\end{equation}
and the symbol $(\rho)$ denotes the class of conjugate elements of $\Hgr$ determined by the numbers $\alpha,\beta,\gamma,\cdots$ of the decomposition (2.), § 1. As with the classes, I imagine the different solutions of equation (3.), where the order of the summands is immaterial, to be numbered quite arbitrarily and denoted by $(0),(1),\cdots$, and in general by $(\lambda)$.

For a fixed $\lambda$, the $k$ numbers
\[
 \chi_\rho^{(\lambda)}\qquad(\rho=0,1,\cdots,k-1)
\]
form, as I now show, a simple character of $\Hgr$, and for different $\lambda$ they form different characters. If $m=n$, there are precisely $k$ systems, each of $n$ different numbers, satisfying the condition
\begin{equation}\tag{4.}
 \kappa_1+\kappa_2+\cdots+\kappa_n=\frac12 n(n+1)
\end{equation}
For if, in a solution, one orders the summands so that
\begin{equation}\tag{5.}
 \kappa_1<\kappa_2<\cdots<\kappa_n
\end{equation}
then to it there corresponds a decomposition of the number $n$ into the parts
\[
 \kappa_1\leqq \kappa_2-1\leqq \kappa_3-2\leqq\cdots\leqq\kappa_n-n+1
\]
and conversely. Hence the alternating function
\begin{equation}\tag{6.}
\begin{aligned}
 &(x_1+x_2+\cdots+x_n)^\alpha(x_1^2+x_2^2+\cdots+x_n^2)^\beta(x_1^3+x_2^3+\cdots+x_n^3)^\gamma\cdots\Delta(x_1,x_2,\cdots,x_n)\\
 &=\sum_{(\kappa)}[\kappa_1,\kappa_2,\cdots,\kappa_n]\chi_\rho^{(\kappa)}x_1^{\kappa_1}x_2^{\kappa_2}\cdots x_n^{\kappa_n}
\end{aligned}
\end{equation}
gives the values of the $k$ characters of the symmetric group $\Hgr$. One already obtains them all by considering, in this sum of $n!k$ terms, only the $k$ terms in which the exponents satisfy condition (5.).

Between the characters $\psi_\rho^{(\kappa)}$ of a group $\Hgr$ there are the relations
\begin{equation}\tag{7.}
 \sum_\rho h_\rho\psi_\rho^{(\kappa)}\psi_\rho^{(\kappa')}=h,
 \qquad
 \sum_\rho h_\rho\psi_\rho^{(\kappa)}\psi_\rho^{(\lambda)}=0,
\end{equation}
provided $(\lambda)$ is different from the character $(\kappa')$ inverse to $(\kappa)$, that is, the conjugate-complex character.  Now let
\[
 \chi_\rho=\sum_\kappa r_\kappa\psi_\rho^{(\kappa)}
\]
be a composite character, so that the coefficients $r_\kappa$ are positive or negative integers. Then, because $\psi_\rho^{(\kappa')}=\psi_\rho^{(\kappa)}$,
\[
 \sum_\rho h_\rho\chi_\rho\chi_\rho'=h\sum_\kappa r_\kappa^2.
\]
If one can therefore show that this sum is equal to $h$, then one of the numbers $r_\kappa$ must be equal to $\pm1$, and the other $k-1$ must be equal to $0$. If moreover $\chi_0$ is a positive integer, where $(0)$ is the principal class, it follows that $\chi_\rho$ is a simple character. If several such characters $\chi_\rho^{(\kappa)},\chi_\rho^{(\lambda)},\cdots$ are known and
\[
 \sum_\rho h_\rho\chi_\rho^{(\kappa)}\chi_\rho^{(\lambda)}=0,
\]
then they are distinct from one another.

For the symmetric group these formulae are further simplified by the fact that the characters $\chi_\rho^{(\lambda)}$ occurring in formula (6.) have only real values, and are therefore equal to their inverse characters.

\section*{§ 3.}

To prove the assertion just stated, I introduce a second system of $m$ independent variables $y_1,y_2,\cdots,y_m$ and form the sum
\begin{equation}\tag{1.}
\begin{aligned}
 &\sum_\rho \frac{h_\rho}{h}(x_1+\cdots+x_m)^\alpha(x_1^2+\cdots+x_m^2)^\beta\cdots
 (y_1+\cdots+y_m)^\alpha(y_1^2+\cdots+y_m^2)^\beta\cdots\\
 &=\sum \frac{1}{1^\alpha\alpha!}(x_1+\cdots+x_m)^\alpha(y_1+\cdots+y_m)^\alpha
 \frac{1}{2^\beta\beta!}(x_1^2+\cdots+x_m^2)^\beta(y_1^2+\cdots+y_m^2)^\beta\cdots .
\end{aligned}
\end{equation}
At first one does not restrict the exponents $\alpha,\beta,\gamma,\cdots$ by the condition
\begin{equation}\tag{2.}
 \alpha+2\beta+3\gamma+\cdots=n,
\end{equation}
but gives each of them the values from $0$ to $\infty$. Then the sum of this series is
\[
 e^{(x_1+\cdots+x_m)(y_1+\cdots+y_m)+\frac12(x_1^2+\cdots+x_m^2)(y_1^2+\cdots+y_m^2)+\cdots}.
\]
If one performs the multiplication in each term, the exponent becomes
\[
 -l(1-x_1y_1)-l(1-x_1y_2)-l(1-x_2y_1)-\cdots-l(1-x_my_m).
\]
Consequently the sum is equal to the reciprocal value of the product
\[
 (1-x_1y_1)(1-x_1y_2)(1-x_2y_1)\cdots(1-x_my_m).
\]
If, therefore, the expression under consideration is further multiplied by
\[
 \Delta(x_1,x_2,\cdots,x_m)\Delta(y_1,y_2,\cdots,y_m),
\]
then, by a well-known formula of \textsc{Cauchy}, it becomes the determinant of degree $m$
\[
 \left|\frac{1}{1-x_\mu y_\nu}\right|\qquad(\mu,\nu=1,2,\cdots,m),
\]
or, written out,
\[
 \sum[\mu_1,\mu_2,\cdots,
\mu_m]\frac{1}{(1-x_{\mu_1}y_1)(1-x_{\mu_2}y_2)\cdots(1-x_{\mu_m}y_m)}.
\]
Now
\[
 \frac{1}{1-xy}=\sum_{\lambda=0}^\infty x^\lambda y^\lambda,
\]
and hence the determinant is equal to
\[
 \sum_\lambda\sum_\mu[\mu_1,\mu_2,\cdots,
\mu_m]x_{\mu_1}^{\lambda_1}y_1^{\lambda_1}x_{\mu_2}^{\lambda_2}y_2^{\lambda_2}\cdots x_{\mu_m}^{\lambda_m}y_m^{\lambda_m},
\]
or, after the factors are suitably rearranged in every term,
\begin{equation}\tag{3.}
 \sum_{\kappa,\lambda}[\kappa_1,\kappa_2,\cdots,
\kappa_m][\lambda_1,\lambda_2,\cdots,
\lambda_m]x_1^{\kappa_1}x_2^{\kappa_2}\cdots x_m^{\kappa_m}y_1^{\lambda_1}y_2^{\lambda_2}\cdots y_m^{\lambda_m}.
\end{equation}
Each of the $m$ exponents $\lambda_1,\lambda_2,\cdots,\lambda_m$ runs through the values from $0$ to $\infty$. If two of these exponents are equal, the term has coefficient $0$. For every system of $m$ distinct exponents, $\kappa_1,\kappa_2,\cdots,\kappa_m$ run through the $m!$ permutations of $\lambda_1,\lambda_2,\cdots,\lambda_m$.

If in the sum (1.) one takes only the terms which satisfy condition (2.), then in the series (3.) one has to restrict oneself to the terms in which
\begin{equation}\tag{4.}
 \lambda_1+\lambda_2+\cdots+\lambda_m=n+\frac12m(m-1)
\end{equation}
Therefore, by (6.), § 2, this finite series is equal to
\[
 \sum_{\kappa,\lambda}\left(\sum_\rho\frac{h_\rho}{h}\chi_\rho^{(\kappa)}\chi_\rho^{(\lambda)}\right)
 [\kappa_1,\cdots,
\kappa_m][\lambda_1,\cdots,
\lambda_m]x_1^{\kappa_1}\cdots x_m^{\kappa_m}y_1^{\lambda_1}\cdots y_m^{\lambda_m},
\]
where the summation is over all systems of values $\kappa,\lambda$ which satisfy (4.) and the analogous condition
\[
 \kappa_1+\kappa_2+\cdots+\kappa_m=n+\frac12m(m-1)
\]
Thus, if $\kappa_1,\kappa_2,\cdots,\kappa_m$ are not obtained from $\lambda_1,\lambda_2,\cdots,\lambda_m$ by a mere change of order, then
\[
 \sum_\rho h_\rho\chi_\rho^{(\kappa)}\chi_\rho^{(\lambda)}=0,
\]
in the other case the sum is equal to $h$. With this the essential part of the proof is completed, and it remains only to show that the numbers $\chi_0^{(\lambda)}=f^{(\lambda)}$ are positive.

For the principal class (0) one has $\alpha=n$, $\beta=0$, $\gamma=0$, $\cdots$, and hence
\begin{equation}\tag{5.}
 (x_1+x_2+\cdots+x_m)^n\Delta(x_1,x_2,\cdots,x_m)=\sum[\lambda_1,\cdots,
\lambda_m]f^{(\lambda)}x_1^{\lambda_1}\cdots x_m^{\lambda_m}.
\end{equation}
The left hand side is equal to
\[
 \left(\sum \frac{n!}{\mu_1!\cdots\mu_m!}x_1^{\mu_1}\cdots x_m^{\mu_m}\right)
 \left(\sum[\kappa_1,\cdots,
\kappa_m]x_1^{\kappa_1}\cdots x_m^{\kappa_m}\right),
\]
where $\mu_1+\cdots+\mu_m=n$ and $\kappa_1,\kappa_2,\cdots,
\kappa_m$ denotes one of the $m!$ permutations of the numbers $0,1,\cdots,m-1$. The function $1/\mu!$ is to be taken as zero when $\mu$ is negative. The product of the two sums is equal to $n!$ times
\[
 \sum[\kappa_1,
\cdots,
\kappa_m]\frac{1}{(\lambda_1-
\kappa_1)!}\cdots\frac{1}{(\lambda_m-
\kappa_m)!}x_1^{\lambda_1}\cdots x_m^{\lambda_m}.
\]
In this sum the coefficient of $x_1^{\lambda_1}\cdots x_m^{\lambda_m}$ is equal to the determinant of degree $m$
\[
 \frac{f^{(\lambda)}}{h}=\left|\frac{1}{(\lambda_\mu-
u+1)!}\right|\qquad(\mu,
\nu=1,2,
\cdots,m).
\]
Multiplying the elements of the $\mu^{\ten}$ row by $\lambda_\mu!$ and dividing those of the $\nu^{\ten}$ column by $(\nu-1)!$, one obtains
\[
 \frac{\lambda_1!\cdots\lambda_m!}{0!\cdots(m-1)!}\frac{f^{(\lambda)}}{h}=\left|\binom{\lambda_\mu}{\nu-1}\right|.
\]
I use the symbol
\[
 \binom{x}{n}=\frac{x(x-1)\cdots(x-n+1)}{1\cdot2\cdots n}
\]
also in the case where $x$ is a variable. The rows of the last determinant are the values of the functions
\[
 \binom{x}{0},\binom{x}{1},\cdots,\binom{x}{m-1}\quad\text{for}\quad x=\lambda_1,\lambda_2,\cdots,
\lambda_m.
\]
The degrees of these functions are $0,1,\cdots,m-1$. If $\lambda_1,\lambda_2,\cdots,\lambda_m$ are independent variables, then this determinant is equal to $\Delta(\lambda_1,\lambda_2,\cdots,\lambda_m)$ up to a constant divisor $\Delta(0,1,\cdots,m-1)$, which one obtains by putting $\lambda_1=0,\lambda_2=1,\cdots,\lambda_m=m-1$. Since
\[
 \Delta(0,1,
\cdots,m-1)=0!1!\cdots(m-1)!
\]
one has
\begin{equation}\tag{6.}
 f^{(\lambda)}=\frac{n!\Delta(\lambda_1,
\lambda_2,
\cdots,
\lambda_m)}{\lambda_1!\lambda_2!\cdots\lambda_m!},
\end{equation}
and this is positive when $\lambda_1<\lambda_2<\cdots<\lambda_m$.


\section*{§ 4.}

In order to obtain the $k$ different characters, we must put $m=n$. Then $\chi^{(\kappa)}$ is characterized by a system of $n$ numbers
\begin{equation}\tag{1.}
 (\kappa)=\kappa_1,
\kappa_2,
\cdots,
\kappa_n
\end{equation}
which satisfy the conditions
\begin{equation}\tag{2.}
 \kappa_1+
\kappa_2+\cdots+
\kappa_n=\frac12n(n+1),\qquad 0\leqq\kappa_1<\kappa_2<\cdots<\kappa_n
\end{equation}
As regards the relation between the two developments (2.) and (6.), § 2, one has
\[
 [\lambda_1,
\lambda_2,
\cdots,
\lambda_m]\chi_\rho^{(\lambda)}=[\kappa_1,
\kappa_2,
\cdots,
\kappa_n]\chi_\rho^{(\kappa)},
\]
if, in the case $m<n$, $(\kappa)$ is the system of the $n$ numbers
\begin{equation}\tag{3.}
 \kappa_1,
\kappa_2,
\cdots,
\kappa_n=0,1,
\cdots,
\mu-1,
\mu+
\lambda_1,
\cdots,
\mu+
\lambda_m\qquad(m+
\mu=n)
\end{equation}
Conversely, in the case $m>n$,
\begin{equation}\tag{4.}
 \lambda_1,
\lambda_2,
\cdots,
\lambda_m=0,1,
\cdots,
\mu-1,
\mu+
\kappa_1,
\cdots,
\mu+
\kappa_n\qquad(n+
\mu=m).
\end{equation}
Thus, in order to be independent of $m$, one could use, for the definition of the character $\chi$, those among the numbers
\[
 \lambda_1\leqq\lambda_2-1\leqq\lambda_3-2\leqq\cdots\leqq\lambda_m-m+1
\]
which are different from zero. It proves more practical, however, to use the following characteristic for this purpose.

The $n$ numbers (1.) are all $\leqq2n-1$, for $\kappa_n$ reaches the maximal value $2n-1$ when the other $n-1$ numbers have the minimal values $\kappa_1=0,
\kappa_2=1,
\cdots,
\kappa_{n-1}=n-2$. Let $a_1,
 a_2,
\cdots,
 a_n$ be the numbers $0,1,
\cdots,n-1$ in some order. Then
\[
 n-1-a_1,
\qquad n-1-a_2,
\cdots,
 n-1-a_n
\]
are the same numbers. Of these, let $n-r$ occur among the numbers (1.), for instance $n-1-a_{r+1},
\cdots,n-1-a_n$. The remaining $r$ of the numbers (1.) are $\geqq n$, and are to be denoted by $n+b_1,
\cdots,n+b_r$. I then call
\begin{equation}\tag{5.}
 (\kappa)=\left(\begin{matrix}a_1&a_2&\cdots&a_r\\ b_1&b_2&\cdots&b_r\end{matrix}\right)
\end{equation}
the characteristic of $\chi^{(\kappa)}$. In order that it be completely determined, suppose that
\begin{equation}\tag{6.}
 a_1<a_2<\cdots<a_r,\qquad b_1<b_2<\cdots<b_r,
\end{equation}
holds. The equation (2.) now reads
\[
 (n-1-a_{r+1})+\cdots+(n-1-a_n)+(n+b_1)+\cdots+(n+b_r)=\frac12n(n+1).
\]
Furthermore
\[
 a_1+\cdots+a_r+a_{r+1}+\cdots+a_n=0+1+\cdots+(n-1)=\frac12n(n-1),
\]
and hence
\begin{equation}\tag{7.}
 a_1+a_2+\cdots+a_r+b_1+b_2+\cdots+b_r=n-r.
\end{equation}
The left-hand side is $\geqq2(0+1+\cdots+(r-1))=r(r-1)$, and consequently the number $r$, which I call the rank of the characteristic (5.), is
\begin{equation}\tag{8.}
 r\leqq\sqrt n.
\end{equation}
The degree of the character $\chi^{(\kappa)}$ now becomes
\begin{equation}\tag{9.}
 f^{(\kappa)}=\frac{n!\Delta(\kappa_1,
\kappa_2,
\cdots,
\kappa_n)}{\kappa_1!\kappa_2!\cdots\kappa_n!}
 =\frac{n!\Delta(a_1,
\cdots,a_r)\Delta(b_1,
\cdots,b_r)}{a_1!\cdots a_r!\,b_1!\cdots b_r!\prod(a_\alpha+b_\beta+1)},
\end{equation}
where in the last product $\alpha$ and $\beta$ run through the values $1,2,\cdots,r$.

In carrying out this transformation we shall, since $f^{(\kappa)}$ is positive, disregard the signs. Then it is to be shown that
\[
\begin{aligned}
 &a_1!\cdots a_r!\,b_1!\cdots b_r!\prod(a_\alpha+b_\beta+1)
 \Delta(-a_{r+1}-1,\cdots,-a_n-1,b_1,\cdots,b_r)\\
 &=\pm (n-a_{r+1}-1)!\cdots(n-a_n-1)!(n+b_1)!\cdots(n+b_r)!
 \Delta(a_1,\cdots,a_r)\Delta(b_1,\cdots,b_r)
\end{aligned}
\]
holds. Now
\[
 \frac{(b_\beta+n)!}{b_\beta!}=\prod_\alpha(a_\alpha+b_\beta+1)\prod_\gamma(a_\gamma+b_\beta+1),
\]
where $\gamma$ runs through the values $r+1,
\cdots,n$, and
\[
 \pm\Delta(-a_{r+1}-1,
\cdots,-a_n-1,b_1,
\cdots,b_r)=\Delta(a_{r+1},
\cdots,a_n)\Delta(b_1,
\cdots,b_r)
 \prod_{\gamma,
\beta}(a_\gamma+b_\beta+1).
\]
It remains, therefore, to prove that
\[
 a_1!\cdots a_r!\Delta(a_{r+1},\cdots,a_n)=(n-1-a_{r+1})!\cdots(n-1-a_n)!\Delta(a_1,
\cdots,a_r)
\]
holds. Now, if in the products the factor $a_\gamma-a_\alpha$ is omitted,
\[
 (a_\alpha-a_1)(a_\alpha-a_2)\cdots(a_\alpha-a_n)=(a_\alpha-0)(a_\alpha-1)\cdots(a_\alpha-n+1)=\pm a_\alpha!(n-1-a_\alpha)!.
\]
If the above equation is therefore multiplied by $(n-1-a_1)!\cdots(n-1-a_r)!$, one obtains
\[
 \Delta(a_1,\cdots,a_r)^2\prod(a_\alpha-a_\gamma)\Delta(a_{r+1},\cdots,a_n)
 =\pm0!1!\cdots(n-1)!\Delta(a_1,\cdots,a_r).
\]
This formula is obtained from the equation
\[
\begin{aligned}
0!1!\cdots(n-1)!&=\Delta(0,1,\cdots,n-1)\\
&=\pm\Delta(a_1,\cdots,a_r,a_{r+1},\cdots,a_n)\\
&=\pm\Delta(a_1,\cdots,a_r)\Delta(a_{r+1},\cdots,a_n)\prod(a_\gamma-a_\alpha).
\end{aligned}
\]
The expression (9.) can be brought into the form of a determinant of degree $r$,
\begin{equation}\tag{9.}
 f^{(\kappa)}=\frac{n!}{a_1!\cdots a_r!\cdots b_1!\cdots b_r!}
 \left|\frac{1}{a_\alpha+b_\beta+1}\right|\qquad(\alpha,
\beta=1,2,
\cdots,r)
\end{equation}
.

The characteristic (5.) can also be derived from the $m$ numbers which above were denoted by $\lambda_1,
\lambda_2,
\cdots,
\lambda_m$. Of these, $r$ are $\geqq m$, and equal to $m+b_1,m+b_2,
\cdots,m+b_r$; and, among the numbers $0,1,
\cdots,m-1$, the $r$ numbers $m-1-a_1,m-1-a_2,m-1-a_r$ are missing from them. In order to compute $\chi_\rho^{(\lambda)}$ by means of the development of the expression (2.), § 2, it is only necessary to choose $m\geqq a_r+1$ there.

The character corresponding to the characteristic (5.) I shall also denote by
\[
 \chi_\rho^{(\kappa)}=\chi\left(\begin{matrix}a_1&a_2&\cdots&a_r\\ b_1&b_2&\cdots&b_r\end{matrix}\right)_\rho
\]
If, for example, in formula (2.), § 2, one puts $m=2$, then one obtains
\begin{equation}\tag{10.}
\begin{aligned}
 &(1-x)(1+x)^\alpha(1+x^2)^\beta(1+x^3)^\gamma\cdots\\
 &=1+x\chi\left(\begin{matrix}1\\ n-2\end{matrix}\right)
 +x^2\chi\left(\begin{matrix}0&1\\0&n-3\end{matrix}\right)
 +x^3\chi\left(\begin{matrix}0&1\\1&n-4\end{matrix}\right)+\cdots\\
 &\quad +x^\lambda\chi\left(\begin{matrix}0&1\\ \lambda-2&n-\lambda-1\end{matrix}\right)+\cdots,
\end{aligned}
\end{equation}
provided that $\lambda<\frac12(n+1)$. If $n+1$ is even and $\lambda=\frac12(n+1)$, then the coefficient of $x^\lambda$ is equal to $0$. Then the same coefficients recur in the reverse order, but with the opposite signs. Hence, for example,
\begin{equation}\tag{11.}
\begin{aligned}
\chi\left(\begin{matrix}1\\ n-2\end{matrix}\right)&=\alpha-1,\qquad
\chi\left(\begin{matrix}0&1\\0&n-3\end{matrix}\right)=\frac12\alpha(\alpha-3)+\beta,\\
\chi\left(\begin{matrix}0&1\\1&n-4\end{matrix}\right)&=\frac16\alpha(\alpha-1)(\alpha-5)+(\alpha-1)\beta+
\gamma,\\
\chi\left(\begin{matrix}0&1\\2&n-5\end{matrix}\right)&=\frac1{24}\alpha(\alpha-1)(\alpha-2)(\alpha-7)
 +\frac12(\alpha-1)(\alpha-2)\beta+\frac12\beta(\beta-3)+(\alpha-1)\gamma+
\delta.
\end{aligned}
\end{equation}
The last expression accordingly has the value $0$ for all solutions of the equation $\alpha+2\beta+3\gamma+4\delta=7$.

\section*{§ 5.}

The characters of first rank can be computed in an especially simple way. To them belongs the character $\chi(R)=\alpha-1$, which, by ($U.$, § 5), belongs to every doubly transitive group of substitutions. Its characteristic function is
\begin{equation}\tag{1.}
 F(x)=(1-x)^{\alpha-1}(1-x^2)^\beta(1-x^3)^\gamma\cdots.
\end{equation}
For the sum of the roots of the equation $F(x)=0$ is $\chi(R)=\alpha-1$, and the sum of their $x^{\ten}$ powers is $\chi(R^x)$, namely equal to the number of symbols which $R^x$ leaves unchanged, diminished by $1$. According to the results of § 2 of my paper \emph{On the composition of the characters of a group}, Sitzungsberichte 1899, the coefficient of $(-x)^{n-1-
\lambda}$ in the development of $F(x)$ according to powers of $x$ is a linear combination of the characters of $\Hgr$ with integral positive numerical coefficients. In order to determine these, I make $F(x)$ homogeneous and form the sum
\[
 \sum_\rho\frac{h_\rho}{h}(u-v)^\alpha(u^2-v^2)^\beta(u^3-v^3)^\gamma\cdots
 (x_1+
\cdots+x_n)^\alpha(x_1^2+
\cdots+x_n^2)^\beta(x_1^3+
\cdots+x_n^3)^\gamma\cdots,
\]
in which
\[
 \frac{h_\rho}{h}=\frac{1}{1^\alpha\alpha!\,2^\beta\beta!\,3^\gamma\gamma!\cdots}
\]
If, at first, one disregards the restriction on the summation
\[
 n=\alpha+2\beta+3\gamma+
\cdots
\]
then, as in § 3, one obtains for the infinite series the expression
\[
 \left(\frac{(1-x_1v)\cdots(1-x_nv)}{(1-x_1u)\cdots(1-x_nu)}-1\right).
\]
Here, in order to avoid extensions, the system of values $\alpha=\beta=\gamma=\cdots=0$ has been excluded from the summation. This expression is to be multiplied by
\[
 \frac{\Delta(x_1,x_2,
\cdots,x_n)}{u-v}
\]
and then, in its development according to ascending powers of $u$ and $v$, the coefficient of $u^\lambda(-v)^{n-1-
\lambda}$ is to be determined. If one replaces $u$ and $v$ by their reciprocal values, and puts
\[
 f(x)=(x-x_1)(x-x_2)\cdots(x-x_n),
\]
then, in the development of the expression
\[
 \frac{\Delta(x_1,x_2,
\cdots,x_n)}{v-u}\left(\frac{f(v)}{f(u)}\frac{v^n}{u^n}\right)
\]
according to descending powers of $u$ and $v$, the coefficient of
\[
 (-1)^{n-1-
\lambda}u^{-n-1-
\lambda}v^\lambda\qquad(\lambda=0,1,
\cdots,n-1)
\]
is to be found. Now
\[
 \frac1{v-u}\left(\frac{f(v)}{f(u)}-\frac{v^n}{u^n}\right)=\frac1{u^n}\frac{u^n-v^n}{u-v}+
 \sum_\nu\frac{f(v)}{f'(x_\nu)(v-x_\nu)(u-x_\nu)},
\]
and in this expression the coefficient of $u^{-n-1-
\lambda}$, multiplied by $\Delta(x_1,x_2,
\cdots,x_n)$, is equal to
\[
 \Delta(x_1,x_2,
\cdots,x_n)\sum_\nu\frac{f(v)x_\nu^{n+\lambda}}{f'(x_\nu)(v-x_\nu)}
 =\sum_\nu\Delta(x_1,
\cdots,x_{\nu-1},v,x_{\nu+1},
\cdots,x_n)x_\nu^{n+
\lambda}.
\]
In this sum the coefficient $(-1)^{n-1-
\lambda}v^\lambda x_1^{\kappa_1}x_2^{\kappa_2}\cdots x_n^{\kappa_n}$ is to be computed, where $\kappa_1,
\kappa_2,
\cdots,
\kappa_n$ satisfy the conditions (2.), § 4. This coefficient is zero, except when $\kappa_1,
\kappa_2,
\cdots,
\kappa_{n-1}\leqq n-1$ and $\kappa_n=n+
\lambda$. Hence this decomposition has rank $r=1$, and
\[
 (\kappa)=\left(\begin{matrix}n-1-
\lambda\\ \lambda\end{matrix}\right).
\]
Finally the coefficient of
\[
 (-1)^{n-1-
\lambda}v^\lambda x_1^0x_2^1\cdots x_\lambda^{\lambda-1}x_{\lambda+1}^{\lambda+1}\cdots x_{n-1}^{n-1}x_n^{n+
\lambda}
\]
in $\Delta(x_1,
\cdots,x_{n-1},v)x_n^{n+
\lambda}$ is equal to $+1$.

Consequently
\begin{equation}\tag{2.}
\begin{aligned}
 F(x)&=\chi\left(\begin{matrix}0\\ n-1\end{matrix}\right)
 -x\chi\left(\begin{matrix}1\\ n-2\end{matrix}\right)
 +x^2\chi\left(\begin{matrix}2\\ n-3\end{matrix}\right)\cdots\\
 &\quad+(-x)^\lambda\chi\left(\begin{matrix}\lambda\\ n-1-
\lambda\end{matrix}\right)+\cdots+(-x)^{n-2}\chi\left(\begin{matrix}n-2\\1\end{matrix}\right)+(-x)^{n-1}\chi\left(\begin{matrix}n-1\\0\end{matrix}\right).
\end{aligned}
\end{equation}
The characters of first rank are therefore
\begin{equation}\tag{3.}
\begin{aligned}
\chi\left(\begin{matrix}0\\ n-1\end{matrix}\right)&=1,\qquad
\chi\left(\begin{matrix}1\\ n-2\end{matrix}\right)=\alpha-1,
\qquad
\chi\left(\begin{matrix}2\\ n-3\end{matrix}\right)=\binom{\alpha-1}{2}-\beta,\\
\chi\left(\begin{matrix}3\\ n-4\end{matrix}\right)&=\binom{\alpha-1}{3}-(\alpha-1)\beta+
\gamma,\\
\chi\left(\begin{matrix}4\\ n-5\end{matrix}\right)&=\binom{\alpha-1}{4}-\binom{\alpha-1}{2}\beta+(
\alpha-1)\gamma+
\binom{\beta}{2}-\delta,
\cdots.
\end{aligned}
\end{equation}
Furthermore
\begin{equation}\tag{4.}
 \chi\left(\begin{matrix}n-1\\0\end{matrix}\right)=(-1)^{\beta+
\delta+\cdots}=(-1)^{n-s},
\end{equation}
where $s=\alpha+
\beta+
\gamma+
\delta+
\cdots$ is the number of cycles of $R$, and in general
\[
 \chi\left(\begin{matrix}b\\a\end{matrix}\right)_\rho=(-1)^{\beta+
\delta+
\cdots}\chi\left(\begin{matrix}a\\b\end{matrix}\right)_\rho.
\]

\section*{§ 6.}

Let $\Hgr$ be the symmetric group and $\Ggr$ the alternating group of degree $n$, and let $T$ be any odd permutation. Then $\Hgr=\Ggr+
\Ggr T$, where the complex $\Ggr$ contains the $\frac12n!$ even permutations, and the complex $\Ggr T=T\Ggr$ contains the $\frac12n!$ odd permutations. These two complexes are the elements of the group $\Hgr/\Ggr$ of order $2$. It has two characters, $1,1$ and $1,-1$. Consequently the group $\Hgr$ has a character
\begin{equation}\tag{1.}
 \chi_\rho^{(1)}=\chi\left(\begin{matrix}n-1\\0\end{matrix}\right)_\rho=(-1)^{\beta+
\delta+
\cdots}=(-1)^{n-s},
\end{equation}
which has the value $+1$ for the elements of $\Ggr$, and the value $-1$ for those of $\Ggr T$. Multiplication by this character of first degree gives from every character $\chi_\rho^{(\lambda)}$ another character $\chi_\rho^{(\Lambda)}$. Two such characters, for which
\begin{equation}\tag{2.}
 \chi^{(\lambda)}=\chi^{(\kappa)}\chi^{(1)},\qquad
 \chi^{(\kappa)}=\chi^{(\lambda)}\chi^{(1)}
\end{equation}
are called by me associated. If $\chi^{(\lambda)}(R)$ has the value $0$ for every odd permutation $R$, then this character is associated with itself. Otherwise two associated characters are distinct. We characterize $\chi^{(\lambda)}$ by a system of $n$ numbers satisfying
\begin{equation}\tag{3.}
 \frac12n(n+1)=\lambda_1+
\lambda_2+
\cdots+
\lambda_n,
\qquad 0\leqq\lambda_1<\lambda_2<\cdots<\lambda_n
\end{equation}
and $\chi^{(\Lambda)}$ by an analogous decomposition (2.), § 4; and I shall also call these two decompositions associated. I will now show that then the numbers
\begin{equation}\tag{4.}
 \kappa_1,
\kappa_2,
\cdots,
\kappa_n,
 2n-1-
\lambda_1,
 2n-1-
\lambda_2,
\cdots,
 2n-1-
\lambda_n
\end{equation}
agree, apart from order, with the numbers $0,1,\cdots,2n-1$. In this way, first, the decompositions (3.) are paired with one another. For $\kappa_1,
\kappa_2,
\cdots,
\kappa_n$ are $n$ of the $2n$ numbers $0,1,
\cdots,2n-1$. If the remaining $n$ are denoted by $2n-1-
\lambda_1,2n-1-
\lambda_2,
\cdots,2n-1-
\lambda_n$, then
\[
 \sum(\kappa)+
\sum(2n-1-
\lambda)=0+1+
\cdots+2n-1=n(2n-1),
\]
and hence $\sum\kappa=\sum\lambda$. That the relations (2.) then hold follows from the equation
\begin{equation}\tag{5.}
 \sum_\rho\frac{h_\rho}{h}\chi_\rho^{(1)}\chi_\rho^{(\kappa)}\chi_\rho^{(\lambda)}=1
\end{equation}
To prove it I form the infinite series
\[
\begin{aligned}
 &\sum (-1)^{\beta+
\delta+
\cdots}\frac1{\alpha!}(x_1+
\cdots+x_n)^\alpha(y_1+
\cdots+y_n)^\alpha
 \frac1{2^\beta\beta!}(x_1^2+
\cdots+x_n^2)^\beta(y_1^2+
\cdots+y_n^2)^\beta\cdots\\
 &=e^{(x_1+
\cdots+x_n)(y_1+
\cdots+y_n)-\frac12(x_1^2+
\cdots+x_n^2)(y_1^2+
\cdots+y_n^2)+
\cdots}\\
 &=\prod(1+x_\mu y_\nu).
\end{aligned}
\]
Then the sum (5.) is equal to the coefficient of $x_1^{\kappa_1}\cdots x_n^{\kappa_n}y_1^{\lambda_1}\cdots y_n^{\lambda_n}$ in the expansion of the product
\[
 \Delta(x_1,
\cdots,x_n)\Delta(y_1,
\cdots,y_n)\prod(1+x_\mu y_\nu).
\]
If $y_1,\cdots,y_n$ are replaced by $-1/y_1,
\cdots,-1/y_n$, this is the coefficient of
\[
\pm x_1^{\kappa_1}\cdots x_n^{\kappa_n}
y_1^{2n-1-\lambda_1}\cdots y_n^{2n-1-\lambda_n}
\]
in the expansion of
\[
 \Delta(x_1,
\cdots,x_n)\Delta(y_1,
\cdots,y_n)\prod(y_\nu-x_\mu)=\Delta(x_1,
\cdots,x_n,y_1,
\cdots,y_n).
\]
This coefficient is $0$ unless the $2n$ numbers (4.) are all different from one another, and so agree, apart from order, with $0,1,\cdots,2n-1$. In that case the coefficient is $\pm1$. Even without entering into the exact determination of the sign, this gives equation (5.), since the sum on the left can have only one of the two values $0$ or $+1$.

The relation between two associated characters can be displayed more conveniently when $\chi^{(\kappa)}$ is defined by the characteristic (5.), § 4. Let $b_1,
\cdots,b_r,b_{r+1},\cdots,b_n$ be the numbers $0,1,
\cdots,n-1$. If then $\kappa_1,
\cdots,
\kappa_n$ are the numbers $n-1-a_{r+1},
\cdots,n-1-a_n,n+b_1,
\cdots,n+b_r$, then $2n-1-
\lambda_1,
\cdots,2n-1-
\lambda_n$ are the numbers $n-1-a_1,
\cdots,n-1-a_r,n+b_{r+1},
\cdots,n+b_n$, and $\lambda_1,
\cdots,
\lambda_n$ themselves are the numbers $n-1-b_{r+1},
\cdots,n-1-b_n,n+a_1,
\cdots,n+a_r$. Therefore
\begin{equation}\tag{6.}
 (\kappa)=\left(\begin{matrix}a_1&\cdots&a_r\\ b_1&\cdots&b_r\end{matrix}\right),
 \qquad
 (\lambda)=\left(\begin{matrix}b_1&\cdots&b_r\\ a_1&\cdots&a_r\end{matrix}\right)
\end{equation}
are associated characteristics, or
\begin{equation}\tag{7.}
 \chi\left(\begin{matrix}b_1&\cdots&b_r\\ a_1&\cdots&a_r\end{matrix}\right)_\rho
 =(-1)^{n-s}\chi\left(\begin{matrix}a_1&\cdots&a_r\\ b_1&\cdots&b_r\end{matrix}\right)_\rho.
\end{equation}
A characteristic is therefore associated with itself always and only when $a_1=b_1,
\cdots,a_r=b_r$. In that case, from equation (7.), § 4, one obtains
\begin{equation}\tag{8.}
 n=(2a_1+1)+
\cdots+(2a_r+1)
\end{equation}
a representation of $n$ as a sum of distinct odd numbers only. Conversely, to every such representation there corresponds a self-associated characteristic
\begin{equation}\tag{9.}
 (\kappa)=\left(\begin{matrix}a_1&\cdots&a_r\\ a_1&\cdots&a_r\end{matrix}\right),
 \qquad
 a_1+a_2+
\cdots+a_r=\frac12(n-r).
\end{equation}
The number $v$ of self-associated characters of $\Hgr$ is therefore equal to the number of representations of $n$ as a sum of distinct odd numbers only. By the well-known formula
\[
 \prod(1+x^{2\alpha-1})=1:\prod(1-(-1)^{\alpha-1}x^\alpha)
\]
this number is equal to the difference between the number of decompositions
\begin{equation}\tag{10.}
 n=c_1+c_2+c_3+
\cdots,
\end{equation}
for which $\sum(c-1)$ is even, and the number of those for which $\sum(c-1)$ is odd; it is therefore equal to the difference between the number of the even and the number of the odd classes. The $k$ classes consequently decompose into
\begin{equation}\tag{11.}
 \frac12(k+v)=u+v\quad\text{even classes and}\quad \frac12(k-v)=u\quad\text{odd classes.}
\end{equation}
To every representation (8.) there corresponds a decomposition
\[
 n=1+(2a_1+1)+
\cdots+(2a_{r-1}+1)+2a_r,
\]
for which $\sum(c-1)$ is odd. Hence $u>v$, except for $n=3$, where $u=v$.

The theorems developed here are of importance for the computation of the characters of the alternating group, which, as I shall show shortly, can also be determined completely from their relations to the characters of subgroups.

\section*{§ 7.}

I single out some special cases in which certain values of characters can be determined in a simple way. Let the substitutions of the $\rho^{\ten}$ class consist of $s$ cycles on respectively $c_1,c_2,\cdots,c_s$ symbols. Then
\begin{equation}\tag{1.}
 n=c_1+c_2+
\cdots+c_s,
\end{equation}
and let
\begin{equation}\tag{2.}
 c_1\leqq c_2\leqq\cdots\leqq c_s.
\end{equation}
Now $[\kappa_1,
\kappa_2,
\cdots,
\kappa_n]\chi_\rho^{(\kappa)}$ is the coefficient of $x_1^{\kappa_1}x_2^{\kappa_2}\cdots x_n^{\kappa_n}$ in the expansion of the product
\begin{equation}\tag{3.}
 (x_1^{c_1}+
\cdots+x_n^{c_1})\cdots(x_1^{c_s}+
\cdots+x_n^{c_s})\Delta(x_1,
\cdots,x_n).
\end{equation}
If the characteristic of $\chi^{(\kappa)}$ is
\begin{equation}\tag{4.}
 (\kappa)=\left(\begin{matrix}a_1&a_2&\cdots&a_r\\ b_1&b_2&\cdots&b_r\end{matrix}\right),
\end{equation}
then
\[
 \kappa_1,
\kappa_2,
\cdots,
\kappa_n=n-1-a_{r+1},
\cdots,n-1-a_n,n+b_1,
\cdots,n+b_r.
\]
Any term in the expansion of $\Delta(x_1,\cdots,x_n)$ is $\pm x_1^{k_1}x_2^{k_2}\cdots x_n^{k_n}$, where $k_1,k_2,\cdots,k_n$ are the numbers $0,1,\cdots,n-1$. Any term in the expansion of the product
\begin{equation}\tag{5.}
 (x_1^{c_1}+
\cdots+x_n^{c_1})\cdots(x_1^{c_s}+
\cdots+x_n^{c_s})
\end{equation}
is $x_1^{l_1}x_2^{l_2}\cdots x_n^{l_n}$, where $l_1=c_{\alpha_1}+c_{\alpha_2}+\cdots$, $l_2=c_{\tau_1}+c_{\tau_2}+\cdots$, and where
\[
 \sigma_1,
\sigma_2,
\cdots\qquad \tau_1,
\tau_2,
\cdots
\]
together are the numbers $1,2,\cdots,s$. Consequently $s'\leq s$, and $s'=s$ only when $l_1,l_2,\cdots,l_s$ agree, apart from order, with $c_1,c_2,\cdots,c_s$. Multiplying these two terms, $s'$ of the exponents $k_1,k_2,\cdots,k_n$ are increased. If the term $x_1^{\kappa_1}x_2^{\kappa_2}\cdots x_n^{\kappa_n}$ is obtained in this way, then $r$ of the $n$ exponents $k_1,k_2,\cdots,k_n$, which are $\leq n-1$, are made larger than $n-1$ by this increase, namely equal to
\[
 n+b_1,
\cdots,n+b_r.
\]
Thus $s'\geq r$. This is impossible when $s<r$, and consequently in that case $\chi_\rho^{(\kappa)}=0$.

If the number $s$ of the cycles of which the substitutions of the $\rho^{\ten}$ class consist is smaller than the rank $r$ of the characteristic $(\kappa)$, then $\chi_\rho^{(\kappa)}=0$.

Now let $s=r$. Then exactly $s$ of the exponents $k_1,k_2,\cdots,k_n$ must be increased so that they become equal to $n+b_1,
\cdots,n+b_r$, and the remaining $n-s$ exponents must agree with $n-1-a_{r+1},\cdots,n-1-a_n$. Further, $s'=s$ must hold and $l_1,l_2,
\cdots,l_s=c_1,c_2,\cdots,c_s$. The $s$ exponents $n-1-a_1,
\cdots,n-1-a_s$ are each increased by one of the numbers $c_1,
\cdots,c_s$, and so give the numbers $n+b_1,
\cdots,n+b_s$. Therefore
\begin{equation}\tag{6.}
 c_1-1=a_\alpha+b_\lambda,
\qquad c_2-1=a_\beta+b_\mu,
\quad\cdots\quad c_s-1=a_\vartheta+b_\tau,
\end{equation}
where $\alpha,\beta,\cdots,
\delta$ and $\lambda,\mu,
\cdots,
\tau$ are two permutations of the numbers $1,2,\cdots,s$. In another order these $s$ equations may read
\[
 c_\xi-1=a_\gamma+b_1,
\qquad c_\eta-1=a_\delta+b_2,
\quad\cdots\quad c_\zeta-1=a_\nu+b_s.
\]
To the term
\[
 [n+b_1,
\cdots,n+b_s,n-1-a_{s+1},
\cdots,n-1-a_n]\chi_\rho^{(\kappa)}x_1^{n+b_1}\cdots x_s^{n+b_s}x_{s+1}^{n-1-a_{s+1}}\cdots x_n^{n-1-a_n}
\]
of (3.) one therefore obtains a contribution by multiplying the term
\[
 x_1^{n-1-a_\gamma}\cdots x_s^{n-1-a_\nu}x_{s+1}^{n-1-a_{s+1}}\cdots x_n^{n-1-a_n}
\]
of $\Delta(x_1,
\cdots,x_n)$ by the term $x_\alpha^{c'_1}\cdots x_\delta^{c'_s}$ of (5.). Hence
\[
 \chi_\rho^{(\kappa)}=[-a_{s+1}-1,
\cdots,-a_n-1,b_1,
\cdots,b_s]
 \sum[-a_{s+1},
\cdots,-a_n,-a_\gamma,
\cdots,-a_\nu]
\]
or, if the order (6.) is again restored,
\[
 \chi_\rho^{(\kappa)}=\sum[-a_{s+1}-1,
\cdots,-a_n-1,b_\lambda,
\cdots,b_\tau][-a_{s+1},
\cdots,-a_n,-a_\alpha,
\cdots,-a_\vartheta],
\]
the sum extending over all solutions of equations (6.). The product of the two signs is equal to
\[
 \sign.(a_{s+1}-a_\alpha)\cdots(a_n-a_\alpha)\cdots(a_{s+1}-a_\vartheta)\cdots(a_n-a_\vartheta)
 [b_\lambda,
\cdots,b_\tau][a_\alpha,
\cdots,a_\vartheta](-1)^{\frac12s(s-1)}.
\]
The latter sign is, since $a_\alpha,
a_\beta,
\cdots,
a_\vartheta$ differ from $a_1,a_2,
\cdots,a_s$ only in their order, equal to
\[
 \sign.(a_{s+1}-a_1)
\cdots(a_n-a_1)
\cdots(a_{s+1}-a_s)
\cdots(a_n-a_s).
\]
Since $a_1<a_2<\cdots<a_s$, and $a_1,
\cdots,a_s,a_{s+1},
\cdots,a_n=0,1,
\cdots,n-1$, the numbers $0,1,
\cdots,a_1-1$ occur among $a_{s+1},
\cdots,a_n$; thus among the $n-s$ differences $a_{s+1}-a_1,
\cdots,a_n-a_1$ precisely $a_1$ are negative. Likewise the numbers $0,1,
\cdots,a_2-1$, with the exception of $a_1$, occur among $a_{s+1},
\cdots,a_n$, and therefore among the $n-s$ differences $a_{s+1}-a_2,
\cdots,a_n-a_2$ precisely $a_2-1$ are negative.

The last sign is therefore equal to
\[
 (-1)^{a_1+(a_2-1)+(a_3-2)+\cdots+(a_s-s+1)}.
\]
Since, finally, $b_1<b_2<\cdots<b_s$, one has $[b_\lambda,b_\mu,
\cdots,b_\tau]=[\lambda,\mu,
\cdots,
\tau]$, and hence
\begin{equation}\tag{7.}
 \chi_\rho^{(\kappa)}=(-1)^{a_1+a_2+
\cdots+a_s}\sum[\alpha,
\beta,
\cdots,
\vartheta][\lambda,
\mu,
\cdots,
\tau],
\end{equation}
where the sum extends over all solutions of equations (6.). A term of this sum has value $+1$ or $-1$ according as $\lambda,
\mu,
\cdots,
\tau$ is a positive or a negative permutation of $\alpha,
\beta,
\cdots,
\vartheta$.

If, for example,
\begin{equation}\tag{8.}
 c_1-1=a_1+b_1,
\qquad c_2-1=a_2+b_2,
\quad\cdots\quad c_s-1=a_s+b_s,
\end{equation}
then a solution of equations (6.) is thereby given. They can have no other. For if $a_1+b_1=c_1-1=a_\alpha+b_\lambda$, then, since $a_\alpha\geq a_1$ and $b_\lambda\geq b_1$, it must be that $a_\alpha=a_1$ and $b_\lambda=b_1$, and hence $\alpha=1$ and $\lambda=1$, and so on. Thus, under condition (8.),
\begin{equation}\tag{9.}
 \chi_\rho^{(\kappa)}=(-1)^{a_1+a_2+
\cdots+a_s}.
\end{equation}
Every character $\chi^{(\kappa)}$ therefore has, for at least one class, the value $\pm1$. If more specially $a_1=b_1,
\cdots,a_s=b_s$, and hence by (8.)
\begin{equation}\tag{10.}
 c_1=2a_1+1,
\qquad c_2=2a_2+1,
\quad\cdots\quad c_s=2a_s+1,
\end{equation}
then
\begin{equation}\tag{11.}
 \chi_\rho^{(\kappa)}=(-1)^{\frac12(n-s)}=(-1)^{\frac12(p-1)},
\end{equation}
where
\begin{equation}\tag{12.}
 p=c_1c_2\cdots c_s
\end{equation}
For this product satisfies
\[
 (1+2a_1)\cdots(1+2a_s)\equiv 1+2(a_1+
\cdots+a_s)\pmod{4}.
\]

But if $r=s$ and $a_1=b_1,
\cdots,a_s=b_s$, without condition (10.) being satisfied, then
\[
 \chi_\rho^{(\kappa)}=(-1)^{\frac12(n-s)}\sum[\alpha,
\beta,
\cdots,
\vartheta][\lambda,
\mu,
\cdots,
\tau],
\]
the sum being extended over all solutions of the equations
\[
 c_1-1=a_\alpha+a_\lambda,
\qquad c_2-1=a_\beta+a_\mu,
\quad\cdots\quad c_s-1=a_\vartheta+a_\tau.
\]
One does not satisfy these equations by putting
\[
 \alpha=\lambda,
\qquad \beta=\mu,
\quad\cdots\quad \vartheta=\tau
\]
for otherwise, by (2.), $\alpha=\lambda=1$, $\beta=\mu=2$, $\cdots$, $\delta=\tau=s$. Hence to every solution there corresponds another,
\[
 c_1-1=a_\lambda+a_\alpha,
\qquad c_2-1=a_\mu+a_\beta,
\quad\cdots\quad c_s-1=a_\tau+a_\vartheta,
\]
which is different from it, and consequently $\chi_\rho^{(\kappa)}$ is an even number. For all other classes also, as I shall show, $\chi_\rho^{(\kappa)}$ is even, with the exception of the one class (10.).

To treat another example, I shall compute $\chi_\rho^{(\kappa)}$ for the case in which the permutations of the $\rho^{\ten}$ class consist of one cycle on $c$ symbols and of $n-c$ cycles on one symbol each, and I do this by the method used in § 3 to determine $f^{(\lambda)}$. Using also the notation employed there, one has
\[
\begin{aligned}
 &\sum[\lambda_1,
\cdots,
\lambda_n]\chi_\rho^{(\lambda)}x_1^{\lambda_1}\cdots x_n^{\lambda_n} \\
 &\quad=(x_1^c+
\cdots+x_n^c)(x_1+
\cdots+x_n)^{n-c}\Delta(x_1,
\cdots,x_n)\\
 &\quad=(x_1^c+
\cdots+x_n^c)
 \left(\sum\frac{(n-c)!}{\mu_1!\cdots\mu_n!}x_1^{\mu_1}\cdots x_n^{\mu_n}\right)\\
 &\qquad\cdot
 \left(\sum[\kappa_1,
\cdots,
\kappa_n]x_1^{\kappa_1}\cdots x_n^{\kappa_n}\right),
\end{aligned}
\]
\[
 \sum[\kappa_1,
\kappa_2,
\cdots,
\kappa_n]\frac1{(\lambda_1-c-
\kappa_1)!}\frac1{(\lambda_2-
\kappa_2)!}\cdots\frac1{(\lambda_n-
\kappa_n)!},
\]
where $\mu_1+\cdots+
\mu_n=n-c$, and $\kappa_1,
\cdots,
\kappa_n$ are the numbers $0,1,
\cdots,n-1$. If $\lambda_1<\lambda_2<\cdots<\lambda_n$, then $\chi_\rho^{(\lambda)}:(n-c)!$ consists of $n$ parts. The first is obtained by taking from the sum $x_1^c+
\cdots+x_n^c$ only the first term. It is equal to
\[
 \frac{\Delta(\lambda_1-c,
\lambda_2,
\cdots,
\lambda_n)}{(\lambda_1-c)!\lambda_2!\cdots\lambda_n!}.
\]
and hence, by § 3, equal to
\[
 f^{(\lambda)}=\frac{n!\Delta(\lambda_1,
\lambda_2,
\cdots,
\lambda_n)}{\lambda_1!\lambda_2!\cdots\lambda_n!}
\]
Dividing by
\[
 h_\rho=\frac{n!}{c(n-c)!},
\]
and multiplying by
\[
 -\frac{c^2h_\rho\chi_\rho^{(\lambda)}}{f^{(\lambda)}}
\]
one sees that $-\dfrac{c^2h_\rho\chi_\rho^{(\lambda)}}{f^{(\lambda)}}$ is a sum of $n$ terms, the first of which is
\[
 \frac{(\lambda_1-c-
\lambda_1)(\lambda_1-c-
\lambda_2)\cdots(\lambda_1-c-
\lambda_n)\lambda_1(\lambda_1-1)\cdots(\lambda_1-c+1)}{(\lambda_1-
\lambda_2)\cdots(\lambda_1-
\lambda_n)}
\]
Putting
\[
 \varphi(x)=(x-
\lambda_1)\cdots(x-
\lambda_n),
\quad
 \psi(x)=(x-c-
\lambda_1)\cdots(x-c-
\lambda_n)x(x-1)\cdots(x-c+1),
\]
one has therefore
\[
 -\frac{c^2h_\rho\chi_\rho^{(\lambda)}}{f^{(\lambda)}}=
 \frac{\psi(\lambda_1)}{\varphi'(\lambda_1)}+
\cdots+
\frac{\psi(\lambda_n)}{\varphi'(\lambda_n)},
\]
or it is equal to the coefficient of $x^{-1}$ in the expansion of $\psi(x)/\varphi(x)$ according to descending powers of $x$. If
\[
 (\lambda)=\left(\begin{matrix}a_1&a_2&\cdots&a_r\\ b_1&b_2&\cdots&b_r\end{matrix}\right),
\]
then this fraction is to be multiplied in numerator and denominator by
\[
 (x-(n-1-a_1))\cdots(x-(n-1-a_r))\,(x-c-(n-1-a_1))\cdots(x-c-(n-1-a_r)).
\]
Now the numbers
\[
 \lambda_1,
\cdots,
\lambda_n,
 n-1-a_1,
\cdots,n-1-a_r
\]
are equal to
\[
 0,1,
\cdots,n-1,
 n+b_1,
\cdots,n+b_r.
\]
If $x$ is now replaced by $x+n$, the fraction becomes
\[
 \frac{(x+a_1+1)\cdots(x+a_r+1)(x-c-b_1)\cdots(x-c-b_r)x(x-1)\cdots(x-c+1)}{(x-b_1)\cdots(x-b_r)(x-c+a_1+1)\cdots(x-c+a_r+1)}.
\]
Thus, if
\begin{equation}\tag{13.}
 f(x)=\frac{(x-b_1)\cdots(x-b_r)}{(x+a_1+1)\cdots(x+a_r+1)},
\end{equation}
then
\begin{equation}\tag{14.}
\begin{aligned}
 -\frac{c^2h_\rho\chi_\rho^{(\lambda)}}{f^{(\lambda)}}
 &=\left[\frac{f(x-c)x(x-1)\cdots(x-c+1)}{f(x)}\right]_{x^{-1}}\\
 &=\left[\frac{\varphi(x-c)x(x-1)\cdots(x-c+1)}{\varphi(x)}\right]_{x^{-1}},
\end{aligned}
\end{equation}
That is, it is equal to the coefficient of $x^{-1}$ in the expansion of this rational function according to descending powers of $x$; and this formula remains valid if, using the notation (3.) or (4.), § 4, one puts
\begin{equation}\tag{15.}
 \varphi(x)=(x-
\lambda_1)(x-
\lambda_2)\cdots(x-
\lambda_m)
\end{equation}
For $c=2$ one therefore obtains
\begin{equation}\tag{16.}
 \frac{h_\rho\chi_\rho^{(\lambda)}}{f^{(\lambda)}}=\frac12\left(\sum b(b+1)-\sum a(a+1)\right).
\end{equation}
I have already proved, in § 2 of my paper \emph{On group characters}, Sitzungsberichte 1896, that this expression is an integer.
% ===== END WORK 060 =====

\clearpage

% ===== BEGIN WORK 061: 061_ber_die_Charaktere_der_alternirenden_Grup =====
\frobtitle{On the characters of the alternating group}

The characters of the alternating group of the permutations of $n$ symbols can be determined by a procedure similar to the one which I used (Sitzungsberichte 1900) for the computation of the characters of the symmetric group.

The alternating group $\Hgr$ of degree $n$ consists of the even, or positive, permutations of the symmetric group $\Hgr'$. A permutation which consists of $s$ cycles $C_1,C_2,\ldots,C_s$ on respectively $c_1,c_2,\ldots,c_s$ symbols is even or odd according as $\sum(c-1)=n-s$ is even or odd. The groups $\Hgr$ and $\Hgr'$ have orders $h=\frac12 n!$ and $h'=2h$; let their elements decompose into $k$ and $k'$ classes. Also in $\Hgr$ two permutations can be conjugate only when they consist of equally many cycles of the same order. But the question arises whether all such permutations also form only one class in $\Hgr$. For two even permutations $R$ and $S$ are always, and only, conjugate in $\Hgr$ when there is an even permutation $P$ satisfying the condition $P^{-1}RP=S$.

If $R$ and $S$ are conjugate in $\Hgr'$, then this condition is always satisfied when $R$ is commutable with a negative permutation $T$. For in $\Hgr'$ there is a permutation $Q$ such that $Q^{-1}RQ=S$. Moreover $(TQ)^{-1}R(TQ)=Q^{-1}(T^{-1}RT)Q=Q^{-1}RQ=S$, and of the two permutations $Q$ and $TQ$ one is always even. In this case, therefore, the class of $\Hgr$ is completely determined by the numbers $c_1,c_2,\ldots,c_s$.

But if $R$ is commutable with no negative permutation, then $R$ and $T^{-1}RT=S$ cannot be conjugate in $\Hgr$ if $T$ is odd. For otherwise there would be a positive permutation $P$ such that $P^{-1}SP=R$, and hence the negative permutation $TP$ would be commutable with $R$.

If $P_1,P_2,\ldots,P_h$ are the $h$ positive permutations, then $TP_1,TP_2,\ldots,TP_h$ are the $h$ negative ones. Therefore
\[
 P_\lambda^{-1}RP_\lambda,\qquad P_\lambda^{-1}T^{-1}RTP_\lambda\qquad(\lambda=1,2,\ldots,h)
\]
are all the permutations which are conjugate with $R$ in $\Hgr'$. They form one class in $\Hgr'$, but two classes in $\Hgr$, which are represented by $R$ and $T^{-1}RT=S$. Two such \emph{conjugate classes} of $\Hgr$ have the same order.

In order, however, that $R$ be commutable only with positive permutations, it is necessary and sufficient that the orders of the cycles of $R$ be all distinct odd numbers. For each individual cycle $C$ of $R$ is a permutation commutable with $R$, and indeed a negative one when its order $c$ is even. But if $c$ is odd, and two of these cycles, say $(1,2,\ldots,c)$ and $(c+1,c+2,\ldots,2c)$, have the same order, then $R$ is commutable with the permutation $(1,c+1)(2,c+2)\cdots(c,2c)$, which consists of the odd number $c$ of transpositions and hence is odd.

But if the orders $c_1,c_2,\ldots,c_s$ of the cycles $C_1,C_2,\ldots,C_s$ of $R$ are all distinct odd numbers, then $R$ is commutable only with positive permutations. These form a group of commutable elements
\begin{equation}\tag{1.}
 C_1^{\gamma_1}C_2^{\gamma_2}\cdots C_s^{\gamma_s}\qquad(\gamma_1=0,1,\ldots,c_1-1,\ldots,\gamma_s=0,1,\ldots,c_s-1)
\end{equation}
of order
\begin{equation}\tag{2.}
 p=c_1c_2\cdots c_s
\end{equation}
and are all even, because the basis-elements $C_1,C_2,\ldots,C_s$ are all even.

In the symmetric group $\Hgr'$, let $u$ be the number of odd classes, and $u+v$ the number of even classes. Then $v$ is the number of classes of $\Hgr'$ whose permutations consist of cycles of distinct odd orders only. In the group $\Hgr$, each of these $v$ classes of $\Hgr'$ decomposes into two conjugate classes; but each of the other $u$ even classes of $\Hgr'$ remains one class also in $\Hgr$. Therefore
\begin{equation}\tag{3.}
 k=u+2v,	\qquad k'=2u+v,
\end{equation}
and since $u>v$, one also has $k'>k$.

\section*{§ 2.}

With the aid of the rules which I developed in § 1 of my paper \emph{On relations between the characters of a group and those of its subgroups}, Sitzungsberichte 1898 (cited in what follows as U.), the $k$ characters of $\Hgr$ can for the most part be derived from those of $\Hgr'$. If, in a character $\chi(S)$ of $\Hgr'$, one lets $S$ run only over the $h$ elements $R$ of $\Hgr$, then $\chi(R)$ is a linear combination of the $k$ characters $\phi^{(\varkappa)}(R)$ of the group $\Hgr$,
\[
 \chi(R)=\sum_{\varkappa} r_\varkappa\phi^{(\varkappa)}(R)
\]
with positive integral coefficients $r_\varkappa$, and hence
\[
 \sum_R \chi(R)\chi(R^{-1})=h\sum_\varkappa r_\varkappa^2.
\]
For the symmetric group the two (in general conjugate complex) values $\chi(R)$ and $\chi(R^{-1})$ are equal. If $S$ runs through the $2h$ elements of $\Hgr'$, then
\[
 \sum_S \chi(S)^2=2h.
\]
If $T$ is a fixed negative permutation, then $R$ and $RT$ together run through the $2h$ elements $S$ of $\Hgr'$, and hence
\[
 \sum_R \chi(R)^2+\sum_R \chi(RT)^2=2h.
\]
The character of $\Hgr'$ associated with $\chi(S)$ has the values $\omega(R)=\chi(R)$, $\omega(RT)=-\chi(RT)$. If now $\chi(S)$ is not associated with itself, then $\chi(S)$ and $\omega(S)$ are distinct characters of $\Hgr'$, and consequently
\[
 \sum_S \chi(S)\omega(S)=0,\qquad \sum_R \chi(R)^2=\sum_R\chi(RT)^2,
\]
so that
\[
 \sum_R\chi(R)\chi(R^{-1})=h,	\qquad \sum_\varkappa r_\varkappa^2=1.
\]
Therefore, of the $k$ positive numbers $r_\varkappa$, one is equal to $1$ and the others are equal to $0$, and hence $\chi(R)$ is a character of $\Hgr$. Thus from each of the $u$ pairs of associated characters of $\Hgr'$ there arises a character of $\Hgr$ satisfying the condition
\begin{equation}\tag{1.}
 \chi(T^{-1}RT)=\chi(R),
\end{equation}
that is, one which is self-conjugate.

But if $\chi(S)$ is associated with itself, then $\chi(RT)=0$,
\[
 \sum_R \chi(R)^2=2h,	\qquad \sum_\varkappa r_\varkappa^2=2.
\]
Consequently two of the numbers $r_\varkappa$ are equal to $1$, the others to $0$, and $\chi(R)$ is equal to the sum of two distinct characters $\phi(R)$ and $\psi(R)$ of $\Hgr$. By U. § 2 these, since $\Hgr$ is an invariant subgroup of $\Hgr'$, are two conjugate characters, hence
\begin{equation}\tag{2.}
 \psi(R)=\phi(T^{-1}RT),\qquad \chi(R)=\phi(R)+\psi(R)=\phi(R)+\phi(T^{-1}RT).
\end{equation}
If $R$ belongs to one of the $u$ non-splitting classes of $\Hgr'$, then $R$ and $T^{-1}RT$ are also conjugate in $\Hgr$, and hence for such an element $R$
\begin{equation}\tag{u.}
 \phi(R)=\psi(R)=\frac12\chi(R).
\end{equation}
In particular this holds for $R=E$. The degrees of the two conjugate characters $\phi$ and $\psi$ are therefore equal to one another and equal to half the degree of $\chi$. But if the orders of the cycles of $R$ are all distinct odd numbers, so that $R$ belongs to one of the $v$ splitting classes of $\Hgr'$, then the values of $\phi(R)$ and $\psi(R)$ may be distinct; more precisely, putting $T^{-1}RT=S$, one has
\begin{equation}\tag{v.}
 \phi(R)=\psi(S),	\qquad \phi(S)=\psi(R).
\end{equation}
From each of the $v$ self-associated characters of $\Hgr'$ there arise, therefore, two distinct, indeed conjugate, characters of $\Hgr$.

The $u+2v=k$ characters of $\Hgr$ so obtained are all distinct from one another. For let $\phi(R)=\chi(R)$ and $\overline{\phi}(R)=\overline{\chi}(R)$ be two of the first $u$ characters of $\Hgr$, and suppose that, when $S$ runs through the $2h$ elements of $\Hgr'$, $\overline{\chi}(S)$ is distinct both from $\chi(S)$ and from the character associated with $\chi(S)$. Then
\[
 \sum_R \chi(R)\overline{\chi}(R)+\sum_R \chi(RT)\overline{\chi}(RT)=0,
\]
and, replacing $\chi(R)$ by the associated character, also
\[
 \sum_R \chi(R)\overline{\chi}(R)-\sum_R\chi(RT)\overline{\chi}(RT)=0,
\]
and consequently
\[
 \sum_R \phi(R)\overline{\phi}(R^{-1})=0,
\]
so that $\overline{\phi}(R)$ is distinct from $\phi(R)$.

If, however, $\overline{\phi}(R)=\overline{\chi}(R)$ is one of the first $u$, and $\phi(R)$ one of the last $2v$ characters of $\Hgr$, arising from $\chi(R)=\phi(R)+\psi(R)$, then $\chi(RT)=0$, and therefore
\[
 \sum_S \overline{\chi}(S^{-1})\chi(S)=0,	\qquad \sum_R\overline{\phi}(R^{-1})(\phi(R)+\psi(R))=0.
\]
Now $\phi$ and $\psi$ are distinct. If therefore $\overline{\phi}$ were equal to one of these two characters, this sum would be equal to $h$. Finally one sees in the same way that any two of the last $2v$ characters of $\Hgr$ are distinct.

Thus, in order to determine the $u+2v=k$ distinct characters of $\Hgr$, it only remains to split each of the $v$ self-associated characters $\chi$ of $\Hgr'$ into two conjugate characters $\phi$ and $\psi$, and this only for the $v$ permutations $R$ which consist of cycles of distinct odd orders only.

The $k'=2u+v$ characters of $\Hgr'$ consist of $u$ pairs of associated characters and $v$ characters which are self-associated. The $k=u+2v$ characters of $\Hgr$ consist of $v$ pairs of conjugate characters and $u$ characters which are self-conjugate.

\section*{§ 3.}

I first refer to the examples $n=4$ and $n=5$ in § 3 of my paper \emph{On group characters}, and $n=6$ in § 4 of my paper \emph{On the composition of the characters of a group}, Sitzungsberichte 1896 and 1899. In these cases $v=1$, likewise for $n=7$. I therefore give here still the example $n=8$, where $v=2$.

In the first column of the adjoining table, the class of a permutation is given by the formula $n=1\cdot\alpha+2\cdot\beta+\cdots$, that is, by the fact that it consists of $\alpha$ cycles of order $1$, $\beta$ cycles of order $2$, and so forth. If only one cycle of a certain order occurs, then the multiplier $\alpha,\beta,\cdots$ is omitted. First come the $u=10$ odd classes, then the $u+v=12$ even ones.

From each of the $u=10$ pairs of associated characters only one is listed. The other arises from it by multiplication with the character of first degree $\chi\ch{7}{0}$, which is associated with the principal character $\chi\ch{0}{7}=1$. In the last two columns stand the $v=2$ self-associated characters.

The characters are given by formulas (11.), § 4 and (3.), § 5 of my paper \emph{On the characters of the symmetric group}, and by the following:
\begin{align}\tag{1.}
 \chi\ch{0\quad 2}{0\quad n-4}&=\frac13\alpha(\alpha-2)(\alpha-4)-\gamma,\\
 \chi\ch{0\quad 2}{1\quad n-5}&=\frac18\alpha(\alpha-1)(\alpha-3)(\alpha-6)+\frac12(\alpha-1)(\alpha-2)\beta-\frac12\beta(\beta-1)-\delta,\\
 \chi\ch{1\quad 2}{0\quad n-5}&=\frac1{12}\alpha(\alpha-1)(\alpha-4)(\alpha-5)+\beta(\beta-2)-(\alpha-1)\gamma,\\
 \chi\ch{0\quad 3}{0\quad n-5}&=\frac18\alpha(\alpha-2)(\alpha-3)(\alpha-5)-\frac12\alpha(\alpha-3)\beta-\frac12\beta(\beta-1)+\delta,\\
 \chi\ch{1\quad 2}{1\quad n-6}&=\frac1{24}\alpha(\alpha-1)(\alpha-2)(\alpha-5)(\alpha-7)+\frac16\alpha(\alpha-1)(\alpha-5)\beta\notag\\
 &\quad+\frac12(\alpha-1)\beta(\beta-1)-\frac12(\alpha-1)(\alpha-2)\gamma+\beta\gamma-(\alpha-1)\delta.
\end{align}
In the alternating group the class $(1+7)$ decomposes into two classes, which I distinguish by an appended sign $+$ or $-$. Likewise the character $\chi\ch{03}{03}$ decomposes into two. The same holds for the class $(3+5)$ and for the corresponding character $\chi\ch{12}{12}$.

\begin{landscape}
\begin{center}
\small
\textbf{Symmetric group.}\\
$h=8!=40320.$\\[0.5em]
\setlength{\tabcolsep}{3.5pt}
\renewcommand{\arraystretch}{1.15}
\begin{tabular}{c r|rrrrrrrrrrrr}
\toprule
$1\cdot\alpha+2\cdot\beta+\cdots$ & $h_\rho$ & $\ch{7}{0}$ & $\ch{1}{6}$ & $\ch{2}{5}$ & $\ch{3}{4}$ & $\ch{01}{05}$ & $\ch{01}{14}$ & $\ch{01}{23}$ & $\ch{02}{04}$ & $\ch{02}{13}$ & $\ch{12}{03}$ & $\ch{03}{03}$ & $\ch{12}{12}$\\
\midrule
$1\cdot6+2$ & 28 & $-1$ & 5 & 9 & 5 & 10 & 10 & 4 & 16 & 10 & 4 & 0 & 0\\
$1\cdot4+4$ & 420 & $-1$ & 3 & 3 & 1 & 2 & $-2$ & $-2$ & 0 & $-4$ & 0 & 0 & 0\\
$1\cdot3+2+3$ & 1120 & $-1$ & 2 & 0 & $-1$ & 1 & 1 & 1 & $-2$ & 1 & $-2$ & 0 & 0\\
$1\cdot2+2\cdot3$ & 420 & $-1$ & 1 & $-3$ & $-3$ & 2 & 2 & 0 & 0 & $-2$ & 4 & 0 & 0\\
$1\cdot2+6$ & 3360 & $-1$ & 1 & 0 & 0 & $-1$ & $-1$ & 0 & 0 & 1 & 1 & 0 & 0\\
$1+2+5$ & 4032 & $-1$ & 0 & $-1$ & 0 & 0 & 0 & $-1$ & 1 & 0 & $-1$ & 0 & 0\\
$1+3+4$ & 3360 & $-1$ & 0 & 0 & 1 & $-1$ & 1 & 1 & 0 & $-1$ & 0 & 0 & 0\\
$2\cdot2+4$ & 1260 & $-1$ & $-1$ & $-1$ & 1 & 2 & $-2$ & 2 & 0 & 0 & 0 & 0 & 0\\
$2+3\cdot2$ & 1120 & $-1$ & $-1$ & 0 & 2 & 1 & 1 & $-2$ & $-2$ & 1 & 1 & 0 & 0\\
$8$ & 5040 & $-1$ & $-1$ & 1 & $-1$ & 0 & 0 & 0 & 0 & 0 & 0 & 0 & 0\\
\midrule
$1\cdot8$ & 1 & 1 & 7 & 21 & 35 & 20 & 28 & 14 & 64 & 70 & 56 & 90 & 42\\
$1\cdot5+3$ & 112 & 1 & 4 & 6 & 5 & 5 & 1 & $-1$ & 4 & $-5$ & $-4$ & 0 & $-6$\\
$1\cdot4+2\cdot2$ & 210 & 1 & 3 & 1 & $-5$ & 4 & 4 & 2 & 0 & 2 & 0 & $-6$ & 2\\
$1\cdot3+5$ & 1344 & 1 & 2 & 1 & 0 & 0 & $-2$ & $-1$ & $-1$ & 0 & 1 & 0 & 2\\
$1\cdot2+2+4$ & 2520 & 1 & 1 & $-1$ & $-1$ & 0 & 0 & 0 & 0 & 0 & 0 & 2 & $-2$\\
$1\cdot2+3\cdot2$ & 1120 & 1 & 1 & 0 & 2 & $-1$ & 1 & 2 & $-2$ & 1 & $-1$ & 0 & 0\\
$1+2\cdot2+3$ & 1680 & 1 & 0 & $-2$ & 1 & 1 & 1 & $-1$ & 0 & $-1$ & 0 & 0 & 2\\
$2\cdot4$ & 105 & 1 & $-1$ & $-3$ & 3 & 4 & $-4$ & 6 & 0 & $-2$ & 8 & $-6$ & $-6$\\
$2+6$ & 3360 & 1 & $-1$ & 0 & 0 & 1 & $-1$ & 0 & 0 & 1 & $-1$ & 0 & 0\\
$4\cdot2$ & 1260 & 1 & $-1$ & 1 & $-1$ & 0 & 0 & 2 & 0 & $-2$ & 0 & 2 & 2\\
\midrule
$1+7$ & 5760 & 1 & 0 & 0 & 0 & $-1$ & 0 & 0 & 1 & 0 & 0 & $-1$ & 0\\
$3+5$ & 2688 & 1 & $-1$ & 1 & 0 & 0 & 1 & $-1$ & $-1$ & 0 & 1 & 0 & $-1$\\
\bottomrule
\end{tabular}
\end{center}
\end{landscape}

\begin{landscape}
\begin{center}
\small
\textbf{Alternating group.}\\
$h=\frac12 8! =20160.$\\[0.5em]
\setlength{\tabcolsep}{2.7pt}
\renewcommand{\arraystretch}{1.15}
\begin{tabular}{c r|rrrrrrrrrrrrrr}
\toprule
$1\cdot\alpha+2\cdot\beta$ & $h_\rho$ & $\ch{0}{7}$ & $\ch{1}{6}$ & $\ch{2}{5}$ & $\ch{3}{4}$ & $\ch{01}{05}$ & $\ch{01}{14}$ & $\ch{01}{23}$ & $\ch{02}{04}$ & $\ch{02}{13}$ & $\ch{12}{03}$ & $\ch{03}{03}_{+}$ & $\ch{03}{03}_{-}$ & $\ch{12}{12}_{+}$ & $\ch{12}{12}_{-}$\\
\midrule
$1\cdot8$ & 1 & 1 & 7 & 21 & 35 & 20 & 28 & 14 & 64 & 70 & 56 & 45 & 45 & 21 & 21\\
$1\cdot5+3$ & 112 & 1 & 4 & 6 & 5 & 5 & 1 & $-1$ & 4 & $-5$ & $-4$ & 0 & 0 & $-3$ & $-3$\\
$1\cdot4+2\cdot2$ & 210 & 1 & 3 & 1 & $-5$ & 4 & 4 & 2 & 0 & 2 & 0 & $-3$ & $-3$ & 1 & 1\\
$1\cdot3+5$ & 1344 & 1 & 2 & 1 & 0 & 0 & $-2$ & $-1$ & $-1$ & 0 & 1 & 0 & 0 & 1 & 1\\
$1\cdot2+2+4$ & 2520 & 1 & 1 & $-1$ & $-1$ & 0 & 0 & 0 & 0 & 0 & 0 & 1 & 1 & $-1$ & $-1$\\
$1\cdot2+3\cdot2$ & 1120 & 1 & 1 & 0 & 2 & $-1$ & 1 & 2 & $-2$ & 1 & $-1$ & 0 & 0 & 0 & 0\\
$1+2\cdot2+3$ & 1680 & 1 & 0 & $-2$ & 1 & 1 & 1 & $-1$ & 0 & $-1$ & 0 & 0 & 0 & 1 & 1\\
$2\cdot4$ & 105 & 1 & $-1$ & $-3$ & 3 & 4 & $-4$ & 6 & 0 & $-2$ & 8 & $-3$ & $-3$ & $-3$ & $-3$\\
$2+6$ & 3360 & 1 & $-1$ & 0 & 0 & 1 & $-1$ & 0 & 0 & 1 & $-1$ & 0 & 0 & 0 & 0\\
$4\cdot2$ & 1260 & 1 & $-1$ & 1 & $-1$ & 0 & 0 & 2 & 0 & $-2$ & 0 & 1 & 1 & 1 & 1\\
\midrule
$(1+7)_{+}$ & 2880 & 1 & 0 & 0 & 0 & $-1$ & 0 & 0 & 1 & 0 & 0 & $\frac12(-1+\sqrt{-7})$ & $\frac12(-1-\sqrt{-7})$ & 0 & 0\\
$(1+7)_{-}$ & 2880 & 1 & 0 & 0 & 0 & $-1$ & 0 & 0 & 1 & 0 & 0 & $\frac12(-1-\sqrt{-7})$ & $\frac12(-1+\sqrt{-7})$ & 0 & 0\\
$(3+5)_{+}$ & 1344 & 1 & $-1$ & 1 & 0 & 0 & 1 & $-1$ & $-1$ & 0 & 1 & 0 & 0 & $\frac12(-1+\sqrt{-15})$ & $\frac12(-1-\sqrt{-15})$\\
$(3+5)_{-}$ & 1344 & 1 & $-1$ & 1 & 0 & 0 & 1 & $-1$ & $-1$ & 0 & 1 & 0 & 0 & $\frac12(-1-\sqrt{-15})$ & $\frac12(-1+\sqrt{-15})$\\
\bottomrule
\end{tabular}
\end{center}
\end{landscape}

In the same way as in these examples one finds, for every value of $n$, the characters of the alternating group $\Hgr$. Let
\begin{equation}\tag{2.}
 c_1=2a_1+1,	\quad c_2=2a_2+1,	\quad \ldots,	\quad c_s=2a_s+1\qquad(c_1<c_2<\cdots<c_s)
\end{equation}
be $s$ distinct odd numbers whose sum is
\begin{equation}\tag{3.}
 c_1+c_2+\cdots+c_s=n.
\end{equation}
To the class (2.) of the symmetric group $\Hgr'$ corresponds the self-associated character
\begin{equation}\tag{4.}
 \chi\ch{a_1\ a_2\ \cdots\ a_s}{a_1\ a_2\ \cdots\ a_s},
\end{equation}
and conversely to this character corresponds the class (2.). In this way the $v$ self-associated characters (4.) of the group $\Hgr'$ and the $v$ classes (2.) are mutually and uniquely assigned to one another. If $R$ belongs to the class (2.), and if $T$ is any negative permutation, then $R$ and $T^{-1}RT=S$ represent the two conjugate classes into which the class (2.) splits in the alternating group $\Hgr$. Likewise the character (4.) of $\Hgr'$ splits into two conjugate characters of $\Hgr$,
\begin{equation}\tag{5.}
 \chi(P)=\phi(P)+\psi(P),
\end{equation}
where
\[
 \psi(P)=\phi(T^{-1}PT).
\]
Then, for every class of $\Hgr$ represented by $Q$,
\begin{equation}\tag{6.}
 \phi(Q)=\psi(Q)=\frac12\chi(Q)
\end{equation}
with the exception of the two classes $(R)$ and $(S)$. I have already shown in § 2 that this equation holds for the $u$ classes of $\Hgr$ whose cycles do not satisfy the conditions (2.). But it also holds for $2v-2$ of the remaining classes. On the other hand,
\begin{equation}\tag{7.}
\begin{aligned}
 \phi(R)&=\frac12(\eps+\sqrt{\eps p}),&\qquad \psi(R)&=\frac12(\eps-\sqrt{\eps p}),\\
 \phi(S)&=\frac12(\eps-\sqrt{\eps p}),&\qquad \psi(S)&=\frac12(\eps+\sqrt{\eps p}),
\end{aligned}
\end{equation}
where
\begin{equation}\tag{8.}
 p=c_1c_2\cdots c_s=\frac{h}{h_R},	\qquad \eps=(-1)^{\frac12(p-1)}=(-1)^{\frac12(n-s)}.
\end{equation}

The cyclic permutation $C=(0,1,2,\ldots,a,a+1,\ldots,2a-1,2a)$ of order $c=2a+1$ is transformed by the permutation
\[
 G=(1,2a)(2,2a-1)\cdots(a,a+1)
\]
into
\[
 G^{-1}CG=(0,2a,2a-1,\ldots,a+1,a,\ldots,2,1)=C^{-1}.
\]
The permutation $G$ consists of $a=\frac12(c-1)$ transpositions. Therefore $R=C_1C_2\cdots C_s$ is transformed into $H^{-1}RH=R^{-1}$ by a permutation $H$ which consists of $\frac12\sum(c-1)=\frac12(n-s)$ transpositions. Now, under the assumption (2.), the permutations which transform $R$ into $R^{-1}$ are either all even or all odd. Therefore $R$ and $R^{-1}$ belong in $\Hgr$ to the same class or to two distinct conjugate classes according as $\eps=+1$ or $-1$. In the first case $\phi(R)=\phi(R^{-1})$ is real; in the other $\phi(R)$ and $\phi(R^{-1})=\psi(R)$ are conjugate complex quantities.

\section*{§ 4.}

Before the proof of the theorem stated above I put two remarks which hold for arbitrary groups.

1. Let $\Hgr_1$ and $\Hgr_2$ be two groups of orders $h_1$ and $h_2$ which have only the principal element in common. Further, let every element of $\Hgr_1$ be commutable with every element of $\Hgr_2$. Then $\Hgr=\Hgr_1\Hgr_2=\Hgr_2\Hgr_1$ is a group of order $h=h_1h_2$. If the elements of $\Hgr_1$ and $\Hgr_2$ split into $k_1$ and $k_2$ classes, then those of $\Hgr$ split into $k=k_1k_2$ classes. If $A_1,B_1,\ldots,R_1,\ldots$ ($A_2,B_2,\ldots,R_2,\ldots$) are representatives of the $k_1$ ($k_2$) classes of $\Hgr_1$ ($\Hgr_2$), then $A_1A_2,A_1B_2,A_2B_1,A_2B_2,\ldots,R_1R_2,\ldots$ represent the $k_1k_2=k$ classes of $\Hgr$. If $\chi_1(R_1)$ ($\chi_2(R_2)$) is a character of $\Hgr_1$ ($\Hgr_2$), and if $R=R_1R_2$, then $\chi(R)=\chi_1(R_1)\chi_2(R_2)$ is a character of $\Hgr$. If one substitutes for $\chi_1$ ($\chi_2$), successively, each of the $k_1$ ($k_2$) distinct characters of $\Hgr_1$ ($\Hgr_2$), then one obtains the $k_1k_2=k$ distinct characters of $\Hgr$. These propositions follow immediately from the properties by which I defined characters in § 1 of my paper \emph{On the representation of finite groups by linear substitutions}, Sitzungsberichte 1897.

2. Between the characters $\chi$ of a group $\Hgr$ and the characters $\psi$ of one of its subgroups $\Ggr$ there is, by U. § 1 (5.), the relation
\begin{equation}\tag{1.}
 \sum_\lambda r_{\varkappa\lambda}\chi^{(\lambda)}(R)=\frac{h}{g h_\rho}\sum_{P\in(\rho)}\psi^{(\varkappa)}(P).
\end{equation}
The positive integral numbers $r_{\varkappa\lambda}$ are independent of the class $(\rho)$ of the element $R$. The $g_\rho$ elements $P$ of the $\rho^{\ten}$ class of $\Hgr$ which are contained in $\Ggr$ are distributed among several, say among $m$, distinct classes of $\Ggr$. Let them be represented by $P_1,P_2,\ldots,P_m$. If the class of $P_\mu$ consists of $l_\mu$ elements, then
\[
 \sum_{(\rho)}\psi(P)=l_1\psi(P_1)+l_2\psi(P_2)+\cdots+l_m\psi(P_m).
\]
The group of the elements of $\Ggr$ which are commutable with $P$ is contained in the group of the elements of $\Hgr$ which are commutable with $P$. In the case which I shall consider in the sequel, these two groups are identical for each of the $g_\rho$ elements $P$ of $\Ggr$ which belong to a fixed class $(\rho)$ of $\Hgr$, or every element of $\Hgr$ commutable with $P$ is contained in $\Ggr$. The number of elements of $\Hgr$ commutable with $P$ is $h/h_\rho$, and the number of elements of $\Ggr$ commutable with $P_\mu$ is $g/l_\mu$. Therefore
\[
 l_1=l_2=\cdots=l_m=\frac{g_\rho}{m},	\qquad \frac{h}{h_\rho}=\frac{g}{l_\mu}=\frac{gm}{g_\rho},
\]
and by (1.)
\begin{equation}\tag{2.}
 \psi^{(\varkappa)}(P_1)+\psi^{(\varkappa)}(P_2)+\cdots+\psi^{(\varkappa)}(P_m)=\sum_\lambda r_{\varkappa\lambda}\chi^{(\lambda)}(R).
\end{equation}

\section*{§ 5.}

I now turn to the proof of the rule which I stated in § 3 for the computation of the characters of the alternating group $\Hgr$ of degree $n$. In doing so I assume that this rule has already been proved for every alternating group $\Hgr_1$ of degree $n_1<n$.

I use the notations (2.) and (3.), § 3, but, if $n$ is odd, I assume $s>1$. I divide the $n$ symbols into the first $c_1$ and the last $n-c_1$, and form the symmetric and the alternating group $\Hgr'_1$, $\Hgr_1$ and $\Hgr'_2$, $\Hgr_2$ for those $c_1$ and for these $n-c_1$ symbols. If $c_1=1$, then the following developments require a modification, which, for the sake of simplicity, I omit.

Then $\Ggr=\Hgr_1\Hgr_2$ is a subgroup of the alternating group $\Hgr$ of degree $n$. Let $R_1$ consist of a cycle on the first $c_1$ symbols, and $R_2$ of $s-1$ cycles on respectively $c_2,\ldots,c_s$ of the last $n-c_1$ symbols. Let $T_1$ ($T_2$) be a fixed negative permutation of $\Hgr'_1$ ($\Hgr'_2$). Further let $T_1^{-1}R_1T_1=S_1$ and $T_2^{-1}R_2T_2=S_2$. If then
\[
 p_1=c_1,	\quad p_2=c_2\cdots c_s,	\quad \eps_1=(-1)^{\frac12(p_1-1)},	\quad \eps_2=(-1)^{\frac12(p_2-1)},
\]
there are in $\Hgr_1$ two conjugate characters $\phi_1(P_1)$ and $\psi_1(P_1)=\phi_1(T_1^{-1}P_1T_1)$, for which
\[
 \phi_1(R_1)=\psi_1(S_1)=\frac12(\eps_1+\sqrt{\eps_1p_1}),	\quad
 \phi_1(S_1)=\psi_1(R_1)=\frac12(\eps_1-\sqrt{\eps_1p_1}),
\]
and in $\Hgr_2$ two conjugate characters $\phi_2(P_2)$ and $\psi_2(P_2)=\phi_2(T_2^{-1}P_2T_2)$, for which
\[
 \phi_2(R_2)=\psi_2(S_2)=\frac12(\eps_2+\sqrt{\eps_2p_2}),	\quad
 \phi_2(S_2)=\psi_2(R_2)=\frac12(\eps_2-\sqrt{\eps_2p_2}),
\]
while for every other class $\phi_1(Q_1)=\psi_1(Q_1)$ and $\phi_2(Q_2)=\psi_2(Q_2)$ is a rational integer. From this, by § 4, there arise four distinct characters of $\Ggr$, $\phi_1\phi_2$, $\phi_1\psi_2$, $\phi_2\psi_1$, $\phi_2\psi_2$. If
\[
 R_1R_2=R,	\qquad R_1S_2=S,	\qquad T_1T_2=T,
\]
then $T$ is a positive permutation, and
\[
 T^{-1}RT=S_1S_2,	\qquad T^{-1}ST=S_1R_2,	\qquad T_2^{-1}RT_2=S.
\]
Therefore the two classes $(R_1R_2)$ and $(S_1S_2)$ of $\Ggr$ unite into one class $(R)$ of $\Hgr$, and the two classes $(R_1S_2)$ and $(S_1R_2)$ into one class $(S)$; but the two classes $(R)$ and $(S)$ of $\Hgr$ are distinct, because $T_2$ is odd.

The class $(R)$ can contain only such elements of $\Ggr$ as belong to one of the two classes $(R_1R_2)$ or $(S_1S_2)$. For every permutation $P$ of the class $(R)$ consists of $s$ cycles $C_1,C_2,\ldots,C_s$ on respectively $c_1,c_2,\ldots,c_s$ symbols. If it is to belong to the group $\Ggr$, then it must be $P=P_1P_2$, where $P_1$ moves only the first $c_1$ symbols and $P_2$ only the last $n-c_1$ symbols. But this is possible only when $P_1=C_1$ and $P_2=C_2\cdots C_s$, because $c_1$ is the smallest of the numbers $c_1,c_2,\ldots,c_s$, and hence cannot be represented as a sum of some of the numbers $c_2,\ldots,c_s$.

This is the reason why I have divided the $n$ symbols precisely in this way into two divisions of $n_1=c_1$ and $n_2=n-c_1$ symbols. The same purpose is also reached if one sets $n_1=c_2$, $n_2=c_1+c_3+\cdots+c_s$, and one then has in addition the advantage that $n_1=1$ cannot occur.

Accordingly $P_1$ belongs to the class $(R_1)$ or $(S_1)$, and $P_2$ to the class $(R_2)$ or $(S_2)$. As shown above, however, only the combinations
\[
 (R)=(R_1R_2)+(S_1S_2),	\qquad (S)=(R_1S_2)+(S_1R_2)
\]
are possible.

The elements of $\Hgr$ commutable with $R$ are, by (1.), § 1, all contained in $\Ggr$. Therefore there is a composite character $\phi(P)$ of $\Hgr$ for which
\[
 \phi(R)=\frac14(\eps_1+\sqrt{\eps_1p_1})(\eps_2+\sqrt{\eps_2p_2})+\frac14(\eps_1-\sqrt{\eps_1p_1})(\eps_2-\sqrt{\eps_2p_2})
\]
holds. Put
\[
 p=p_1p_2=c_1c_2\cdots c_s,	\qquad \eps=\eps_1\eps_2=(-1)^{\frac12(p-1)}=(-1)^{\frac12(n-s)},
\]
and let $\psi(P)=\phi(T_2^{-1}PT_2)$ be the character conjugate with $\phi(P)$. Then
\[
 \phi(R)=\psi(S)=\frac12(\eps+\sqrt{\eps p}),	\qquad \phi(S)=\psi(R)=\frac12(\eps-\sqrt{\eps p}).
\]
For every element $Q_2$ of $\Hgr_2$ which does not belong to the class $(R_2)$, one has $\phi_2(Q_2)=\phi_2(T_2^{-1}Q_2T_2)$. If one now applies the formula (1.), § 4 to the element $Q=R_1Q_2$ of $\Hgr$, then on the right, beside every term $\phi_1(R_1)\phi_2(Q_2)$, there also stands the term $\phi_1(T_1^{-1}R_1T_1)\phi_2(T_2^{-1}Q_2T_2)$, and the sum of these two terms is a rational number which does not change if one replaces $Q$ by $T_2^{-1}QT_2$. Therefore $\phi(Q)=\psi(Q)$ is a rational integer, and the same holds for an element $Q=Q_1R_2$ or $Q=Q_1Q_2$.

Consequently $\Hgr$ has a composite character $\phi-\psi=\vartheta$, which has the values
\begin{equation}\tag{1.}
 \vartheta(R)=\sqrt{\eps p},\qquad \vartheta(S)=-\sqrt{\eps p},\qquad \vartheta(Q)=0
\end{equation}
if $(Q)$ is any class distinct from the two fixed conjugate classes $(R)$ and $(S)$. For each of these two classes $h_\rho=h/p$, and hence
\[
 \sum_\rho h_\rho\vartheta_\rho\overline{\vartheta}_\rho=2h,
\]
where $\overline{\vartheta}_\rho$ denotes the complex quantity conjugate with $\vartheta_\rho$. Thus, if, as in § 2, $\vartheta(P)=\sum r_\varkappa\phi^{(\varkappa)}(P)$, then $\sum r_\varkappa^2=2$; hence two of the numbers $r_\varkappa$ are equal to $\pm1$, and the others are equal to $0$. Now $\phi^{(\varkappa)}(E)$ is a positive integer and $\vartheta(E)=0$. Consequently $\vartheta(P)$ is the difference of two simple characters of $\Hgr$. The two characters $\phi$ and $\psi$ of $\Hgr$ as whose difference $\vartheta$ has been obtained above may be composite. From now on I denote by $\phi$ and $\psi$ the two simple characters of $\Hgr$ whose difference is $\phi-\psi=\vartheta$.

If $\chi(P)$ is one of the $u$ characters of $\Hgr$ which arise from a pair of associated characters of $\Hgr'$, then $\chi(R)=\chi(S)$, and by (1.) consequently $\sum h_\rho\vartheta_\rho\chi_\rho=0$. Hence neither of the two characters $\phi$ or $\psi$ is equal to $\chi$; rather, $\phi$ and $\psi$ belong to the $2v$ characters of $\Hgr$ which arise from the $v$ self-associated characters of $\Hgr'$, and which split into $v$ pairs of conjugate characters. The character $\omega$ conjugate with $\phi$ is therefore distinct from $\phi$. If $\chi=\phi+\omega$, then $\chi(R)=\chi(S)$ and hence $\sum h_\rho\vartheta_\rho\chi_\rho=0$, or $\sum h_\rho(\phi_\rho-\psi_\rho)(\phi_\rho+\omega_\rho)=0$, and because $\phi$ is distinct from $\psi$ and from $\omega$, $\sum h_\rho\psi_\rho\omega_\rho=h$. Consequently $\omega=\psi$, or $\phi$ and $\psi$ are conjugate characters.

I assign the pair of conjugate classes $(R)$ and $(S)$ and the pair of conjugate characters $\phi$ and $\psi$ to one another. Then by (1.) to distinct pairs $(R),(S)$ correspond distinct pairs $\phi,\psi$, and conversely. The values of a character are algebraic integers. Therefore $\chi(R)=\chi(S)=2\phi(R)-\sqrt{\eps p}$ is an odd (rational) number, whereas $\chi(Q)=2\phi(Q)=2\psi(Q)$ is an even number.

In order to be able to use the induction step, $s>1$ has so far been assumed. Thus, if $n$ is odd, there remains a pair of conjugate classes $(R),(S)$ whose permutations consist of a single cycle on $n$ symbols, and hence also a pair of conjugate characters $\phi,\psi$, whose sum and difference I denote by $\chi$ and $\vartheta$. If then $(\overline{R}),(\overline{S})$ is any other pair of conjugate classes, and $\overline{\phi},\overline{\psi}$ the corresponding pair of conjugate characters, and if $\overline{\vartheta}=\overline{\phi}-\overline{\psi}$, then $\sum h_\rho\overline{\vartheta}_\rho\phi_\rho=0$, and hence $\phi(\overline{R})=\phi(\overline{S})=\psi(\overline{R})=\psi(\overline{S})$. Thus, if $(Q)$ is any class distinct from $(R)$ and $(S)$ and also distinct from its conjugate, then $\vartheta(Q)=0$, $\phi(Q)=\psi(Q)=\frac12\chi(Q)$. By § 2 one has $\phi(R)=\psi(S)$ and $\phi(S)=\psi(R)$. The elements $R$ and $R^{-1}$ belong to the same class when $n\equiv1\pmod 4$, and to two conjugate classes when $n\equiv -1\pmod 4$. In the first case, therefore, $\phi(R)$ and $\phi(S)$ are real quantities; in the second, conjugate complex quantities. Now $\sum h_\rho\phi_\rho\vartheta_\rho=h$, hence if $\eps=(-1)^{\frac12(n-1)}=+1$, then $h_R(\phi(R)\vartheta(R)+\phi(S)\vartheta(S))=h$, and therefore
\[
 \vartheta(R)^2=\frac{h}{h_R}=p=n.
\]
But if $\eps=-1$, then $h_R(\phi(R)\vartheta(S)+\phi(S)\vartheta(R))=h$, and therefore $\vartheta(R)^2=-p$. Hence $\vartheta(R)=\sqrt{\eps p}$, $\vartheta(S)=-\sqrt{\eps p}$, and consequently the same laws hold for this assignment as above.

By the condition that $\chi(R)=\chi(S)$ is odd, but that otherwise $\chi(Q)$ is even, the assignment between a pair of conjugate characters whose sum is $\chi$ and the corresponding pair of conjugate classes $(R),(S)$ is completely determined; and, as asserted above, $\chi$ must have the characteristic
\[
 \ch{a_1\ a_2\ \cdots\ a_s}{a_1\ a_2\ \cdots\ a_s}
\]
when
\[
 c_1=2a_1+1,	\quad c_2=2a_2+1,	\quad \cdots,	\quad c_s=2a_s+1
\]
are the orders of the cycles of $R$. For by formula (11.) § 7 of my paper \emph{On the characters of the symmetric group},
\[
 \chi\ch{a_1\ \cdots\ a_s}{a_1\ \cdots\ a_s}(R)=\eps=(-1)^{\frac12(p-1)}=(-1)^{\frac12(n-s)}
\]
is odd. With this the $k$ characters of $\Hgr$ are completely determined.

If the class $(\varkappa)$ and its conjugate correspond to the character $\chi^{(\varkappa)}$ and its conjugate, then by formula (9.), § 4 of the just-cited paper and by formula (6.), § 3 of this paper,
\[
 f^{(\varkappa)}=\frac12\frac{n!\,\Delta(a_1,\ldots,a_s)^2}{(a_1!\cdots a_s!)^2\prod(a_\alpha+a_\beta+1)}.
\]
Now
\[
 \frac{h}{h_\varkappa}=(2a_1+1)\cdots(2a_s+1),
\]
and consequently
\begin{equation}\tag{2.}
 \frac{h_\varkappa}{f^{(\varkappa)}}=\left(\frac{a_1!\cdots a_s!\prod'(a_\alpha+a_\beta+1)}{\Delta(a_1,\ldots,a_s)}\right)^2
\end{equation}
is the square of an integer. The product $\prod'$ is to be extended only over the $\frac12s(s-1)$ pairs of distinct indices. The square root occurring in $\chi^{(\varkappa)}_\varkappa$ can therefore be drawn from
\[
 \frac{\eps h}{h_\varkappa}\quad\hbox{or}\quad \frac{\eps h}{f^{(\varkappa)}}.
\]
It can be rational in exceptional cases, for example when $n=9$ and $R$ consists of a cycle on $9$ symbols.
% ===== END WORK 061 =====

\clearpage

% ===== BEGIN WORK 070: 070_Theorie_der_hyperkomplexen_Gr_en =====
\frobtitle{Theory of hypercomplex quantities}

The methods which I developed in my paper \emph{On the prime factors of the group determinant}, Sitzungsberichte 1896 (cited below as \emph{Gr.}), for the investigation of the properties of the determinant of a finite group, also suffice for the investigation of an arbitrary system of hypercomplex quantities formed from $n$ fundamental quantities. Such quantities are the subject of Mr. \textsc{Molien}'s fundamental memoir \emph{On systems of higher complex numbers}, Math. Ann. vol. 41 (cited below as \emph{Mol.}).

He has made the understanding of his work somewhat more difficult by suppressing the calculation as far as possible and by trying to present the conceptual content of the proofs in a pure form. Thus I have not become clear whether the auxiliary means developed by him suffice for a rigorous proof of Theorem 25. Conversely, he corrected a fallacy made in § 5 in a note in the 42nd volume of the Math. Ann. by a rather lengthy calculation. By using, in place of the equation which he calls the \textsc{Killing} equation, a linear combination $S(x)+T(y)$ of the matrices of the two antistrophic groups expressed in different variables, I succeed in simplifying precisely this proof considerably.

In addition, as I recently observed, there is in his proof of Theorem 19 a not insignificant gap. But notwithstanding these small defects, his pioneering and thoughtful work, which he carried through despite a rather imperfect apparatus with unremitting perseverance and penetrating acuteness, marks one of the most important advances in the part of algebra which one calls group theory. For the reasons stated, I consider it appropriate to take up the investigation anew with the auxiliary means which I developed especially in my paper \emph{On commutable matrices}, Sitzungsberichte 1896.

The methods of \textsc{Lie} are not used at all in this purely algebraic work. On the other hand, my presentation has many points of contact with that of \textsc{Dedekind} in his memoir \emph{On the theory of complex quantities formed from $n$ principal units}, Göttinger Nachrichten 1885 (cited below as \emph{Ded.}).

The basis of my investigation is formula (2.) § 5. With its help I show directly in § 6 that the variables occurring in the different prime factors of the group determinant are all independent of one another, and thereby avoid Theorems 3, 4, and 5 of \textsc{Molien}, as well as the proof of Theorem 25. And with the same formula I prove in § 7 that the elementary divisors of the determinant of a simple group are all linear. The whole development rests on these two results.

For the classification of groups, for which only insufficient beginnings have hitherto been made, the most suitable invariants turn out to be not so much the exponents and degrees of the elementary divisors into which the determinants of the two antistrophic groups decompose, but rather (§ 9) the same numbers for the parastrophic matrix $R(\xi)$, which has hitherto hardly been noticed, especially the rank of $R(\xi)$; and, for non-commutative groups, above all the elementary invariants of the matrix $uR(\xi)+vR'(\xi)$ and of the matrices $S(x)+T(y)$ and $R(\xi)+R'(\eta)$ depending on $2n$ variables. These invariants are significant even when, for example, the determinant of $R(\xi)$ or of $R(\xi)+R'(\eta)$ vanishes identically.

Only after completing this investigation did I become aware of Mr. \textsc{Cartan}'s excellent memoir \emph{Sur les groupes bilinéaires et les systèmes de nombres complexes}, Ann. de Toulouse, tome XII, 1898, in which he derives the results of \textsc{Molien}, apparently without knowing his work. The path taken by Mr. \textsc{Cartan} has nothing at all in common either with \textsc{Molien}'s methods or with those used here. The transformation of the basis, the starting point and the final goal of his investigation, has been avoided by me as long as possible (§ 9). The invariant properties, independent of every representation of the group, with which I begin, appear in his work only at the end through interpretation of a normal form of the group obtained by a long series of transformations, whose aim becomes clear only at the end of the development. Thus the difference between the two methods is the same as that between the method of \textsc{Weierstrass} and that of \textsc{Kronecker} in the theory of pencils of bilinear forms. A particularly noteworthy formula of Mr. \textsc{Cartan} (§ 65, (37.)), which is not found in \textsc{Molien}, I had obtained (§ 12, (5.)) in the simplest way by decomposing the determinant $|S(x)+T(y)|$ into prime factors.

\section*{§ 1.}

The coordinates $x_1,x_2,\ldots,x_n$ of the hypercomplex quantity
\[
\fnum{1} x=\eps_1x_1+\eps_2x_2+\cdots+\eps_nx_n
\]
formed from $n$ fundamental quantities $\eps_1,\eps_2,\ldots,\eps_n$ can assume all real and complex values. I call the totality of these quantities, which reproduce themselves under addition and multiplication by ordinary quantities, a group $(\eps)$, if, in addition, the product of any two again belongs to the system, and if for the multiplication of any three the associative law $(xy)z=x(yz)$ holds.

The fundamental quantities may be connected by linear relations. I call the number of independent ones among them the order of the group. This case can easily be reduced to the one in which the fundamental quantities are independent. I therefore assume that $n$ is the order of the group $(\eps)$ under consideration.

If the multiplication of the fundamental quantities is carried out according to the law
\[
\fnum{2} \eps_\beta\eps_\gamma=\sum_\alpha a_{\alpha\beta\gamma}\eps_\alpha,
\]
then from the associative principle and from the independence of the fundamental quantities there result, between the coefficients occurring here, the relations
\[
\fnum{3} \sum_\kappa a_{\kappa\alpha\beta}a_{\gamma\kappa\delta}
      =\sum_\kappa a_{\kappa\beta\delta}a_{\gamma\alpha\kappa}.
\]
If then $x=yz$ is the product of the two quantities
\[
 y=\sum_\lambda\eps_\lambda y_\lambda,\qquad
 z=\sum_\lambda\eps_\lambda z_\lambda,
\]
then
\[
\fnum{4} x_\alpha=\sum_{\beta,\gamma}a_{\alpha\beta\gamma}y_\beta z_\gamma,
\]
or, if $\xi_1,\xi_2,\ldots,\xi_n$ are variables whose system I denote by $\xi$ and call a parameter,
\[
\fnum{5} \xi(x)=x(\xi)=\sum_\alpha \xi_\alpha x_\alpha
  =F(\xi,y,z)=\sum_{\alpha,\beta,\gamma}a_{\alpha\beta\gamma}\xi_\alpha y_\beta z_\gamma.
\]
Putting
\[
\fnum{6} r_{\alpha\beta}(\xi)=\sum_\kappa a_{\kappa\alpha\beta}\xi_\kappa,\quad
s_{\alpha\beta}(y)=\sum_\lambda a_{\alpha\lambda\beta}y_\lambda,\quad
t_{\alpha\beta}(z)=\sum_\mu a_{\alpha\beta\mu}z_\mu,
\]
therefore
\[
s_{\alpha\beta}(y)=\frac{\partial x_\alpha}{\partial z_\beta},\qquad
t_{\alpha\beta}(z)=\frac{\partial x_\alpha}{\partial y_\beta},
\]
the trilinear form is consequently
\[
\fnum{7} F(\xi,y,z)=\sum_{\beta,\gamma}r_{\beta\gamma}(\xi)y_\beta z_\gamma
=\sum_{\alpha,\gamma}s_{\alpha\gamma}(y)\xi_\alpha z_\gamma
=\sum_{\alpha,\beta}t_{\alpha\beta}(z)\xi_\alpha y_\beta .
\]
The matrices of degree $n$ formed from those three systems of linear functions I denote by
\[
 R(\xi)=R_\xi=(r_{\alpha\beta}(\xi)),\qquad
 S(y)=S_y=(s_{\alpha\beta}(y)),\qquad
 T(z)=T_z=(t_{\alpha\beta}(z)).
\]

By $f(x)$ I always understand a function of the coordinates $x_1,x_2,\ldots,x_n$ of the quantity $x$; by $S(x)$ or $S_x$ I understand a matrix whose elements are (linear) functions of these coordinates. When the matter is a function of the hypercomplex quantity $x$ itself, I shall use the notation $f((x))$.

The conditions (3.) of the associative principle are more conveniently discussed when one combines them by multiplying them by variables and adding. Thus, when $x$ and $y$ are two arbitrary quantities, one obtains (\emph{Ded.} (36.))
\[
\fnum{8} \sum_\lambda s_{\alpha\lambda}(x)t_{\lambda\beta}(y)
      =\sum_\lambda t_{\alpha\lambda}(y)s_{\lambda\beta}(x),
\]
or more simply
\[
\fnum{9} S_xT_y=T_yS_x,
\]
and
\[
\fnum{10} \sum_\lambda r_{\alpha\lambda}(\xi)s_{\lambda\beta}(x)
       =\sum_\lambda t_{\lambda\alpha}(x)r_{\lambda\beta}(\xi),
\]
or
\[
\fnum{11} R_\xi S_x=T'_xR_\xi,
\]
where $T'$ is the matrix conjugate to $T$.

I call $S(x)$ and $|S(x)|$ the group matrix and the group determinant, $T(x)$ and $|T(x)|$ the antistrophic group matrix and group determinant, and $R(\xi)$ and $|R(\xi)|$ the parastrophic matrix and determinant. The theorem contained in formula (9.) can be stated and conversely interpreted as follows:

\emph{The group matrix $S(x)$ is commutable with the antistrophic matrix $T(y)$. If the parastrophic determinant does not vanish identically, then every matrix independent of $x$ which, for every value of $x$, is commutable with $S(x)$ ($T(x)$), can be brought to the form $T(y)$ ($S(y)$).}

For let the parameter $\xi$ be chosen so that the determinant of the matrix $R(\xi)=R=(r_{\alpha\beta})$ does not vanish, and let $U=(u_{\alpha\beta})$ be any matrix of degree $n$ independent of $x$ which is commutable with $S(x)$ for every $x$. Then one can determine $y$ so that for $\lambda=1,2,\ldots,n$
\[
 \sum_\kappa \xi_\kappa u_{\kappa\lambda}
 =\sum_\alpha r_{\lambda\alpha}y_\alpha
 =\sum_\kappa \xi_\kappa t_{\kappa\lambda}(y)
\]
(by (7.)) holds. Now $SU=US$, hence
\[
 \sum_\lambda s_{\alpha\lambda}(x)u_{\lambda\beta}
 =\sum_\lambda u_{\alpha\lambda}s_{\lambda\beta}(x),
\]
and consequently
\[
 \sum_\lambda\xi_\lambda s_{\lambda\alpha}(x)u_{\alpha\beta}
 =\sum_{\lambda,\mu}\xi_\lambda u_{\lambda\mu}s_{\mu\beta}(x)
 =\sum_{\lambda,\mu}\xi_\lambda t_{\lambda\mu}(y)s_{\mu\beta}(x).
\]
Thus, by (8.),
\[
 \sum_{\mu,\lambda}\xi_\lambda s_{\lambda\alpha}(x)u_{\alpha\beta}
 =\sum_{\mu,\lambda}\xi_\lambda s_{\lambda\alpha}(x)t_{\alpha\beta}(y).
\]
But by (7.)
\[
 \sum_\kappa \xi_\kappa s_{\kappa\lambda}(x)=\sum_\alpha r_{\alpha\lambda}x_\alpha.
\]
Substituting this, comparison of the coefficients of $x_\alpha$ gives
\[
 \sum_\lambda r_{\alpha\lambda}u_{\lambda\beta}
 =\sum_\lambda r_{\alpha\lambda}t_{\lambda\beta}(y),
\]
and hence, since $|R|$ is different from zero, $u_{\alpha\beta}=t_{\alpha\beta}(y)$, $U=T(y)$.

If in particular $U=E=(e_{\alpha\beta})$ is the principal matrix, then one can determine a quantity $y=e$ so that $T(e)=E$ (\emph{Ded.} (44.)). By (11.) one then also has $S(e)=E$. This quantity
\[
\fnum{12} e=e_1\eps_1+e_2\eps_2+\cdots+e_n\eps_n
\]
is called the principal unit. Its coordinates satisfy the conditions
\[
\fnum{13} s_{\alpha\beta}(e)=t_{\alpha\beta}(e)=e_{\alpha\beta}
\]
or
\[
\fnum{14} \sum_\lambda a_{\alpha\lambda\beta}e_\lambda=e_{\alpha\beta},\qquad
\sum_\mu a_{\alpha\beta\mu}e_\mu=e_{\alpha\beta}.
\]
Thus if $S(x)=S(y)$ or $T(x)=T(y)$ or $S(x)=T(y)$, then $x=y$; if $R(\xi)=R(\eta)$ or $R(\xi)=R'(\eta)$, then $\xi=\eta$.

\section*{§ 2.}

If $x=yz$, hence
\[
\fnum{1} x_\alpha=\sum_{\lambda,\mu}a_{\alpha\lambda\mu}y_\lambda z_\mu
=\sum_\mu s_{\alpha\mu}(y)z_\mu=\sum_\lambda t_{\alpha\lambda}(z)y_\lambda,
\]
then
\[
s_{\alpha\beta}(x)=\sum_\lambda a_{\alpha\lambda\beta}x_\lambda
 =\sum_{\kappa,\lambda,\mu}a_{\lambda\kappa\mu}a_{\alpha\lambda\beta}y_\kappa z_\mu
 =\sum_\lambda s_{\alpha\lambda}(y)s_{\lambda\beta}(z)
\]
and
\[
t_{\alpha\beta}(x)=\sum_\lambda a_{\alpha\beta\lambda}x_\lambda
 =\sum_{\kappa,\lambda,\mu}a_{\lambda\kappa\mu}a_{\alpha\beta\lambda}y_\kappa z_\mu
 =\sum_\lambda t_{\lambda\beta}(z)t_{\alpha\lambda}(y).
\]
By $f(yz)$ or $S(yz)$ I denote what is obtained from $f(x)$ or $S(x)$ when the $n$ variables $x_\alpha$ in it are replaced by the coordinates (1.) of the product $yz$. Then one obtains
\[
\fnum{2} S(yz)=S(y)S(z),\qquad T(yz)=T(z)T(y).
\]

If $S(x)=x_1E_1+x_2E_2+\cdots+x_nE_n$, then consequently
\[
\fnum{3} E_\beta E_\gamma=\sum_\alpha a_{\alpha\beta\gamma}E_\alpha,
\]
and since $S(x)=0$ only when $x=0$, the $n$ constant matrices $E_1,E_2,\ldots,E_n$ are linearly independent. They therefore form a representation (cf. § 16) of the group $(\eps)$, and from this it follows that the assumption of the independence of the fundamental quantities is compatible with the quadratic equations (2.) § 1, whose coefficients satisfy equations (3.) § 1 and the inequalities (1.) § 3, or that no linear relations between the fundamental quantities flow from these conditions. Only on the basis of this certainty are proofs admissible, such as the one linked to formula (3.) § 3, in which the hypercomplex quantities themselves are used.

If further $T(x)=x_1F_1+x_2F_2+\cdots+x_nF_n$, then
\[
\fnum{4} F_\beta F_\gamma=\sum_\alpha a_{\alpha\gamma\beta}F_\alpha.
\]
Since the associative law holds for matrices, $F_1,F_2,\ldots,F_n$ form the basis of a group, the antistrophic group $(\eps')$. Another representation of $(\eps)$ is furnished by the matrix $T'$ conjugate to $T$. One could therefore also call $T'$ the matrix of the group $(\eps)$ and $S'$ that of the antistrophic group $(\eps')$.

Let $\xi$ be a parameter, and put
\[
\fnum{5^*} \zeta_\mu=\sum_{\kappa,\lambda}a_{\kappa\lambda\mu}\xi_\kappa y_\lambda
=\sum_\lambda r_{\lambda\mu}(\xi)y_\lambda
=\sum_\kappa s_{\kappa\mu}(y)\xi_\kappa.
\]
Then
\[
 r_{\alpha\beta}(\zeta)=\sum_\mu a_{\mu\alpha\beta}\zeta_\mu
 =\sum_{\kappa,\mu}\xi_\kappa s_{\kappa\mu}(y)a_{\mu\alpha\beta}.
\]
This is the coefficient of $z_\beta$ in
\[
 \sum_\beta r_{\alpha\beta}(\zeta)z_\beta
 =\sum_{\kappa,\mu}\xi_\kappa s_{\kappa\mu}(y)t_{\mu\alpha}(z)
 =\sum_{\kappa,\mu}\xi_\kappa t_{\kappa\mu}(z)s_{\mu\alpha}(y)
 =\sum_{\mu,\beta}s_{\mu\alpha}(y)r_{\mu\beta}(\xi)z_\beta .
\]
Therefore
\[
\fnum{5} r_{\alpha\beta}(\zeta)=\sum_\mu s_{\mu\alpha}(y)r_{\mu\beta}(\xi)
\]
or
\[
\fnum{5} R(\zeta)=S'(y)R(\xi).
\]
Putting finally
\[
\fnum{6^*} \eta_\lambda=\sum_{\kappa,\mu}a_{\kappa\lambda\mu}\xi_\kappa z_\mu
=\sum_\mu r_{\lambda\mu}(\xi)z_\mu
=\sum_\kappa t_{\kappa\lambda}(z)\xi_\kappa,
\]
one obtains
\[
 r_{\alpha\beta}(\eta)=\sum_\lambda a_{\lambda\alpha\beta}\eta_\lambda
 =\sum_{\kappa,\lambda}\xi_\kappa t_{\kappa\lambda}(z)a_{\lambda\alpha\beta}
\]
and therefore
\[
 \sum_\alpha r_{\alpha\beta}(\eta)y_\alpha
 =\sum_{\kappa,\lambda}\xi_\kappa t_{\kappa\lambda}(z)s_{\lambda\beta}(y)
 =\sum_{\kappa,\lambda}\xi_\kappa s_{\kappa\lambda}(y)t_{\lambda\beta}(z)
 =\sum_{\alpha,\lambda}r_{\alpha\lambda}(\xi)t_{\lambda\beta}(z)y_\alpha,
\]
hence
\[
\fnum{6} r_{\alpha\beta}(\eta)=\sum_\lambda r_{\alpha\lambda}(\xi)t_{\lambda\beta}(z)
\]
or
\[
\fnum{6} R(\eta)=R(\xi)T(z).
\]

With suitably changed notation the following results from this:

\emph{If the parastrophic determinant does not vanish identically, and if for the value $\xi$ of the parameter the determinant of the matrix $R_\xi=R=(r_{\alpha\beta})$ is different from zero, then $R_\xi$ is transformed by the substitution $\xi_\alpha=\sum_\beta r_{\alpha\beta}x_\beta$ into $RT_x$, and by the conjugate substitution $\xi_\beta=\sum_\alpha r_{\alpha\beta}x_\alpha$ into $S'_xR$.}

\emph{Under this condition the group determinant is identically equal to the antistrophic group determinant, and, up to the factor $|R|$, also equal to the function into which the parastrophic determinant passes under any one of those two conjugate substitutions; and all three determinants agree in their elementary divisors.}

If $\xi,\eta,\zeta,\tau$ are arbitrary parameters, it follows from (9.) § 1 that $R_\eta^{-1}R_\xi$ is commutable with $R_\tau^{-1}R'_\zeta$, or that the matrix
\[
 R_\zeta R_\eta^{-1}R_\xi R_\tau^{\prime -1}
\]
is independent of $\zeta$.

\section*{§ 3.}

If $|R(\xi)|$ does not vanish identically, then neither of the two determinants $|S(x)|$ or $|T(x)|$ vanishes identically, and there is a number $e$ for which $S(e)=T(e)=E$. If, however, $|R(\xi)|=0$ identically, then I add to the previous assumptions the further one that neither of the two determinants
\[
\fnum{1} |S(x)|,\qquad |T(x)|
\]
is identically zero. Choose $z$ so that $|S(z)|$ and $|T(z)|$ are both different from zero. Then (cf. for example \emph{Mol.} Theorem 1), since $|S(z)|$ is different from zero, one can determine a number $e$ satisfying the equation $ze=z$. If furthermore $x$ is an arbitrary quantity, then, since $|T(z)|$ is different from zero, one can determine $y$ so that $yz=x$. Then also $xe=x$ and $zex=zx$, hence, since $|S(z)|$ is different from zero, $ex=x$. Since moreover by (2.) § 2 $S(z)=S(z)S(e)$ and $T(z)=T(e)T(z)$, one has $S(e)=T(e)=E$.

The theorem of § 1 remains correct when neither $|S(x)|$ nor $|T(x)|$ vanishes identically, even if $|R(\xi)|=0$ identically. For if one puts $\sum_\kappa u_{\alpha\kappa}e_\kappa=y_\alpha$, then
\[
 \sum_{\beta,\kappa}u_{\alpha\beta}s_{\beta\kappa}(x)e_\kappa
 =\sum_{\gamma,\kappa}s_{\alpha\gamma}(x)u_{\gamma\kappa}e_\kappa,
\]
so
\[
 \sum_\beta u_{\alpha\beta}x_\beta=\sum_\gamma s_{\alpha\gamma}(x)y_\gamma=\sum_\beta t_{\alpha\beta}(y)x_\beta,
\]
and consequently $u_{\alpha\beta}=t_{\alpha\beta}(y)$.

If $|S(x)|$ is different from zero for a definite quantity $x$, then one can determine $y$ so that $xy=e$. Then by (2.) § 2 $T(y)T(x)=T(e)=E$. Thus, if $|S(x)|$ is different from zero for the value $x$, then $|T(x)|$ is also different from zero, and conversely. Thus if $|S(x)|=0$, then also $|T(x)|=0$. Every prime factor $\Phi(x)$ of the group determinant $|S(x)|$ is therefore also contained in the antistrophic determinant $|T(x)|$, and conversely. In the two decompositions
\[
\fnum{2} |S(x)|=\prod \Phi(x)^s,\qquad |T(x)|=\prod \Phi(x)^t
\]
exactly the same prime factors $\Phi(x)$ occur, though the exponents $s$ and $t$ can be different, but only if $|R(\xi)|=0$ identically.

The simplest example for this possibility is supplied, as \textsc{Molien} has informed me, by the unique non-commutative group of order $n=3$, for which
\[
 x_1=y_1z_1,\qquad x_2=y_2z_2,\qquad x_3=y_1z_3+y_3z_2
\]
holds. Accordingly
\[
 R(\xi)=
 \begin{pmatrix}\xi_1&0&\xi_3\\0&\xi_2&0\\0&\xi_3&0\end{pmatrix},\quad
 S(y)=
 \begin{pmatrix}y_1&0&0\\0&y_2&0\\0&y_3&y_1\end{pmatrix},\quad
 T(z)=
 \begin{pmatrix}z_1&0&0\\0&z_2&0\\z_3&0&z_2\end{pmatrix},
\]
and therefore, decomposed into elementary divisors,
\[
 |S(x)|=x_1^2x_2,\qquad |T(x)|=x_1x_2^2,\qquad
 |uR(\xi)+vR'(\xi)|=-(\xi_1+\xi_2)\xi_3^2uv(u+v).
\]

From equation (2.) § 2 it follows, by repeated application, that
\[
 S_x^\nu=S(x^\nu),\qquad T_x^\nu=T(x^\nu).
\]
By $u,v,w$ I here always understand ordinary, not hypercomplex, quantities. If $g(u)$ is an integral function of the variable $u$, so that $g((x))$ is an integral function of the hypercomplex quantity $x$ itself, then
\[
\fnum{3} g(S_x)=S(g((x))),\qquad g(T_x)=T(g((x))).
\]
If therefore one of the expressions $g((x))$, $g(S_x)$, or $g(T_x)$ vanishes for all values of $x$, then the other two also vanish (\emph{Mol.} § 4). The equation of lowest degree which a matrix $S_x=S$ satisfies, $g(S)=0$, is called its reduced equation. The function $g(u)$ is obtained by dividing the characteristic determinant of $S$, the determinant of the matrix
\[
\fnum{4} uE-S(x)=S(ue-x),
\]
by the greatest common divisor of its minors of degree $(n-1)$. By (3.), $S$ and $T$ satisfy the same reduced equation $g(S)=0$ and $g(T)=0$, and this is at the same time the equation of lowest degree which a variable quantity $x$ satisfies, $g((x))=0$. Thus, even if the determinants of the two antistrophic groups can be different, their first elementary divisors always agree.

The reduced function $g(u)$ vanishes for each of the characteristic roots of $S$, that is, for the roots of the characteristic equation $|uE-S|=0$. Thus it is proved anew that the two antistrophic group determinants contain exactly the same prime factors.

\section*{§ 4.}

If $\Theta(x)=|S(x)|$ is the group determinant, then by (2.) § 2 $\Theta(yz)=\Theta(y)\Theta(z)$. In each prime factor $\Phi(x)$ of $\Theta(x)$ I suppose the constant factor to have been chosen so that $\Phi(e)=1$. If then $\Phi(x)$ is such a factor, or also a product of several, then also
\[
\fnum{1} \Phi(yz)=\Phi(y)\Phi(z).
\]
Conversely, every homogeneous function $\Phi(x)$ which possesses this property must be a product of prime factors of $\Theta(x)$ (\emph{Gr.} § 1). For determine $y$ so that $xy=\Theta(x)e$. Then $y_1,y_2,\ldots,y_n$ are found from $n$ linear equations with determinant $\Theta(x)$ and are therefore integral functions of $x_1,x_2,\ldots,x_n$. If $r$ is the degree of $\Phi(x)$, then by (1.)
\[
 \Phi(x)\Phi(y)=\Theta(x)^r.
\]
The roots $u_1,u_2,\ldots,u_r$ of the equation $\Phi(ue-x)=0$ I call the characteristic roots of $\Phi(x)$. Then
\[
\fnum{2} \Phi(ue-x)=u^r-\Phi_1(x)u^{r-1}+\Phi_2(x)u^{r-2}-\cdots\pm\Phi_r(x)
=(u-u_1)(u-u_2)\cdots(u-u_r).
\]
The sum of the $r$ roots, $\Phi_1(x)$, or
\[
\fnum{3} u_1+u_2+\cdots+u_r=\chi(x)=\sum_\alpha\chi_\alpha x_\alpha,
\]
I call the trace of $\Phi(x)$. By the system of its $n$ coefficients, which I likewise denote by $\chi$ and call the parameter characteristic for $\Phi(x)$, the function $\Phi(x)$ is completely determined (\emph{Mol.} § 3; \emph{Gr.} § 3). For if
\[
 g(u)=a(u+v_1)(u+v_2)\cdots(u+v_m)
\]
is an integral function of $u$, then also
\[
 y=g((x))=a(x+v_1e)(x+v_2e)\cdots(x+v_me)
\]
and hence by (1.)
\[
 \Phi(y)=a^r\Phi(x+v_1e)\Phi(x+v_2e)\cdots\Phi(x+v_me).
\]
Putting here
\[
 x=u_ie,
\]
one obtains
\[
 \Phi(y)=g(u_1)g(u_2)\cdots g(u_r).
\]
If now $g(u)$ is replaced by $g(u)-v$, one sees that $g(u_1),g(u_2),\ldots,g(u_r)$ are the characteristic roots of $\Phi(y)$. Therefore
\[
\fnum{4} \chi(g((x)))=g(u_1)+g(u_2)+\cdots+g(u_r)
\]
and in particular
\[
\fnum{5} \chi(x^\nu)=u_1^\nu+u_2^\nu+\cdots+u_r^\nu,
\]
so for $\nu=0$
\[
\fnum{6} r=\chi(e)=\sum_\alpha\chi_\alpha e_\alpha.
\]
By the well-known relations between the coefficients of an equation and the power sums of its roots, the following is therefore an identity for prime functions of first degree:
\[
\fnum{7} \chi(x)\chi(y)-\chi(xy)=0,
\]
and for those of second degree:
\[
\fnum{8} \chi(x)\chi(y)\chi(z)-\chi(x)\chi(yz)-\chi(y)\chi(zx)-\chi(z)\chi(xy)+\chi(xyz)+\chi(xzy)=0.
\]
The general law for the formation of these equations I set forth in \emph{Gr.} § 3.

In particular, the number of different linear prime functions is equal to the number of solutions of equations (2.) § 1 by ordinary numbers (\emph{Ded.} (19.)). As I shall show in § 6, these solutions are all linearly independent.

If $\Theta(x)$ is different from zero for a definite value $x$, then by (1.) § 4
\[
 \Phi(x^{-1}yx+ue)=\Phi(x^{-1})\Phi(y+ue)\Phi(x)=\Phi(y+ue)
\]
and hence $\chi(x^{-1}yx)=\chi(y)$; or, replacing $y$ by $xy$ (\emph{Mol.} Theorem 14; \emph{Gr.} § 3, (13.)),
\[
\fnum{9} \chi(xy)=\chi(yx).
\]
The matrix of this symmetric bilinear form $F(\chi,x,y)$ is
\[
\fnum{10} R_\chi=R'_\chi.
\]
Thus I call $\xi$ a symmetric parameter when it satisfies the equations $r_{\alpha\beta}(\xi)=r_{\beta\alpha}(\xi)$ or $R_\xi=R'_\xi$; then $\chi$ is such a parameter.

A quantity $x$ is called an invariant quantity of the group $(\eps)$ if it is commutable with every quantity $y$ of this group, $xy=yx$. Accordingly
\[
 \sum_\gamma a_{\alpha\gamma\beta}x_\gamma=\sum_\gamma a_{\alpha\beta\gamma}x_\gamma,\qquad
 t_{\alpha\beta}(x)=s_{\alpha\beta}(x).
\]
Conversely, if the variability of $x$ is restricted by the equation $S(x)=T(x)$, then $x$ is an invariant quantity of $(\eps)$.

Also the more general equation $S(x)=T(y)$, or
\[
 x_1E_1+\cdots+x_nE_n=y_1F_1+\cdots+y_nF_n,
\]
can, by (14.) § 1, hold only when $x=y$ is an invariant quantity. The group formed by the invariant quantities is therefore the greatest common divisor of the two antistrophic groups (3.) and (4.) § 2.

If $y$ varies without restriction but $x$ is an invariant variable, then $S(x)=T(x)$ is commutable with $S(y)$, and therefore the characteristic roots $u_1,u_2,\ldots$ of the matrix $S(x)$ can be assigned to the characteristic roots $v_1,v_2,\ldots$ of the matrix $S(y)$ in such a way that $u_1v_1,u_2v_2,\ldots$ are the characteristic roots of $S(x)S(y)=S(xy)$. Hence also
\[
 \Phi(ue-xy)=(u-u_1v_1)(u-u_2v_2)\cdots(u-u_rv_r),
\]
where $u_1,u_2,\ldots,u_r$ depend only on $x$, and $v_1,v_2,\ldots,v_r$ only on $y$. Putting $y=e$, one sees that $u_1,u_2,\ldots,u_r$ are the characteristic roots of $\Phi(x)$; putting $x=e$, that $v_1,v_2,\ldots,v_r$ are those of $\Phi(y)$.

In determining the way in which these roots are assigned to each other in the above decomposition, I restrict myself to the case in which $\Phi(y)$ is a prime function. Then $\Phi(ue-y)$, as a function of $u$, is irreducible, that is, it cannot be represented as a product of integral functions of $u$ whose coefficients are rational functions of the $n$ unrestricted variable quantities $y_1,y_2,\ldots,y_n$. Since $u_1$ is independent of these, the function $\Phi(ue-u_1y)$ of the variable $u$ is irreducible in the same sense. The function $\Phi(ue-xy)$ has the factor $u-u_1v_1$ in common with it and consequently is identical with it: $\Phi(ue-xy)=\Phi(ue-u_1y)$. Hence $\chi(xy)=u_1\chi(y)$, and therefore, for $y=e$, $\chi(x)=ru_1$. Putting $u=0$, $y=e$, one obtains $\Phi(x)=u_1^r$. Thus the theorem results (\emph{Gr.} § 6):

\emph{If $\chi(x)$ is the trace of the prime factor $\Phi(x)$ of degree $r$ of the group determinant, if $y$ is unrestrictedly variable and $x$ is an invariant quantity, then}
\[
\fnum{11} \chi(x)\chi(y)=r\chi(xy)=\chi(e)\chi(xy),
\]
\emph{and}
\[
\fnum{12} \Phi(x)=\left(\frac1r\chi(x)\right)^r,\qquad
 \Phi(ue-x)=\left(u-\frac1r\chi(x)\right)^r
\]
\emph{is the $r$th power of a linear function, and the $r$ characteristic roots of $\Phi(x)$ are all equal to $\frac1r\chi(x)$.}

\section*{§ 5.}

If $t=xy$, then $t+ux=x(y+ue)$ and consequently, if $\Phi(x)$ is a factor of a power of $\Theta(x)$, $\Phi(t+ux)=\Phi(x)\Phi(y+ue)$. Comparing the coefficients of $u^{r-1}$ in this, one obtains
\[
 \sum_\alpha \frac{\partial\Phi(x)}{\partial x_\alpha}t_\alpha=\Phi(x)\chi(y),
\]
and therefore, since $t_\alpha=\sum_\beta s_{\alpha\beta}(x)y_\beta$,
\[
 \sum_\alpha\frac{\partial\Phi(x)}{\partial x_\alpha}s_{\alpha\beta}(x)=\Phi(x)\chi_\beta .
\]
In this way, and by the decomposition $yx+ux=(y+ue)x$, one obtains (\emph{Gr.} § 5, (3.))
\[
\fnum{1} \sum_\kappa\frac{\partial\Phi(x)}{\partial x_\kappa}s_{\kappa\alpha}(x)
=\sum_\kappa\frac{\partial\Phi(x)}{\partial x_\kappa}t_{\kappa\alpha}(x)
=\Phi(x)\chi_\alpha .
\]
Thus if in formula $(5^*.)$ § 2 one puts $\xi_\kappa=\frac{\partial\Phi}{\partial x_\kappa}$, then $\zeta_\alpha=\Phi(x)\chi_\alpha$, and by (5.) and (6.) § 2 one therefore has (\emph{Gr.} § 4, (1.))
\[
\fnum{2} S'(x)R\!\left(\frac{\partial\Phi}{\partial x}\right)
=R\!\left(\frac{\partial\Phi}{\partial x}\right)T(x)
=\Phi(x)R(\chi),
\]
where the matrix
\[
 R\!\left(\frac{\partial\Phi}{\partial x}\right)
 =\left(\sum_\kappa a_{\kappa\alpha\beta}\frac{\partial\Phi}{\partial x_\kappa}\right)
\]
is meant. If, when $x$ is replaced by $ue-x$, this passes into
\[
 R_1u^{r-1}+R_2u^{r-2}+\cdots+R_r,
\]
then by (2.) § 4
\[
 (R_1u^{r-1}+\cdots+R_r)(Eu-T)
=R_\chi(u^r-\Phi_1u^{r-1}+\Phi_2u^{r-2}-\cdots\pm\Phi_r).
\]
Comparing the coefficients of $u^r,u^{r-1},\ldots,u^0$, one obtains $r+1$ equations. Multiplying these on the right by $T^r,T^{r-1},\ldots,T^0$ and adding, one obtains (\emph{Gr.} § 4, (5.))
\[
\fnum{3} R_\chi(T^r-\Phi_1T^{r-1}+\Phi_2T^{r-2}-\cdots\pm\Phi_r)=0
\]
or more simply
\[
 R_\chi\Phi(xE-eT)=0.
\]
In the same way one obtains $\Phi(xE-eS')R_\chi=0$, hence, taking the conjugate matrices, by (10.) § 4
\[
\fnum{3} R_\chi\Phi(xE-eS)=0,\qquad R_\chi\Phi(xE-eT)=0.
\]
By (3.) § 3, for $\Phi(xE-eS)$ one can also write $S(\Phi(x-((x))e))$, that is, $S(y)$, where $y=g((x))$ and $g(u)=\Phi(x-ue)$.

One also reaches formula (2.), on which the following development essentially rests, by putting $t=xyz$ (or $yzx$), and by comparing, in the equation
\[
 \sum_\kappa\frac{\partial\Phi(x)}{\partial x_\kappa}t_\kappa=\Phi(x)\chi(yz),
\]
where
\[
 t_\kappa=\sum_{\lambda,\mu,\alpha,\beta}a_{\kappa\lambda\mu}x_\lambda a_{\mu\alpha\beta}y_\alpha z_\beta,\qquad
 \chi(yz)=\sum_{\alpha,\beta}r_{\alpha\beta}(\chi)y_\alpha z_\beta,
\]
the coefficients of $y_\alpha z_\beta$.

\section*{§ 6.}

The smallest number of independent linear combinations of the variables by means of which a function or a system of functions can be expressed I call its \emph{linear rank}, or also only its \emph{rank}, as soon as confusion with the concept of the rank of a determinant or matrix is excluded. The linear rank $m$ of a quadratic function
\[
\fnum{1} \sum_{\alpha,\beta}r_{\alpha\beta}(\chi)x_\alpha x_\beta
=F(\chi,x,x)=\chi(x^2)=\Phi_1^2-2\Phi_2
\]
is equal to the rank of the symmetric matrix $R_\chi$ formed from its coefficients, and equal to the linear rank of the system of the $n$ linear functions
\[
\fnum{2} \sum_\beta r_{\alpha\beta}(\chi)x_\beta,
\]
the half derivatives of the quadratic function. By any $m$ precisely of these linear functions which are independent of one another, $\chi(x^2)$ can be represented. Likewise its covariant, the symmetric bilinear form
\[
\fnum{3} \sum_{\alpha,\beta}r_{\alpha\beta}(\chi)x_\alpha y_\beta
=F(\chi,x,y)=\chi(xy)=\chi(yx)
\]
can be expressed by those $m$ variables and by the $m$ among the $n$ variables $\sum_\beta r_{\alpha\beta}(\chi)y_\beta$ which are independent of one another.

The trilinear function $\chi(xyz)=\chi((xy)z)$ is a bilinear function of the coordinates of $xy$ and of $z$, and can therefore be expressed by the $m$ variables which are independent among the $n$ variables $\sum_\beta r_{\alpha\beta}z_\beta$. It is equal to $\chi(x(yz))=\chi(z(xy))=\chi((zx)y)$, and therefore also depends only on the variables $\sum_\beta r_{\alpha\beta}x_\beta$ and the variables $\sum_\beta r_{\alpha\beta}y_\beta$. The same holds for $\chi(xyzt)=\chi(yztx)=\chi(ztxy)=\chi(txyz)$. Therefore also
\[
 \chi(x^\nu)=u_1^\nu+u_2^\nu+\cdots+u_r^\nu
\]
depends only on the variables (2.), and hence also the product $\Phi(x)=u_1u_2\cdots u_r$, which is an integral function of those power sums. Conversely, by the linear combinations of the variables $x_1,x_2,\ldots,x_n$ on which $\Phi(x)$ depends, the functions $\Phi_1(x),\Phi_2(x),\ldots$ occurring in the expansion of $\Phi(ue-x)$ can also be expressed, and therefore also $\chi(x^2)$.

\emph{I. If $\Phi(x)$ is a product of prime factors of the group determinant, and $\chi(x)$ the trace of $\Phi(x)$, then the linear rank of $\Phi(x)$ is equal to that of the quadratic function $\chi(x^2)$, hence equal to the rank of the matrix $R_\chi$, and the function $\Phi(x)$ can be expressed by the derivatives of $\chi(x^2)$.}

Now let $\Phi,\Phi',\Phi'',\ldots$ be the $k$ different prime factors of $\Theta$, and $\chi,\chi',\chi'',\ldots$ their traces. If then $C,C',C'',\ldots$ are any matrices of degree $n$, I shall show that a relation
\[
\fnum{4} CR_\chi+C'R_{\chi'}+C''R_{\chi''}+\cdots=0
\]
can exist only when separately
\[
 CR_\chi=C'R_{\chi'}=C''R_{\chi''}=\cdots=0
\]
holds. For by (2.) § 5
\[
 R_\chi=R\!\left(\frac{\partial l(\Phi)}{\partial x}\right)T(x)
\]
and hence, since $|T(x)|$ is different from zero,
\[
 CR\!\left(\frac{\partial l(\Phi)}{\partial x}\right)
+C'R\!\left(\frac{\partial l(\Phi')}{\partial x}\right)
+C''R\!\left(\frac{\partial l(\Phi'')}{\partial x}\right)+\cdots=0.
\]
Now let $\Psi=\Phi'\Phi''\cdots$ be the product of the $k-1$ prime functions different from $\Phi$. Multiplying by $\Phi\Psi$, one sees that the elements of the matrix $\Psi(x)CR\!\left(\frac{\partial\Phi}{\partial x}\right)$ are all divisible by the function $\Phi(x)$ of degree $r$; hence, since $\Phi$ and $\Psi$ are relatively prime, so are the elements of the matrix $CR\!\left(\frac{\partial\Phi}{\partial x}\right)$. But since these are only of degree $(r-1)$, they must be zero. Hence
\[
 0=CR\!\left(\frac{\partial l(\Phi)}{\partial x}\right)T(x)=CR_\chi.
\]
One can also reach this result by means of formula (3.) § 5. According to it, $R_\chi\Phi(xE-eS)=0$, and since any two integral functions of the matrix $S$ are commutable with one another, also
\[
 R_\chi\Psi(xE-eS)=R_\chi\Phi''(xE-eS)\Phi'(xE-eS)\cdots=0.
\]
Consequently also
\[
 CR_\chi\bigl(\Psi(xE-eS)+\Phi(xE-eS)\bigr)=0.
\]
The determinant of the matrix in parentheses, a function of $S$, is different from zero because the function $\Psi(x-ue)+\Phi(x-ue)$ of the variable $u$ vanishes for no root of the characteristic equation $\Theta(x-ue)=0$. Consequently $CR_\chi=0$.

Let $m,m',m'',\ldots$ be the ranks of the matrices $R_\chi,R_{\chi'},R_{\chi''},\ldots$. Then among the variables $\sum_\alpha r_{\alpha\beta}(\chi)x_\beta$ exactly $m$ are independent, among the variables $\sum_\alpha r_{\alpha\beta}(\chi')x_\beta$ exactly $m'$ are independent, among the variables $\sum_\alpha r_{\alpha\beta}(\chi'')x_\beta$ exactly $m''$, and so on. The $m+m'+m''+\cdots$ linear combinations of the $n$ variables $x_1,x_2,\ldots,x_n$ obtained in this way are all independent of one another. For an identity in $x_1,x_2,\ldots,x_n$ of the form
\[
 \sum_\beta c_\beta\left(\sum_\alpha r_{\alpha\beta}(\chi)x_\alpha\right)
+\sum_\beta c'_\beta\left(\sum_\alpha r_{\alpha\beta}(\chi')x_\alpha\right)
+\sum_\beta c''_\beta\left(\sum_\alpha r_{\alpha\beta}(\chi'')x_\alpha\right)+\cdots=0
\]
can exist only if the $k$ partial sums vanish separately. The $n$ equations
\[
 \sum_\alpha c_\alpha r_{\alpha\beta}(\chi)
+\sum_\alpha c'_\alpha r_{\alpha\beta}(\chi')
+\sum_\alpha c''_\alpha r_{\alpha\beta}(\chi'')+\cdots=0
\]
namely have the form (4.) if $C$ is a matrix in which one row consists of the elements $c_1,c_2,\ldots,c_n$, while the elements of the other rows vanish.

Since $\sum_\alpha x_\alpha\chi_\alpha=\sum_\alpha e_\alpha\bigl(\sum_\beta r_{\alpha\beta}(\chi)x_\beta\bigr)$ is a linear combination of the $n$ variables (2.), the $k$ functions
\[
\fnum{5} \chi(x),\quad \chi'(x),\quad \chi''(x),\quad\ldots
\]
are also linearly independent.

Furthermore, $\chi(x^2)$ can be expressed as a quadratic function of $m$ independent variables with non-vanishing determinant, $\chi'(x^2)$ as a function of $m'$ variables of rank $m'$, and so on; and these $m+m'+m''+\cdots$ variables are all independent of one another. If therefore $c,c',c'',\ldots$ are constants, the rank of the quadratic function
\[
 c\chi(x^2)+c'\chi'(x^2)+c''\chi''(x^2)+\cdots
\]
is equal to the sum of the ranks of the individual summands. The rank of $c\chi(x^2)$ is $0$ or $m$ according as $c=0$ or not.

\emph{II. The rank of the matrix $cR_\chi+c'R_{\chi'}+c''R_{\chi''}+\cdots$ is equal to the sum of the ranks of the matrices $cR_\chi,c'R_{\chi'},c''R_{\chi''},\ldots$.}

\emph{The linear rank of a product of prime factors of the group determinant is equal to the sum of the ranks of its different prime factors.}

\section*{§ 7.}

I denote the traces of the determinants $|S(x)|$ and $|T(x)|$ by
\[
\fnum{1} \sigma(x)=\sum_\alpha \sigma_\alpha x_\alpha=\sum_\kappa s_{\kappa\kappa}(x),\qquad
\tau(x)=\sum_\alpha \tau_\alpha x_\alpha=\sum_\kappa t_{\kappa\kappa}(x).
\]
Their coefficients are
\[
\fnum{2} \sigma_\alpha=\sum_\kappa a_{\kappa\alpha\kappa},\qquad
\tau_\alpha=\sum_\kappa a_{\kappa\kappa\alpha},
\]
and by (2.) § 3 one has
\[
\fnum{3} \sigma=s\chi+s'\chi'+s''\chi''+\cdots=\sum s\chi,
\qquad
\tau=t\chi+t'\chi'+t''\chi''+\cdots=\sum t\chi.
\]
Therefore, by Theorem II of § 6, the matrices $R_\sigma$ and $R_\tau$ both have rank $m+m'+m''+\cdots$.

A group $(\eps)$ for which the determinant of the symmetric matrix
\[
\fnum{4} R_\sigma=P=(p_{\alpha\beta})
\]
is different from zero I call a \textsc{Dedekind} group (\textsc{Ded}. (27.)), since \textsc{Dedekind} first recognized, at least for commutative groups, the significance of this condition. For such a group, by the theorem of § 2, $\Theta(x)=|S(x)|=|T(x)|$, hence $\sigma=\tau$, and the rank of $R_\sigma$ is
\[
 n=m+m'+m''+\cdots=\sum m.
\]
A \textsc{Dedekind} group can therefore also be defined as one whose order is equal to the linear rank of its determinant. Let
\[
 \psi=\chi+\chi'+\chi''+\cdots=\sum\chi
\]
be the trace of the product of the $k$ different prime factors of $\Theta$,
\[
 \Psi(x)=\Phi\Phi'\Phi''\cdots=\prod\Phi(x).
\]
Then, again by Theorem II of § 6, $R_\psi$ has rank $\sum m=n$. Consequently $|R_\psi|$ is different from zero. But by (3.) § 5 one has
\[
 R_\psi\Psi(Se-Ex)=0,
\]
and consequently
\[
\fnum{5} \Psi(Se-Ex)=0.
\]
Since the function
\[
\fnum{6} g(u)=\Psi(ue-x)=\prod\Phi(ue-x)
\]
has no repeated factors, it follows that $g(S)=0$ is the reduced equation for the matrix $S$ (cf. \textsc{Mol}. Theorem 24).

This can also be seen as follows. By (2.) § 5,
\[
 S'R\!\left(\frac{\partial\Theta}{\partial x}\right)=\Theta R_\sigma.
\]
Now $\Theta S'^{-1}=\overline{S'}$ is the adjugate matrix of $S'$, and hence (\textsc{Ded}. (63.))
\[
\fnum{7} \overline{S'}=R\!\left(\frac{\partial\Theta}{\partial x}\right)R_\sigma^{-1}.
\]
The elements of the matrix $\overline{S'}$ are the minors of degree $n-1$ of $\Theta$. They are therefore linear combinations of the derivatives of $\Theta$ with respect to the $n$ variables $x_1,x_2,\ldots,x_n$. Conversely, the derivatives of a determinant are linear combinations of its minors. Therefore the greatest common divisor of the minors is equal to that of the derivatives, hence, when $x$ is replaced by $ue-x$, equal to
\[
 \prod\bigl(\Phi(ue-x)\bigr)^{s-1}.
\]
But by the rule mentioned in § 3 one obtains the reduced equation $g(S)=0$ by dividing
\[
 \Theta(ue-x)=\prod\bigl(\Phi(ue-x)\bigr)^s
\]
by that expression; consequently $g(u)$ is the function (6.). The roots $u_1,u_2,\ldots,u_p$ of the equation $g(u)=0$ are the $p=r+r'+r''+\cdots$ distinct roots among the $n$ roots $u_1,u_2,\ldots,u_n$ of the equation $\Theta(ue-x)=0$.

\section*{§ 8.}

If $A$ and $B$ are two commuting matrices of degree $n$, then their characteristic roots $a_1,a_2,\ldots,a_n$ and $b_1,b_2,\ldots,b_n$ can be assigned to one another in such a way that $f(a_1,b_1),f(a_2,b_2),\ldots,f(a_n,b_n)$ are the characteristic roots of the matrix $f(A,B)$, and this assignment is independent of the choice of the integral function $f(u,v)$. In § 3 of the preceding paper I further proved:

\emph{I. If the characteristic determinants of two (or more) mutually commuting matrices $A$ and $B$ decompose into only linear elementary divisors, then every matrix $f(A,B)$ has the same property.}

\emph{II. If, moreover, always $b_\kappa=b_\lambda$ whenever $a_\kappa=a_\lambda$, then $B$ is an integral function of $A$.}

Now
\[
\fnum{1} g(u,x)=\prod\Phi(ue-x)
\]
is an integral function of degree $p$ in $u$, whose coefficients are integral functions of the $n$ variables $x_1,x_2,\ldots,x_n$. The $p$ roots $u_1,u_2,\ldots,u_p$ of the equation $g(u,x)=0$ are all mutually distinct. If $X=S(x)$, then the elementary divisors of $|uE-X|$ are all linear and $X$ satisfies the equation $g(X,x)=0$.

This equation, being an identity in $x_1,x_2,\ldots,x_n$, is therefore also satisfied by every special system of values of the variables. For such a system the numbers $u_1,u_2,\ldots,u_p$ need not all be distinct, the elementary divisors of $|uE-X|$ need not all be linear, and $g(X,x)=0$ need not be the reduced equation for $X$. But if $x$ is so chosen that $u_1,u_2,\ldots,u_p$ are distinct, then $g(X,x)=0$ is the reduced equation for the matrix $X$. For the equation $|uE-X|=0$ cannot have more than $p$ distinct roots, and if $\psi(X)=0$, then $\psi(u)$ must vanish for every root of this characteristic equation. Since then the equation $g(u,x)=0$ has no repeated root, the elementary divisors of $|uE-X|$ are all linear.

Let $x$ be chosen as just described, and let $y$ be a quantity commuting with $x$. Then the matrix $Y=S(y)$ also commutes with $X=S(x)$. Therefore the determinant
\[
\fnum{2} |uX+vY+wE|=|S(ux+vy+we)|=\prod\bigl(\Phi(ux+vy+we)\bigr)^s
\]
is a product of linear factors $uu_\alpha+vv_\alpha+w$, where $u_\alpha$ and $v_\alpha$ are corresponding characteristic roots of the two commuting matrices $X$ and $Y$; hence also
\[
\fnum{3} \Phi(ux+vy+we)=\prod_\alpha^{r}(uu_\alpha+vv_\alpha+w)
\]
is a product of linear functions of $u,v,w$. Setting $v=0$ or $u=0$, one sees that to the $r$ characteristic roots $u_1,\ldots,u_r$ of $\Phi(x)$ there are assigned the $r$ roots $v_1,\ldots,v_r$ of $\Phi(y)$ in a definite order, and consequently also to the $p$ roots $u_1,\ldots,u_p$ of $g(u,x)=0$ the $p$ roots $v_1,\ldots,v_p$ of $g(v,y)=0$. Since, further, $u_1,\ldots,u_p$ are distinct, formulas (2.) and (3.) show that equal roots $u_\kappa=u_\lambda$ among the $n$ characteristic roots $u_1,\ldots,u_n$ of $X$ have assigned to them equal roots $v_\kappa=v_\lambda$ of $Y$.

If $l$ is a constant, then $v_1-lu_1,\ldots,v_p-lu_p$ are the roots of the equation $g(w,y-lx)=0$. If $l$ is not equal to one of the fractions $\frac{v_\alpha-v_\beta}{u_\alpha-u_\beta}$, then those $p$ roots are all distinct, and consequently the elementary divisors of the characteristic determinant of the matrix $Z=Y-lX$ are all linear. Since $X$ and $Z$ commute, and $Y=lX+Z$ is a function of $X$ and $Z$, by Theorem I the elementary divisors of $|vE-Y|$ are also all linear. By Theorem II, therefore, $Y$ is an integral function of $X$ (cf. \textsc{Mol}. Theorem 27).

The equation of lowest degree satisfied by $x$ has degree $p$. Therefore $x^0,x^1,\ldots,x^{p-1}$ are linearly independent, while every integral function of $x$ depends on these $p$. The linear equations $xy-yx=0$ among the $n$ unknowns $y_1,y_2,\ldots,y_n$ therefore have exactly $p$ independent solutions. The matrix of their coefficients, $S(x)-T(x)$, therefore has rank $n-p$.

\emph{III. If, for a definite quantity $x$ of a \textsc{Dedekind} group, the $r+r'+r''+\cdots$ roots of the equation $\prod\Phi(ue-x)=0$ are all distinct, then the matrix $S(x)-T(x)$ has rank $n-(r+r'+r''+\cdots)$. Then every quantity $y$ commuting with $x$ is an integral function of $x$, and for every such quantity $y$ the elementary divisors of $|vE-S(y)|$ are all linear.}

\section*{§ 9.}

I shall bring the theory of \textsc{Dedekind} groups begun in §§ 7 and 8 to completion in §§ 13 to 16. I now return to the general theory. In place of the basis $\eps_1,\eps_2,\ldots,\eps_n$ of the group $(\eps)$ one may introduce a new basis by a linear substitution
\[
\fnum{1} \eps_\beta=\sum_\alpha c_{\alpha\beta}\bar\eps_\alpha
\]
where the determinant of the matrix $C=(c_{\alpha\beta})$ is different from zero. If then
\[
\fnum{2} \sum \eps_\alpha x_\alpha=\sum\bar\eps_\alpha\bar x_\alpha,
\]
one has
\[
\fnum{3} \bar x_\alpha=\sum_\beta c_{\alpha\beta}x_\beta.
\]
The quantities $\frac{\partial\Phi}{\partial x_\alpha},\xi_\alpha,\chi_\alpha,\sigma_\alpha$ are cogredient with the basic numbers $\eps_\alpha$; the quantities $y_\alpha,z_\alpha$ are cogredient with the coordinates $x_\alpha$ and contragredient with the basic numbers $\eps_\alpha$. Further,
\[
\fnum{4}
\begin{aligned}
S(x)&=C^{-1}\bar S(\bar x)C,
& T(x)&=C^{-1}\bar T(\bar x)C,\\
R(\xi)&=C'\bar R(\bar\xi)C,
& R'(\xi)&=C'\bar R'(\bar\xi)C.
\end{aligned}
\]
Therefore the exponents and the degrees of the elementary divisors into which the determinants of the matrices $R(\xi),S(x),T(x)$ decompose are invariants of the group. The same also holds for the determinants of the matrices depending on $2n$ variables,
\[
\fnum{5} S(x)+T(y),\qquad R(\xi)+R'(\eta),
\]
and in particular also for the elementary divisors of the determinant of the matrix of the bilinear form $uF(\xi,y,z)+vF(\xi,z,y)$ in $y$ and $z$,
\[
\fnum{6} uR(\xi)+vR'(\xi).
\]
To illustrate this remark I consider some examples from the work of Mr. \textsc{Study}, \emph{Über Systeme von complexen Zahlen}, Göttinger Nachrichten 1889. If (\textsc{St}. IX)
\[
\fnum{7} R(\xi)=
\begin{pmatrix}
\xi_1&\xi_2&\xi_3&\xi_4\\
\xi_2&\xi_4&\xi_4&0\\
\xi_3&-\xi_4&c\xi_4&0\\
\xi_4&0&0&0
\end{pmatrix},
\]
then those elementary divisors are
\[
\fnum{8} \xi_4^3\xi_4(u+v)(u+v)\bigl((u-v)^2+c(u+v)^2\bigr),
\]
where the last factor, for $c=0$, is a quadratic elementary divisor, and otherwise decomposes into two distinct linear elementary divisors. This formula shows that $c$ is a (non-numerical) invariant of the group.

If further (\textsc{St}. XIV)
\[
\fnum{9} R(\xi)=
\begin{pmatrix}
\xi_1&\xi_2&\xi_3&\xi_4\\
\xi_2&0&\xi_4&0\\
\xi_3&-\xi_4&0&0\\
\xi_4&0&0&0
\end{pmatrix},
\]
then the elementary divisors of $|uR(\xi)+vR'(\xi)|$ are
\[
\fnum{10} \xi_4^2\xi_4^2(u+v)(u+v)(u-v)(u-v),
\]
thus essentially different from the preceding ones (for $c=0$). I shall not, however, enter here more closely into the classification of the groups.

It is possible that the quantities (2.) still form a group $(\mathfrak S)$ even when the coordinates are subjected to certain linear relations. Then $(\mathfrak S)$ is called a subgroup of $(\eps)$. Its basic numbers $\mathfrak s_1,\mathfrak s_2,\ldots$ are linear combinations of the $n$ basic numbers $\eps_\alpha$, and in order that $(\mathfrak S)$ be a group, each of the products $\mathfrak s_\alpha\mathfrak s_\beta$ must be a linear combination of $\mathfrak s_1,\mathfrak s_2,\ldots$.

If also each of the products $\eps_\alpha\mathfrak s_\beta$ is a linear combination of $\mathfrak s_1,\mathfrak s_2,\ldots$, then $(\mathfrak S)$ is called an invariant subgroup\footnote{If $\mathfrak H$ is a finite group of order $n$ and $\mathfrak G$ is a subgroup of order $n/m$, if $R$ is an element of $\mathfrak H$ and $P$ an element of $\mathfrak G$, then the hypercomplex quantities $\sum R x_R$ form a group $(\eps)$ of order $n$, and the quantities $\sum P x_P$ form a subgroup $(\mathfrak S)$. It is not, for that reason, an invariant subgroup of $(\eps)$ when $\mathfrak G$ is such a subgroup of $\mathfrak H$. But if $\mathfrak H=\mathfrak G+\mathfrak G A+\mathfrak G B+\cdots$, then one obtains an invariant subgroup of order $n-m$ of $(\eps)$ when the variability of the $n$ coordinates $x_R$ is restricted by the $n$ equations $\sum_P x_P=0$, $\sum_P x_{AP}=0$, $\sum_P x_{BP}=0$, $\cdots$. The group $(\eta)$ complementary to it and homomorphic with $(\eps)$ stands to the group $\mathfrak H/\mathfrak G$ in the same relation as $(\eps)$ stands to $\mathfrak H$.} of $(\eps)$. Then the formulas (2.) § 1 define a group $(\eta)$ even when one subjects the $n$ basic numbers $\eps_\alpha$ to the linear relations $\mathfrak s_1=0,\mathfrak s_2=0,\ldots$, or when any two quantities of $(\eps)$ are regarded as equal mod. $(\mathfrak S)$ if their difference belongs to the group $(\mathfrak S)$. It is called the group complementary to the invariant subgroup $(\mathfrak S)$ and homomorphic with $(\eps)$ (accompanying number system in \textsc{Molien}).

To see this, transform the basis so that $\mathfrak s_1,\mathfrak s_2,\ldots$ become the basic numbers $\eps_{m+1},\ldots,\eps_n$. Then the subgroup $(\mathfrak S)$ is determined by the $m$ linear equations $x_1=0,\ldots,x_m=0$ among the coordinates. In order that $(\mathfrak S)$ be a subgroup it is necessary and sufficient that $a_{\alpha\beta\gamma}=0$ if $\alpha\leq m$, $\beta>m$ and $\gamma>m$; in order that $(\mathfrak S)$ be an invariant subgroup, that $a_{\alpha\beta\gamma}=0$ if $\alpha\leq m$ and even only one of the two numbers $\beta$ or $\gamma$ is $>m$ (\textsc{Mol}. Theorem 2). Then the equations (2.) § 1 still define a group when $\eps_{m+1}=\cdots=\eps_n=0$ is put. For in formulas (4.) § 1, $x_1,\ldots,x_m$ depend only on $y_1,\ldots,y_m$ and $z_1,\ldots,z_m$. In the two matrices $S(x)$ and $T(x)$ all elements therefore vanish which the first $m$ rows have in common with the last $n-m$ columns. Accordingly
\[
 S=\begin{pmatrix}S_1&0\\ S_{21}&S_2\end{pmatrix},
 \qquad
 T=\begin{pmatrix}T_1&0\\ T_{21}&T_2\end{pmatrix},
\]
where, for example, the submatrix $S_1$ contains only the first $m$ rows and columns of $S$. From $S(y)S(z)=S(yz)$ it follows that $S_1(y)S_1(z)=S_1(yz)$, and since $S_1(x)=x_1E_1+\cdots+x_mE_m$ depends only on $x_1,\ldots,x_m$, the equations
\[
 \eps_\beta\eps_\gamma=\sum_{\alpha=1}^{m}a_{\alpha\beta\gamma}\eps_\alpha\qquad(\beta,\gamma=1,2,\ldots,m)
\]
define a group $(\eta)$ which arises from $(\eps)$ by putting $\eps_{m+1}=\cdots=\eps_n=0$. Since $|S|=|S_1||S_2|$ and $|T|=|T_1||T_2|$, $(\eta)$ also satisfies the conditions (1.) § 3. Further, the determinant of the group $(\eps)$ is divisible by the determinant of every group $(\eta)$ homomorphic with $(\eps)$, and the same holds for the groups antistrophic to $(\eps)$ and $(\eta)$.

If, moreover, $a_{\alpha\beta\gamma}=0$ whenever $\alpha>m$ and one of the two numbers $\beta$ or $\gamma$ is $\leq m$, then $(\eps)$ decomposes into two groups, each of which is both an invariant subgroup of $(\eps)$ and a group homomorphic with $(\eps)$.

\section*{§ 10.}

Let $m,m',m'',\ldots$ be the ranks of the $k$ matrices $R_\chi,R_{\chi'},R_{\chi''},\ldots$. Then $\Phi$ depends only on the $m$ independent ones among the $n$ linear functions $\sum_\beta r_{\alpha\beta}(\chi)x_\beta$, $\Phi'$ only on the $m'$ independent ones among the functions $\sum_\beta r_{\alpha\beta}(\chi')x_\beta$, and so on, and these $m+m'+m''+\cdots$ variables are all independent of one another. One may therefore choose the basic numbers so that $\Phi$ depends only on $x_1,\ldots,x_m$, $\Phi'$ only on $x_{m+1},\ldots,x_{m+m'}$, and so on. Then $\chi(x^2)=\sum r_{\alpha\beta}(\chi)x_\alpha x_\beta$ also depends only on $x_1,\ldots,x_m$. Hence $r_{\alpha\beta}(\chi)=0$ if $\alpha$ or $\beta>m$, and since $R_\chi$ has rank $m$, the determinant of degree $m$
\[
 |r_{\kappa\lambda}(\chi)|\qquad(\kappa,\lambda=1,2,\ldots,m)
\]
is different from zero. By § 6, further, the trilinear function
\[
 \chi(xyz)=\sum r_{\alpha\delta}(\chi)a_{\delta\beta\gamma}x_\alpha y_\beta z_\gamma
 =\sum r_{\kappa\lambda}(\chi)a_{\lambda\beta\gamma}x_\kappa y_\beta z_\gamma
\]
also depends only on $x_1,\ldots,x_m$, $y_1,\ldots,y_m$, $z_1,\ldots,z_m$. If therefore $\beta$ or $\gamma>m$, then $\sum_{\lambda=1}^m r_{\kappa\lambda}(\chi)a_{\lambda\beta\gamma}=0$ and hence $a_{\lambda\beta\gamma}=0$. Thus to the $k$ different prime factors of $\Theta$ there correspond $k$ groups homomorphic with $(\eps)$.

Or still more simply: all quantities $x$ whose coordinates satisfy the linear equations $\sum_\beta r_{\alpha\beta}(\chi)x_\beta=0$ form an invariant subgroup $(\mathfrak S)$ of $(\eps)$. For these equations express that $\chi(tx)=0$ for every quantity $t$. If $y$ is any quantity of $(\eps)$, and if one replaces $t$ by $yt$ or $ty$, one obtains $\chi(t(xy))=0$ or $\chi(t(yx))=0$. Thus if $x$ belongs to the complex $(\mathfrak S)$, then $xy$ and $yx$ also belong to it. Consequently $(\mathfrak S)$ is an invariant subgroup of $(\eps)$.

A group $(\eps)$ of order $r$ is called simple (\textsc{Molien}: original number system) if it has no invariant subgroup (except $(\eps)$), that is, if it has no homomorphic group whose order is $<n$.

By the preceding development, the determinant $\Theta$ of a simple group may contain no prime factor $\Phi$ whose linear rank $m<n$. By Theorem II of § 6, therefore, $\Theta$ may not contain two different prime factors, and must therefore be $\Theta=\Phi^s$, a power of a prime function, and hence $\sigma=s\chi$. Further, the rank $m$ of $\Phi$, and hence the rank of the matrix $R_\sigma=sR_\chi$, must be equal to $m=n=rs$; therefore $(\eps)$ must be a \textsc{Dedekind} group.

Conversely, if $|R_\sigma|$ is different from zero, and $\Theta=\Phi^s$ is a power of a prime function, then $(\eps)$ is a simple group, that is, it is not homomorphic with any group $(\eta)$ whose order $n'<n$ (\textsc{Mol}. Theorem 23). For the determinant of $(\eta)$ is a divisor of $\Theta$, hence equal to $\Theta'=\Phi^{s'}$, and therefore has the same rank as $\Phi$, namely $n$. But the rank $n$ of $\Theta'$ cannot be greater than the order $n'$ of $(\eta)$.

Since $S(x)$ and $T(y)$ are commuting matrices, the determinant (\emph{Gr.} § 10, (1.))
\[
 |uS(x)+vT(y)+wE|
\]
decomposes into a product of linear functions of $u,v,w$. If $uu_1+vv_1+w$ is one of them, then, by setting $v=0$ ($u=0$), one sees that $u_1$ ($v_1$) is a characteristic root of $\Phi(x)$ ($\Phi(y)$). If one regards $y$ as constant and $x$ as variable, then $\Phi(ux+we)$ is irreducible as a function of $w$, and so is $\Phi(ux+(vv_1+w)e)$. With this function the determinant has the factor $uu_1+vv_1+w$ in common, hence all $r$ factors $uu_\alpha+vv_1+w$. If one then regards $x$ as constant and $y$ as variable, one sees that the determinant contains each of the $r^2$ factors
\[
 uu_\alpha+vv_\beta+w\qquad(\alpha,\beta=1,2,\ldots,r)
\]
and contains each equally often. Hence
\[
 |uS(x)+vT(y)+wE|=\left(\prod(uu_\alpha+vv_\beta+w)\right)^c
\]
and therefore $s=rc$. In particular (\emph{Gr.} § 10, (8.))
\[
\fnum{1} |S(x)-T(x)+wE|=w^s\prod_{\beta>\alpha}\bigl(w^2-(u_\beta-u_\alpha)^2\bigr)^c.
\]
The elementary divisors of the characteristic determinants of the two commuting matrices $S(x)$ and $T(x)$ are all linear. By Theorem I, § 8, the elementary divisors of the determinant (1.) are therefore also linear, and in particular the $s$ of them which vanish for $w=0$. Consequently the matrix $S(x)-T(x)$ has rank $n-s$.

By Theorem III, § 8, however, it has rank $n-r$. Therefore (\textsc{Mol}. Theorem 29)
\[
\fnum{2} r=s,
\qquad m=n=r^2,
\]
hence $c=1$ and
\[
\fnum{3} |uS(x)+vT(y)+wE|=\prod_{\alpha,\beta}^{r}(uu_\alpha+vv_\beta+w).
\]
\emph{The determinant of a simple group is a power of a prime function whose exponent is equal to the degree of the prime function. The linear rank of this function, which agrees with the order of the group, is equal to the square of its degree.}

\section*{§ 11.}

Suppose that, for a definite quantity $h$ of the simple group $(\eps)$ of order $n=r^2$, the $r$ roots of the equation $\Phi(ue-h)=0$ are all distinct, and let $a$ be one definite one of these $r$ roots. Then the determinant $|S(ue-h)|$ has only linear elementary divisors, and consequently the matrix $S(ae-h)$ has rank $n-r$. Hence (\textsc{Mol}. § 8) the linear equations $(ae-h)t=0$ among the coordinates $t_1,\ldots,t_n$ of the unknown quantity $t$ have exactly $r$ independent solutions, $t=t^{(1)},t^{(2)},\ldots,t^{(r)}$.

If $x$ is a variable quantity, then $t^{(\kappa)}x$ is also a solution, and hence
\[
 t^{(\kappa)}x=\sum_\lambda x_{\kappa\lambda}t^{(\lambda)}.
\]
Since the $r$ solutions $t^{(\lambda)}$ are independent, the quantities $x_{\kappa\lambda}=f_{\kappa\lambda}(x)$ are linear functions of the coordinates $x_1,\ldots,x_n$. If $y$ and $z$ are two other quantities, and if $y_{\kappa\lambda}=f_{\kappa\lambda}(y)$, $z_{\kappa\lambda}=f_{\kappa\lambda}(z)$, then also
\[
 t^{(\kappa)}y=\sum_\lambda y_{\kappa\lambda}t^{(\lambda)},\qquad
 t^{(\kappa)}z=\sum_\lambda z_{\kappa\lambda}t^{(\lambda)},
\]
and therefore
\[
 t^{(\kappa)}yz=\sum_\mu y_{\kappa\mu}t^{(\mu)}z
 =\sum_{\lambda,\mu}y_{\kappa\mu}z_{\mu\lambda}t^{(\lambda)}.
\]
Thus if $x=yz$, then
\[
\fnum{1} x_{\kappa\lambda}=\sum_\mu y_{\kappa\mu}z_{\mu\lambda}.
\]
Among the $n=r^2$ linear functions $x_{\kappa\lambda}=f_{\kappa\lambda}(x)$ let $m$ be independent. If then $x=yz$, these $m$ combinations of $x_1,\ldots,x_n$ depend, by (1.), on the same $m$ combinations of $y_1,\ldots,y_n$ and on those of $z_1,\ldots,z_n$. Consequently $(\eps)$ has a homomorphic group of order $m$. But since $(\eps)$ is simple, $m=n=r^2$, and hence one can introduce $r^2$ new basic numbers $\eps_{\kappa\lambda}$ such that
\[
 \sum \eps_\alpha x_\alpha=\sum \eps_{\kappa\lambda}x_{\kappa\lambda}.
\]
For these, formulas (4.) § 1 pass into (1.) and the relations (2.) § 1 pass into
\[
\fnum{2} \eps_{\alpha\beta}\eps_{\delta\gamma}=\eps_{\alpha\gamma},
\qquad
\eps_{\alpha\delta}\eps_{\beta\gamma}=0\quad(\beta\ne\delta).
\]

To the $k$ different prime factors of $\Theta$ there correspond, as was shown at the beginning of § 10, $k$ groups homomorphic with $(\eps)$. The first has order $m$, its determinant is a divisor of $\Theta$, and it depends only on $x_1,\ldots,x_m$. Since $\Phi',\Phi'',\ldots$ are independent of $x_1,\ldots,x_m$, its determinant is a power of $\Phi$, and therefore, like $\Phi$, has linear rank $m$. Since its order is likewise equal to $m$, by § 10 it is a simple group. Consequently its determinant is equal to $\Phi^r$, and the rank of $\Phi$ is $m=r^2$.

\emph{I. The linear rank of every prime factor of the group determinant is equal to the square of its degree.}

I now denote the $k$ different functions $\Phi$ by $\Phi_1,\Phi_2,\ldots,\Phi_k$. In the group determinant $S$ all elements $s_{\alpha\beta}$ vanish which the first $m_1$ rows have in common with the last $n-m_1$ columns, likewise all elements which the following $m_2$ rows have in common with the first $m_1$ and the last $n-m_1-m_2$ columns, further those which the following $m_3$ rows have in common with the first $m_1+m_2$ and the last $n-m_1-m_2-m_3$ columns, and so on, whereas nothing is known about the last $n-m_1-\cdots-m_k$ rows. The same holds for the antistrophic matrix $T$, hence also for $uS(x)+vT(y)+wE$. If the first $m_1$ rows and columns of $S$ ($T$) form the matrix $S_1$ ($T_1$), the following $m_2$ rows and columns the matrix $S_2$ ($T_2$), $\ldots$, let
\[
 U=\begin{pmatrix}
 S_1&0&\cdots&0\\
 0&S_2&\cdots&0\\
 \cdot&\cdot&\cdots&\cdot\\
 0&0&\cdots&S_k
 \end{pmatrix},
 \qquad
 V=\begin{pmatrix}
 T_1&0&\cdots&0\\
 0&T_2&\cdots&0\\
 \cdot&\cdot&\cdots&\cdot\\
 0&0&\cdots&T_k
 \end{pmatrix}.
\]
Then
\[
 S=\begin{pmatrix}U&0\\ U_0&S_0\end{pmatrix},
 \qquad
 T=\begin{pmatrix}V&0\\ V_0&T_0\end{pmatrix}.
\]
From this follows the theorem of \textsc{Molien}:

\emph{II. Every group is homomorphic to a \textsc{Dedekind} group whose determinant contains every prime factor of its determinant $\Theta$, and whose order is equal to the linear rank of $\Theta$.}

This rank $m$ is equal to the rank of the quadratic form $\sigma(x^2)=\sum p_{\alpha\beta}x_\alpha x_\beta$, and to obtain that group one need only introduce the $m$ independent ones among the $n$ functions $\sum_\beta p_{\alpha\beta}x_\beta$ as new variables; that is, in formula (4.) § 9 one chooses the substitution $C$ so that it transforms the form $\sigma(x^2)$ into a function of $m$ variables, for instance into a sum of $m$ squares. I have not succeeded, however, in proving this theorem directly, that is, without using the prime functions $\Phi_\kappa$.

Further,
\[
 |S|=|S_1|\cdots|S_k||S_0|,
 \qquad
 |T|=|T_1|\cdots|T_k||T_0|
\]
and
\[
 |uS(x)+vT(y)+wE|=|uS_1+vT_1+wE_1|\cdots|uS_k+vT_k+wE_k|\,|uS_0+vT_0+wE_0|.
\]
Since $S_\kappa$ is the matrix of a simple group, by (3.) § 10
\[
 \Psi_\kappa=|uS_\kappa(x)+vT_\kappa(y)+wE_\kappa|
 =\prod(uu_\alpha+vv_\beta+w),
\]
where $u_1,\ldots,u_r$ are the characteristic roots of $\Phi_\kappa(x)$ and $v_1,\ldots,v_r$ those of $\Phi_\kappa(y)$. These results remain unchanged if the basic numbers of $(\eps)$ are chosen not in the way described above, but arbitrarily.

Also $\Psi_0=|uS_0(x)+vT_0(y)+wE_0|$ decomposes into linear factors $uu_0+vv_0+w$. Here $u_0$ is a characteristic root of a prime factor $\Phi_\kappa(x)$, and $v_0$ is such a root of a prime factor $\Phi_\lambda(y)$. From the irreducibility of $\Phi_\kappa(ue-x)$ and $\Phi_\lambda(ve-y)$ it follows that $\Psi_0$ also contains all $r_\kappa r_\lambda$ factors of the product $\prod(uu_\alpha+vv_\beta+w)$, where $u_\alpha$ runs through the $r_\kappa$ characteristic roots of $\Phi_\kappa(x)$ and $v_\beta$ through the $r_\lambda$ characteristic roots of $\Phi_\lambda(y)$. If $\kappa=\lambda$, the factor $\Psi_\kappa$ therefore occurs repeatedly. But, as the example of § 3 shows, $\kappa$ may also be different from $\lambda$.

If one chooses for $x$ an invariant quantity, then by (12.) § 4 the $r_\kappa$ quantities $u_\alpha$ are all equal to $\frac1{r_\kappa}\chi^{(\kappa)}(x)$. Formulas (2.) and (3.) § 8 show, however, that then to the root $v_\beta$ there is assigned the root $\frac1{r_\lambda}\chi^{(\lambda)}(x)$, and only this root. Hence $\kappa$ can be different from $\lambda$ only if, for an invariant quantity $x$,
\[
\fnum{1} \frac1{r_\kappa}\chi^{(\kappa)}(x)=\frac1{r_\lambda}\chi^{(\lambda)}(x)\qquad(S_x=T_x)
\]
holds.

Now let $w=0$, $u=v=1$, and
\[
\fnum{2} \Psi_{\kappa\lambda}(x,y)=\prod(u_\alpha+v_\beta),
\]
where $u_\alpha$ runs through the $r_\kappa$ characteristic roots of $\Phi_\kappa(x)$, and $v_\beta$ through the $r_\lambda$ roots of $\Phi_\lambda(y)$. Then $\Psi_{\kappa\lambda}$ is an irreducible function of $x_1,\ldots,x_n,y_1,\ldots,y_n$, of degree $r_\kappa$ with respect to the former and of degree $r_\lambda$ with respect to the latter, and
\[
\fnum{3} \Psi(x,y)=|S(x)+T(y)|=\prod\bigl(\Psi_{\kappa\lambda}(x,y)\bigr)^{c_{\kappa\lambda}},
\]
where the integer $c_{\kappa\kappa}>0$, but $c_{\kappa\lambda}\geq0$, and $c_{\kappa\lambda}>0$ can occur only under condition (1.). Since
\[
 \Psi_{\kappa\lambda}(x,0)=\Phi_\kappa(x)^{r_\lambda},
 \qquad
 \Psi_{\kappa\lambda}(0,y)=\Phi_\lambda(y)^{r_\kappa},
\]
it follows that for the exponents occurring in the decompositions
\[
\fnum{4} |S(x)|=\prod\Phi_\lambda(x)^{s_\lambda},
\qquad
|T(x)|=\prod\Phi_\lambda(x)^{t_\lambda},
\]
one has the expressions
\[
\fnum{5} s_\lambda=c_{\lambda1}r_1+c_{\lambda2}r_2+\cdots+c_{\lambda k}r_k,
\qquad
 t_\lambda=c_{1\lambda}r_1+c_{2\lambda}r_2+\cdots+c_{k\lambda}r_k,
\]
so that always
\[
\fnum{6} s_\lambda\geq r_\lambda,
\qquad t_\lambda\geq r_\lambda.
\]
The two functions $|S(x)|$ and $|T(x)|$ have the same linear rank
\[
\fnum{7} m=r_1^2+r_2^2+\cdots+r_k^2,
\]
this is equal to the rank of the matrix $R_\sigma$ and that of $R_\tau$, and equal to the sum of the ranks of the $k$ matrices $R(\chi^{(\kappa)})$ or of the $k$ prime functions $\Phi_\kappa(x)$. On the other hand, the order of the group $(\eps)$ is
\[
\fnum{8} n=r_1s_1+\cdots+r_ks_k=r_1t_1+\cdots+r_kt_k=\sum c_{\kappa\lambda}r_\kappa r_\lambda,
\]
where $r_\kappa r_\lambda$, if $\kappa$ is different from $\lambda$, has the coefficient $c_{\kappa\lambda}+c_{\lambda\kappa}$.

\section*{§ 13.}

I now resume the theory of \textsc{Dedekind} groups begun in §§ 7 and 8. For such a group $m=n$, hence by (6.), (7.) and (8.) § 12 $s_\kappa=t_\kappa=r_\kappa$.

\emph{I. The exponent of the power of a prime function occurring in the determinant of a \textsc{Dedekind} group is equal to the degree of that function.}

The group $(\eps)$ decomposes into $k$ simple groups of orders $r^2,r'^2,\ldots$; in each of them $S(x)$ decomposes into $r$ identical matrices of degree $r$, and the $r^2$ elements $x_{\kappa\lambda}$ of such a matrix are independent of one another and of the variables of the other simple groups. In my paper \emph{Über die Darstellung der endlichen Gruppen durch lineare Substitutionen} II., Sitzungsberichte 1899, § 5, I derived from this the theorem:

\emph{II. If $a$ and $b$ are two definite quantities of a \textsc{Dedekind} group, and if the elementary invariants of the two matrices $S(ue-a)$ and $S(ue-b)$ agree, then there exists in the group a quantity $c$ for which $|S(c)|$ is different from zero and which satisfies the condition $c^{-1}ac=b$.}

The determinant of a \textsc{Dedekind} group is
\[
\fnum{1} \Theta(x)=|S(x)|=|T(x)|=\prod\Phi^r,
\]
its trace is
\[
\fnum{2} \sigma=\tau=r\chi+r'\chi'+r''\chi''+\cdots=\sum r\chi,
\]
and therefore also
\[
\fnum{3} P=R_\sigma=rR_\chi+r'R_{\chi'}+r''R_{\chi''}+\cdots=\sum rR_\chi.
\]
The matrix $R_\chi$ has rank $r^2$, while the matrix $r'R_{\chi'}+r''R_{\chi''}+\cdots=P-rR_\chi$ has, by Theorem II of § 6, rank $r'^2+r''^2+\cdots=n-r^2$. Therefore the determinant $|uP-rR_\chi|$ vanishes for $u=0$ at least to order $n-r^2$, for $u=1$ at least to order $r^2$, and since it does not vanish for more than $r$ values, not only is
\[
\fnum{4} |uP-rR_\chi|=|P|u^{n-r^2}(u-1)^{r^2},
\]
but the elementary divisors of this determinant are also all linear (since otherwise it would vanish for $u=0$ or $1$ to higher order). Consequently the matrix $H=rP^{-1}R_\chi$ satisfies the equation $H^2=H$, and therefore
\[
\fnum{5} R_\chi P^{-1}R_\chi=\frac1r R_\chi.
\]
If one therefore appends on the left in (3.) the factor $R_\chi P^{-1}$, one obtains
\[
 0=r'(R_\chi P^{-1})R_{\chi'}+r''(R_\chi P^{-1})R_{\chi''}+\cdots,
\]
and hence by (4.) § 6
\[
\fnum{6} R_\chi P^{-1}R_{\chi'}=0.
\]
If between $\xi_1,\ldots,\xi_n$ and $x_1,\ldots,x_n$ there are the equations
\[
\fnum{7} \xi_\alpha=\sum_\beta p_{\alpha\beta}x_\beta,
\]
then I call the parameter $\xi$ and the quantity $x$ conjugate. Then
\[
\fnum{8} R_\xi=S_x'P=PT_x,
\qquad
R_\xi'=PS_x=T_x'P.
\]
Thus if $R_\xi=R_\xi'$, then $S_x=T_x$; to a symmetric parameter there is conjugate an invariant quantity, and conversely.

The invariant quantity $x=h$ conjugate to the characteristic parameter $\xi=\chi$ I call a characteristic quantity for the prime function $\Phi$. Then
\[
\fnum{9} \chi_\alpha=\sum_\beta p_{\alpha\beta}h_\beta
\]
and
\[
\fnum{10} R_\chi=PS_h=PT_h.
\]
From (5.) and (10.) it follows that $\frac1r S(h)=S(h)^2=S(h^2)$ and $0=S(h)S(h')=S(hh')$, hence (\emph{Gr.} § 5 (9.))
\[
\fnum{11} h^2=\frac1r h,
\qquad hh'=0,
\]
and from (3.)
\[
\fnum{12} e=rh+r'h'+r''h''+\cdots=\sum rh,
\]
so that the parameter $\sigma$ and the principal unit $e$ are conjugate. Further,
\[
 rR_\chi S_h=R_\chi,
 \qquad
 r\sum_\lambda r_{\alpha\lambda}(\chi)s_{\lambda\beta}(h)=r_{\alpha\beta}(\chi).
\]
Multiplying by $e_\alpha e_\beta$ and summing over $\alpha$ and $\beta$, one obtains (cf. \textsc{Ded}. (13.))
\[
\fnum{13} \sum_\alpha \chi_\alpha h_\alpha=1,
\qquad
\sum_\alpha \chi'_{\alpha}h_\alpha=0,
\]
or, if one puts
\[
 h(\xi)=h_1\xi_1+\cdots+h_n\xi_n,
\]
then
\[
\fnum{14} \chi(h)=h(\chi)=1,
\qquad
\chi(h')=h(\chi')=0,
\]
or finally, if $P^{-1}=Q=(q_{\alpha\beta})$,
\[
\fnum{15} \sum p_{\alpha\beta}h_\alpha h_\beta=\sum q_{\alpha\beta}\chi_\alpha\chi_\beta=1,
\qquad
\sum p_{\alpha\beta}h_\alpha h'_{\beta}=\sum q_{\alpha\beta}\chi_\alpha\chi'_{\beta}=0.
\]

\section*{§ 14.}

Let $u_1,\ldots,u_r$ be the characteristic roots of $\Phi(x)$, and $v_1,\ldots,v_r$ those of $\Phi(y)$. If $g(u,v)$ is an integral function of $u$ and $v$, then by § 12 the $n$ characteristic roots of the matrix $g(S_x,T_y)$ are the $r^2$ quantities $g(u_\alpha,v_\beta)$ and the $r'^2,r''^2,\ldots$ quantities obtained from $\Phi',\Phi''$ in the same way.

If $x$ is an invariant quantity, then by the theorem of § 4
\[
 u_1=\cdots=u_r=\frac1r\chi(x)=c.
\]
By (13.) § 13 the characteristic determinant of $S_h=T_h$ is therefore equal to $(u-\frac1r)^{r^2}u^{n-r^2}$, and the characteristic roots of $(S(x)-cE)T(h)=S((x-ce)h)$ are all zero. Thus a power of this matrix vanishes. But since it is commutable with every group matrix $S(y)$, by Theorem III, § 8 the elementary divisors of its characteristic determinant are all linear. Therefore the first power vanishes, and hence
\[
\fnum{1} xh=ch=h\frac1r\chi(x).
\]
Multiplying (12.) § 13 by $x$, one therefore obtains
\[
\fnum{2} x=rch+r'c'h'+r''c''h''+\cdots
=h\chi(x)+h'\chi'(x)+h''\chi''(x)+\cdots,
\]
where the coefficients $rc=\chi(x)$, $r'c'=\chi'(x)$, $\ldots$ are ordinary quantities.

The $k$ quantities $h,h',h'',\ldots$, which by (5.) § 6 or by (11.) § 13 are linearly independent of one another, therefore form a complete system of solutions of the linear equations $S(x)=T(x)$. Among the $n^2$ equations $s_{\alpha\beta}(x)=t_{\alpha\beta}(x)$ there are consequently $n-k$ independent ones. If, in addition, the quantity $x$ satisfies the equations $\chi(x)=\chi'(x)=\cdots=0$, then by (2.) $x=0$.

This can also be seen as follows. If $xy=yx$, then by (11.) § 4, $r\chi(xy)=\chi(x)\chi(y)$. Thus if $\chi(x)=0$, then also $\chi(xy)=0$, or $\sum_\beta r_{\alpha\beta}(\chi)x_\beta=0$. If this also holds for $\chi',\chi'',\ldots$, then $r^2+r'^2+\cdots=n$ independent linear equations result, and hence $x_\beta=0$.

\emph{I. Among the $n^2$ linear functions $s_{\alpha\beta}(x)-t_{\alpha\beta}(x)$ of the $n$ variables $x_1,\ldots,x_n$, $n-k$ are independent. Together with the $k$ functions $\chi(x),\chi'(x),\ldots$ they form $n$ independent functions. The $k$ characteristic quantities $h,h',\ldots$ form a complete system of independent solutions of the equations $s_{\alpha\beta}(x)=t_{\alpha\beta}(x)$.}

\emph{II. Among the linear functions $r_{\alpha\beta}(\xi)-r_{\beta\alpha}(\xi)$ of the $n$ variables $\xi_1,\ldots,\xi_n$, $n-k$ are independent. Together with the $k$ functions $h(\xi),h'(\xi),\ldots$ they form $n$ independent functions. The $k$ characteristic parameters $\chi,\chi',\ldots$ form a complete system of independent solutions of the equations $r_{\alpha\beta}(\xi)=r_{\beta\alpha}(\xi)$.}

A special case of these theorems is the elegant criterion which Mr. \textsc{Molien} gives in Theorems 9 and 10, § 3 for the simplicity of a group. Every simple group is by § 10 a \textsc{Dedekind} group, and in order that such a group be simple it is necessary and sufficient that $k=1$.

\emph{III. In order that a group be simple, it is necessary and sufficient that the determinant of degree $n$, $\left|\sum_\kappa \sigma_\kappa a_{\kappa\alpha\beta}\right|$, be different from zero, and that $\xi_\kappa=\sigma_\kappa$ be the only solution of the linear equations}
\[
 \sum_\kappa(a_{\kappa\alpha\beta}-a_{\kappa\beta\alpha})\xi_\kappa=0
\]
\emph{.}

Or also:

\emph{IV. In order that a group be simple, it is necessary and sufficient that the determinant of degree $n$, $\left|\sum_\kappa \sigma_\kappa a_{\kappa\alpha\beta}\right|$, be different from zero, and that $x_\lambda=e_\lambda$ be the only solution of the linear equations}
\[
 \sum_\lambda(a_{\alpha\lambda\beta}-a_{\alpha\beta\lambda})x_\lambda=0
\]
\emph{.}

That each theorem appears here in two different forms has its reason in the fact that, for the group itself, two conjugate forms can be placed side by side. If to the parameters $\xi,\eta,\zeta$ there are conjugate the quantities $x,y,z$, let
\[
 F(\xi,y,z)=G(x,\eta,\zeta)=\sum b_{\alpha\beta\gamma}x_\alpha\eta_\beta\zeta_\gamma.
\]
Then
\[
\fnum{3} a_{\alpha\beta\gamma}=\sum_\lambda b_{\beta\lambda\alpha}p_{\lambda\gamma}
=\sum_\lambda b_{\gamma\alpha\lambda}p_{\lambda\beta}.
\]
By this transformation the group $(\eps)$ thus passes into an equivalent group $(\bar\eps)$, the conjugate group. If one uses for it the notations (4.) § 9, then, if in those formulas $C=P$ is put, one obtains
\[
\fnum{4} R_x=\bar T_xP,
\qquad R'_x=\bar S_xP,
\qquad S_x=\bar R'_xP,
\qquad T_x=\bar R_xP.
\]
Further, $\bar\sigma_\alpha=e_\alpha$. Just as $\bar a_{\alpha\beta\gamma}=b_{\alpha\beta\gamma}$ has been put, let $\bar p_{\alpha\beta}=q_{\alpha\beta}$ be put. Then $q_{\alpha\beta}=\sum_\kappa e_\kappa b_{\kappa\alpha\beta}$ and hence
\[
 \sum_\lambda p_{\alpha\lambda}q_{\lambda\beta}
 =\sum_{\kappa,\lambda}p_{\alpha\lambda}e_\kappa b_{\kappa\lambda\beta}
 =\sum_\kappa e_\kappa a_{\beta\kappa\alpha}=e_{\alpha\beta},
\]
so that $Q=P^{-1}$. Consequently the conjugate group of $(\bar\eps)$ is again the original group,
\[
\fnum{5} b_{\alpha\beta\gamma}=\sum_\lambda a_{\beta\lambda\alpha}q_{\lambda\gamma}
=\sum_\lambda a_{\gamma\alpha\lambda}q_{\lambda\beta}.
\]

\section*{§ 15.}

In § 6, (5.) it was shown that the $k$ functions $\chi(x),\chi'(x),\ldots$ are linearly independent. By Theorem I, § 14 this remains the case even when $x$ is not allowed to vary without restriction, but runs only through the invariant quantities of $(\eps)$. These form a subgroup $(\eta)$ of $(\eps)$, whose order is $k$ (in general not an invariant subgroup). As basic numbers $\eta_1,\eta_2,\ldots,\eta_k$ of $(\eta)$ one can take the $k$ quantities $rh,r'h',\ldots$. Then
\[
\fnum{1} \eta_\kappa^2=\eta_\kappa,
\qquad \eta_\kappa\eta_\lambda=0.
\]
Consequently $(\eta)$ is a commutative \textsc{Dedekind} group, whose determinant, if the quantities $\eta_1,\ldots,\eta_k$ are expressed by $\eps_1,\ldots,\eps_n$, is equal to $\prod \frac1r\chi(x)$. In fact these $k$ linear factors, in which the variability of the coordinates is restricted by the conditions $s_{\alpha\beta}(x)=t_{\alpha\beta}(x)$, are linearly independent.

\emph{I. If $k$ is the number of distinct prime factors of the determinant of a \textsc{Dedekind} group, then its invariant quantities form a commutative group of order $k$, which is also a \textsc{Dedekind} group. Its determinant is $\prod\frac1r\chi(x)$, while the determinant of the given group for an invariant variable $x$ is equal to $\prod\left(\frac1r\chi(x)\right)^{r^2}$.}

If the degrees $r=r'=\cdots=1$, then $k=n$, and hence by Theorem I, § 14 identically $s_{\alpha\beta}(x)=t_{\alpha\beta}(x)$, $a_{\alpha\beta\gamma}=a_{\alpha\gamma\beta}$.

\emph{II. The determinant of a \textsc{Dedekind} group always decomposes, but also only then decomposes, into purely linear factors when the group is commutative.}

In order that the determinant of an arbitrary group decompose into purely linear factors, it is necessary and sufficient, by Theorem II, § 12, that the \textsc{Dedekind} group homomorphic with it and of order $m=\sum r^2$ have this property, and therefore be a commutative group. This group is obtained by introducing the $m$ independent ones among the functions $\sum_\beta p_{\alpha\beta}t_\beta$ as coordinates of a quantity $t$. Thus if one puts $t=yz$, these bilinear functions of $y_1,\ldots,y_n,z_1,\ldots,z_n$ must remain unchanged under interchange of $y$ and $z$. If $x$ is a third variable, then the function
\[
 \sum_{\alpha,\beta}p_{\alpha\beta}x_\alpha t_\beta=\sigma(xt)=\sigma(xyz)
\]
has the same property. This function, by § 6, always remains unchanged under a cyclic permutation of $x,y,z$, and in the case under consideration therefore under every permutation. It is the trace of the matrix $S(xyz)=S(x)S(y)S(z)$, and is therefore equal to
\[
\fnum{2} \sigma(xyz)=\sum_{\kappa,\lambda,\mu}s_{\lambda\mu}(x)s_{\mu\kappa}(y)s_{\kappa\lambda}(z)
=\sum a_{\lambda\alpha\mu}a_{\mu\beta\kappa}a_{\kappa\gamma\lambda}x_\alpha y_\beta z_\gamma.
\]
If one uses the antistrophic group, then $\tau(xyz)$ takes the place of $\sigma(xyz)$. Thus the following theorem results (cf. \textsc{Cartan}, \emph{I. Thèse, Sur la structure des groupes de transformations finis et continus}, Paris 1894, p. 48, (5.)):

\emph{III. In order that the determinant of a group decompose into purely linear factors, it is necessary and sufficient that the trilinear function $\sigma(xyz)$ (or $\tau(xyz)$) remain unchanged under interchange of $y$ and $z$, and hence be a symmetric function of the three rows of variables; that is, that the expressions}
\[
\fnum{3} \sum_{\kappa,\lambda,\mu}a_{\lambda\alpha\mu}a_{\mu\beta\kappa}a_{\kappa\gamma\lambda}
\quad\text{or}\quad
\sum_{\kappa,\lambda,\mu}a_{\lambda\mu\alpha}a_{\mu\kappa\beta}a_{\kappa\lambda\gamma}
\]
\emph{also remain unchanged under a transposition of $\alpha,\beta,\gamma$.}

Since $\sigma(xyz)$ is also the trace of $S(x)S(yz)$, in place of the sums (3.) one may also take the expressions
\[
\fnum{4} \sum_{\kappa,\lambda,\mu}a_{\kappa\lambda\mu}a_{\mu\alpha\kappa}a_{\lambda\beta\gamma}
\quad\text{or}\quad
\sum_{\kappa,\lambda,\mu}a_{\kappa\lambda\mu}a_{\lambda\kappa\alpha}a_{\mu\beta\gamma}.
\]
Finally, by means of formulas (1.) and (2.) § 7, one obtains the expressions
\[
\fnum{5}\begin{gathered}
 \sum_{\kappa,\lambda,\mu}a_{\kappa\lambda\kappa}a_{\lambda\alpha\mu}a_{\mu\beta\gamma}
 \quad\text{or}\quad
 \sum_{\kappa,\lambda,\mu}a_{\kappa\kappa\lambda}a_{\lambda\alpha\mu}a_{\mu\beta\gamma},\\
 \sum_{\kappa,\lambda,\mu}a_{\kappa\lambda\kappa}a_{\lambda\mu\alpha}a_{\mu\beta\gamma}
 \quad\text{or}\quad
 \sum_{\kappa,\lambda,\mu}a_{\kappa\kappa\lambda}a_{\lambda\mu\alpha}a_{\mu\beta\gamma}.
\end{gathered}
\]

With the aid of the preceding theorems, the decomposition of the determinant of a \textsc{Dedekind} group into its prime factors can be carried out. The linear equations $s_{\alpha\beta}(x)-t_{\alpha\beta}(x)=0$ have $k$ independent solutions. One now sets up the quadratic equations among the unknowns $\chi_1,\chi_2,\ldots,\chi_n$ which are obtained from
\[
\fnum{6} \chi(x)\chi(y)=\chi(e)\chi(xy)
\]
by substituting for $x$ successively those $k$ solutions, and for $y$ the $n$ basic numbers. In combination with the linear equations $r_{\alpha\beta}(\chi)=r_{\beta\alpha}(\chi)$ they provide, for the ratios $\chi_1,\chi_2,\ldots,\chi_n$, $k$ different systems of values. The constant factor is chosen each time so that $\sum q_{\alpha\beta}\chi_\alpha\chi_\beta=1$ and $\sum\chi_\alpha e_\alpha$ becomes positive. It is then equal to a positive integer $r$, the degree of the prime function $\Phi$ determined by the characteristic parameter $\chi$.

Or one denotes $k$ independent solutions of the equations $s_{\alpha\beta}(x)=t_{\alpha\beta}(x)$ by $g,g',\ldots$, or also by $\eta_1,\eta_2,\ldots,\eta_k$. These form the basis of a commutative \textsc{Dedekind} group of order $k$, so that $\eta_\beta\eta_\gamma=\sum_\alpha c_{\alpha\beta\gamma}\eta_\alpha$, where $c_{\alpha\beta\gamma}=c_{\alpha\gamma\beta}$. Its determinant is a product of $k$ independent linear functions and, when the basis $\eps_1,\eps_2,\ldots,\eps_n$ is reintroduced, passes into $\prod\psi(x)$, where $\psi(x)=\sum\psi_\alpha x_\alpha$. Then, since $\psi_1,\psi_2,\ldots,\psi_n$ contain an arbitrary constant factor, one determines the ratios of the unknowns $\chi_1,\chi_2,\ldots,\chi_n$ from the linear equations
\[
 r_{\alpha\beta}(\chi)-r_{\beta\alpha}(\chi)=0,
 \qquad g(\chi)=g(\psi),
 \qquad g'(\chi)=g'(\psi),\ldots,
\]
and from these the actual values as above.

\section*{§ 16.}

If $E_1,E_2,\ldots,E_n$ are $n$ matrices of degree $m$ which satisfy the conditions
\[
\fnum{1} E_\beta E_\gamma=\sum_\alpha a_{\alpha\beta\gamma}E_\alpha
\]
then they form a representation of the group $(\eps)$. If these $n$ matrices are not linearly independent (compare, however, \textsc{Mol}. p. 126, condition 1), then they represent a group homomorphic with $(\eps)$.

If then
\[
\fnum{2} (x_{\kappa\lambda})=\sum x_\alpha E_\alpha,
\]
then, if $x=yz$,
\[
 (x_{\kappa\lambda})=(y_{\kappa\lambda})(z_{\kappa\lambda}).
\]
Therefore I call the matrix (2.) a matrix belonging to the group $(\eps)$. If its determinant is different from zero, then by (1.) § 4 it is
\[
\fnum{3} |x_{\kappa\lambda}|=\prod\Phi(x)^s
\]
a product of prime factors of the group determinant $\Theta$.

If $C$ is a constant matrix of degree $m$ whose determinant is different from zero, then the $n$ matrices $C^{-1}E_1C,C^{-1}E_2C,\ldots,C^{-1}E_nC$ also form a representation of $(\eps)$, which is called equivalent to the first. If $C$ cannot be chosen so that these $n$ matrices assume the form
\[
 \begin{pmatrix}E_1'&0\\ E_1^{(0)}&E_1''\end{pmatrix},
 \quad
 \begin{pmatrix}E_2'&0\\ E_2^{(0)}&E_2''\end{pmatrix},
 \quad \ldots \quad
 \begin{pmatrix}E_n'&0\\ E_n^{(0)}&E_n''\end{pmatrix},
\]
then I call the representation primitive or irreducible; if this is possible, I call it imprimitive or reducible; and if also $E_1^{(0)},\ldots,E_n^{(0)}=0$, decomposed or decomposable.

The methods which I developed in my paper \emph{Über die Darstellung der endlichen Gruppen durch lineare Substitutionen}, II., Sitzungsberichte 1899, can be transferred directly (with the modification indicated at the end of § 3) to arbitrary \textsc{Dedekind} groups.

The $r^2$ linear functions of $x_1,\ldots,x_n$ which I denoted in § 11 by $x_{\kappa\lambda}$ form a matrix of degree $r$ belonging to the group $(\eps)$, whose determinant $|x_{\kappa\lambda}|=\Phi(x)$. I denote the corresponding representation of $(\eps)$ by $[\Phi]$. The representations so obtained, corresponding to the $k$ different prime factors of the group determinant $\Theta$, are all the primitive representations of the group $(\eps)$.

A representation equivalent to (2.) then decomposes into $s$ primitive representations $[\Phi]$, $s'$ representations $[\Phi']$, and so on. The coefficient of $u^{m-1}$ in the determinant $|x_{\kappa\lambda}+ue_{\kappa\lambda}|$,
\[
 \varphi(x)=\sum s\chi(x),
\]
is called the trace of the representation (2.). In order that two representations of a \textsc{Dedekind} group be equivalent, it is necessary and sufficient that their traces agree.
% ===== END WORK 070 =====

\clearpage

% ===== BEGIN WORK 071: 071_Theorie_der_hyperkomplexen_Gr_en_II =====
\frobtitle{Theory of Hypercomplex Quantities II}

In my paper \emph{Theory of Hypercomplex Quantities}, cited in what follows as \emph{H.}, I treat, in connection with my investigations on the determinants of finite groups, chiefly the properties of the \textsc{Dedekind} groups. As Mr. \textsc{Molien} has shown (\emph{H.} § 11, II), every group $(\eps)$ with principal unit is homomorphic to a \textsc{Dedekind} group $(\grS)$ whose determinant is divisible by every prime factor of the determinant of the whole group. If $(\eta)$ is the largest invariant subgroup of $(\eps)$ which consists entirely of roots of zero -- I call it the \emph{radical} of $(\eps)$ -- then $(\eps)$ mod. $(\eta)=(\grS)$; that is, $(\grS)$ is the group $(\eps)$ when in it any two quantities are regarded as equal whose difference is contained in $(\eta)$.

If, among the $n$ independent fundamental units $\eps_1,\ldots,\eps_n$ of a group $(\eps)$, the first $m$ form the basis of a group $(\grS)$ and the last $n-m$ form that of a group $(\eta)$, then I call $(\eps)=(\grS)+(\eta)$ the \emph{sum} of these two subgroups. If one of them, $(\eta)$, is an invariant subgroup of $(\eps)$, then the other, $(\grS)$, is a group homomorphic to $(\eps)$, because
\[
 \eps_1x_1+\cdots+\eps_mx_m+\eps_{m+1}x_{m+1}+\cdots+\eps_nx_n
 \equiv \eps_1x_1+\cdots+\eps_mx_m \pmod{\eta}
\]
holds. If both subgroups are invariant, then the product of every quantity of $(\grS)$ and every quantity of $(\eta)$ vanishes, and $(\eps)$ \emph{splits} into the two groups $(\grS)$ and $(\eta)$. (\emph{H.} § 9.)

A group $(\eta)$ contained in $(\eps)$ is called an \emph{invariant subgroup} of $(\eps)$ if $xy$ and $yx$ are quantities of $(\eta)$ whenever $y$ is any quantity of $(\eta)$ and $x$ is any quantity of $(\eps)$. If $n-m$ is the order of $(\eta)$, then $(\eps)$ (mod. $\eta$) is homomorphic to a group $(\grS)$ of order $m$. There need not, however, always exist a subgroup $(\grS)$ of $(\eps)$ of such a kind that $(\eps)=(\grS)+(\eta)$.

Mr. \textsc{Cartan}, in his paper \emph{Sur les groupes bilinéaires et les systèmes de nombres complexes}, Ann. de Toulouse, tome XII, 1898, cited in what follows as \emph{C.}, has given a more precise formulation of \textsc{Molien}'s theorem, which I shall derive from the results of my first paper in order to complete my presentation:

\begin{quote}\itshape
Every group with principal unit is the sum of its radical and of a {\normalfont\textsc{Dedekind}} group whose determinant is divisible by every prime factor of the determinant of the whole group.
\end{quote}

But whereas the invariant subgroup $(\eta)$, which I call the radical of $(\eps)$, is completely determined, the subgroup $(\grS)$ can usually be chosen in infinitely many ways.

If a power of a number vanishes, I call it a \emph{root of zero} (pseudo-nul in \textsc{Cartan}, no. 21; nilpotent in \textsc{Peirce}). A group which consists entirely of roots of zero I call a \emph{root group}.

If a group $(\eps)$ is the sum of a \textsc{Dedekind} group $(\grS)$ of order $m$ and a root group $(\eta)$ of order $n-m$, and if this root group is an invariant subgroup of $(\eps)$, then $(\eta)$ is the radical of $(\eps)$. For since $(\eps)$ and $(\grS)$ are homomorphic, the determinant of the group $(\eps)$ is divisible by that of the group $(\grS)$ (\emph{H.} § 9). The linear rank of the determinant of a \textsc{Dedekind} group $(\grS)$ is equal to its order $m$. Hence the linear rank of the determinant of $(\eps)$, $\bar m$, is $\geq m$. The radical $(\bar\eta)$ of $(\eps)$ has order $n-\bar m$. Since $(\eta)$ is contained in $(\bar\eta)$, one has $n-m\leq n-\bar m$. Thus $m=\bar m$ and $(\eta)=(\bar\eta)$.

Before I pass to the proof of \textsc{Cartan}'s theorem (§ 5), I shall briefly derive the principal properties of root groups.

\section*{§ 2.}

If neither of the two determinants $|S(x)|$ and $|T(x)|$ vanishes identically, then the group $(\eps)$ possesses a \emph{principal unit} $e$, and for every quantity $x$ in $(\eps)$ one has
\begin{equation}\tag{1.}
 ex=xe=x.
\end{equation}

Conversely, Mr. \textsc{Cartan} (\emph{C.} 15) has shown: if for every quantity $x$ one has $ex=x$, then $S(e)=E$, hence $|S(e)|=1$, and if $xe=x$, then $T(e)=E$, hence $|T(e)|=1$. For, if one seeks to determine $x$ so that $ex=0$, one obtains for the coordinates $x_1,\ldots,x_n$ $n$ homogeneous linear equations with the matrix $S(e)$. But then $x=ex=0$, so only the value system $x_1=\cdots=x_n=0$ satisfies these equations, and consequently their determinant $|S(e)|$ is different from zero. Since now $S(e)^2=S(e^2)=S(e)$, it follows that $S(e)=E$.

If $|S(x)|=\Theta$ and $|T(x)|=\Theta'$ are the two antistrophic determinants of the group $\eps$, then one can determine two quantities $y$ and $z$, whose coordinates are integral functions of $x_1,\ldots,x_n$, so that $xy=\Theta e$ and $zx=\Theta'e$. Then $T(y)T(x)=T(xy)=\Theta T(e)$, and hence $|T(y)|\Theta'=\Theta^n$; likewise $|S(z)|\Theta=\Theta'{}^n$. Consequently every prime factor of one of the two determinants $\Theta$ and $\Theta'$ is also contained in the other (\emph{H.} § 3).

Since I also frequently have to consider groups without principal unit, I first make the following remarks about them. Suppose that between the $(n-1)^3$ quantities $a_{\alpha\beta\gamma}$ $(\alpha,\beta,\gamma=1,2,\ldots,n-1)$ there exist the relations
\begin{equation}\tag{2.}
 \sum_{\kappa} a_{\kappa\alpha\beta}a_{\gamma\kappa\delta}
 =\sum_{\kappa} a_{\kappa\beta\delta}a_{\gamma\alpha\kappa}.
\end{equation}
Then there is always a group $(\bar\eps)$ with $n-1$ linearly independent fundamental units $\eps_1,\ldots,\eps_{n-1}$, between which the relations
\begin{equation}\tag{3.}
 \eps_{\beta}\eps_{\gamma}=\sum_{\alpha}a_{\alpha\beta\gamma}\eps_{\alpha}
\end{equation}
hold. If $(\bar\eps)$ possesses a principal unit, I proved this in \emph{H.} § 2 (compare also \textsc{Study} in the encyclopedia). In the other case set
\begin{equation}\tag{4.}
 a_{\alpha0\alpha}=a_{\alpha\alpha0}=1\qquad(\alpha=0,1,\ldots,n-1),
\end{equation}
and otherwise set $a_{\alpha\beta\gamma}=0$ if one of the indices is $0$. Then equations (2.) also hold for the $n^3$ quantities $a_{\alpha\beta\gamma}$ $(\alpha,\beta,\gamma=0,1,\ldots,n-1)$. If, for a fixed $\lambda$, one denotes the matrix $(a_{\alpha\lambda\beta})$ by $E_{\lambda}$, then $E_0=E$ and
\begin{equation}\tag{5.}
 E_{\beta}E_{\gamma}=\sum_{\alpha}a_{\alpha\beta\gamma}E_{\alpha}.
\end{equation}
If further $\sum c_{\lambda}E_{\lambda}=0$, then $\sum c_{\lambda}a_{\alpha\lambda0}=0$, and so $c_{\alpha}=0$. Therefore $E_1,\ldots,E_{n-1}$ form a representation of the group $(\bar\eps)$ of order $n-1$ by matrices of degree $n$.

This group $(\bar\eps)$ is an invariant subgroup of a group $(\eps)$ with basis $\eps_0,\eps_1,\ldots,\eps_{n-1}$, for which $\eps_0$ is the principal unit. In its two antistrophic matrices $S(x)$ and $T(x)$ one has $s_{00}(x)=t_{00}(x)=x_0$, but $s_{0\beta}(x)=t_{0\beta}(x)=0$ (and $s_{\alpha0}(x)=t_{\alpha0}(x)=x_{\alpha}$). If one removes the first row and column, and in the resulting matrices of degree $(n-1)$ sets $x_0=0$, then one obtains the two antistrophic matrices $\bar S(x)$ and $\bar T(x)$ of $(\bar\eps)$. Since at least one of their determinants vanishes identically, one of the two determinants $|S(x)|$ or $|T(x)|$ is divisible by $x_0^2$.

Conversely, let $(\eps)$ be a group of order $n$ with principal unit $e$, and containing an invariant subgroup $(\bar\eps)$ of order $n-1$. Let $\eps_1,\ldots,\eps_{n-1}$ be a basis of $(\bar\eps)$, and let $\eps_0$ be a quantity of $(\eps)$ not contained in $(\bar\eps)$. Then $|S(x)|$ and $|T(x)|$ are both divisible by $x_0$. According as neither of these two determinants, or at least one of them, is divisible by $x_0^2$, the group $(\bar\eps)$ possesses a principal unit $(\bar e)$ or does not. If $e$ were contained in $(\bar\eps)$, then, by the definition of an invariant subgroup, one would also have $e\eps_0=\eps_0$. Therefore one can always choose the principal unit $e$ for $\eps_0$.

\section*{§ 3.}

If $A$ and $B$ are two matrices of degree $n$, then the characteristic functions of $AB$ and $BA$ are equal. For if $|B|$ is different from zero, then $B^{-1}(BA)B=AB$ is similar to the matrix $BA$. Hence, if the elements of $A$ and $B$ are variable quantities, the equation $|uE-AB|=|uE-BA|$ holds for all values of these variables for which $|B|$ is not zero, and therefore it holds identically.

If $y$ and $z$ are two quantities of the group $(\eps)$, then the coordinates of their product $x=yz$ are
\begin{equation}\tag{1.}
 x_{\alpha}=\sum_{\beta,\gamma}a_{\alpha\beta\gamma}y_{\beta}z_{\gamma}
 =\sum_{\gamma}s_{\alpha\gamma}(y)z_{\gamma}.
\end{equation}
If now $S(y)=0$, so that $y_1,\ldots,y_n$ satisfy the linear equations $s_{\alpha\gamma}(y)=0$, then $yz=0$ for every quantity $z$ of $(\eps)$. Thus, if $(\eps)$ has a principal unit, then $y=0$; in every case, however, $y^2=0$.

More generally, if $x$ is a quantity of $(\eps)$ for which $S(x)^m=S(x^m)=0$, then $x^mz=0$ and $x^{m+1}=0$. Thus, in order that $x$ be a root of zero, it is necessary and sufficient that the matrix $S(x)$ (or $T(x)$) be a root of zero; equivalently, that all characteristic roots of $S(x)$ vanish. If $zy$ is a root of zero, then $yz$ is also one. For all characteristic roots of the matrix $S(zy)=S(z)S(y)$ are zero, and consequently so are those of the matrix $S(y)S(z)=S(yz)$.

In a group $(\eps)$ a quantity $y$ is called a \emph{root quantity} if $yz$ is always a root of zero whenever $z$ is any quantity of $(\eps)$. If $x$ and $z$ are any two quantities of $(\eps)$, then $y(zx)=(yz)x$ is also a root of zero, and hence also $x(yz)$. Thus, if $y$ is a root quantity of $(\eps)$, then $xy$, $yz$, and $xyz$ are also root quantities.

Then the characteristic roots of the matrix $S(yz)$ vanish, and consequently so does their sum $\sigma(yz)$. Conversely, suppose that $\sigma(yz)=0$ for every quantity $z$ of $(\eps)$. Replacing $z$ by $z(yz)^{\kappa-1}$, one obtains $\sigma((yz)^\kappa)=0$. Therefore (\emph{H.} § 4, (5.)) the sum of the $\kappa^{\ten}$ powers of the characteristic roots of the matrix $S(yz)$ vanishes, so that these roots themselves also vanish, and hence $yz$ is a root of zero. The coordinates $y_1,\ldots,y_n$ of the root quantities $y$ are therefore found by solving the homogeneous linear equations obtained by setting equal to zero the derivatives of the bilinear form $\sigma(yz)$ with respect to $z_1,\ldots,z_n$ (or those of the quadratic form $\sigma(y^2)$ with respect to $y_1,\ldots,y_n$). Hence the root quantities of $(\eps)$ reproduce themselves under addition and under multiplication by ordinary quantities. Moreover $yz$ and $zy$ are root quantities whenever $y$ is one. Consequently the root quantities of $(\eps)$ form an invariant subgroup, which I call the radical of $(\eps)$.

\begin{quote}\itshape
I. The radical of a group $(\eps)$ is formed by all quantities $y$ whose coordinates, for every quantity $z$, satisfy the linear equations $\sigma(yz)=0$ (or $\tau(yz)=0$).
\end{quote}

That also for groups without principal unit (\emph{H.} § 4, (9.))
\begin{equation}\tag{2.}
 \sigma(xy)=\sigma(yx)
\end{equation}
holds follows from (2.), § 2, by setting $\gamma=\delta$ and summing over $\gamma$ (\textsc{Mol}. § 3, (4.)).

A root group can also be defined as a group which is equal to its radical. If a nonzero quantity $x$ satisfying the equation $x^2=x$ is called a \emph{unit}, then \textsc{Peirce}, \emph{Linear Associative Algebra}, no. 51 (American Journ. tome IV), has shown:

\begin{quote}\itshape
II. In order that a group be a root group, it is necessary and sufficient that it contain no unit.
\end{quote}

For, if $m$ is the smallest exponent for which $x^m=0$, then one cannot have $x=x^2$, since otherwise $x^{m-1}=x^m=0$ would hold.

Let further $\psi((x))=0$ be the \emph{reduced} equation satisfied by the quantity $x$ of the group $(\eps)$. This is the equation of lowest degree of the form $px^0+qx+rx^2+\cdots=0$ or $px+qx^2+rx^3+\cdots=0$, according as $(\eps)$ has a principal unit $x^0=e$ or not. Let $\psi(u)=(u-a)^\alpha\chi(u)$ and let $\chi(a)$ be nonzero. Determine (\emph{On Commuting Matrices}, Sitzungsberichte 1896, § 3) an integral function $f(u)$ such that $f(u)-1$ is divisible by $(u-a)^\alpha$. If $(\eps)$ possesses a principal unit, let $f(u)$ also be divisible by $\chi(u)$; if this is not the case and $a$ is nonzero, let $f(u)$ be divisible by $u\chi(u)$. Then $f(u)^2-f(u)$ is divisible by $\psi(u)$, but $f(u)$ itself is not. Thus, if $y=f((x))$, then $y^2=y$ and $y$ is nonzero. Accordingly, if $(\eps)$ is not a root group, it contains a unit.

If $(\eps)$ has a principal unit, and if
\[
 \psi(u)=(u-a)^\alpha(u-b)^\beta(u-c)^\gamma\cdots,
\]
then to the different roots $a,b,c,\ldots$ of $\psi(u)$ there correspond, in the indicated way, the integral functions $f(u),g(u),h(u),\ldots$. Then $f(u)g(u)$ and $f(u)+g(u)+h(u)+\cdots-1$ are divisible by $\psi(u)$, and hence
\begin{equation}\tag{3.}
 f((x))^2=f((x)),\qquad f((x))g((x))=0
\end{equation}
and
\begin{equation}\tag{4.}
 f((x))+g((x))+h((x))+\cdots=e.
\end{equation}

In this way (\emph{C.} 17) the principal unit $e$ can be decomposed into as many independent units as the equation $\psi(u)=0$ has distinct roots.

\begin{quote}\itshape
III. If the prime factors of the determinant of a group with principal unit are all linear, then its radical consists of all roots of zero contained in the group.
\end{quote}

For if $x$ and $y$ are two quantities of the group $(\eps)$, then
\[
 |S(ux+vy+we)|=\prod (uu_\alpha+vv_\alpha+w)
\]
is a product of linear factors. Setting $v=0$ or $u=0$, one recognizes that $u_1,\ldots,u_n$ are the characteristic roots of $S(x)$ and $v_1,\ldots,v_n$ those of $S(y)$. If $y$ is a root of zero, then $v_1=\cdots=v_n=0$. If $z=g((y))$, the characteristic roots of $S(z)$ are therefore all equal to $g(0)$. Thus, if $v$ is nonzero and $z=u(y+ve)^{-1}$, they are all equal to $\frac{u}{v}$. Consequently
\[
 |S(x+u(y+ve)^{-1})|=\prod\left(u_\alpha+\frac{u}{v}\right),\qquad |S(y+ve)|=v^n,
\]
and, since
\[
 S(x+u(y+ve)^{-1})S(y+ve)=S(xy+vx+ue),
\]
one has
\[
 |S(xy+vx+ue)|=\prod(u_\alpha v+u).
\]
Since both sides of this equation are integral functions of $v$, it remains valid for the previously excluded value $v=0$. Thus $|S(xy+ue)|=u^n$, and consequently (\emph{C.} 28) $xy$ is a root of zero and $y$ is a root quantity of $(\eps)$.

The converse of this theorem follows from the results obtained in § 5, according to which, when $r>1$ there, $\eps_{12}^2=0$, but $\eps_{12}$ does not belong to the radical of $(\eps)$.

\section*{§ 4.}

\begin{quote}\itshape
I. In every root group there is a nonzero quantity $x$ which satisfies the equations $xy=yx=0$ whenever $y$ is an arbitrary quantity of the group.
\end{quote}

For this theorem Mr. \textsc{Cartan} (nos. 31--34) gives a very acute but somewhat circumstantial proof, which I shall replace here by an essentially simpler one. The theorem can also be expressed in the following way:

\begin{quote}\itshape
II. Every root group contains an invariant subgroup of order 1.
\end{quote}

That the two theorems are identical follows from a lemma, much used by Mr. \textsc{Cartan} (\emph{C.} 29), in which $x$ and $y$ denote two quantities of a group with or without principal unit:

\begin{quote}\itshape
III. If $xy=ax$ (or $yx=ax$), where $a$ is an ordinary quantity, and if $y$ is a root of zero, then either $x=0$ or $a=0$.
\end{quote}

By hypothesis there is a number $k$ such that $y^k=0$, and hence also a number $l$ such that $xy^l=0$. Suppose $x\neq0$, and let $m$ be the smallest number of this kind. Then $m>0$ and $xy^m=0$, but $xy^{m-1}$ (which, for $m=1$, is to be understood as $x$) is different from zero. From $xy^m=axy^{m-1}$ it follows that $a=0$. One obtains the same result by multiplying the equation $x(y-a)=0$ on the right by $y^{m-1}+y^{m-2}a+\cdots+a^{m-1}$. The theorem can be generalized as follows:

\begin{quote}\itshape
IV. If $x_1,x_2,\ldots,x_m$ are quantities of a root group, and if the product $x_1x_2\cdots x_m$ is nonzero, then the $m$ quantities
\[
 x_1,
 \quad x_1x_2,
 \quad x_1x_2x_3,
 \ldots
 \quad x_1x_2\cdots x_m
\]
are linearly independent.
\end{quote}

For suppose
\[
 ax_1\cdots x_l+bx_1\cdots x_lx_{l+1}+cx_1\cdots x_lx_{l+1}x_{l+2}+\cdots=0,
\]
where $a$ is the first coefficient different from zero. Then
\[
 bx_{l+1}+cx_{l+1}x_{l+2}+\cdots=-y
\]
is a quantity of the root group $(\eta)$, hence a root of zero, and
\[
 x_1\cdots x_l(a-y)=0.
\]
Therefore $x_1\cdots x_l=0$ and hence also $x_1\cdots x_m=0$.

\begin{quote}\itshape
V. In a root group of order $n-1$, every product of $n$ quantities vanishes.
\end{quote}

For if $x_1\cdots x_n$ were nonzero, then the $n$ quantities
\[
 x_1,
 \quad x_1x_2,
 \quad x_1x_2x_3,
 \ldots
 \quad x_1x_2\cdots x_n
\]
would be linearly independent, and so the order of the group would be $\geq n$.

By this theorem, for a root group $(\eta)$ there is an invariant number $m$ such that every product of $m$ quantities of $(\eta)$ vanishes, but not every product of $m-1$ quantities does. If $x$ is a nonvanishing product of $m-1$ quantities of $(\eta)$, then $xy$ and $yx$ are products of $m$ quantities, and therefore zero. Repeated application of this theorem gives (\emph{C.} 31):

\begin{quote}\itshape
VI. The fundamental units $\eta_1,\eta_2,\ldots,\eta_{n-1}$ of a root group of order $n-1$ can be chosen so that the product $\eta_\alpha\eta_\beta$ is a linear combination of fundamental units whose index is $>\alpha$ and $>\beta$.
\end{quote}

By theorem V there is furthermore, for a root group $(\eta)$, an invariant number $l$ $(\leq m)$, the smallest one such that the $l^{\ten}$ power of every quantity of the group vanishes. For a commutative group $l=m$. The following theorem also holds:

\begin{quote}\itshape
VII. If in a root group whose order is $n-1$ there are $n-1$ quantities whose product $x_1x_2\cdots x_{n-1}$ is nonzero, then $x_1^{n-1}$ is also nonzero, and the group is the commutative group formed from all integral functions of $x_1$.
\end{quote}

For by theorem IV the $n-1$ quantities
\[
 x_1,
 \quad x_1x_2,
 \quad x_1x_2x_3,
 \ldots
 \quad x_1x_2\cdots x_{n-1}
\]
are linearly independent and therefore form a basis of the group $(\eta)$. Hence
\[
 x_\kappa=a_{\kappa1}x_1+a_{\kappa2}x_1x_2+\cdots+a_{\kappa,n-1}x_1x_2\cdots x_{n-1}\quad(\kappa=1,2,\ldots,n-1).
\]
Since every product of $n$ quantities of the group $(\eta)$ vanishes, multiplying these $n-1$ equations yields
\[
 x_1x_2\cdots x_{n-1}=a_{11}a_{21}\cdots a_{n-1,1}x_1^{n-1},
\]
and thus $x_1^{n-1}$ is nonzero. Therefore, if one sets $x_1=x$, the quantities
\[
 x,
\quad x^2,
\ldots,
\quad x^{n-1}
\]
also form a basis of $(\eta)$. Adjoining to them $x^0=e$, one obtains a group $(\eps)$ of order $n$, consisting of the quantities $z=zx^0+z_1x+\cdots+z_{n-1}x^{n-1}$, whose determinant $|S(z)|=z^n$ contains only one elementary divisor.

\section*{§ 5.}

I now turn to the proof of \textsc{Cartan}'s theorem. Let $(\eps)$ be a group of order $n$ with principal unit, let $(\eta)$ be its radical, and let $n-m$ be its order. Then $m$ is the linear rank of the determinant of $(\eps)$
\begin{equation}\tag{1.}
 \Theta(x)=|S(x)|=\prod \Phi^{s}.
\end{equation}
Each of its $k$ prime factors can, by a suitable choice of the fundamental units, be brought to the form of a determinant
\begin{equation}\tag{2.}
 \Phi(x)=|x_{\alpha\beta}|\qquad(\alpha,\beta=1,2,\ldots,r).
\end{equation}
The
\begin{equation}\tag{3.}
 m=r^2+r'{}^2+r''{}^2+\cdots
\end{equation}
elements
\begin{equation}\tag{4.}
 x_{\alpha\beta},\qquad x'_{\alpha\beta},\qquad x''_{\alpha\beta},\ldots
\end{equation}
of the $k$ different determinants $\Phi,\Phi',\Phi'',\ldots$ of degrees $r,r',r'',\ldots$ are all independent variables. The determinant of the \textsc{Dedekind} group $(\grS)$ homomorphic to $(\eps)$ is $\prod \Phi^r$; the exponent $r$ of $\Phi$ is equal to the degree of the prime factor.

One can choose the $n$ fundamental units of $(\eps)$ so that $\eps_{m+1},\ldots,\eps_n$ form the basis of the radical $(\eta)$, and so that the quantities (4.) become the coordinates of $\eps_1,\ldots,\eps_m$. I therefore denote these $m$ fundamental units by
\begin{equation}\tag{5.}
 \eps_{\alpha\beta},\qquad \eps'_{\alpha\beta},\qquad \eps''_{\alpha\beta},\ldots,
\end{equation}
and the remaining $n-m$ by $\eta_1,\ldots,\eta_{n-m}$, so that
\begin{equation}\tag{6.}
 x=\sum x_{\alpha\beta}\eps_{\alpha\beta}
 +\sum x'_{\alpha\beta}\eps'_{\alpha\beta}
 +\sum x''_{\alpha\beta}\eps''_{\alpha\beta}
 +\cdots+\sum y_\rho\eta_\rho
\end{equation}
becomes an arbitrary quantity of $(\eps)$. The characteristic determinant of $S(x)$ is then
\begin{equation}\tag{7.}
 |S(ue-x)|=\prod |u e_{\alpha\beta}-x_{\alpha\beta}|^{s}.
\end{equation}
Between the fundamental units (5.) there exist the relations
\begin{equation}\tag{8.}
 \eps_{\alpha\beta}\eps_{\beta\gamma}\equiv\eps_{\alpha\gamma},
 \qquad
 \eps_{\alpha\delta}\eps_{\beta\gamma}\equiv0,
 \qquad
 \eps_{\alpha\delta}\eps'_{\beta\gamma}\equiv0
 \pmod{\eta},
\end{equation}
where in the second relation $\beta$ and $\delta$ are assumed to be different. It remains to show that the fundamental units (5.) can be modified modulo $(\eta)$ so that these equations hold absolutely. Let
\[
 z=a_1\eps_{11}+\cdots+a_r\eps_{rr}+a_{r+1}\eps'_{11}+\cdots+a_{r+r'}\eps'_{r'r'}+\cdots,
\]
or, if one sets
\begin{equation}\tag{9.}
 \eps_{11}=\eps_1,
 \ldots,
 \eps_{rr}=\eps_r,
 \qquad
 \eps'_{11}=\eps_{r+1},
 \ldots,
 \eps'_{r'r'}=\eps_{r+r'},\ldots,
\end{equation}
one has $z=\sum \eps_\lambda a_\lambda$. The
\begin{equation}\tag{10.}
 p=r+r'+r''+\cdots
\end{equation}
coordinates $a_\lambda$ may be chosen in any fixed way; only they must all be distinct from one another. Then
\[
 \varphi(u)=|S(ue-z)|=(u-a_1)^s\cdots(u-a_r)^s(u-a_{r+1})^{s'}\cdots(u-a_{r+r'})^{s'}\cdots,
\]
and, if $f(u)$ is an integral function of $u$, it follows from (8.) that
\[
 f((z))\equiv \sum \eps_\lambda f(a_\lambda)\pmod{\eta}.
\]
As in § 3, (3.), I now determine $p$ integral functions $f_\lambda(u)$ so that $(u-a_\lambda)^s f_\lambda(u)$ is divisible by $\varphi(u)$ and $f_\lambda(u)-1$ is divisible by $(u-a_\lambda)^s$. Then
\begin{equation}\tag{11$^{*}$.}
 f_\lambda((z))^2=f_\lambda((z)),
 \qquad
 f_\kappa((z))f_\lambda((z))=0
\end{equation}
and
\begin{equation}\tag{12$^{*}$.}
 \sum f_\lambda((z))=e.
\end{equation}
Furthermore $f_\lambda((z))\equiv \eps_\lambda$ (mod. $\eta$), and therefore one can change $\eps_\lambda$ (mod. $\eta$) so that $f_\lambda(z)=\eps_\lambda$. Then
\begin{equation}\tag{11.}
 \eps_{\alpha\alpha}^2=\eps_{\alpha\alpha},
 \qquad
 \eps_{\alpha\alpha}\eps_{\beta\beta}=0,
 \qquad
 \eps_{\alpha\alpha}\eps'_{\beta\beta}=0
\end{equation}
and
\begin{equation}\tag{12.}
 \sum \eps_\lambda=e.
\end{equation}

Now, by (8.),
\[
 \eps_{\alpha\beta}\equiv\eps_{\alpha\alpha}\eps_{\alpha\beta}\eps_{\beta\beta},
 \qquad
 \eps'_{\alpha\beta}\equiv\eps'_{\alpha\alpha}\eps'_{\alpha\beta}\eps'_{\beta\beta},\ldots
 \pmod{\eta}.
\]
If $\alpha=\beta$, then the equations hold absolutely. If $\alpha$ is different from $\beta$, replace $\eps_{\alpha\beta}$ by $\eps_{\alpha\alpha}\eps_{\alpha\beta}\eps_{\beta\beta}$. This changes that fundamental unit (mod. $\eta$) so that, by (11.), now
\begin{equation}\tag{13.}
 \eps_{\alpha\alpha}\eps_{\alpha\beta}=\eps_{\alpha\beta},
 \qquad
 \eps_{\alpha\beta}\eps_{\beta\beta}=\eps_{\alpha\beta}
\end{equation}
and
\begin{equation}\tag{14.}
 \eps_{\alpha\beta}\eps_{\delta\gamma}=0,
 \qquad
 \eps_{\alpha\beta}\eps'_{\delta\gamma}=0
\end{equation}
hold, provided in the first equation $\delta$ is different from $\beta$.

I do not further alter the fundamental units $\eps_{\alpha\alpha},\eps'_{\alpha\alpha},\ldots$, but I do alter the fundamental units $\eps_{\alpha\beta},\eps'_{\alpha\beta},\ldots$, in such a way that the relations already obtained remain in force. If $\eta$ is an arbitrary quantity of $(\eps)$, I call
\begin{equation}\tag{15.}
 \eps_\kappa\eta\eps_\lambda=\eta_{\kappa\lambda}
\end{equation}
a quantity of type $(\kappa,\lambda)$. It satisfies the equations
\begin{equation}\tag{16.}
 \eps_\kappa\eta_{\kappa\lambda}=\eta_{\kappa\lambda},
 \qquad
 \eta_{\kappa\lambda}\eps_\lambda=\eta_{\kappa\lambda},
\end{equation}
and, if $\mu$ is different from $\kappa$ and $\nu$ from $\lambda$,
\begin{equation}\tag{17.}
 \eps_\mu\eta_{\kappa\lambda}=0,
 \qquad
 \eta_{\kappa\lambda}\eps_\nu=0.
\end{equation}
If $\eta$ is a quantity of the invariant subgroup $(\eta)$, then $\eta_{\kappa\lambda}$ also belongs to the radical. Therefore one can change $\eps_{\alpha\beta}$ by an arbitrary root quantity $\eta_{\alpha\beta}$ of type $(\alpha,\beta)$ without changing equations (13.) and (14.).

Now $\eps_{\alpha\beta}\eps_{\beta\alpha}=\eps_{\alpha\alpha}$, and hence
\begin{equation}\tag{18.}
 \eps_{\alpha\beta}\eps_{\beta\alpha}=\eps_{\alpha\alpha}-\eta_{\alpha\alpha},
\end{equation}
where $\eta_{\alpha\alpha}$ is a quantity of $(\eta)$, and indeed, by (13.), one of type $(\alpha,\alpha)$. Hence the right hand side is equal to $\eps_{\alpha\alpha}(e-\eta_{\alpha\alpha})$. From the equation
\[
 e-x^l=(e-x)(x^0+x+\cdots+x^{l-1})
\]
it follows that
\[
 e=(e-x)(x^0+x+\cdots+x^{l-1}),
\]
if $x^l=0$. If now $\eta_{\alpha\alpha}^l=0$, multiply equation (18.) on the right by
\[
 \eta_{\alpha\alpha}^{0}+\eta_{\alpha\alpha}+\cdots+\eta_{\alpha\alpha}^{l-1}.
\]
One obtains
\[
 \eps_{\alpha\beta}(\eps_{\beta\alpha}+\eta_{\beta\alpha})=\eps_{\alpha\alpha},
\]
where
\[
 \eta_{\beta\alpha}=\eps_{\beta\alpha}(\eta_{\alpha\alpha}+\eta_{\alpha\alpha}^2+\cdots+\eta_{\alpha\alpha}^{l-1})
\]
is a root quantity of type $(\beta,\alpha)$. Thus, if one changes $\eps_{\beta\alpha}$ by $\eta_{\beta\alpha}$, one obtains
\begin{equation}\tag{19.}
 \eps_{\alpha\beta}\eps_{\beta\alpha}=\eps_{\alpha\alpha}.
\end{equation}
But then one always also has (\emph{C.} 52.)
\begin{equation}\tag{20.}
 \eps_{\beta\alpha}\eps_{\alpha\beta}=\eps_{\beta\beta}.
\end{equation}
For at first $\eps_{\beta\alpha}\eps_{\alpha\beta}=\eps_{\beta\beta}-\eta_{\beta\beta}$, and from this, as before, one obtains $\eps_{\beta\alpha}(\eps_{\alpha\beta}+\eta_{\alpha\beta})=\eps_{\beta\beta}$; multiplying on the left by $\eps_{\alpha\beta}$ gives $\eps_{\alpha\beta}+\eta_{\alpha\beta}=\eps_{\alpha\beta}$, and therefore $\eta_{\alpha\beta}=0$.

Now choose $\eps_{12},\ldots,\eps_{1r}$ in any manner (\emph{C.} 67) consistent with the relations established so far, and then choose $\eps_{21},\ldots,\eps_{r1}$ so that $\eps_{1\alpha}\eps_{\alpha1}=\eps_{11}$, whence also $\eps_{\alpha1}\eps_{1\alpha}=\eps_{\alpha\alpha}$. Then the equation
\begin{equation}\tag{21.}
 \eps_{\alpha\beta}=\eps_{\alpha1}\eps_{1\beta}
\end{equation}
holds if $\alpha=1$, or $\beta=1$, or $\alpha=\beta$, that is, in the cases where $\eps_{\alpha\beta}$ has already been disposed of. But if $\alpha$ and $\beta$ are different and both are $>1$, then $\eps_{\alpha\beta}=\eps_{\alpha1}\eps_{1\beta}+\eta_{\alpha\beta}$, where $\eta_{\alpha\beta}$ is a root quantity of type $(\alpha,\beta)$. Therefore one can change $\eps_{\alpha\beta}$ so that equation (21.) is satisfied. Then
\[
 \eps_{\alpha\beta}\eps_{\beta\gamma}
 =\eps_{\alpha1}(\eps_{1\beta}\eps_{\beta1})\eps_{1\gamma}
 =\eps_{\alpha1}\eps_{11}\eps_{1\gamma}
 =\eps_{\alpha1}\eps_{1\gamma}
 =\eps_{\alpha\gamma}.
\]
The equations (8.) now all hold absolutely, and hence the $m$ quantities (5.) are the fundamental units of a group $(\grS)$.

By (12.) one has
\[
 \eta=(\eps_1+\cdots+\eps_p)\eta(\eps_1+\cdots+\eps_p)=\sum(\eps_\kappa\eta\eps_\lambda).
\]
Therefore all root quantities can be assembled from the $(n-m)p^2$ quantities $\eps_\kappa\eta_i\eps_\lambda$. If from these one chooses an independent system, then each of the fundamental units of the radical also belongs to a definite type $(\kappa,\lambda)$. In the subgroup $(\grS)$, $\eps_{\alpha\beta}$ has type $(\alpha,\beta)$ if $\alpha$ and $\beta$ are two of the numbers from 1 to $r$. By contrast it contains, for example, no quantity of type $(1,r+1)$.

If the $n$ fundamental units of $(\eps)$ are chosen in the indicated way, then by (17.) all quantities $\xi_{\kappa\lambda}$ of type $(\kappa,\lambda)$ can be expressed only by the fundamental units which have the same type. The fundamental units of type $(1,1)$ form the basis of a group for which $\eps_1=\eps_{11}$ is the principal unit and at the same time the sole unit. The remaining fundamental units of type $(1,1)$ form the basis of its radical. If $\alpha$ and $\beta$ are two of the numbers $1,2,\ldots,r$, then of the two equations
\[
 \eps_{\alpha\beta}\xi_{\beta\lambda}=\xi_{\alpha\lambda},
 \qquad
 \eps_{\beta\alpha}\xi_{\alpha\lambda}=\xi_{\beta\lambda}
\]
each is a consequence of the other. Therefore (\emph{C.} 57) one obtains all quantities of type $(\beta,\lambda)$ by multiplying $\eps_{\beta\alpha}$ with all quantities of type $(\alpha,\lambda)$.

Furthermore, all fundamental units of the types
\[
\begin{gathered}
 (1,1),(1,r+1),(1,r+r'+1),\ldots\ (r+1,1)(r+1,r+r),(r+1,r+r'+1),\ldots \\
 (r+r'+1,1),(r+r'+1,r+1),(r+r'+1,r+r'+1),\ldots
\end{gathered}
\]
together form the basis of a group $(\eps')$; the $k$ quantities among them which are contained in $(\grS)$, namely $\eps_1,\eps_{r+1},\eps_{r+r'+1},\ldots$, form the basis of a commutative \textsc{Dedekind} group $(\grS')$, and the remaining ones form the basis of a root group $(\eta')$, which is an invariant subgroup of $(\eps')$. By § 1, therefore, $(\eta')$ is the radical of $(\eps')$. Since $(\grS')$ is a commutative group, the determinant of $(\eps')$ splits into linear factors only. Mr. \textsc{Cartan} (\emph{C.} 57--60) has described so fully how the relations between the fundamental units of this group $(\eps')$ allow one to derive those between the fundamental units of $(\eps)$ that I do not wish to return to it here.
% ===== END WORK 071 =====

\end{document}
